% AUTO-GENERATED by scripts/build_topics.py — do not edit.
% Edit topics/bodies/*.tex or topics/preamble.tex instead.
\documentclass[11pt,openany]{report}
\usepackage[a4paper, margin=1in, top=1.1in, bottom=1.1in]{geometry}
% ---------------------------------------------------------------------------
%  Shared preamble for the "Вероятности и Статистика" lecture notes.
%
%  Used by both builds:
%    * lecture_NN.tex   — standalone article, one lecture   (\section = top level)
%    * lectures_full.tex — collected book, report class     (\chapter = per lecture)
%
%  The driver sets \numberwithin{thm}{section|chapter} after \input-ing this.
% ---------------------------------------------------------------------------

\usepackage{fontspec}
% The whole book is Bulgarian, so the main face has to carry Cyrillic properly
% rather than merely contain the glyphs. PT Serif was designed at ParaType for
% Cyrillic and is the right first choice; the fallbacks are ordered by how well
% they set Bulgarian, not by how common they are.
% Ligatures=TeX is not optional here. Without it fontspec leaves `` '' as
% literal backticks and --- as three hyphens on the page, which is what the
% whole book had been doing. It is set on every text family below for the same
% reason; the MONOSPACE family is deliberately left without it, so that code in
% \texttt and verbatim keeps -- and quotes exactly as typed.
\IfFontExistsTF{PT Serif}
  {\setmainfont{PT Serif}[Ligatures=TeX]}
  {\IfFontExistsTF{Times New Roman}
    {\setmainfont{Times New Roman}[Ligatures=TeX]}
    {\setmainfont{Liberation Serif}[Ligatures=TeX]}}

% polyglossia will not typeset a single Cyrillic word until it is told which
% font to use for the script — it does not infer it from \setmainfont, and the
% error it raises ("does not contain the Cyrillic script") is about the
% declaration being absent, not about the font lacking the glyphs.
% All three families need declaring, not just the roman one: \texttt on a
% Cyrillic word fails exactly as \textrm does, and folder names and identifiers
% in this book are Bulgarian. PT Serif/Sans/Mono are one designed family.
\IfFontExistsTF{PT Serif}
  {\newfontfamily\cyrillicfont{PT Serif}[Ligatures=TeX]
   \newfontfamily\cyrillicfontsf{PT Sans}[Ligatures=TeX]
   \newfontfamily\cyrillicfonttt{PT Mono}[Scale=MatchLowercase]
   \setmonofont{PT Mono}[Scale=MatchLowercase]}
  {\newfontfamily\cyrillicfont{Times New Roman}[Ligatures=TeX]
   \newfontfamily\cyrillicfonttt{Menlo}[Scale=MatchLowercase]
   \setmonofont{Menlo}[Scale=MatchLowercase]}

% Bulgarian: loads hyphenation patterns (without these TeX cannot break a single
% Cyrillic word, which is what produced the loose lines and overfull boxes) and
% translates the auto-generated captions ("Contents" -> "Съдържание").
\usepackage{polyglossia}
\setmainlanguage{bulgarian}
\setotherlanguage{english}

\usepackage{amsmath}
\usepackage{amssymb}
\usepackage{amsfonts}
\usepackage{amsthm}
% Semantic interpretation brackets used in the predicate-logic chapters.
% Keep these local instead of pulling a second symbol font into the PDF.
\newcommand{\llbracket}{[\![}
\newcommand{\rrbracket}{]\!]}
% Shared probability notation. Bodies use these semantic commands exclusively.
\newcommand{\E}{\mathbb{E}}
\newcommand{\Var}{\mathbb{D}}
\newcommand{\ind}{\mathbf{1}}
\newcommand{\given}{\mid}
\DeclareMathOperator{\Bin}{Bi}
\DeclareMathOperator{\Poi}{Po}
\DeclareMathOperator{\Ber}{Ber}
\DeclareMathOperator{\Geo}{Ge}
\DeclareMathOperator{\NegBin}{NB}
\DeclareMathOperator{\Hyp}{HG}
\DeclareMathOperator{\Exponential}{Exp}
\DeclareMathOperator{\Unif}{U}
\DeclareMathOperator{\Cov}{cov}
% Standalone lectures do not load labels from other lecture bodies. The book
% driver upgrades cross-lecture mentions to links after all shared setup.
\newcommand{\crosslecture}[2]{#2}
\usepackage[most]{tcolorbox}
\usepackage{microtype}
\usepackage{fancyhdr}
% Long Bulgarian chapter titles wrap in the running head. Reserve enough room
% so fancyhdr does not silently change \topmargin from page to page.
\setlength{\headheight}{24pt}

% "г.", "т.е.", "п.с." are abbreviations, not sentence ends — suppress the
% extra inter-sentence space LaTeX would otherwise insert after them.
\frenchspacing

\linespread{1.08}
\setlength{\parskip}{0.4em}
\setlength{\parindent}{0pt}
\clubpenalty=10000       % no orphans
\widowpenalty=10000      % no widows
\displaywidowpenalty=10000

% Long unbreakable inline-math runs (\{1,2,\dots,n\}, \text{...} spans) cannot be
% hyphenated. Let TeX stretch a line as a last resort rather than let it stick
% out into the margin.
\setlength{\emergencystretch}{3em}

% --- palette ---------------------------------------------------------------
% Three reading tiers, so the page can be skimmed:
%   gold + thick bar + star = landmark result, memorise it
%   blue / green boxes      = ordinary definition / result
%   hairline rule, no fill  = example or proof, skippable on a first pass
% Each tier is distinguished by RULE WEIGHT and FILL as well as by hue, so the
% hierarchy still reads when the document is printed in greyscale.
\definecolor{defframe}{RGB}{38,90,160}
\definecolor{defback}{RGB}{234,241,250}
\definecolor{thmframe}{RGB}{34,120,70}
\definecolor{thmback}{RGB}{233,246,238}
\definecolor{keyframe}{RGB}{176,124,18}
\definecolor{keyback}{RGB}{253,247,232}
\definecolor{exrule}{RGB}{118,130,146}
\definecolor{proofrule}{RGB}{198,203,212}
\definecolor{suppframe}{RGB}{124,96,150}

% --- figures ---------------------------------------------------------------
% Native TikZ: vector, no external assets, and it can reuse the palette above so
% figures read as part of the same system. Figures are NON-FLOATING: where the
% picture is the argument (the reflection principle), a float that drifts to the
% next page breaks the proof. Meaning is carried by line weight and dash pattern
% as well as colour, so the figures survive greyscale like the boxes do.
\usepackage{tikz}
\usetikzlibrary{arrows.meta, decorations.pathreplacing}
\usepackage{booktabs}   % \toprule/\midrule/\bottomrule in the statistical tables
\usepackage{longtable}  % the Phi table is taller than a page can hold unbroken
\usepackage{caption}
\captionsetup{font=small, labelfont=bf, skip=6pt}
% NB: put \label{} INSIDE the caption argument -- \captionof steps the figure
% counter, so a label placed before it would capture the wrong number.
\newenvironment{lecfigure}[1][]
  {\par\medskip\noindent\begin{minipage}{\linewidth}\centering
   \def\lecfigcap{#1}}
  {\ifx\lecfigcap\empty\else\captionof{figure}{\lecfigcap}\fi
   \end{minipage}\par\medskip}
\tikzset{
  axisline/.style={gray!70, -{Stealth[length=2mm]}, line width=0.4pt},
  walk/.style={thmframe, line width=1pt, line join=round},
  walkrefl/.style={keyframe, line width=1pt, line join=round,
                   dash pattern=on 3pt off 2pt},
  cdfstep/.style={defframe, line width=1pt},
  guide/.style={gray!55, line width=0.4pt, dash pattern=on 2pt off 2pt},
}

% --- numbered statement environments ---------------------------------------
% One shared counter for all five, so a reader scanning for "Твърдение 4.7"
% never has to guess which of five sequences it belongs to.
\newtheoremstyle{lecstyle}
  {\topsep}{\topsep}     % space above / below
  {\normalfont}          % body font — upright; italic Cyrillic at 11pt is hard to read
  {0pt}                  % indent
  {\bfseries}            % head font
  {.}                    % punctuation after head
  {0.5em}                % space after head
  {}                     % default head spec
\theoremstyle{lecstyle}

\newtheorem{thm}{Теорема}
\newtheorem{defn}[thm]{Дефиниция}
\newtheorem{prop}[thm]{Твърдение}
\newtheorem{lem}[thm]{Лема}
\newtheorem{cor}[thm]{Следствие}

% Landmark variants. Same counter, same printed name — the mathematical status
% is identical; only the study weight differs. \keythm{...} reads exactly like
% \thm{...} to a reader who ignores the colour.
\newtheoremstyle{leckey}
  {\topsep}{\topsep}{\normalfont}{0pt}{\bfseries}{.}{0.5em}
  {\textcolor{keyframe}{\ensuremath{\bigstar}}\ \thmname{#1}\thmnumber{ #2}\thmnote{ (#3)}}
\theoremstyle{leckey}
\newtheorem{keythm}[thm]{Теорема}
\newtheorem{keydefn}[thm]{Дефиниция}
\newtheorem{keylem}[thm]{Лема}
\newtheorem{keyaxiom}[thm]{Аксиома}
\newtheorem{keyscheme}[thm]{Схема}

\newtheoremstyle{lecremark}
  {\topsep}{\topsep}{\normalfont}{0pt}{\itshape}{.}{0.5em}{}
\theoremstyle{lecremark}
\newtheorem{remark}[thm]{Забележка}

% Material that is NOT part of the recorded 2021/22 lectures — taken from the
% official course documents. Dashed violet frame and a dagger, so it can never
% be mistaken for something the lecturer said.
\newtheoremstyle{lecsupp}
  {\topsep}{\topsep}{\normalfont}{0pt}{\bfseries\color{suppframe}}{.}{0.5em}
  {\color{suppframe}\ensuremath{\dagger}\ \thmname{#1}\thmnumber{ #2}\thmnote{ (#3)}}
\theoremstyle{lecsupp}
\newtheorem{supp}[thm]{Допълнение}

% Examples are marked to be SKIPPED on a first pass: recessive slate head with a
% ▷ lead-in, no fill, thin left rule.
\newtheoremstyle{lecexample}
  {\topsep}{\topsep}{\normalfont}{0pt}{\bfseries\color{exrule}}{.}{0.5em}
  {\color{exrule}\ensuremath{\triangleright}\ \thmname{#1}\thmnumber{ #2}\thmnote{ (#3)}}
\theoremstyle{lecexample}
\newtheorem{example}[thm]{Пример}

% Definitions in blue, everything provable in green. Proofs stay outside.
% `lines before break` is what keeps a statement whole. tcolorbox first asks
% whether the box fits in the space left on the page; only if it does not does
% this threshold apply, and then it ejects the page and re-asks on a fresh one.
% Set just under a full text height (684pt / 14.7pt per line = 46.6 lines), so
% "does not fit here" always becomes "start it at the top of the next page".
% A box that is genuinely taller than a page still breaks -- but from the top of
% a page, never with a two-line stub stranded under the previous paragraph.
\newcommand{\lecbreakminlines}{45}
% ...but only when ejecting the page would actually achieve something. A box
% that is taller than the text block breaks no matter which page it starts on,
% so pushing it forward buys nothing and strands up to three quarters of a page.
% tcolorbox has just measured the box into \tcb@h@total when this runs, so the
% question can be answered exactly: over-tall boxes keep the default threshold
% and start where they stand, filling the page to its foot before they continue.
\usepackage{etoolbox}
\makeatletter
\apptocmd{\tcb@comp@h@total@standalone}
  {\ifdim\tcb@h@total>\textheight\relax\def\kvtcb@breakminlines{2}\fi}
  {}{\PackageError{preamble}{could not patch tcolorbox height hook}{}}
\makeatother
\tcbset{lecbox/.style={enhanced, breakable, lines before break=\lecbreakminlines,
    arc=2.4mm, boxrule=0.7pt,
    left=2.6mm, right=2.6mm, top=1.8mm, bottom=1.8mm, before skip=0.6em,
    after skip=0.6em}}
\tcolorboxenvironment{defn}{lecbox, colback=defback, colframe=defframe}
\tcolorboxenvironment{thm} {lecbox, colback=thmback, colframe=thmframe}
\tcolorboxenvironment{prop}{lecbox, colback=thmback, colframe=thmframe}
\tcolorboxenvironment{lem} {lecbox, colback=thmback, colframe=thmframe}
\tcolorboxenvironment{cor} {lecbox, colback=thmback, colframe=thmframe}

% Landmark tier: gold, plus a heavy left bar so it is still unmistakable in
% greyscale and when flicking through at speed.
\tcbset{leckeybox/.style={lecbox, colback=keyback, colframe=keyframe,
    leftrule=2.2mm}}
\tcolorboxenvironment{keythm} {leckeybox}
\tcolorboxenvironment{keydefn}{leckeybox}
\tcolorboxenvironment{keylem} {leckeybox}
\tcolorboxenvironment{keyaxiom}{leckeybox}
\tcolorboxenvironment{keyscheme}{leckeybox}

% Supplement tier: dashed violet frame, no fill — visibly outside the lecture record.
\tcolorboxenvironment{supp}{enhanced, breakable, lines before break=\lecbreakminlines,
  colback=white,
  colframe=suppframe, boxrule=0.5pt, frame style={dash pattern=on 3pt off 2pt},
  arc=2.4mm, left=2.6mm, right=2.6mm, top=1.8mm, bottom=1.8mm,
  before skip=0.6em, after skip=0.6em}

% Skippable tier: no fill, no frame, just a marginal rule the eye can run past.
% Examples get a visible slate rule; proofs get a fainter one still.
\tcolorboxenvironment{example}{blanker, breakable,
  borderline west={1.2pt}{0pt}{exrule},
  left=4mm, top=1mm, bottom=1mm, before skip=0.7em, after skip=0.7em}
\tcolorboxenvironment{proof}{enhanced, breakable,
  colback=white, colframe=proofrule, boxrule=0.6pt, arc=1.5mm,
  left=2.6mm, right=2.6mm, top=1.4mm, bottom=1.4mm,
  before skip=0.7em, after skip=0.7em}

% Proofs have a complete frame and an automatic flush-right QED.
% The driver localizes \proofname after \begin{document}, once polyglossia has
% installed its Bulgarian captions.
\renewcommand{\qedsymbol}{$\square$}

% --- "По конспекта" opening panel -------------------------------------------
% The colour tiers above say what KIND of mathematical object a box holds. They
% cannot say what the official annotation ASKS FOR — a definition, a statement
% with no proof, a full proof, a construction, pseudocode, an implementation, a
% complexity analysis, a worked task type or a drawing. That is a different
% axis, and it is the one a student revising against the конспект needs first,
% so it is spelled out in words at the top of every chapter rather than encoded
% in a colour. Slate, deliberately quieter than the three content tiers: this
% panel is a map, not material to be learnt.
\definecolor{konframe}{RGB}{88,96,110}
\definecolor{konback}{RGB}{244,245,247}
\newtcolorbox{konspekt}{enhanced, breakable, lines before break=6,
  colback=konback, colframe=konframe, boxrule=0.7pt, arc=2.4mm,
  left=2.6mm, right=2.6mm, top=1.4mm, bottom=1.6mm,
  before skip=0.4em, after skip=1.1em,
  title={\normalfont\bfseries\small По конспекта: какво трябва да мога},
  coltitle=white, colbacktitle=konframe}

% One obligation per line: the required ACTION first, then where the chapter
% discharges it. "--- липсва" is a legitimate right-hand side; a panel that
% cannot honestly point somewhere must say so rather than quietly omit the row.
\newcommand{\kreq}[2]{\par\noindent\textbf{#1:} #2}

% Panel targets are LINKS, not printed strings: the reader taps "§3.7, стр. 60"
% and lands there. \ref/\pageref rather than a typed number, so a section that
% moves cannot leave the panel pointing at the wrong place — and check_refs.py
% fails the build if a target ever disappears. The starred forms keep the whole
% phrase inside ONE \hyperref instead of nesting links.
\newcommand{\kref}[1]{\hyperref[s:#1]{\S\ref*{s:#1}, стр.~\pageref*{s:#1}}}
\newcommand{\kr}[1]{\hyperref[s:#1]{\S\ref*{s:#1}}}
\newcommand{\krange}[2]{\hyperref[s:#1]{\S\ref*{s:#1}--\S\ref*{s:#2}, стр.~\pageref*{s:#1}}}
\newcommand{\kfig}[1]{\hyperref[f:#1]{фиг.~\ref*{f:#1}, стр.~\pageref*{f:#1}}}

% The конспект's own exam-selection notes (questions 6, 7, 8, 18). Gold, like
% the landmark tier, because missing one changes what a student walks in with.
% It says what may be DRAWN on the exam, not what may be skipped in revision.
\newcommand{\kselect}[1]{\par\medskip\noindent
  \begin{tcolorbox}[enhanced, breakable, colback=keyback, colframe=keyframe,
    boxrule=0.7pt, leftrule=2.2mm, arc=1.5mm,
    left=2.2mm, right=2.2mm, top=1.2mm, bottom=1.2mm,
    before skip=0.2em, after skip=0.4em]
  \textcolor{keyframe}{\ensuremath{\bigstar}}\ \textbf{Бележка на конспекта:} #1
  \end{tcolorbox}}

% Material the chapter carries that the annotation does NOT require at this
% depth. Same violet dagger as the \supp tier, so the two agree on sight.
\newcommand{\kextra}[1]{\par\noindent
  \textcolor{suppframe}{\ensuremath{\dagger}\ \textbf{Допълнително:}}\ #1}

% Prerequisites: recalled here, examined elsewhere or not at all.
\newcommand{\kprereq}[1]{\par\noindent\textbf{Предварително (не се изпитва тук):}\ #1}

% Honest coverage status. Never a "complete" badge earned by mere presence.
\newcommand{\kstatus}[1]{\par\medskip\noindent\textit{Статус на покритието:}\ #1}

% --- links ------------------------------------------------------------------
\usepackage{hyperref}   % load last
\hypersetup{
  colorlinks=true, linkcolor=defframe, urlcolor=defframe, citecolor=defframe,
  linktoc=all, bookmarksnumbered=true, bookmarksopen=true, unicode=true,
  bookmarksdepth=3,   % printed TOC stays lean; the PDF outline keeps full depth
  pdfauthor={}}
% A long unbroken URL in the bibliography sticks out into the margin, because
% \url may only break after a slash by default and a 32-character hash offers
% no slash at all. xurl lifts that restriction; the \Urlmuskip lets the glue
% around the break stretch instead of producing an overfull line.
\usepackage{xurl}
\Urlmuskip=0mu plus 1mu\relax

\numberwithin{thm}{chapter}
\renewcommand{\chaptername}{Въпрос}
\hypersetup{pdftitle={Държавен изпит по Компютърни науки — записки}}
\pagestyle{fancy}\fancyhf{}
\fancyhead[L]{\small\itshape\leftmark}
\fancyhead[R]{\small\thepage}
\renewcommand{\headrulewidth}{0.4pt}
\begin{document}
\renewcommand{\proofname}{Доказателство}
\begin{titlepage}
\centering\vspace*{3cm}
{\Huge\bfseries Държавен изпит\par}\vspace{0.6cm}
{\LARGE Компютърни науки\par}\vspace{2cm}
{\large Записки по конспекта от 30.06.2025\par}
\vspace{0.4cm}
{\small ФМИ, Софийски университет „Св. Климент Охридски“\par}
\vfill
\end{titlepage}
\tableofcontents
\cleardoublepage
\part{Основи на компютърните науки}

\chapter{Множества. Декартово произведение. Релации. Функции.}
%% Drafted from the mapped corpus by scripts/draft_topics.py.
% Model: gpt-5.6-terra; source characters: 17554.
\section{Множества и аксиоматичен подход}

Понятието \emph{множество} се приема за първично: използва се без дефиниция. Елементите на множество $A$ се означават с $x\in A$, а твърдението, че $x$ не е негов елемент, с $x\notin A$. \begin{defn}[Подмножество]
Множеството $A$ е подмножество на $B$, означено с $A\subseteq B$, когато всеки елемент на $A$ е елемент и на $B$:
\[
A\subseteq B\quad\Longleftrightarrow\quad
(\forall x)\,(x\in A\Rightarrow x\in B).
\]
\end{defn}

\begin{keyaxiom}[За обема]
Две множества $A$ и $B$ са равни, ако и само ако имат едни и същи елементи:
\[
A=B\quad\Longleftrightarrow\quad(\forall x)\,(x\in A\Leftrightarrow x\in B).
\]
Формалният запис е
\[
(\forall X)(\forall Y)\left(
(\forall z)\,(z\in X\Leftrightarrow z\in Y)\Rightarrow X=Y
\right).
\]
\end{keyaxiom}

\subsection{Аксиоми за построяване на множества}

\begin{keyscheme}[За отделянето]
Нека $A$ е множество и $\varphi(x,u_1,\ldots,u_n)$ е свойство. Тогава съществува множество, състоящо се точно от онези елементи на $A$, които удовлетворяват това свойство:
\[
B=\{x\in A\mid \varphi(x,u_1,\ldots,u_n)\}.
\]
Формално:
\[
(\forall u_1)\cdots(\forall u_n)(\forall A)(\exists B)(\forall x)
\left(x\in B\Leftrightarrow
\bigl(x\in A\land\varphi(x,u_1,\ldots,u_n)\bigr)\right).
\]
\end{keyscheme}

Схемата за отделянето позволява да се образуват подмножества на вече дадено множество чрез зададено условие. Например множеството на четните естествени числа, които са не по-големи от $20$, може да се запише като
\[
\{x\in\mathbb N\mid x\leq 20\text{ и }x\text{ е четно}\}.
\]

\begin{keyaxiom}[За степенното множество]
За всяко множество $A$ съществува множество от всички негови подмножества, наречено \emph{степенно множество} на $A$ и означавано с $\mathcal P(A)$:
\[
\mathcal P(A)=\{X\mid X\subseteq A\}.
\]
\end{keyaxiom}

Например
\[
\mathcal P(\varnothing)=\{\varnothing\},
\qquad
\mathcal P(\{a,b\})=
\{\varnothing,\{a\},\{b\},\{a,b\}\}.
\]
За крайно множество $A$ с $|A|=n$ е валидно
\[
|\mathcal P(A)|=2^{|A|}=2^n.
\]

\begin{keyaxiom}[За индуктивно генерираните множества]
Нека $U$ е множество, $\varnothing\ne M_0\subseteq U$, а $F$ е множество от операции $f:U^{k_f}\to U$ с крайна положителна арност $k_f$. Тогава множеството $M$ се строи по следния начин:

\begin{enumerate}
\item \emph{База:} $M_0\subseteq M$, т.е. в $M$ са всички елементи на $M_0$.
\item \emph{Индуктивно предположение:} нека $M_i$ е вече построеното на $i$-тата стъпка множество. То е непразно, защото $\varnothing\neq M_0\subseteq M_i$.
\item \emph{Индуктивна стъпка:} в $M_{i+1}$ включваме всички досегашни елементи и всички елементи, които се получават от тях чрез прилагане на операциите от $F$.
\item \emph{Заключение:} в $M$ няма други елементи освен базовите и получените чрез краен брой индуктивни стъпки.
\end{enumerate}

При едноместни операции конструкцията се записва формално като
\[
M_{i+1}=M_i\cup\{f(m)\mid f\in F,\ m\in M_i\},
\qquad
M=\bigcup_{i=0}^{\infty}M_i.
\]
При многоместна операция $f$ се прилагат всички допустими наредени групи от елементи на $M_i$. Построеното множество се означава с
\[
M:=\langle M_0,F\rangle.
\]
\end{keyaxiom}

Заключението изразява минималността на построеното множество: елементите му са точно базовите елементи и елементите, достижими чрез краен брой приложения на допустимите операции.

\begin{supp}[Бележка]
Обединението с $M_i$ в индуктивната стъпка е съществено: то запазва вече получените елементи при преминаването към следващото ниво.
\end{supp}

\section{Математическа индукция}

Математическата индукция е метод за доказване на твърдения за естествени числа. Нека $m\in\mathbb N$ и $\varphi(n)$ е твърдение, зависещо от $n$.

\begin{keythm}[Принцип на слабата математическа индукция]
Ако са изпълнени:

\begin{enumerate}
\item базата: $\varphi(m)$ е вярно;
\item индуктивната стъпка:
\[
(\forall k\geq m)\,\bigl(\varphi(k)\Rightarrow\varphi(k+1)\bigr),
\]
\end{enumerate}

то
\[
(\forall n\in\mathbb N,\ n\geq m)\,\varphi(n).
\]
\end{keythm}

\begin{proof}
Използваме принципа за най-малкия елемент на естествените числа. Ако има $n\ge m$, за което $\varphi(n)$ е невярно, нека $n_0$ е най-малкото такова число. Базата дава $n_0>m$. От минималността $\varphi(n_0-1)$ е вярно, а индуктивната стъпка тогава дава $\varphi(n_0)$ --- противоречие. Следователно няма контрапример. Принципът за най-малкия елемент и индукцията са еквивалентни основни свойства на $\mathbb N$; тук първият е приет за изходен.
\end{proof}

Прилагането на принципа има ясен строеж. Първо се доказва твърдението за началната стойност. След това се приема, че е вярно за произволно, но фиксирано $k$, и при това допускане се доказва за $k+1$. Така от верността за $m$ следва верността за $m+1$, после за $m+2$ и т.н.

\begin{keythm}[Принцип на пълната математическа индукция]
Ако $\varphi(m)$ е вярно и за всяко $k>m$ от
\[
\varphi(m)\land\varphi(m+1)\land\cdots\land\varphi(k-1)
\]
следва $\varphi(k)$, то $\varphi(n)$ е вярно за всяко естествено число $n\geq m$.
\end{keythm}

\begin{proof}
Прилагаме слабата индукция към $\psi(k)\equiv\bigwedge_{j=m}^{k}\varphi(j)$. Базата $\psi(m)$ е дадена. Ако $\psi(k)$ е вярно, предпоставката на пълната индукция дава $\varphi(k+1)$, следователно и $\psi(k+1)$. Получаваме $\psi(k)$ за всяко $k\ge m$, откъдето следва $\varphi(k)$.
\end{proof}

Разликата е в индуктивното предположение: при слабата индукция се използва само $\varphi(k)$, а при пълната индукция може да се използва верността на твърдението за всички предишни стойности.

\begin{example}
Да се докаже, че за всяко $n\in\mathbb N$:
\[
0+1+\cdots+n=\frac{n(n+1)}{2}.
\]

За $n=0$ равенството е вярно. Нека е вярно за $n=k$. Тогава
\[
0+1+\cdots+k+(k+1)
=
\frac{k(k+1)}{2}+(k+1)
=
\frac{(k+1)(k+2)}{2}.
\]
Следователно формулата е вярна и за $k+1$, а по математическа индукция --- за всяко $n\in\mathbb N$.
\end{example}

\section{Операции върху множества}

\begin{defn}[Операции върху множества]
Нека $A$ и $B$ са множества. Основните операции са следните:
\[
A\cup B=\{x\mid x\in A\lor x\in B\},
\]
\[
A\cap B=\{x\mid x\in A\land x\in B\},
\]
\[
A\setminus B=\{x\mid x\in A\land x\notin B\},
\]
\[
A\triangle B=(A\setminus B)\cup(B\setminus A).
\]
\end{defn}

Обединението съдържа елементите, които принадлежат поне на едно от множествата; сечението съдържа общите им елементи; разликата $A\setminus B$ съдържа елементите на $A$, които не са в $B$. Симетричната разлика съдържа елементите, които принадлежат точно на едно от двете множества.

\begin{defn}[Допълнение]
При фиксиран универсум $U$, съдържащ разглежданите множества, \emph{допълнение} на $A$ до $U$ е
\[
\overline A=U\setminus A=\{x\in U\mid x\notin A\}.
\]
\end{defn}

\begin{defn}[Обобщено обединение и сечение]
Нека $A_i\subseteq U$ за всяко $i\in I$, където $U$ е фиксиран универсум. Използват се обобщенията
\[
\bigcup_{i\in I}A_i
=
\{x\mid (\exists i\in I)\ x\in A_i\},
\]
\[
\bigcap_{i\in I}A_i
=
\{x\in U\mid (\forall i\in I)\ x\in A_i\}.
\]
При $I=\varnothing$ обединението е $\varnothing$, а сечението е $U$.
\end{defn}

\begin{prop}[Основни свойства]
За произволни множества $A$, $B$ и $C$ са валидни:

\begin{align*}
A\cup A&=A, & A\cap A&=A;\\
A\cup B&=B\cup A, & A\cap B&=B\cap A;\\
(A\cup B)\cup C&=A\cup(B\cup C), &
(A\cap B)\cap C&=A\cap(B\cap C);\\
A\cup(B\cap C)&=(A\cup B)\cap(A\cup C),&
A\cap(B\cup C)&=(A\cap B)\cup(A\cap C);\\
A\cup\varnothing&=A, & A\cap U&=A;\\
A\cup U&=U, & A\cap\varnothing&=\varnothing;\\
A\cup\overline A&=U, & A\cap\overline A&=\varnothing;\\
\overline{\overline A}&=A. &&
\end{align*}

Освен това са в сила законите на Де Морган:
\[
\overline{A\cup B}=\overline A\cap\overline B,
\qquad
\overline{A\cap B}=\overline A\cup\overline B.
\]
\end{prop}

\begin{proof}
Фиксираме произволен $x\in U$ и означаваме $p\equiv x\in A$, $q\equiv x\in B$, $r\equiv x\in C$. Принадлежността на обединение, сечение и допълнение се превежда съответно като $\lor$, $\land$ и $\neg$. Идемпотентността, комутативността, асоциативността и дистрибутивността следват от едноименните булеви тъждества (проверими за всички осем оценки на $p,q,r$). Принадлежността на $\varnothing$ е неистина, а на $U$ --- истина, което дава четирите равенства с тези множества. Законите за допълнението следват от $p\lor\neg p=1$, $p\land\neg p=0$ и $\neg\neg p=p$. Накрая $\neg(p\lor q)\equiv\neg p\land\neg q$ и $\neg(p\land q)\equiv\neg p\lor\neg q$ дават двата закона на Де Морган. Във всеки случай двете страни имат едни и същи елементи и аксиомата за обема дава равенството.
\end{proof}

\begin{example}
Нека
\[
A=\{x\in\mathbb N\mid x>1\},
\qquad
B=\{x\in\mathbb N\mid x>3\}.
\]
Тогава $B\subseteq A$ и следователно
\[
A\cup B=A,\qquad A\cap B=B,
\]
\[
A\setminus B=\{x\in\mathbb N\mid 1<x\leq 3\},
\qquad
B\setminus A=\varnothing,
\]
\[
A\triangle B=\{x\in\mathbb N\mid 1<x\leq 3\}.
\]
\end{example}

\section{Наредени двойки и декартови произведения}

\begin{keydefn}[Наредена двойка]
Наредена двойка от елементите $a$ и $b$ се означава с $(a,b)$. Тя се определя чрез характеристичното свойство
\[
(a,b)=(c,d)\quad\Longleftrightarrow\quad a=c\land b=d.
\]
\end{keydefn}

Редът е съществен: по принцип $(a,b)\neq(b,a)$. Една теоретико-множествена реализация на наредена двойка е конструкцията на Куратовски:
\[
(a,b)=\bigl\{\{a\},\{a,b\}\bigr\}.
\]

\begin{keydefn}[Наредена $n$-орка]
Полагаме $(a_1)=a_1$ като база на рекурсивния запис. За $n\geq 2$ наредената $n$-орка от $a_1,\ldots,a_n$ се означава с
\[
(a_1,\ldots,a_n)
\]
и може рекурсивно да се разбира като
\[
(a_1,(a_2,\ldots,a_n)).
\]
\end{keydefn}

\begin{keydefn}[Декартово произведение]
За множества $A$ и $B$ тяхното декартово произведение е
\[
A\times B=\{(a,b)\mid a\in A\land b\in B\}.
\]
\end{keydefn}

По принцип $A\times B\neq B\times A$, защото се променя редът на компонентите. Например
\[
\{1,2\}\times\{a\}=\{(1,a),(2,a)\},
\]
докато
\[
\{a\}\times\{1,2\}=\{(a,1),(a,2)\}.
\]

\begin{keydefn}[Обобщено декартово произведение]
За множества $A_1,\ldots,A_n$:
\[
\prod_{i=1}^{n}A_i
=
A_1\times\cdots\times A_n
=
\{(a_1,\ldots,a_n)\mid a_i\in A_i,\ 1\leq i\leq n\}.
\]
\end{keydefn}

\section{Релации}

\begin{keydefn}[$n$-местна релация]
Нека $A_1,\ldots,A_n$ са множества. Всяко подмножество
\[
R\subseteq A_1\times A_2\times\cdots\times A_n
\]
се нарича $n$-местна, или $n$-арна, релация над домейните $A_1,\ldots,A_n$.
\end{keydefn}

При $n=2$ говорим за бинарна релация. Ако домейните са еднакви, например $R\subseteq A\times A$, релацията е хомогенна. Вместо $(a,b)\in R$ често се пише $aRb$.

\begin{defn}[Домейн и множество от стойности на релация]
За $R\subseteq A\times B$ се дефинират:
\[
\operatorname{dom}(R)=
\{a\in A\mid(\exists b\in B)\ (a,b)\in R\},
\]
\[
\operatorname{rng}(R)=
\{b\in B\mid(\exists a\in A)\ (a,b)\in R\}.
\]
Това са съответно домейнът и множеството от стойности на релацията.
\end{defn}

\subsection{Свойства на бинарните релации}

\begin{defn}[Свойства на бинарна релация]
Нека $R\subseteq A\times A$. Казваме, че $R$ е:

\begin{itemize}
\item \emph{рефлексивна}, ако
\[
(\forall x\in A)\ xRx;
\]

\item \emph{антирефлексивна}, или \emph{ирефлексивна}, ако
\[
(\forall x\in A)\ \neg xRx;
\]

\item \emph{симетрична}, ако
\[
(\forall x,y\in A)\ (xRy\Rightarrow yRx);
\]

\item \emph{антисиметрична}, ако
\[
(\forall x,y\in A)\ (xRy\land yRx\Rightarrow x=y);
\]

\item \emph{силно антисиметрична}, ако за всеки две различни $x,y\in A$ е изпълнено точно едно от $xRy$ и $yRx$;

\item \emph{асиметрична}, ако
\[
(\forall x,y\in A)\ (xRy\Rightarrow\neg yRx);
\]

\item \emph{транзитивна}, ако
\[
(\forall x,y,z\in A)\ (xRy\land yRz\Rightarrow xRz).
\]
\end{itemize}
\end{defn}

Антисиметричността не означава, че за различни елементи винаги съществува сравнение. Тя забранява едновременното наличие на $xRy$ и $yRx$ при $x\neq y$. Силната антисиметричност добавя изискването за сравнимост на всяка двойка различни елементи.

\subsection{Еквивалентност и класове на еквивалентност}

\begin{keydefn}[Релация на еквивалентност]
Релация $R\subseteq A\times A$ е релация на еквивалентност, ако е рефлексивна, симетрична и транзитивна.
\end{keydefn}

\begin{keydefn}[Клас на еквивалентност]
Нека $R$ е релация на еквивалентност над $A$ и $a\in A$. Класът на еквивалентност на $a$ е
\[
[a]_R=\{b\in A\mid aRb\}.
\]
Когато релацията е ясна от контекста, се пише просто $[a]$.
\end{keydefn}

\begin{thm}
Всяка релация на еквивалентност над множество $A$ поражда разбиване на $A$ на класове на еквивалентност.
\end{thm}

\begin{proof}
По рефлексивност $a\in[a]_R$, затова класовете са непразни и покриват $A$. Ако $c\in[a]_R\cap[b]_R$, то $aRc$ и $bRc$, откъдето по симетричност и транзитивност $aRb$. За всяко $x\in[a]_R$ имаме $bRaRx$, следователно $x\in[b]_R$. Обратното включване е аналогично. Значи два класа или съвпадат, или са непресичащи се, както изисква определението за разбиване.
\end{proof}

\begin{defn}[Разбиване на множество]
Фамилията $\{A_i\mid i\in I\}$ е \emph{разбиване} на $A$, ако всички $A_i$ са непразни, обединението им е $A$ и различните части нямат общи елементи:
\[
\bigcup_{i\in I}A_i=A,
\]
\[
(\forall i,j\in I)\,
(A_i\cap A_j\neq\varnothing\Rightarrow A_i=A_j).
\]
\end{defn}

\begin{example}
В $\mathbb Z$ задаваме $aRb$, ако $a$ и $b$ дават еднакъв остатък при деление на $3$. Това е релация на еквивалентност. Класовете са числата с остатък $0$, $1$ и $2$ при деление на $3$.
\end{example}

\subsection{Частични и пълни наредби}

\begin{keydefn}[Частична наредба]
Релация $R\subseteq A\times A$ е нестрога частична наредба, ако е рефлексивна, антисиметрична и транзитивна.
\end{keydefn}

\begin{defn}[Строга частична наредба]
Строга частична наредба върху $A$ е антирефлексивна и транзитивна релация $R\subseteq A\times A$. Тези условия влекат асиметричност: от $xRy$ и $yRx$ по транзитивност би следвало $xRx$, противоречие. Допълнителната антисиметричност в изходната формулировка е излишна.
\end{defn}

\begin{keydefn}[Пълна наредба]
Частична наредба $R$ над $A$ е пълна, ако всеки два различни елемента са сравними:
\[
(\forall a,b\in A)\,
\left(a\neq b\Rightarrow(aRb\lor bRa)\right).
\]
\end{keydefn}

Релациите $\leq$ и $<$ над естествените числа са примери за пълни наредби в нестрогия и строгия вариант.

\begin{keydefn}[Минимален и максимален елемент]
Нека $R$ е частична наредба над $A$. Елементът $a\in A$ е минимален, ако
\[
(\forall b\in A)\,(b\neq a\Rightarrow\neg bRa).
\]
Той е максимален, ако
\[
(\forall b\in A)\,(b\neq a\Rightarrow\neg aRb).
\]
\end{keydefn}

\begin{thm}
Всяка частична наредба над непразно крайно множество притежава минимален и максимален елемент.
\end{thm}

\begin{proof}
Избираме $a_0\in A$. Ако той не е минимален, има различен $a_1$ с $a_1Ra_0$. Продължаваме, докато достигнем минимален елемент. Повторение $a_i=a_j$ с $i<j$ е невъзможно: транзитивността и антисиметричността биха дали $a_i=a_{i+1}$, въпреки строгия избор. Понеже $A$ е крайно, процесът спира. Прилагаме същото разсъждение към обратната наредба $R^{-1}$ и получаваме максимален елемент.
\end{proof}

\begin{keydefn}[Диаграма на Хассе]
За крайно частично наредено множество $(A,\le)$ диаграмата на Хассе има върхове елементите на $A$. Ребро от $x$ нагоре към $y$ се чертае точно когато $x<y$ и няма $z$ с $x<z<y$. Тогава $y$ покрива $x$. Рефлексивните примки и ребрата, следващи по транзитивност, не се рисуват. Това уточнява общото описание на граф на релация в източника.
\end{keydefn}

При частична наредба тя позволява удобно да се видят несравнимите елементи, както и минималните и максималните елементи.

\subsection{Линейно разширение и топологично сортиране}

\begin{keythm}[Влагане в пълна наредба]
Нека $A$ е непразно крайно множество и $R$ е частична наредба над $A$. Съществува пълна наредба $S$ над $A$, за която
\[
R\subseteq S.
\]
\end{keythm}

\begin{proof}
Индукция по $|A|$. За едноелементно множество твърдението е ясно. Избираме минимален елемент $a$; той съществува по предходната теорема. По индукция ограничението на $R$ върху $A\setminus\{a\}$ има линейно разширение. Поставяме $a$ пред всички елементи на това разширение. Така получаваме пълна наредба. Няма отношение $bRa$ за $b\ne a$ по минималността, а всички останали отношения на $R$ се запазват по индукционната хипотеза. Следователно новата наредба съдържа $R$.
\end{proof}

\begin{defn}[Линейно разширение]
Релацията $S$ се нарича линейно разширение на $R$. Построяването му е известно като топологично сортиране.
\end{defn}

\begin{verbatim}
Вход: крайно множество A и частична наредба R над A
Изход: редица B, задаваща линейно разширение

i <- 1
докато A не е празно:
    избери минимален елемент a на A спрямо R
    B[i] <- a
    изтрий a от A и ограничи R върху оставащите елементи
    i <- i + 1
върни B
\end{verbatim}

Идеята за коректност е следната. Във всяка непразна стъпка оставащото крайно частично наредено множество има минимален елемент, следователно алгоритъмът може да продължи. Ако елемент $a$ предшества $b$ по $R$, $b$ не може да бъде избран преди $a$, защото тогава $b$ не би бил минимален сред оставащите елементи. Следователно редицата запазва всички отношения от $R$ и задава пълна наредба, съдържаща $R$.

\section{Функции}

\begin{keydefn}[Частична функция]
Нека $A$ и $B$ са множества. Релация $f\subseteq A\times B$ е частична функция от $A$ в $B$, ако за всеки $a\in A$ има най-много един $b\in B$, за който $(a,b)\in f$:
\[
(\forall a\in A)(\forall b_1,b_2\in B)\,
\bigl((a,b_1)\in f\land(a,b_2)\in f\Rightarrow b_1=b_2\bigr).
\]
\end{keydefn}

При частична функция не е задължително всеки елемент на $A$ да има стойност.

\begin{keydefn}[Тотална функция]
Частичната функция $f\subseteq A\times B$ е тотална функция от $A$ в $B$, ако
\[
\operatorname{dom}(f)=A.
\]
Еквивалентно:
\[
(\forall a\in A)(\exists!b\in B)\ (a,b)\in f.
\]
Пише се $f:A\to B$, а вместо $(a,b)\in f$ се пише $f(a)=b$.
\end{keydefn}

\begin{keydefn}[Инекция, сюрекция и биекция]
Нека $f:A\to B$ е тотална функция.

\begin{itemize}
\item $f$ е \emph{инекция}, ако
\[
(\forall a_1,a_2\in A)\,
(f(a_1)=f(a_2)\Rightarrow a_1=a_2).
\]

\item $f$ е \emph{сюрекция}, ако
\[
(\forall b\in B)(\exists a\in A)\ f(a)=b.
\]

\item $f$ е \emph{биекция}, ако е едновременно инекция и сюрекция.
\end{itemize}
\end{keydefn}

Инекцията не допуска различни аргументи с една и съща стойност. Сюрекцията изисква всяка стойност от $B$ да бъде достигната. Биекцията установява съответствие един към един между елементите на $A$ и $B$.

\section{Крайни и изброими множества. Принцип на Дирихле}

\begin{keydefn}[Крайно множество и кардиналност]
Множество $A$ е крайно, ако $A=\varnothing$ или съществува $n\in\mathbb N$, за което има биекция
\[
A\longrightarrow \{0,1,\ldots,n-1\}.
\]
Числото $n$ се нарича кардиналност на $A$ и се означава с $|A|=n$. За празното множество $|\varnothing|=0$.
\end{keydefn}

\begin{keydefn}[Изброимо безкрайно множество]
Множество $A$ е изброимо безкрайно, ако съществува биекция между $A$ и $\mathbb N$.
\end{keydefn}

\begin{keythm}[Принцип на Дирихле]
Нека $X$ и $Y$ са крайни множества и
\[
|X|>|Y|.
\]
Тогава за всяка тотална функция $f:X\to Y$ съществуват различни $a,b\in X$, за които
\[
f(a)=f(b).
\]
Следователно не съществува инекция от $X$ в $Y$.
\end{keythm}

\begin{proof}
Ако никои два различни аргумента нямат един и същ образ, всяко влакно $f^{-1}(\{y\})$ има най-много един елемент. Тези влакна са непресичащи се и покриват $X$, понеже $f$ е тотална. Следователно
\[|X|=\sum_{y\in Y}|f^{-1}(\{y\})|\le |Y|,\]
което противоречи на хипотезата. По същия начин, ако всяка от $n$ кутии съдържа най-много $k$ обекта, общият брой е най-много $kn$, което доказва и обобщената форма.
\end{proof}

Обобщената форма гласи: ако $kn+1$ обекта се разпределят в $n$ кутии, то поне в една кутия има повече от $k$ обекта.

\begin{example}
При разпределяне на $13$ студенти в $12$ месеца според месеца им на раждане поне двама студенти са родени в един и същ месец. Това е пряко приложение на принципа на Дирихле.
\end{example}

\section{Примерни задачи}

\begin{example}
Нека
\[
A=\{1,2,3\},\qquad B=\{2,3,4\}.
\]
Тогава
\[
A\cup B=\{1,2,3,4\},\qquad
A\cap B=\{2,3\},
\]
\[
A\setminus B=\{1\},\qquad
A\triangle B=\{1,4\}.
\]
\end{example}

\begin{example}
Нека $A=\{1,2\}$ и $B=\{a,b\}$. Тогава
\[
A\times B=\{(1,a),(1,b),(2,a),(2,b)\}.
\]
Релацията
\[
R=\{(1,a),(2,b)\}\subseteq A\times B
\]
е тотална функция от $A$ към $B$. Тя е биекция, защото различните елементи на $A$ имат различни образи и всеки елемент на $B$ е образ на елемент от $A$.
\end{example}

\begin{example}
Нека $A=\{1,2,3,6\}$ и над $A$ е зададена релацията делимост:
\[
aRb\quad\Longleftrightarrow\quad a\mid b.
\]
Тя е частична наредба. Елементите $2$ и $3$ са несравними, защото нито $2\mid 3$, нито $3\mid 2$. Минимален елемент е $1$, а максимален елемент е $6$.
\end{example}

\begin{example}
Нека $R$ над $\mathbb Z$ е зададена с
\[
aRb\quad\Longleftrightarrow\quad a-b\text{ е четно}.
\]
Релацията е рефлексивна, симетрична и транзитивна, следователно е релация на еквивалентност. Тя разделя $\mathbb Z$ на два класа: клас на четните и клас на нечетните цели числа.
\end{example}

\section{Изпитен фокус}

\begin{itemize}
\item Да се формулират аксиомата за обема, схемата за отделянето, аксиомата за степенното множество и идеята за индуктивно генерирано множество.
\item Да се обяснят слабата и пълната математическа индукция и да се проведе индукционно доказателство с база и индуктивна стъпка.
\item Да се дефинират обединение, сечение, разлика, симетрична разлика, допълнение и степенно множество; да се използват законите на Де Морган.
\item Да се дефинират наредена двойка, наредена $n$-орка, декартово и обобщено декартово произведение.
\item Да се дефинира $n$-местна релация, както и домейн и множество от стойности на бинарна релация.
\item Да се формулират рефлексивност, антирефлексивност, симетричност, антисиметричност, силна антисиметричност, асиметричност и транзитивност.
\item Да се дефинират релация на еквивалентност, клас на еквивалентност и разбиване на множество.
\item Да се разграничават частична и пълна наредба; да се определят минимални и максимални елементи и да се интерпретира диаграма на Хассе.
\item Да се възпроизведе идеята на топологичното сортиране: многократно избиране и премахване на минимален елемент, докато се получи линейно разширение.
\item Да се дефинират частична и тотална функция, инекция, сюрекция и биекция.
\item Да се дефинират крайно множество, кардиналност и изброимо безкрайно множество.
\item Да се формулира принципът на Дирихле и да се разпознава приложението му като доказателство, че две различни стойности имат еднакъв образ.
\end{itemize}


\chapter{Основни комбинаторни принципи и конфигурации. Рекурентни уравнения.}
%% Drafted from the mapped corpus by scripts/draft_topics.py.
% Model: gpt-5.6-terra; source characters: 20514.
\begin{konspekt}
\kreq{Формулировки}{седемте принципа --- Дирихле, биекция, събиране, изваждане, умножение, деление, включване и изключване --- \kref{2.1}.}
\kreq{Доказателство}{принцип за включване и изключване --- \kref{2.1}.}
\kreq{Дефиниции и извеждане на формулите}{четирите основни конфигурации --- с или без наредба, с или без повтаряне --- \kref{2.2}.}
\kreq{Формулировка и доказателства}{биномни коефициенти, теорема на Нютон, комбинаторни тъждества чрез двукратно броене --- \kref{2.3}.}
\kreq{Алгоритъм}{решаване на линейни рекурентни уравнения с константни коефициенти, хомогенни и нехомогенни --- \kref{2.4}.}
\kreq{Задача 1}{броене на частични функции, тотални функции, инекции и сюрекции --- \kref{2.5.1}.}
\kreq{Задача 2}{целочислени решения на $x_1+\dots+x_n=M$ с ограничения --- \kref{2.5.2}.}
\kreq{Задача 3}{решаване на зададено рекурентно уравнение до край --- \kref{2.5.3}.}
\kprereq{свойствата на бинарните релации, включително доказването и опровергаването им, принадлежат на Въпрос 1.}
\kreq{Литература}{[12], [33], [40].}
\end{konspekt}

\section{Комбинаторни принципи}\label{s:2.1}

В тази глава всички разглеждани множества са крайни, освен ако изрично не е посочено друго. Основната задача на изброителната комбинаторика е да се намери броят на обектите в дадена крайна конфигурация, като вместо непосредствено изброяване се използват структурни свойства на множествата и функциите.

\begin{keythm}[Принцип на Дирихле]
Нека $X$ и $Y$ са крайни множества и $|X|>|Y|$. Тогава не съществува инекция
\[
f:X\to Y.
\]
Еквивалентно: за всяка тотална функция $f:X\to Y$ съществуват различни $a,b\in X$, за които $f(a)=f(b)$.
\end{keythm}

\begin{proof}
Ако $f$ е инекция, всеки образ има най-много един прообраз. Непресичащите се влакна $f^{-1}(\{y\})$, $y\in Y$, покриват $X$, затова $|X|=\sum_{y\in Y}|f^{-1}(\{y\})|\le|Y|$, противоречие. Това е същият принцип, доказан във Въпрос 1.
\end{proof}

Принципът изразява, че ако повече обекти се разпределят в по-малко на брой „кутии“, поне една кутия съдържа поне два обекта.

\begin{keythm}[Принцип на биекцията]
За крайни множества $A$ и $B$ е изпълнено
\[
|A|=|B|
\]
тогава и само тогава, когато съществува биекция $f:A\to B$.
\end{keythm}

\begin{proof}
Ако $|A|=|B|=n$, избираме номерации $a:\{1,\ldots,n\}\to A$ и $b:\{1,\ldots,n\}\to B$; тогава $b\circ a^{-1}$ е биекция. Обратно, биекция $f:A\to B$ пренася всяка номерация на $A$ в номерация на $B$ със същия брой елементи. При празни множества се използва празната биекция.
\end{proof}

Принципът на биекцията е особено полезен, когато разглежданата конфигурация се съпостави еднозначно с друга, чиято мощност вече е известна.

\begin{keythm}[Принцип на събирането]
Нека
\[
A=S_1\cup S_2\cup\cdots\cup S_k,
\]
където множествата $S_1,\ldots,S_k$ са две по две непресичащи се. Тогава
\[
|A|=\sum_{i=1}^{k}|S_i|.
\]
\end{keythm}

\begin{proof}
Номерираме последователно елементите на $S_1$, после на $S_2$ и така до $S_k$. Непресичането гарантира, че няма повторения, а обединението гарантира, че не пропускаме елемент на $A$. Получената номерация има точно $\sum_i|S_i|$ позиции.
\end{proof}

Условието за непресичане е съществено: всеки елемент на $A$ трябва да бъде преброен точно веднъж.

\begin{keythm}[Принцип на изваждането]
Нека $A\subseteq U$, където $U$ е универсумът на разглежданите обекти. Тогава
\[
|A|=|U|-|U\setminus A|.
\]
\end{keythm}

\begin{proof}
Множествата $A$ и $U\setminus A$ са непресичащи се и обединението им е $U$. Принципът на събирането дава $|U|=|A|+|U\setminus A|$. Пренасяме последното събираемо от другата страна.
\end{proof}

\begin{keythm}[Принцип на умножението]
За крайни множества $A$ и $B$ е в сила
\[
|A\times B|=|A|\cdot |B|.
\]
По-общо,
\[
|X_1\times X_2\times\cdots\times X_r|
 =\prod_{i=1}^{r}|X_i|.
\]
\end{keythm}

\begin{proof}
Разбиваме $A\times B$ на непресичащите се множества $\{a\}\times B$ за $a\in A$. Всяко е в биекция с $B$ чрез $(a,b)\mapsto b$. Принципът на събирането дава $|A\times B|=\sum_{a\in A}|B|=|A||B|$. Индукция по положителния брой $r$ дава общата формула. Ако някой множител е празен, и произведението е празно.
\end{proof}

Този принцип се прилага, когато обектът се изгражда последователно чрез независим избор на компоненти.

\begin{keythm}[Принцип на делението]
Нека $A$ е непразно крайно множество и $R\subseteq A^2$ е релация на еквивалентност върху него. Ако $A$ има $k$ класа на еквивалентност и всеки клас има по $m$ елемента, то
\[
|A|=km,
\qquad\text{следователно}\qquad
m=\frac{|A|}{k}.
\]
\end{keythm}

\begin{proof}
Класовете на еквивалентност образуват разбиване. По принципа на събирането $|A|$ е сумата на размерите им, тоест $k$ еднакви събираеми $m$. Следователно $|A|=km$. Понеже разглеждаме непразно $A$, имаме $k,m>0$ и двете формули с деление са допустими.
\end{proof}

При използване за преброяване обикновено са известни $|A|$ и размерът $m$ на всеки клас; тогава броят на класовете е $|A|/m$.

\begin{keythm}[Принцип на включването и изключването]
За крайни множества $A_1,\ldots,A_n$ е изпълнено
\[
\left|\bigcup_{i=1}^{n}A_i\right|
=
\sum_{i=1}^{n}|A_i|
-\sum_{1\leq i<j\leq n}|A_i\cap A_j|
+\sum_{1\leq i<j<k\leq n}|A_i\cap A_j\cap A_k|
-\cdots
+(-1)^{n-1}|A_1\cap\cdots\cap A_n|.
\]
\end{keythm}

\begin{proof}
Доказателството е с индукция по $n$. При $n=1$ равенството е очевидно. За индукционната стъпка записваме
\[
\left|\bigcup_{i=1}^{n}A_i\right|
=
\left|\left(\bigcup_{i=1}^{n-1}A_i\right)\cup A_n\right|.
\]
За две множества е в сила
\[
|B\cup C|=|B|+|C|-|B\cap C|.
\]
Следователно
\[
\left|\bigcup_{i=1}^{n}A_i\right|
=
\left|\bigcup_{i=1}^{n-1}A_i\right|
+|A_n|
-\left|\left(\bigcup_{i=1}^{n-1}A_i\right)\cap A_n\right|.
\]
От дистрибутивния закон
\[
\left(\bigcup_{i=1}^{n-1}A_i\right)\cap A_n
=
\bigcup_{i=1}^{n-1}(A_i\cap A_n).
\]
Прилагат се индукционната хипотеза към двете обединения. След събиране на членовете се получават всички пресичания на $A_1,\ldots,A_n$ със съответните редуващи се знаци.
\end{proof}

\section{Основни комбинаторни конфигурации}\label{s:2.2}

Нека основното множество има $n$ елемента. Разглеждаме избиране на $m$ елемента, като са възможни четири случая според това дали редът е важен и дали са разрешени повторения.

\subsection{Конфигурации с наредба и повторение}\label{s:2.2.1}

\begin{defn}
Множеството от конфигурациите с наредба и повторение с дължина $m$ над множество $A$, $|A|=n$, се означава с $K_{\mathrm{нп}}(n,m)$. Неговите елементи са наредени $m$-торки
\[
(a_1,\ldots,a_m),\qquad a_i\in A.
\]
\end{defn}

\begin{thm}
\[
|K_{\mathrm{нп}}(n,m)|=n^m.
\]
\end{thm}

\begin{proof}
За всяка от $m$-те позиции има независимо по $n$ възможности. Принципът на умножението дава
\[
\underbrace{n\cdot n\cdots n}_{m\text{ пъти}}=n^m.
\]
Еквивалентно, конфигурациите са функциите от $m$-елементно множество към $n$-елементно множество.
\end{proof}

\begin{example}
Броят на думите с дължина $4$, образувани от азбука с $n$ символа, е $n^4$, ако един и същ символ може да се появява многократно.
\end{example}

\subsection{Конфигурации с наредба без повторение}\label{s:2.2.2}

\begin{defn}
Конфигурация с наредба без повторение е наредена $m$-торка от различни елементи на $n$-елементно множество. Множеството на тези конфигурации се означава с $K_{\mathrm{н}}(n,m)$.
\end{defn}

\begin{thm}
\[
|K_{\mathrm{н}}(n,m)|
=
n(n-1)\cdots(n-m+1).
\]
При $m>n$ броят е $0$.
\end{thm}

\begin{proof}
За първата позиция има $n$ възможности. След избора й остават $n-1$ възможности за втората позиция, после $n-2$ за третата и т.н. Принципът на умножението дава формулата. Ако $m>n$, инекция от $m$-елементно множество в $n$-елементно множество не съществува по принципа на Дирихле.
\end{proof}

При $m=n$ конфигурациите се наричат \emph{пермутации} и броят им е
\[
n!=n(n-1)\cdots 2\cdot 1.
\]

\subsection{Конфигурации без наредба и без повторение}\label{s:2.2.3}

\begin{defn}
Конфигурация без наредба и без повторение е $m$-елементно подмножество на $n$-елементно множество. Броят им се означава с
\[
\binom{n}{m}.
\]
\end{defn}

\begin{keythm}
За $0\leq m\leq n$ е изпълнено
\[
\binom{n}{m}
=
\frac{n(n-1)\cdots(n-m+1)}{m!}
=
\frac{n!}{m!(n-m)!}.
\]
Ако $m>n$, то
\[
\binom{n}{m}=0.
\]
\end{keythm}



\begin{proof}
Нека разгледаме множеството $K_{\mathrm{н}}(n,m)$ на наредените конфигурации без повторение. Въвеждаме релация $R$ чрез
\[
XRY
\quad\Longleftrightarrow\quad
X\text{ и }Y\text{ съдържат едни и същи елементи}.
\]
Това е релация на еквивалентност. Всеки клас на еквивалентност съдържа всички наредби на едно и също $m$-елементно множество, следователно има $m!$ елемента. Класовете са точно конфигурациите без наредба и без повторение. По принципа на делението
\[
\binom{n}{m}
=
\frac{|K_{\mathrm{н}}(n,m)|}{m!}
=
\frac{n(n-1)\cdots(n-m+1)}{m!}.
\]
\end{proof}

\subsection{Конфигурации без наредба с повторение}\label{s:2.2.4}

\begin{defn}
Конфигурация без наредба с повторение с големина $m$ над $n$-елементно основно множество е мултимножество с обща кардиналност $m$. Множеството от тези конфигурации се означава с $K_a(n,m)$.
\end{defn}

\begin{thm}
\[
|K_a(n,m)|
=
\binom{n+m-1}{m}
=
\binom{n+m-1}{n-1}.
\]
\end{thm}

\begin{proof}
Нека $x_i$ е броят на избраните обекти от $i$-тия вид. Тогава
\[
x_1+x_2+\cdots+x_n=m,
\qquad x_i\geq 0.
\]
Представяме решението чрез $m$ звездички и $n-1$ разделителя. Например броят звездички преди първия разделител е $x_1$, между първия и втория е $x_2$ и т.н. Така получаваме биекция между избора на мултимножество и избора на позициите на $m$ звездички сред общо $m+n-1$ позиции. Следователно броят е
\[
\binom{m+n-1}{m}.
\]
\end{proof}

\section{Биномни коефициенти и двукратно броене}\label{s:2.3}

\begin{prop}
За цели $0\leq m\leq n$ са в сила следните равенства; тези с избор на $1$ или $n-1$ елемента се отнасят за $n\ge1$:
\[
\binom{n}{0}=\binom{n}{n}=1,
\qquad
\binom{n}{1}=\binom{n}{n-1}=n,
\]
\[
\binom{n}{m}=\binom{n}{n-m},
\]
\[
\sum_{i=0}^{n}\binom{n}{i}=2^n.
\]
\end{prop}

\begin{proof}
Празното множество и цялото множество са единствените подмножества с размер съответно $0$ и $n$. За $n\ge1$ едноелементните подмножества са $n$, както и допълненията им. Съответствието $B\mapsto A\setminus B$ е биекция между подмножествата с размер $m$ и $n-m$. Накрая всяко от $n$-те различни елемента може независимо да бъде включено или не, затова всички подмножества са $2^n$; разбиването им по размер дава сумата.
\end{proof}

Равенството
\[
\binom{n}{m}=\binom{n}{n-m}
\]
следва от биекцията, която на всяко $m$-елементно подмножество $B$ на $A$ съпоставя допълнението $A\setminus B$, имащо $n-m$ елемента.

Равенството
\[
\sum_{i=0}^{n}\binom{n}{i}=2^n
\]
се получава, като всички подмножества на $n$-елементно множество се преброят по два начина. От една страна те са $2^n$; от друга страна ги разбиваме според кардиналността им.

\begin{keythm}[Формула на Паскал]
За $n,m\geq 1$ е в сила
\[
\binom{n}{m}
=
\binom{n-1}{m}
+
\binom{n-1}{m-1}.
\]
\end{keythm}

\begin{proof}
Нека $A$ е $n$-елементно множество и фиксираме $a\in A$. Разбиваме $m$-елементните подмножества на $A$ на две непресичащи се групи: тези, които не съдържат $a$, и тези, които съдържат $a$. В първата група има
\[
\binom{n-1}{m}
\]
подмножества. За всяко подмножество от втората група, след премахване на $a$, остава $(m-1)$-елементно подмножество на $A\setminus\{a\}$, следователно такива са
\[
\binom{n-1}{m-1}.
\]
Принципът на събирането дава търсеното равенство.
\end{proof}

\begin{keythm}[Биномна теорема на Нютон]
За $x,y\in\mathbb{R}$ и $n\in\mathbb{N}$ е изпълнено
\[
(x+y)^n
=
\sum_{k=0}^{n}\binom{n}{k}x^k y^{n-k}.
\]
\end{keythm}

\begin{proof}
Произведението $(x+y)^n$ съдържа $n$ множителя от вида $(x+y)$. При разкриване на скобите за получаване на член $x^k y^{n-k}$ трябва да изберем точно $k$ от $n$-те множителя, от които да вземем $x$; от останалите се взема $y$. Броят на тези избори е $\binom{n}{k}$. Сумирането по всички възможни $k$ дава формулата.
\end{proof}

\begin{example}
Чрез двукратно броене се получава тъждеството
\[
\binom{2n}{n}
=
\sum_{k=0}^{n}\binom{n}{k}^2.
\]
Лявата страна брои $n$-елементните подмножества на множество, разделено на две $n$-елементни части. Ако от първата част са избрани $k$ елемента, то от втората трябва да са избрани $n-k$ елемента. Броят е
\[
\binom{n}{k}\binom{n}{n-k}
=
\binom{n}{k}^2.
\]
Сумираме по всички $k$.
\end{example}

\section{Рекурентни уравнения}\label{s:2.4}

\begin{defn}
Рекурентно отношение от ред $r$ е зависимост, при която за $n\geq r$ членът $a_n$ на редица се определя чрез предходните $r$ члена:
\[
a_n=f(a_{n-1},\ldots,a_{n-r},n).
\]
Заедно с началните условия то задава редицата.
\end{defn}

\begin{example}
Отношението
\[
l(n)=
\begin{cases}
l(n-1)+n,& n\in\mathbb{N}^{*},\\
1,& n=0
\end{cases}
\]
задава числова редица. Решаването на рекурентно уравнение означава да се намери явна формула за $n$-тия член на същата редица.
\end{example}

\subsection{Хомогенни линейни рекурентни уравнения}\label{s:2.4.1}

Разглеждаме линейно рекурентно уравнение от ред $k$ с постоянни коефициенти:
\[
a_n=c_1a_{n-1}+c_2a_{n-2}+\cdots+c_ka_{n-k},
\]
където $c_k\neq 0$, а $a_0,\ldots,a_{k-1}$ са дадени начални условия.

\begin{keydefn}[Характеристично уравнение]
На горното рекурентно уравнение съответства характеристичното уравнение
\[
x^k-c_1x^{k-1}-c_2x^{k-2}-\cdots-c_k=0.
\]
\end{keydefn}

То се получава формално чрез замяна на $a_n$ с $x^n$:
\[
x^n=c_1x^{n-1}+\cdots+c_kx^{n-k},
\]
след което се дели на $x^{n-k}$.

\begin{thm}
Ако характеристичното уравнение има $k$ различни корена
\[
\alpha_1,\ldots,\alpha_k,
\]
то общото решение е
\[
a_n=A_1\alpha_1^n+\cdots+A_k\alpha_k^n,
\]
където константите $A_1,\ldots,A_k$ се определят от началните условия.
\end{thm}

\begin{proof}
Работим над $\mathbb C$ (за реална редица константите се избират така, че решението да е реално). Всяка редица $\alpha_i^n$ удовлетворява рекурентното уравнение чрез пряко заместване. Пространството на решенията е $k$-мерно, защото първите $k$ стойности се избират произволно и определят еднозначно всички следващи. Матрицата на първите $k$ стойности на тези редици е матрица на Вандермонд с детерминанта $\prod_{i<j}(\alpha_j-\alpha_i)\ne0$. Следователно редиците образуват базис и константите се определят еднозначно от началните условия.
\end{proof}

\begin{thm}
Нека различните корени на характеристичното уравнение са
\[
\beta_1,\ldots,\beta_l,
\]
като коренът $\beta_i$ има кратност $r_i$ и
\[
r_1+\cdots+r_l=k.
\]
Тогава към корена $\beta_i$ съответстват членовете
\[
A_{i,1}\beta_i^n+
A_{i,2}n\beta_i^n+
\cdots+
A_{i,r_i}n^{r_i-1}\beta_i^n.
\]
Общото решение е сумата на тези изрази за всички различни корени.
\end{thm}

\begin{proof}
Нека $E$ е операторът на изместване $(Ea)_n=a_{n+1}$ и $P(x)=\prod_i(x-\beta_i)^{r_i}$ е характеристичният полином. Уравнението е $P(E)a=0$. Понеже $c_k\ne0$, всички $\beta_i$ са ненулеви. За полином $Q$ имаме
\[(E-\beta)(\beta^nQ(n))=\beta^{n+1}\bigl(Q(n+1)-Q(n)\bigr).\]
Разликата намалява степента, затова $(E-\beta_i)^{r_i}$ анулира $\beta_i^nQ(n)$ при $\deg Q<r_i$. Тези $r_i$ редици са независими, защото ненулев полином не се анулира върху всички естествени числа. За различни $i$ пространствата са независими: при предполагаема зависимост прилагаме $\prod_{j\ne i}(E-\beta_j)^{r_j}$. Той анулира всички други слагаеми, а върху $\beta_i^nQ(n)$ е обратим, защото всяко $E-\beta_j$ има ненулев диагонален коефициент $\beta_i-\beta_j$ в полиномиалния базис. Следователно и $i$-тото слагаемо е нула. Получените общо $k$ независими решения образуват базис на $k$-мерното пространство на решенията.
\end{proof}

След намиране на общия вид се заместват дадените начални условия. Получава се система от $k$ линейни уравнения относно неизвестните константи, която се решава.

\subsection{Нехомогенни линейни рекурентни уравнения}\label{s:2.4.2}

Нехомогенно линейно рекурентно уравнение с постоянни коефициенти има вида
\[
a_n=c_1a_{n-1}+\cdots+c_ka_{n-k}
+p_1(n)b_1^n+\cdots+p_l(n)b_l^n,
\]
където $p_1(n),\ldots,p_l(n)$ са полиноми, а $b_1,\ldots,b_l$ са константи.

Първо се разглежда хомогенната част
\[
a_n=c_1a_{n-1}+\cdots+c_ka_{n-k}
\]
и се намират корените на нейното характеристично уравнение. Нехомогенната част е представена като сума от изрази от вида
\[
b^nP(n),
\]
където $P(n)$ е полином.

За частно решение на член $b^nP_d(n)$ се търси израз
\[
a_n^{(p)}=n^s b^nQ_d(n),
\]
където $Q_d$ е полином от същата степен $d$, а $s$ е кратността на $b$ като корен на характеристичния полином на хомогенната част. Ако $b$ не е корен, тогава $s=0$. Неизвестните коефициенти на $Q_d$ се намират чрез заместване и сравняване на коефициентите. За сума от такива членове се събират съответните частни решения. Общото решение е
\[
a_n=a_n^{(h)}+a_n^{(p)},
\]
а свободните константи се определят от началните условия.

\section{Примерни задачи по анотацията}\label{s:2.5}

Анотацията към този въпрос назовава три типа задачи. Методите за тях са
изложени по-горе; тук всеки тип е решен докрай, защото „методът е обяснен“ и
„задачата е решена“ са различни неща на изпита.

\subsection{Задача 1: броене на функции между крайни множества}\label{s:2.5.1}

Нека $|A|=m$ и $|B|=n$. Търсим броя на частичните функции, на тоталните
функции, на инекциите и на сюрекциите от $A$ в $B$.

\begin{prop}[Частични и тотални функции]
Броят на тоталните функции $f:A\to B$ е $n^m$, а броят на частичните функции от
$A$ в $B$ е $(n+1)^m$.
\end{prop}
\begin{proof}
Тотална функция се задава чрез независим избор на образ за всеки от $m$-те
елемента на $A$; по принципа на умножението изборите са $n^m$. Това е точно
броят на конфигурациите с наредба и повторение.

Частична функция се задава чрез избор за всеки елемент на $A$ на една от $n+1$
възможности: някой от $n$-те елемента на $B$ или „недефинирана“. Изборите отново
са независими, откъдето $(n+1)^m$. Еквивалентно: разширяваме $B$ с един нов
елемент $\bot$ и броим тоталните функции $A\to B\cup\{\bot\}$.
\end{proof}

\begin{prop}[Инекции]
Броят на инекциите $f:A\to B$ е
\[
n(n-1)\cdots(n-m+1)=\frac{n!}{(n-m)!},
\]
и е $0$ при $m>n$.
\end{prop}
\begin{proof}
Подреждаме елементите на $A$ и избираме образите последователно. Първият има $n$
възможности, вторият --- $n-1$, тъй като образът на първия е зает, и така
нататък. Това е точно броят на конфигурациите с наредба без повторение. При
$m>n$ инекция не съществува по принципа на Дирихле.
\end{proof}

\begin{prop}[Сюрекции]
Броят на сюрекциите $f:A\to B$ е
\[
\sum_{k=0}^{n}(-1)^k\binom{n}{k}(n-k)^m .
\]
\end{prop}
\begin{proof}
За $b\in B$ нека $A_b$ е множеството от тоталните функции $A\to B$, които
\emph{не} приемат стойността $b$. Сюрекциите са точно функциите извън
$\bigcup_{b\in B}A_b$. За фиксирано $k$-елементно $K\subseteq B$ пресичането
$\bigcap_{b\in K}A_b$ се състои от функциите $A\to B\setminus K$, а те са
$(n-k)^m$ на брой. Подмножествата $K$ с $k$ елемента са $\binom{n}{k}$.
Принципът за включване и изключване дава
\[
\left|\bigcup_{b\in B}A_b\right|=\sum_{k=1}^{n}(-1)^{k-1}\binom{n}{k}(n-k)^m,
\]
а изваждането от общия брой $n^m$ дава сумата от твърдението, чийто член при
$k=0$ е именно $n^m$.
\end{proof}

\begin{example}
При $m=n$ сюрекциите съвпадат с биекциите. Формулата дава
$\sum_{k=0}^{n}(-1)^k\binom{n}{k}(n-k)^n=n!$, а при $m<n$ тя дава $0$, което се
съгласува с невъзможността да се покрие по-голямо множество.
\end{example}

\subsection{Задача 2: целочислени решения с ограничения}\label{s:2.5.2}

\begin{prop}[Неотрицателни решения]
Броят на решенията на
\[
x_1+x_2+\cdots+x_n=M
\]
в неотрицателни цели числа е $\binom{M+n-1}{n-1}$.
\end{prop}
\begin{proof}
Всяко решение се кодира еднозначно с редица от $M$ еднакви звездички и $n-1$
прегради: звездичките между $(i-1)$-вата и $i$-тата преграда са стойността
$x_i$. Различните решения дават различни редици и всяка редица от $M+n-1$
позиции с избрани $n-1$ позиции за прегради дава решение. По принципа на
биекцията броят е броят на подмножествата с $n-1$ елемента, тоест
$\binom{M+n-1}{n-1}$. Това е конфигурация без наредба с повторение.
\end{proof}

\begin{prop}[Долни граници]
Нека $c_1,\ldots,c_n$ са цели числа. Броят на решенията на
$x_1+\cdots+x_n=M$ с $x_i\geq c_i$ за всяко $i$ е
\[
\binom{M-\sum_{i=1}^{n}c_i+n-1}{n-1},
\]
когато $M\geq\sum_i c_i$, и $0$ в противен случай.
\end{prop}
\begin{proof}
Полагаме $y_i=x_i-c_i$. Изображението $(x_i)\mapsto(y_i)$ е биекция между
решенията с $x_i\geq c_i$ и неотрицателните решения на
\[
y_1+\cdots+y_n=M-\sum_{i=1}^{n}c_i .
\]
Прилагаме предходното твърдение. Ако дясната страна е отрицателна, решения няма.
\end{proof}

\begin{example}[Горни граници]
Ограничение отгоре се обработва с включване и изключване, а не с изместване. За
\[
x_1+x_2+x_3=10,\qquad 0\leq x_i\leq 5,
\]
общият брой неотрицателни решения е $\binom{12}{2}=66$. Нека $B_i$ е множеството
на решенията с $x_i\geq 6$. Изместването $x_i\mapsto x_i-6$ свежда $B_i$ до
неотрицателните решения на сума $10-6=4$ с три неизвестни, тоест
$|B_i|=\binom{4+2}{2}=\binom{6}{2}=15$, а $|B_i\cap B_j|=0$, защото $6+6=12>10$.
Следователно търсеният брой е
\[
66-3\cdot 15=21 .
\]
\end{example}

\begin{example}[Смесени условия]
За $x_1+x_2+x_3+x_4=20$ с $x_1\geq 2$, $x_2\geq 3$, $x_3,x_4\geq 0$ изместваме
първите две неизвестни и получаваме $\binom{20-5+3}{3}=\binom{18}{3}=816$.
\end{example}

\subsection{Задача 3: решаване на зададено рекурентно уравнение}\label{s:2.5.3}

\begin{example}[Хомогенно уравнение с начални условия]
Да се реши
\[
a_n=5a_{n-1}-6a_{n-2},\qquad a_0=1,\quad a_1=0 .
\]
Характеристичното уравнение е $x^2-5x+6=0$ с различни корени $2$ и $3$, затова
$a_n=A\cdot 2^n+B\cdot 3^n$. Началните условия дават
\[
A+B=1,\qquad 2A+3B=0,
\]
откъдето $B=-2$ и $A=3$. Следователно
\[
a_n=3\cdot 2^n-2\cdot 3^n .
\]
Проверка: $a_0=3-2=1$, $a_1=6-6=0$ и $a_2=12-18=-6=5\cdot 0-6\cdot 1$.
\end{example}

\begin{example}[Нехомогенно уравнение с резонанс]
Да се реши
\[
a_n=3a_{n-1}-2a_{n-2}+2^n,\qquad a_0=0,\quad a_1=1 .
\]
Хомогенната част има характеристично уравнение $x^2-3x+2=0$ с корени $1$ и $2$,
тоест $a_n^{(h)}=A+B\cdot 2^n$. Основата $2$ на нехомогенния член е корен с
кратност $s=1$, а полиномът пред него е от степен $0$. Затова търсим
$a_n^{(p)}=C\,n\,2^n$. Заместваме и делим на $2^{n-2}$:
\[
4Cn=6C(n-1)-2C(n-2)+4 ,
\]
което се свежда до $0=-2C+4$, тоест $C=2$ и $a_n^{(p)}=n\cdot 2^{n+1}$. Общото
решение е
\[
a_n=A+B\cdot 2^n+n\cdot 2^{n+1}.
\]
Началните условия дават $A+B=0$ и $A+2B+4=1$, откъдето $B=-3$ и $A=3$:
\[
a_n=3-3\cdot 2^n+n\cdot 2^{n+1}.
\]
Проверка: $a_2=3-12+16=7=3\cdot 1-2\cdot 0+4$ и $a_3=3-24+48=27=3\cdot 7-2\cdot 1+8$.
\end{example}

\begin{remark}
Типична грешка е спирането при общото решение. Анотацията иска \emph{решено}
уравнение, тоест константите да са определени от началните условия и отговорът
да е проверен поне за първите няколко члена.
\end{remark}

\section{Изпитен фокус}\label{s:2.6}

\begin{itemize}
\item Формулировки на принципите на Дирихле, биекцията, събирането, изваждането, умножението, делението и включването и изключването.
\item Доказателствената схема с индукция за принципа на включването и изключването.
\item Различаване между четирите вида комбинаторни конфигурации и формулите
\[
n^m,\qquad n(n-1)\cdots(n-m+1),\qquad
\binom{n}{m},\qquad \binom{n+m-1}{m}.
\]
\item Доказване на формулата за комбинации чрез релация на еквивалентност и принципа на делението.
\item Комбинаторни доказателства на симетрията на биномните коефициенти, формулата на Паскал, равенството
\[
\sum_{i=0}^{n}\binom{n}{i}=2^n
\]
и биномната теорема на Нютон.
\item Алгоритъм за хомогенно линейно рекурентно уравнение: характеристично уравнение, различни и кратни корени, общ вид на решението и определяне на константите от началните условия.
\item Разделяне на нехомогенно рекурентно уравнение на хомогенна и нехомогенна част; за член $b^nP_d(n)$ да се търси частно решение $n^sb^nQ_d(n)$, където $s$ е кратността на $b$ като характеристичен корен.
\item Задача 1 от анотацията: броене на частични функции $(n+1)^m$, тотални функции $n^m$, инекции $n!/(n-m)!$ и сюрекции $\sum_k(-1)^k\binom{n}{k}(n-k)^m$ --- с обосновката за всяка от четирите.
\item Задача 2 от анотацията: неотрицателни целочислени решения чрез звездички и прегради; изместване $y_i=x_i-c_i$ при долни граници $x_i\geq c_i$; включване и изключване при горни граници.
\item Задача 3 от анотацията: пълно решаване на зададено хомогенно и на зададено нехомогенно уравнение --- характеристични корени, частно решение с резонансен множител $n^s$, определяне на константите от началните условия и числена проверка.
\item Свойствата на бинарните релации (рефлексивност, симетричност, транзитивност и опровергаването им) принадлежат на Въпрос 1, а не на този въпрос.
\end{itemize}


\chapter{Графи. Дървета. Обхождания на графи.}
%% Drafted from the mapped corpus by scripts/draft_topics.py.
% Model: gpt-5.6-terra; source characters: 23818.
\begin{konspekt}
\kreq{Дефиниции}{краен ориентиран и краен неориентиран (мулти)граф --- \kref{3.2}.}
\kreq{Дефиниции}{път и цикъл в ориентиран и неориентиран мултиграф; свързаност и свързани компоненти --- \kref{3.3}.}
\kreq{Дефиниции}{дърво, кореново дърво, покриващо дърво --- \kref{3.4}.}
\kreq{Дефиниции}{Ойлеров път и Ойлеров цикъл --- \kref{3.6}.}
\kreq{Доказателства}{всяко кореново дърво е дърво и $|V|=|E|+1$ --- \kref{3.4}.}
\kreq{Доказателство}{теорема за съществуване на Ойлеров цикъл --- \kref{3.6}.}
\kreq{Формулировка}{критерий за съществуване на Ойлеров път --- \kref{3.6}.}
\kreq{Алгоритми}{обхождане в широчина и в дълбочина --- \kref{3.5}.}
\kreq{Задача 1}{покриващо дърво в широчина или в дълбочина --- \kref{3.7.1}.}
\kreq{Задача 2}{Ойлеров цикъл или път, или доказателство, че няма --- \kref{3.7.2}.}
\kreq{Задача 3}{разбиване на ребрата на минимален брой пътища без общо ребро --- \kref{3.7.3}.}
\kprereq{\kr{3.1} повтаря релациите на еквивалентност от Въпрос 1; тук от тях се използва само връзката „достижимост $=$ еквивалентност, класовете $=$ компоненти“.}
\kreq{Литература}{[12], [33], [40].}
\end{konspekt}

\section{Релации на еквивалентност и класове на еквивалентност}\label{s:3.1}

\textbf{Предварителен раздел, който конспектът не изисква по този въпрос.}
Анотацията към Въпрос 3 започва от дефинициите за граф; релациите на
еквивалентност принадлежат на Въпрос 1. Разделът е тук само защото
достижимостта в неориентиран граф е релация на еквивалентност, а класовете й са
свързаните компоненти --- това е използвано по-нататък. При повторение той може
да се прескочи и да се започне направо от \S3.2.

\begin{defn}
Нека $A$ е множество. Релация $\sim\subseteq A\times A$ се нарича \emph{релация на еквивалентност}, ако за всички $x,y,z\in A$ са изпълнени:
\begin{enumerate}
    \item \emph{рефлексивност}: $x\sim x$;
    \item \emph{симетричност}: ако $x\sim y$, то $y\sim x$;
    \item \emph{транзитивност}: ако $x\sim y$ и $y\sim z$, то $x\sim z$.
\end{enumerate}
\end{defn}

\begin{defn}
Нека $\sim$ е релация на еквивалентност върху $A$ и $a\in A$. \emph{Класът на еквивалентност} на $a$ е множеството
\[
[a]_{\sim}=\{x\in A\mid x\sim a\}.
\]
Когато релацията е ясна от контекста, пишем просто $[a]$.
\end{defn}

За да се определят класовете на еквивалентност, за всеки още необработен елемент $a\in A$ се намират всички елементи $x$, за които $x\sim a$. Полученото множество е $[a]$. След това всички негови елементи се отбелязват като обработени и процедурата се повтаря с елемент извън вече намерените класове.

\begin{prop}
Ако $\sim$ е релация на еквивалентност върху $A$, то за произволни $a,b\in A$ класовете $[a]$ и $[b]$ или съвпадат, или са непресичащи се.
\end{prop}

\begin{proof}
Да допуснем, че съществува $x\in [a]\cap[b]$. Тогава $x\sim a$ и $x\sim b$. От симетричността следва $a\sim x$, а от транзитивността с $x\sim b$ получаваме $a\sim b$.

Нека $y\in[a]$. Тогава $y\sim a$ и $a\sim b$, следователно $y\sim b$, т.е. $y\in[b]$. Значи $[a]\subseteq[b]$. Аналогично $[b]\subseteq[a]$, следователно $[a]=[b]$.
\end{proof}

\begin{cor}
Класовете на еквивалентност на релация на еквивалентност образуват разбиение на множеството $A$: всеки елемент принадлежи на точно един клас.
\end{cor}

\begin{proof}
Рефлексивността дава $a\in[a]$, така че класовете покриват $A$. Ако $c\in[a]\cap[b]$, симетричността и транзитивността дават $a\sim b$, а после $x\sim a\Longleftrightarrow x\sim b$. Значи $[a]=[b]$. Следователно различните класове са непресичащи се и всеки елемент е в точно един от тях.
\end{proof}

\begin{example}
Върху $\mathbb{Z}$ да е зададена релацията
\[
a\sim b\Longleftrightarrow 3\mid(a-b).
\]
Тя разделя целите числа на три класа:
\[
[0]=\{\ldots,-6,-3,0,3,6,\ldots\},
\]
\[
[1]=\{\ldots,-5,-2,1,4,7,\ldots\},
\]
\[
[2]=\{\ldots,-4,-1,2,5,8,\ldots\}.
\]
За да се намери класът на дадено число, е достатъчно да се определи остатъкът му при деление на $3$.
\end{example}

\section{Графи и основни понятия}\label{s:3.2}

\subsection{Неориентирани графи}\label{s:3.2.1}

\begin{defn}
Краен неориентиран граф е наредена двойка
\[
G=(V,E),
\]
където $V$ е непразно крайно множество, чиито елементи се наричат \emph{върхове}, а
\[
E\subseteq\{X\subseteq V\mid |X|=2\}
\]
е множество от \emph{ребра}.
\end{defn}

Тъй като всяко ребро е двуелементно множество, в неориентиран граф реброто между $u$ и $v$ се записва като $\{u,v\}$; редът на върховете няма значение. В тази дефиниция няма примки и две еднакви ребра не могат да се срещат.

\begin{defn}
Нека $G=(V,E)$ е граф.
\begin{enumerate}
    \item Ако $E=\varnothing$, графът се нарича \emph{празен}.
    \item Ако $|V|=1$, графът се нарича \emph{тривиален}.
    \item Върховете $u,v\in V$ са \emph{съседни}, ако $\{u,v\}\in E$.
    \item Ребрата $e_1,e_2\in E$ са \emph{инцидентни}, ако
    \[
    e_1\cap e_2\neq\varnothing.
    \]
    \item Реброто $e\in E$ е инцидентно с върха $u\in V$, ако $u\in e$.
\end{enumerate}
\end{defn}

\begin{defn}
За връх $u\in V$ дефинираме:
\[
N(u)=\{v\in V\mid u\text{ и }v\text{ са съседни}\},
\]
\[
N[u]=N(u)\cup\{u\},
\]
\[
I(u)=\{e\in E\mid u\in e\}.
\]
Числото
\[
d(u)=|I(u)|
\]
се нарича \emph{степен} на върха $u$.
\end{defn}

Максималната и минималната степен на графа се означават съответно с
\[
\Delta(G)=\max\{d(u)\mid u\in V\},
\qquad
\delta(G)=\min\{d(u)\mid u\in V\}.
\]

\subsection{Ориентирани графи и мултиграфи}\label{s:3.2.2}

\begin{defn}
Краен ориентиран граф е наредена двойка
\[
G=(V,E),
\]
където $V$ е непразно крайно множество от върхове и
\[
E\subseteq (V\times V)\setminus\{(u,u)\mid u\in V\}.
\]
Елементите на $E$ се наричат ориентирани ребра или дъги.
\end{defn}

Дъгата $(u,v)$ има начало $u$ и край $v$. За разлика от неориентираното ребро, като цяло $(u,v)\neq(v,u)$.

\begin{defn}
Неориентиран мултиграф е наредена тройка
\[
G=(V,E,f_G),
\]
където $V$ е непразно множество от върхове, $E$ е множество от ребра, $V\cap E=\varnothing$, и
\[
f_G:E\longrightarrow\{X\subseteq V\mid |X|=2\}
\]
е свързваща функция.
\end{defn}

Свързващата функция указва краищата на всяко ребро. Различни елементи на $E$ могат да имат един и същ образ чрез $f_G$, поради което мултиграфът позволява успоредни ребра.

\begin{defn}
Ориентиран мултиграф е наредена тройка
\[
G=(V,E,f_G),
\]
където $V$ е непразно множество от върхове, $E$ е множество от ребра, $V\cap E=\varnothing$, и
\[
f_G:E\longrightarrow V\times V
\]
е свързваща функция.
\end{defn}

\begin{supp}[Бележка]
В представените материали за ориентиран мултиграф се допуска $f_G(e)=(u,u)$, тоест примки. Това се различава от дадената по-горе дефиниция за обикновен ориентиран граф, в която примките са изключени.
\end{supp}

\section{Пътища, цикли и свързаност}\label{s:3.3}

\subsection{Пътища и цикли}\label{s:3.3.1}

\begin{defn}
Нека $G=(V,E,f_G)$ е мултиграф. \emph{Път} в $G$ е алтернираща редица от върхове и ребра
\[
p=(u_0,e_{k_0},u_1,e_{k_1},\ldots,e_{k_{l-1}},u_l),
\]
където $l\geq 0$, $u_i\in V$ за $0\leq i\leq l$, $e_{k_j}\in E$ за $0\leq j\leq l-1$, и всяко ребро $e_{k_j}$ свързва последователните върхове $u_j$ и $u_{j+1}$.
\end{defn}

При ориентиран мултиграф изискването е преминаването по всяко ребро да е в посоката, зададена от свързващата функция.

Върховете $u_0$ и $u_l$ се наричат съответно начало и край на пътя, а останалите върхове са вътрешни. Дължината на пътя е броят ребра в него:
\[
\lVert p\rVert=l.
\]

\begin{defn}
Пътят $p$ се нарича \emph{прост}, ако всички негови върхове и ребра са различни.
\end{defn}

\begin{defn}
Пътят
\[
p=(u_0,e_{k_0},u_1,\ldots,e_{k_{l-1}},u_l)
\]
се нарича \emph{цикъл}, ако $u_0=u_l$. Той е \emph{прост цикъл}, ако има поне едно ребро, върховете $u_0,u_1,\ldots,u_{l-1}$ са по двойки различни и ребрата му са по двойки различни. С други думи, единственото допустимо повторение на връх е $u_0=u_l$.
\end{defn}

\subsection{Свързаност}\label{s:3.3.2}

\begin{defn}
Нека $G=(V,E)$ е неориентиран граф. Върховете $u,v\in V$ са \emph{свързани}, ако съществува път от $u$ до $v$. Графът $G$ е \emph{свързан}, ако всеки два негови върха са свързани.
\end{defn}

\begin{defn}
Ориентиран граф е \emph{слабо свързан}, ако неориентираният граф, получен чрез премахване на посоките на дъгите, е свързан. Той е \emph{полусвързан}, ако за всеки два върха $u$ и $v$ има ориентиран път поне в едната посока: от $u$ до $v$ или от $v$ до $u$. Той е \emph{силно свързан}, ако за всеки два върха има ориентиран път и от първия до втория, и от втория до първия.
\end{defn}

\begin{defn}
Нека $G=(V,E)$ е неориентиран граф. Релацията на достижимост $R_G$ върху $V$ е зададена чрез
\[
R_G\subseteq V\times V,
\qquad
uR_Gv\Longleftrightarrow u\text{ и }v\text{ са свързани}.
\]
\end{defn}

\begin{prop}
Релацията $R_G$ на достижимост в неориентиран граф е релация на еквивалентност.
\end{prop}

\begin{proof}
За всеки връх $u$ съществува тривиален път с дължина $0$ от $u$ до $u$, следователно $uR_Gu$.

Ако има път от $u$ до $v$, същата последователност от ребра, обходена в обратен ред, е път от $v$ до $u$. Следователно $uR_Gv$ предполага $vR_Gu$.

Накрая, ако има път от $u$ до $v$ и път от $v$ до $w$, те могат да се съединят в път от $u$ до $w$. Следователно релацията е транзитивна.
\end{proof}

\begin{defn}
Подграфите на $G$, индуцирани от класовете на еквивалентност на $R_G$, се наричат \emph{свързани компоненти} на $G$.
\end{defn}

Еквивалентно, свързаните компоненти са максималните по включване свързани подграфи. Следователно намирането на свързаните компоненти е конкретен случай на определяне на класовете на еквивалентност на релация.

\section{Дървета}\label{s:3.4}

\begin{defn}
\emph{Дърво} е неориентиран граф, който е свързан и ацикличен.
\end{defn}

Тривиалният граф е дърво. В дървото между два върха няма цикъл; при добавяне на нов връх към дърво и свързването му с точно едно ребро към стар връх не се създава цикъл.

\begin{defn}
\emph{Кореново дърво} определяме индуктивно. Базата е тривиалният граф
\[
(\{r\},\varnothing),
\]
който има корен $r$. В индуктивната стъпка към вече построено кореново дърво се добавя нов връх $u$ и едно ребро, което го свързва с връх $v$ от старото дърво. Коренът се запазва.
\end{defn}

\begin{thm}
Всяко кореново дърво е дърво.
\end{thm}

\begin{proof}
Доказателството е чрез индукция по построението. Базовият едновърхов граф е свързан и няма цикъл.

Нека $T=(V,E)$ е свързано и ациклично дърво. Добавяме нов връх $u\notin V$ и едно ребро $\{v,u\}$, където $v\in V$. Полученият граф е свързан, защото от всеки стар връх има път до $v$, а след това реброто $\{v,u\}$ води до $u$. Новият връх е инцидентен само с едно ребро, затова не може да участва в цикъл. Старите ребра не образуват цикъл по индукционното предположение. Следователно полученият граф също е дърво.
\end{proof}

Обратно, всяко крайно дърво с избран корен може да се получи чрез това построение. Ако има повече от един връх, дървото има листо, различно от корена. Премахването на листото и единственото му инцидентно ребро оставя дърво със същия корен; твърдението следва с индукция по броя върхове, като накрая листото и реброто се добавят обратно.

Следователно кореновото дърво може еквивалентно да се разглежда като дърво с избран връх $r$, наречен \emph{корен}. Изборът на корен задава йерархия: всеки връх, различен от корена, има единствен непосредствен предшественик по пътя към корена.

\begin{thm}
Нека $T=(V,E)$ е кореново дърво. Тогава
\[
|V|=|E|+1.
\]
\end{thm}

\begin{proof}
Използваме индукция по построението на кореновото дърво. За базата
\[
T=(\{r\},\varnothing)
\]
имаме $|V|=1=0+1=|E|+1$.

При индуктивната стъпка се добавят едновременно един нов връх и едно ново ребро. Ако за старото дърво е изпълнено $|V|=|E|+1$, то за новото дърво $T'$ е в сила
\[
|V(T')|=|V|+1=|E|+2=|E(T')|+1.
\]
\end{proof}

\begin{defn}
Нека $G=(V,E)$ е граф. \emph{Покриващо дърво} на $G$ е дърво
\[
T=(V,E'),
\qquad E'\subseteq E.
\]
\end{defn}

Покриващото дърво съдържа всички върхове на графа, но само част от неговите ребра.

\begin{thm}
Краен непразен неориентиран граф $G=(V,E)$ има покриващо дърво тогава и само тогава, когато е свързан.
\end{thm}

\begin{proof}
Ако има покриващо дърво, пътят в него между всяка двойка върхове е път и в $G$, следователно $G$ е свързан. Обратно, за краен непразен свързан неориентиран граф започваме с един връх. Докато не сме включили всички върхове, свързаността осигурява ребро от включен връх към невключен. Добавяме това ребро и новия връх. Свързаността на построеното подмножество се запазва и цикъл не възниква, защото новият връх има само едно добавено ребро. След $|V|-1$ стъпки получаваме покриващо дърво.
\end{proof}

\section{Обхождания на графи}\label{s:3.5}

Обхождането систематично посещава върховете, достижими от начален връх. При откриването на нов връх се запомня реброто, по което е достигнат; така получените ребра образуват покриващо дърво на достигнатата свързана компонента.

\subsection{Стек и опашка}\label{s:3.5.1}

\begin{defn}
\emph{Стек} е структура от тип LIFO: последно добавеният елемент се извежда първи. Празният стек означаваме с $\{\}$. Ако $S$ е стек и $x$ е елемент, то $\operatorname{push}(S,x)$ добавя $x$ на върха на стека, $\operatorname{top}(S)$ връща горния елемент, а $\operatorname{pop}(S)$ премахва горния елемент.
\end{defn}

\begin{defn}
\emph{Опашка} е структура от тип FIFO: първо добавеният елемент се извежда първи. Операцията $\operatorname{push}(Q,x)$ добавя елемент в края, $\operatorname{top}(Q)$ връща първия елемент, а $\operatorname{pop}(Q)$ го премахва.
\end{defn}

\subsection{Схема за обхождане}\label{s:3.5.2}

Нека върховете на ориентиран мултиграф са означени с $1,\ldots,n$, а началният връх е $1$. Поддържа се булев масив $\operatorname{visited}[1,n]$. В началото всички стойности са \texttt{False}; началният връх се поставя в помощна структура $S$ и се маркира като посетен.

Докато $S$ не е празна, от нея се извежда връх $x$. За всяка дъга $(x,y)$, ако $y$ не е посещаван, той се маркира като посетен и се добавя в $S$. Изборът на конкретна структура за $S$ определя вида обхождане.

\subsection{Обхождане в широчина}\label{s:3.5.3}

\begin{keydefn}
\emph{Обхождане в широчина} (BFS) е обхождане по общата схема, при което помощната структура е опашка.
\end{keydefn}

BFS първо разглежда всички върхове на разстояние $1$ от началния връх, след това върховете на разстояние $2$ и т.н. Ако при откриването на $y$ чрез $x$ запомним $\pi[y]=x$, стойностите $\pi$ задават дърво на обхождането с корен началния връх.

\begin{verbatim}
BFS(G=(V,E))
for i <- 1 to n
    colour[i] <- white
    pi[i] <- NIL
    d[i] <- infinity
colour[1] <- gray
d[1] <- 0
create queue Q
Enqueue(Q,1)
while not isempty(Q)
    x <- Dequeue(Q)
    for y in adj(x)
        if colour[y] = white
            colour[y] <- gray
            pi[y] <- x
            d[y] <- d[x] + 1
            Enqueue(Q,y)
    colour[x] <- black
return colour,pi,d
\end{verbatim}

\begin{keythm}
Обхождането в широчина намира най-късите пътища от началния връх до всички достижими върхове в графа.
\end{keythm}

\begin{proof}
Разстоянието тук е броят на ребрата. Началният връх $s$ получава $d[s]=0$. FIFO опашката обработва изцяло върховете на слой $k$, преди върховете на слой $k+1$. По индукция нека всички върхове с действително разстояние най-много $k$ са открити с правилни стойности. Връх $v$ на разстояние $k+1$ има предшественик на разстояние $k$ по най-къс път; при обработката му $v$ се открива най-късно със стойност $k+1$. Всяка присвоена стойност е дължина на реален път и не може да е по-малка от най-късото разстояние. Така $d[v]=k+1$. Всички достижими върхове попадат в някой слой, което завършва доказателството.
\end{proof}

Стойността $d[y]$, присвоена при първото откриване на $y$, е дължината на най-къс път от началния връх до $y$. Това следва от факта, че опашката обработва върховете по нарастване на разстоянието от началния връх.

\subsection{Обхождане в дълбочина}\label{s:3.5.4}

\begin{keydefn}
\emph{Обхождане в дълбочина} (DFS) е обхождане, при което се продължава към необходен съсед, докато това е възможно; когато няма такъв съсед, алгоритъмът се връща назад. Итеративната реализация използва стек, тоест структура LIFO.
\end{keydefn}

Следната рекурсивна реализация изразява същата идея:

\begin{verbatim}
DFS(G=(V,E))
for i <- 1 to n
    colour[i] <- white
    N[i] <- NIL
DFS-Visit(G,1)
return colour,N

DFS-Visit(G=(V,E), x in V)
colour[x] <- gray
foreach y in adj[x]
    if colour[y] = white
        N[y] <- x
        DFS-Visit(G,y)
colour[x] <- black
\end{verbatim}

При достигане на необходен връх $y$ от $x$, стойността $N[y]=x$ запомня непосредствения предшественик на $y$ в дървото на DFS. Ако графът не е свързан, едно обхождане от избран начален връх посещава само неговата компонента; за обхождане на всички върхове алгоритъмът се стартира от всеки все още необходен връх.

И DFS, и BFS могат да бъдат изменени за намиране на път между два върха: единият се избира за начален, а обхождането се прекратява при достигане на другия. Пътят се възстановява чрез запомнените предшественици.

\section{Ойлерови пътища и цикли}\label{s:3.6}

\begin{defn}
Нека $G$ е свързан мултиграф. \emph{Ойлеров път} е път, който преминава точно веднъж през всяко ребро на графа. \emph{Ойлеров цикъл} е Ойлеров път, чийто начален и краен връх съвпадат. Граф, който притежава Ойлеров цикъл, се нарича \emph{Ойлеров граф}.
\end{defn}

\begin{keythm}[Критерий за Ойлеров цикъл в свързан мултиграф]
Нека $G$ е свързан мултиграф (в частност --- свързан граф) с поне едно ребро. Тогава $G$ е Ойлеров тогава и само тогава, когато всеки негов връх има четна степен. Кратните ребра и примките не изискват отделно разглеждане: всяка примка добавя $2$ към степента на върха си, а кратните ребра се броят поотделно.
\end{keythm}

\begin{proof}
Нека $G$ има Ойлеров цикъл. Всеки път, когато цикълът влиза във връх, той трябва да го напусне по друго ребро. Следователно ребрата, инцидентни с произволен връх, се групират по двойки: едно за влизане и едно за излизане. Значи степента на всеки връх е четна.

Обратно, нека $G$ е свързан и всеки връх има четна степен. Избираме начален връх и вървим по неизползвани ребра, като всяко използвано ребро се маркира. Пътят не може да завърши в връх, различен от началния: при такъв краен връх броят използвани инцидентни ребра би бил нечетен, което противоречи на четността на степента му. Така получаваме цикъл.

Ако цикълът съдържа всички ребра, той е Ойлеров. В противен случай поради свързаността съществува връх от получения цикъл, инцидентен с неизползвано ребро. От този връх се построява нов цикъл по неизползвани ребра и се вмъква в стария. Процесът се повтаря. При всяка стъпка броят неизползвани ребра намалява, затова след краен брой стъпки се получава цикъл, съдържащ всяко ребро точно веднъж.
\end{proof}

\begin{cor}
Свързан мултиграф има Ойлеров път, който не е цикъл, тогава и само тогава, когато точно два негови върха са с нечетна степен, а всички останали са с четна степен.
\end{cor}

\begin{proof}
При Ойлеров път всяко посещение на вътрешен връх използва двойка ребра: за влизане и излизане. Само началото и краят имат по едно несдвоено ребро, затова са точно двата върха с нечетна степен. Обратно, ако нечетните върхове са $u,v$, добавяме още едно ребро между тях. Всички степени стават четни, а графът остава свързан, затова по теоремата за Ойлеров цикъл има такъв цикъл. Премахваме добавеното ребро и прочитаме останалата част от цикъла от единия му край: получаваме Ойлеров път от $u$ до $v$.
\end{proof}

\section{Примерни задачи по анотацията}\label{s:3.7}

Анотацията назовава три типа задачи. Общите алгоритми са дадени по-горе; тук и
трите са изпълнени докрай върху конкретни входове.

Навсякъде в този раздел работим с графа
\[
G:\quad V=\{1,2,3,4,5,6\},\qquad
E=\{12,\,13,\,23,\,24,\,35,\,45,\,46,\,56\},
\]
където $uv$ означава реброто между $u$ и $v$. Списъците на съседите са подредени
по нарастване, за да е обхождането еднозначно:
\[
1:2,3\quad 2:1,3,4\quad 3:1,2,5\quad 4:2,5,6\quad 5:3,4,6\quad 6:4,5 .
\]
Степените са
\[
\deg 1=2,\quad \deg 2=3,\quad \deg 3=3,\quad
\deg 4=3,\quad \deg 5=3,\quad \deg 6=2 .
\]

\subsection{Задача 1: покриващо дърво в широчина и в дълбочина}\label{s:3.7.1}

\begin{example}[BFS от връх $1$]
Опашката се обработва отляво надясно; при всеки връх се обхождат необходените му
съседи по реда от списъка.

\begin{center}
\begin{tabular}{llll}
\toprule
Изваден връх & Нови върхове & Ребра на дървото & Опашка след стъпката\\
\midrule
$1$ & $2,3$ & $12,\,13$ & $2,3$\\
$2$ & $4$   & $24$      & $3,4$\\
$3$ & $5$   & $35$      & $4,5$\\
$4$ & $6$   & $46$      & $5,6$\\
$5$ & ---   & ---       & $6$\\
$6$ & ---   & ---       & ---\\
\bottomrule
\end{tabular}
\end{center}

Покриващото дърво е $\{12,13,24,35,46\}$ --- пет ребра при шест върха, както
изисква $|V|=|E|+1$. Разстоянията от началния връх са
\[
d(1)=0,\quad d(2)=d(3)=1,\quad d(4)=d(5)=2,\quad d(6)=3,
\]
и това са дължините на най-късите пътища, защото BFS обхожда върховете по
нарастващо разстояние.
\end{example}

\begin{example}[DFS от връх $1$]
Рекурсията влиза колкото може по-дълбоко, преди да се върне:
\[
1\to 2\to 3\to 5\to 4\to 6,
\]
след което от $6$ няма необходени съседи и обхождането се връща по стека до $1$.
Записаните предшественици са
\[
N[2]=1,\quad N[3]=2,\quad N[5]=3,\quad N[4]=5,\quad N[6]=4,
\]
тоест покриващото дърво е $\{12,23,35,45,46\}$ и е път. Ребрата $13$, $24$ и
$56$ остават извън дървото --- те са обратни ребра към вече посетени върхове.
\end{example}

\begin{remark}
Двете дървета са различни, макар графът и началният връх да са едни и същи. При
такава задача винаги се посочва и редът на съседите: без него отговорът не е
еднозначен.
\end{remark}

\subsection{Задача 2: Ойлеров цикъл или доказателство, че такъв няма}\label{s:3.7.2}

\begin{example}[Отрицателен отговор за $G$]
В $G$ върховете $2,3,4,5$ са с нечетна степен. По критерия Ойлеров цикъл няма,
защото не всички степени са четни; по следствието няма и Ойлеров път, защото
върховете с нечетна степен не са точно два, а четири. Това е пълен отговор ---
позоваване на критерия, не търсене на цикъл „на ръка“.
\end{example}

\begin{example}[Построяване по Hierholzer]
Нека $K_5$ е пълният граф върху $\{1,2,3,4,5\}$. Всички степени са $4$, графът е
свързан, следователно има Ойлеров цикъл. Строим го на три стъпки.

Първо тръгваме от $1$ по неизползвани ребра, докато се върнем в $1$:
\[
C_0:\ 1\to 2\to 3\to 1
\qquad(\text{ребра } 12,23,13).
\]
Връх $2$ още има неизползвани ребра. Строим цикъл от $2$:
\[
2\to 4\to 5\to 2
\qquad(\text{ребра } 24,45,25),
\]
и го вмъкваме в $C_0$ на мястото на $2$:
\[
C_1:\ 1\to 2\to 4\to 5\to 2\to 3\to 1 .
\]
Остават $14,15,34,35$. Строим цикъл от $3$:
\[
3\to 4\to 1\to 5\to 3
\qquad(\text{ребра } 34,14,15,35),
\]
и го вмъкваме на мястото на $3$:
\[
C_2:\ 1\to 2\to 4\to 5\to 2\to 3\to 4\to 1\to 5\to 3\to 1 .
\]
Използвани са всички $\binom{5}{2}=10$ ребра, всяко по веднъж, а началото
съвпада с края. Това е Ойлеров цикъл.
\end{example}

\subsection{Задача 3: разбиване на ребрата на минимален брой пътища без общо ребро}\label{s:3.7.3}

Тук „път“ се разбира в смисъла на \S3.3: върхове могат да се повтарят. Иска се
разбиване на $E$, тоест всяко ребро принадлежи на точно един от пътищата и се
използва в него точно веднъж.

\begin{keythm}[Минимален брой пътища]
Нека $G$ е свързан мултиграф с поне едно ребро и нека точно $2k$ от върховете му
са с нечетна степен (броят на нечетните върхове винаги е четен). Тогава
минималният брой пътища, на които може да се разбие $E$ така, че никои два да
нямат общо ребро, е
\[
\max(k,1).
\]
\end{keythm}

\begin{proof}
\emph{Долна граница.} Разглеждаме произволно разбиване на $E$ на пътища и
фиксираме връх $v$. Всяко преминаване на път през $v$ като вътрешен връх
използва две от инцидентните на $v$ ребра; същото важи и когато път започва и
завършва в $v$. Само път, за който $v$ е точно единият от двата края, използва
нечетен брой инцидентни на $v$ ребра. Тъй като ребрата са разпределени между
пътищата без повторение,
\[
\deg v\equiv\#\{\text{пътища, за които } v \text{ е точно единият край}\}\pmod 2 .
\]
Следователно всеки връх с нечетна степен е край на поне един път. Всеки път има
най-много два края, а нечетните върхове са $2k$, откъдето пътищата са поне $k$.
Освен това $E\ne\varnothing$, тоест поне един път е нужен винаги.

\emph{Конструкция, достигаща границата.} Ако $k=0$, всички степени са четни и по
критерия за Ойлеров цикъл целият граф е един-единствен път (при това затворен),
тоест $1$ път е достатъчен.

Нека $k\geq 1$ и нека нечетните върхове са $u_1,v_1,\ldots,u_k,v_k$ --- $2k$
различни върха, разбити произволно на $k$ двойки. Добавяме $k$ нови ребра
$u_iv_i$. Всяка степен нараства с точно $1$ при нечетните върхове и остава
непроменена при останалите, затова в получения мултиграф $G'$ всички степени са
четни; $G'$ остава свързан. По критерия $G'$ има Ойлеров цикъл $C$.

Премахваме от $C$ добавените $k$ ребра. Понеже $2k$-те върха са различни, никои
две добавени ребра нямат общ връх, следователно никои две не стоят една до друга
в $C$; значи между всеки две последователни премахнати ребра остава поне едно
първоначално ребро. Циклично $C$ се разпада на точно $k$ непразни дъги. Всяка
дъга е път, дъгите нямат общо ребро, а обединението им е точно $E$. Получихме
$k$ пътя.
\end{proof}

\begin{cor}
За граф с няколко свързани компоненти минималният брой пътища е сумата от
$\max(k_i,1)$ по компонентите $i$, съдържащи поне едно ребро. Изолираните
върхове не изискват път.
\end{cor}

\begin{proof}
Път е свързан подграф, затова никой път не пресича две компоненти. Задачата се
разпада независимо по компонентите и се прилага теоремата за всяка от тях.
\end{proof}

\begin{example}[Разбиване на $G$]
В $G$ нечетните върхове са $2,3,4,5$, тоест $2k=4$ и $k=2$: два пътя са
необходими и достатъчни.

Двойките избираме кръстосано --- $(2,5)$ и $(3,4)$ --- защото тези ребра още не
съществуват. Добавяме $e_a=25$ и $e_b=34$. Сега всички степени са четни:
\[
\deg 1=2,\ \deg 2=4,\ \deg 3=4,\ \deg 4=4,\ \deg 5=4,\ \deg 6=2 .
\]
Един Ойлеров цикъл в разширения мултиграф е
\[
2\xrightarrow{12}1\xrightarrow{13}3\xrightarrow{e_b}4\xrightarrow{24}2
\xrightarrow{23}3\xrightarrow{35}5\xrightarrow{56}6\xrightarrow{46}4
\xrightarrow{45}5\xrightarrow{e_a}2 .
\]
Премахваме $e_a$ и $e_b$ и оставащите две дъги са търсените пътища:
\[
P_1:\ 2-1-3,
\qquad
P_2:\ 4-2-3-5-6-4-5 .
\]
Проверка: $P_1$ използва $12,13$; $P_2$ използва $24,23,35,56,46,45$. Заедно това
са всичките $8$ ребра, всяко по веднъж, и общо ребро няма. Краищата на двата
пътя са $2,3$ и $4,5$ --- точно четирите нечетни върха, както изисква долната
граница.
\end{example}

\begin{remark}
Отговорът не е единствен: друг избор на двойки или друг Ойлеров цикъл дава друга
двойка пътища. Единствен е \emph{броят} им. На изпита се иска и обосновка на
минималността, тоест аргументът за долната граница чрез нечетните върхове.
\end{remark}

\section{Изпитен фокус}\label{s:3.8}

\begin{itemize}
    \item Да се дадат определения за неориентиран граф, ориентиран граф, мултиграф, съседство, инцидентност, степен на връх, път, прост път и цикъл.
    \item Да се дефинират свързан граф, слаба и силна свързаност, релация на достижимост и свързани компоненти.
    \item Да се възпроизведе доказателствената идея, че достижимостта в неориентиран граф е релация на еквивалентност; класовете $\,$ѝ са свързаните компоненти.
    \item Да се дефинират дърво, кореново дърво и покриващо дърво.
    \item Да се докаже чрез индукция, че за кореново дърво е изпълнено $|V|=|E|+1$.
    \item Да се формулира критерият: граф има покриващо дърво точно когато е свързан.
    \item Да се опишат BFS и DFS, ролята на опашката при BFS и на стека или рекурсията при DFS, както и дървото на обхождането.
    \item Да се формулира, че BFS намира най-къси пътища от началния връх.
    \item Да се дефинират Ойлеров път и Ойлеров цикъл.
    \item Да се формулира и обясни критерият за Ойлеров цикъл: свързан граф е Ойлеров точно когато всички върхове са с четна степен.
    \item Да се опише идеята на алгоритъма на Hierholzer: строене на цикъл по неизползвани ребра и вмъкване на допълнителни цикли.
    \item Да се формулира условието за Ойлеров път, който не е цикъл: точно два върха са с нечетна степен.
    \item \textbf{Задача 1 от анотацията:} да се построи покриващо дърво на зададен граф в широчина или в дълбочина, с изписан ред на обхождането и на ребрата на дървото (\S3.7.1).
    \item \textbf{Задача 2 от анотацията:} да се построи Ойлеров цикъл или път в зададен мултиграф чрез вмъкване на цикли, или да се докаже, че такъв няма, чрез броя на върховете с нечетна степен (\S3.7.2).
    \item \textbf{Задача 3 от анотацията:} да се разбият ребрата на минимален брой пътища без общо ребро --- долна граница $k$ от $2k$-те нечетни върха, конструкция чрез сдвояване, Ойлеров цикъл и премахване на добавените ребра (\S3.7.3).
    \item Релациите на еквивалентност в \S3.1 са предварителен материал от Въпрос 1; по този въпрос от тях се използва само, че достижимостта е еквивалентност и класовете й са свързаните компоненти.
\end{itemize}


\chapter{Характеризация на регулярните езици. Теорема на Майхил-Нероуд.}
% AUTO-GENERATED by scripts/build_topics.py — do not edit.
% Edit topics/bodies/topic_4.tex or topics/preamble.tex instead.
%% Placeholder. Replace with the written chapter.
\textit{Източници:} PDF \texttt{4.pdf}, снимки \texttt{4та тема}

\begin{supp}[Бележка към картографирането]
4.pdf is image-only (Tier B); it covers automata and regular languages — confirm it reaches Myhill-Nerode
\end{supp}

\subsection*{Анотация от конспекта}
Да се докаже по индукция дадено твърдение върху естествените числа.


\chapter{Лема за разрастването за контекстносвободни езици. Незатвореност относно сечение и допълнение.}
%% Drafted from the mapped corpus by scripts/draft_topics.py.
% Model: gpt-5.6-terra; source characters: 19574.
\begin{konspekt}
\kreq{Дефиниции}{контекстносвободна граматика, извод, контекстносвободен език, дърво на извод --- \kref{5.1}.}
\kreq{Конструкции, без доказателства}{затвореност относно регулярните операции --- обединение, конкатенация, звезда на Клини --- \kref{5.4}.}
\kreq{Формулировка и доказателство}{лема за разрастването за КСЕ --- \kref{5.2}.}
\kreq{Пример с доказателство}{език, който не е контекстносвободен --- \kref{5.3}.}
\kreq{Доказателство}{незатвореност относно сечение --- \kref{5.5}.}
\kreq{Доказателство}{незатвореност относно допълнение --- \kref{5.6}.}
\kreq{Типове задачи}{по даден КСЕ да се построи КСГ с обосновка; да се определи с обосновка дали даден език е контекстносвободен.}
\kreq{Литература}{[34], [38], [43].}
\end{konspekt}

\section{Контекстносвободни граматики и дървета на извода}\label{s:5.1}

\begin{defn}
Контекстносвободна граматика (КСГ) е четворка
\[
G=\langle V,\Sigma,R,S\rangle,
\]
където $V$ е крайно множество от нетерминали, $\Sigma$ е крайна азбука от терминали, $V\cap\Sigma=\varnothing$, $S\in V$ е начален нетерминал и
\[
R\subseteq V\times(V\cup\Sigma)^*
\]
е крайно множество от правила. Правилото $(A,\alpha)\in R$ се записва като $A\to\alpha$.
\end{defn}

\begin{defn}[Едностъпков и краен извод]
Едностъпковият извод, породен от $G$, се означава с $\Rightarrow_G$. По определение
\[
\lambda A\rho\Rightarrow_G\lambda\alpha\rho,
\]
когато $A\to\alpha$ е правило на $G$ и $\lambda,\rho\in(V\cup\Sigma)^*$.
С $\Rightarrow_G^*$ означаваме рефлексивно-транзитивното затваряне на тази релация.
\end{defn}

\begin{defn}
Езикът, породен от граматиката $G$, е
\[
L(G)=\{\omega\in\Sigma^*\mid S\Rightarrow_G^*\omega\}.
\]
Език $L\subseteq\Sigma^*$ се нарича \emph{контекстносвободен език} (КСЕ), ако съществува КСГ $G$, за която $L=L(G)$.
\end{defn}

\begin{defn}[Дърво на извод]
Дърво на извод в $G$ е крайно наредено кореново дърво. Коренът е означен с начален нетерминал. Ако вътрешен връх е означен с $A$, последователните му деца са означени със символите от дясната страна на някое правило $A\to\alpha$. При $A\to\varepsilon$ използваме едно листо $\varepsilon$. Резултатът е редицата от означенията на листата отляво надясно, като $\varepsilon$ не добавя символ. За завършено дърво всички листа са терминали или $\varepsilon$.
\end{defn}

\begin{defn}[Най-ляв извод]
Извод е най-ляв, когато във всяка стъпка се заменя най-левият нетерминал в текущата дума.
\end{defn}

\begin{prop}[Изводи и дървета] За дума $\omega\in\Sigma^*$ са еквивалентни следните твърдения:

\begin{enumerate}
\item $S\Rightarrow_G^*\omega$, тоест $\omega\in L(G)$;
\item съществува дърво на извод с корен $S$ и резултат $\omega$;
\item съществува най-ляв извод на $\omega$ от $S$.
\end{enumerate}

\end{prop}
\begin{proof}
От извод строим дървото по индукция по броя стъпки: началното дърво е един лист $S$, а замяната на нетерминал разширява съответния лист с децата от приложеното правило. Резултатът на частичното дърво е текущата дума. Обратно, за дадено завършено дърво разширяваме винаги най-левия останал нетерминален лист според неговите деца в дървото. Крайността гарантира завършване и получаваме най-ляв извод на резултата. Всеки най-ляв извод е частен случай на извод. Това доказва трите импликации.
\end{proof}

\section{Лема за разрастването}\label{s:5.2}

Лемата за разрастването, наричана още лема за покачването, дава необходимо условие един език да е контекстносвободен. Тя е особено полезна за доказване, че даден език \emph{не} е контекстносвободен.

\begin{keythm}[Лема за разрастването за контекстносвободни езици]
Нека $L\subseteq\Sigma^*$ е контекстносвободен език. Тогава съществува число $p>0$, наречено \emph{константа на разрастването}, такова че за всяка дума $\alpha\in L$ с $|\alpha|\geq p$ съществува разлагане
\[
\alpha=xyuvw,
\]
за което са изпълнени:

\begin{enumerate}
\item $|yv|\geq 1$;
\item $|yuv|\leq p$;
\item за всяко $i\in\mathbb{N}$ е изпълнено
\[
xy^iuv^iw\in L.
\]
\end{enumerate}
\end{keythm}

Идеята е, че достатъчно дълга дума има достатъчно високо дърво на извод. По един път в това дърво неизбежно се повтаря нетерминал. Частта от дървото между двете срещания на същия нетерминал може да се повтаря произволен брой пъти или да се премахва.

\begin{proof}
Нека $G=\langle V,\Sigma,R,S\rangle$ е КСГ, за която $L=L(G)$. Без ограничение на общността приемаме, че $G$ е в нормална форма на Чомски. За непразните думи, които разглеждаме, всяко използвано правило е от един от видовете
\[
A\to BC
\qquad\text{или}\qquad
A\to a,
\]
където $A,B,C\in V$ и $a\in\Sigma$.

Нека
\[
k=|V|
\qquad\text{и}\qquad
p=2^k.
\]
Ще покажем, че $p$ е подходяща константа.

Използваме следното елементарно наблюдение за двоични дървета: ако във всяко корен--листо пътче има по-малко от $k$ ребра, то дървото има по-малко от $2^k$ листа. Действително, на всяко ниво броят на върховете може най-много да се удвои, а листата са не повече от броя на възможните пътища с дължина, по-малка от $k$.

Нека сега $\alpha\in L$ и $|\alpha|\geq p=2^k$. Разглеждаме дърво на извод за $\alpha$ и в него избираме най-дълъг корен--листо път. Тъй като граматиката е в нормална форма на Чомски, дървото е двоично: всеки вътрешен връх има най-много две деца. Броят на листата е $|\alpha|$, следователно дървото има поне $2^k$ листа. Ако всеки път имаше най-много $k$ нетерминални върха, щеше да има най-много $k-1$ двоични разклонявания и най-много $2^{k-1}$ терминални листа. Понеже листата са поне $2^k$, избраният най-дълъг път има поне $k+1$ нетерминални върха.

По този път се срещат поне $k+1$ нетерминала, а нетерминалите в $V$ са само $k$. Следователно, по принципа на Дирихле, по пътя има два върха, означени с един и същ нетерминал $A$.

Сред последните $k+1$ срещания на нетерминали по този най-дълъг път избираме две срещания на един и същ нетерминал $A$. Съответстващата част от дървото дава разлагане
\[
\alpha=xyuvw
\]
и изводи
\[
S\Rightarrow_G^* xAw,
\qquad
A\Rightarrow_G^* yAv,
\qquad
A\Rightarrow_G^* u.
\]
Оттук за всяко $i\in\mathbb{N}$ получаваме
\[
A\Rightarrow_G^* y^iAv^i
\Rightarrow_G^* y^iuv^i,
\]
а следователно
\[
S\Rightarrow_G^* xy^iuv^iw.
\]
Значи
\[
xy^iuv^iw\in L
\]
за всяко $i\in\mathbb{N}$.

Остава да проверим първите две условия. Изводът $A\Rightarrow_G^* yAv$ между двете различни срещания на $A$ съдържа поне една стъпка. При първото разклоняване по избрания път другото дете е корен на поддърво, което в нормална форма на Чомски извежда поне един терминален символ. Този символ попада в $y$ или във $v$. Следователно
\[
|yv|\geq 1.
\]

Горното срещане на $A$ е сред последните $k+1$ нетерминални върха на избрания най-дълъг път. Затова от него до листото има най-много $k$ разклонявания. Понеже избраният път е най-дълъг в цялото дърво, никой страничен път в поддървото на това срещане не е по-дълъг. Следователно височината на поддървото е най-много $k$. То е двоично, затова има най-много $2^k=p$ терминални листа. Тези листа дават точно думата $yuv$, откъдето
\[
|yuv|\leq p.
\]
Така трите условия са доказани.
\end{proof}

\begin{remark}
В лемата е съществено разлагането да съществува за всяка достатъчно дълга дума, но то не е предварително зададено. При приложение за доказване на неконтекстносвободност трябва да се разгледат \emph{всички} разлагания, удовлетворяващи условията на лемата, и за всяко от тях да се намери стойност на $i$, при която получената дума не принадлежи на разглеждания език.
\end{remark}

\section{Приложение: език, който не е контекстносвободен}\label{s:5.3}

\begin{example}
Езикът
\[
L=\{a^n b^n c^n\mid n\in\mathbb{N}\}
\]
не е контекстносвободен.
\end{example}

\begin{proof}
Да допуснем противното: нека $L$ е контекстносвободен. Нека $p$ е константата от лемата за разрастването. Избираме думата
\[
\alpha=a^p b^p c^p\in L.
\]
Тя има дължина поне $p$, следователно според лемата има разлагане
\[
\alpha=xyuvw,
\]
за което
\[
|yuv|\leq p,
\qquad
|yv|\geq 1,
\]
и
\[
xy^iuv^iw\in L
\]
за всяко $i\in\mathbb{N}$.

Тъй като всяка от трите еднородни части $a^p$, $b^p$, $c^p$ има дължина $p$, поддумата $yuv$ с дължина най-много $p$ не може да съдържа едновременно буквите $a$, $b$ и $c$. Следователно $y$ и $v$ засягат най-много две съседни групи от буквите.

Разглеждаме думата при $i=2$:
\[
xy^2uv^2w.
\]
Поне едно от $y,v$ е непразно.

Ако $y$ и $v$ са съставени само от една буква, повторението увеличава броя на поне една от буквите $a,b,c$, без да увеличава броя и на трите по един и същ начин. Следователно в получената дума броят на $a$, $b$ и $c$ не е еднакъв, така че тя не принадлежи на $L$.

Ако някое от $y$ и $v$ съдържа две различни букви, то то пресича границата между две от групите $a^p$, $b^p$, $c^p$. При повторяването му се нарушава формата
\[
a^*b^*c^*.
\]
Следователно и в този случай
\[
xy^2uv^2w\notin L.
\]

Получаваме противоречие с третото условие на лемата за разрастването. Следователно $L$ не е контекстносвободен.
\end{proof}

\section{Затвореност относно регулярните операции}\label{s:5.4}

Конспектът иска \emph{конструкции} за затвореността относно регулярните
операции --- обединение, конкатенация и звезда на Клини --- без доказателства.
Конструкциите по-долу работят за произволни входни граматики, а не само за
конкретните примери от следващите раздели.

Нека
\[
G_1=\langle V_1,\Sigma,R_1,S_1\rangle
\qquad\text{и}\qquad
G_2=\langle V_2,\Sigma,R_2,S_2\rangle
\]
са контекстносвободни граматики над обща азбука $\Sigma$. Преди всяка от
конструкциите преименуваме нетерминалите така, че
\[
V_1\cap V_2=\varnothing,
\]
и избираме нов начален нетерминал $S\notin V_1\cup V_2$. Преименуването не
променя породения език: то е просто биекция върху множеството от нетерминали,
приложена едновременно към лявата и дясната страна на всяко правило. Без него
двете граматики биха си споделили нетерминал и изводите им биха се смесили.

\begin{prop}[Конструкция за обединение]
Езикът $L(G_1)\cup L(G_2)$ се поражда от граматиката
\[
G_\cup=\langle V_1\cup V_2\cup\{S\},\ \Sigma,\ R_1\cup R_2\cup\{S\to S_1,\ S\to S_2\},\ S\rangle .
\]
\end{prop}
\begin{proof}
Първата стъпка на всеки извод от $S$ е $S\to S_1$ или $S\to S_2$; други правила
за $S$ няма. Тъй като $V_1\cap V_2=\varnothing$, след избора на $S_1$ всички
достижими нетерминали са от $V_1$ и приложимите правила са точно тези на $R_1$,
тоест изводът продължава като извод в $G_1$; аналогично за $S_2$ и $G_2$.
Следователно $S\Rightarrow^*_{G_\cup}\omega$ точно когато $S_1\Rightarrow^*_{G_1}\omega$
или $S_2\Rightarrow^*_{G_2}\omega$.
\end{proof}

\begin{prop}[Конструкция за конкатенация]
Езикът $L(G_1)L(G_2)=\{\omega_1\omega_2\mid \omega_1\in L(G_1),\ \omega_2\in L(G_2)\}$
се поражда от граматиката
\[
G_\cdot=\langle V_1\cup V_2\cup\{S\},\ \Sigma,\ R_1\cup R_2\cup\{S\to S_1S_2\},\ S\rangle .
\]
\end{prop}
\begin{proof}
Единственото правило за $S$ дава $S\Rightarrow S_1S_2$. Двата нетерминала
разполагат с непресичащи се множества нетерминали, затова поддървото над $S_1$
използва само правила от $R_1$, а поддървото над $S_2$ --- само правила от $R_2$.
Резултатът е конкатенация на дума, изводима от $S_1$ в $G_1$, и дума, изводима
от $S_2$ в $G_2$; обратно, всяка такава двойка дава извод в $G_\cdot$.
\end{proof}

\begin{prop}[Конструкция за звезда на Клини]
Езикът $L(G_1)^*$ се поражда от граматиката
\[
G_*=\langle V_1\cup\{S\},\ \Sigma,\ R_1\cup\{S\to S_1S,\ S\to\varepsilon\},\ S\rangle .
\]
\end{prop}
\begin{proof}
Правилата за $S$ пораждат точно думите $S_1^k$ за $k\in\mathbb{N}$: $k$ пъти се
прилага $S\to S_1S$ и веднъж $S\to\varepsilon$. Всяко от $k$-те срещания на
$S_1$ извежда независимо дума от $L(G_1)$, тъй като $S\notin V_1$ и правилата на
$R_1$ не могат да въведат $S$. При $k=0$ получаваме $\varepsilon$, което
принадлежи на $L(G_1)^*$ по определение.
\end{proof}

\begin{remark}
Трите конструкции обясняват и защо аналогичен подход не работи за сечение: няма
как с едно правило да се наложи \emph{една и съща} дума да бъде изведена
едновременно в двете граматики. Следващите два раздела показват, че такава
конструкция изобщо не съществува.
\end{remark}

\section{Незатвореност относно сечение}\label{s:5.5}

Класът на контекстносвободните езици е затворен относно обединение, конкатенация и звездата на Клини --- това дадоха конструкциите от предходния раздел. В частност, ако $L_1$ и $L_2$ са КСЕ над една и съща азбука, то $L_1\cup L_2$ е КСЕ. За сечение обаче съответното свойство не е вярно.

\begin{thm}
Класът на контекстносвободните езици не е затворен относно сечение.
\end{thm}

\begin{proof}
Разглеждаме езиците
\[
L_1=\{a^n b^n c^k\mid n,k\in\mathbb{N}\}
\]
и
\[
L_2=\{a^n b^k c^k\mid n,k\in\mathbb{N}\}.
\]
И двата езика са контекстносвободни. Например за $L_1$ може да се използва граматиката с правила
\[
S\to AC,\qquad
A\to aAb\mid\varepsilon,\qquad
C\to cC\mid\varepsilon.
\]
Нетерминалът $A$ поражда еднакъв брой букви $a$ и $b$, а $C$ поражда произволен брой букви $c$.

Аналогично за $L_2$ може да се използва граматиката
\[
S\to AD,\qquad
A\to aA\mid\varepsilon,\qquad
D\to bDc\mid\varepsilon.
\]
Тук $A$ поражда произволен брой букви $a$, а $D$ --- еднакъв брой букви $b$ и $c$.

В дума от сечението $L_1\cap L_2$ броят на $a$ и $b$ трябва да е еднакъв поради принадлежността към $L_1$, а броят на $b$ и $c$ трябва да е еднакъв поради принадлежността към $L_2$. Следователно
\[
L_1\cap L_2
=
\{a^n b^n c^n\mid n\in\mathbb{N}\}.
\]
Но този език не е контекстносвободен според предходния пример. Следователно сечението на два контекстносвободни езика не е непременно контекстносвободно.
\end{proof}

\section{Незатвореност относно допълнение}\label{s:5.6}

Допълнението на език $L\subseteq\Sigma^*$ винаги се разбира относно фиксираната азбука $\Sigma$:
\[
\overline{L}=\Sigma^*\setminus L.
\]

\begin{thm}
Класът на контекстносвободните езици не е затворен относно допълнение.
\end{thm}

\begin{proof}
Да допуснем противното: нека за всеки контекстносвободен език $L$ и неговото допълнение $\overline{L}$ също е контекстносвободно.

Нека $L_1$ и $L_2$ са контекстносвободните езици от доказателството за незатвореност относно сечение. По допускане $\overline{L_1}$ и $\overline{L_2}$ са КСЕ. Тъй като КСЕ са затворени относно обединение,
\[
\overline{L_1}\cup\overline{L_2}
\]
е КСЕ. Прилагаме отново допускането за затвореност относно допълнение и получаваме, че
\[
\overline{\overline{L_1}\cup\overline{L_2}}
\]
също е КСЕ.

По закона на де Морган обаче
\[
\overline{\overline{L_1}\cup\overline{L_2}}
=
L_1\cap L_2.
\]
Това противоречи на вече доказаната незатвореност относно сечение. Следователно класът на контекстносвободните езици не е затворен относно допълнение.
\end{proof}

\begin{supp}[Бележка]
От незатвореността относно сечение сама по себе си не следва непосредствено незатвореност относно допълнение. В доказателството е необходимо и известното свойство, че контекстносвободните езици са затворени относно обединение, за да се приложи законът на де Морган.
\end{supp}

\section{Изпитен фокус}\label{s:5.7}

\begin{itemize}
\item Да се дадат определенията за контекстносвободна граматика, извод, език $L(G)$ и контекстносвободен език.
\item Да се обясни връзката между извод, най-ляв извод и дърво на синтактичен анализ.
\item Да се дадат конструкциите за затвореност относно обединение, конкатенация и звезда на Клини за \emph{произволни} входни граматики: преименуване на нетерминалите до непресичащи се множества, нов начален нетерминал и правилата $S\to S_1\mid S_2$, $S\to S_1S_2$, $S\to S_1S\mid\varepsilon$. Конспектът иска конструкциите, но не и доказателства към тях.
\item Да се формулира точно лемата за разрастването: разлагането $\alpha=xyuvw$, условията $|yv|\geq 1$, $|yuv|\leq p$ и $xy^iuv^iw\in L$ за всяко $i\in\mathbb{N}$.
\item Да се възпроизведе доказателствената идея чрез дърво в нормална форма на Чомски: двоично дърво, дълъг път, принцип на Дирихле и повторение на нетерминал.
\item Да се обясни защо повтореният нетерминал позволява „напомпване“ на частите $y$ и $v$.
\item Да се докаже с лемата, че $\{a^n b^n c^n\mid n\in\mathbb{N}\}$ не е КСЕ, като се разгледат всички допустими разлагания.
\item Да се дадат КСЕ $L_1=\{a^n b^n c^k\}$ и $L_2=\{a^n b^k c^k\}$, чието сечение е $\{a^n b^n c^n\}$.
\item Да се изведе незатвореността относно допълнение от незатвореността относно сечение чрез затвореността относно обединение и закона на де Морган.
\end{itemize}


\chapter{Сортиране чрез сравнения във време O(n lg n).}
% AUTO-GENERATED by scripts/build_topics.py — do not edit.
% Edit topics/bodies/topic_6.tex or topics/preamble.tex instead.
%% Placeholder. Replace with the written chapter.
\textit{Източници:} снимки \texttt{6та тема}

\begin{supp}[Бележка към картографирането]
6.pdf is STALE — it holds 'Модели на изчисленията: машина на Тюринг, RAM', a question from an earlier конспект. The photographs (двоична пирамида, HEAPSORT) are correct. Do not use 6.pdf for this chapter.
\end{supp}

\subsection*{Анотация от конспекта}
Сортиране чрез сравнения във време O(n lg n). Двоична пирамида (binary heap): определение и свойства. Наивно построяване на двоична пирамида: доказателство за коректност и изследване на сложността по време. Бързо построяване на двоична пирамида, използващо функция Heapify: доказателство за коректност и изследване на сложността по време. Сортиращ алгоритъм HEAPSORT: псевдокод, доказателство за коректност и изследване на сложността по време. Сортиращ алгоритъм MERGESORT като пример за алгоритъм по схемата Разделяй и Владей. Псевдокод и доказателство за коректност на MERGESORT. Изследване на сложността по време на MERGESORT. Приложение на MERGESORT за намиране на броя на инверсиите в масив от числа, с доказателство за коректност. Забележка: На изпита ще се падне или HEAPSORT, или MERGESORT. Литература: [31].


\chapter{Минимални покриващи дървета.}
%% Drafted from the mapped corpus by scripts/draft_topics.py.
% Model: gpt-5.6-terra; source characters: 21362.
\begin{konspekt}
\kselect{На изпита ще се падне \textbf{или} алгоритъмът на Прим, \textbf{или} алгоритъмът на Крускал.}
\kreq{Дефиниция на задачата}{какъв вид графи се разглеждат и какво се търси --- \kref{7.1}.}
\kreq{Теорема с доказателство}{за съгласуваното множество --- условия за нарастване на подмножество на МПД --- \kref{7.2}.}
\kreq{Псевдокод, коректност, сложност при различни структури от данни}{алгоритъм на Прим --- \kref{7.3}.}
\kreq{Псевдокод и коректност}{алгоритъм на Крускал --- \kref{7.4}.}
\kreq{Структура от данни}{Union--Find и приложението му за ефикасна реализация --- \kref{7.5}.}
\kreq{Сложност по време}{на алгоритъма на Крускал --- \kref{7.6}.}
\kreq{Литература}{[31].}
\end{konspekt}

\section{Задача за минимално покриващо дърво}\label{s:7.1}

\begin{defn}
Тегловен граф е наредената двойка $(G,\omega)$, където $G=(V,E)$ е граф, а
\[
\omega:E\to\mathbb{R}
\]
е теглова функция, която задава тегло, цена или дължина на всяко ребро.
\end{defn}

В задачата за минимално покриващо дърво се разглежда неориентиран свързан тегловен граф. Ориентацията не е част от задачата, защото понятието „дърво“ тук описва свързан ацикличен неориентиран подграф.

\begin{defn}
Нека $G=(V,E)$ е граф. Покриващ подграф на $G$ е всеки подграф от вида
\[
G'=(V,E'), \qquad E'\subseteq E.
\]
Покриващо дърво на $G$ е покриващ подграф, който е дърво.
\end{defn}

Тъй като графът $G$ е свързан, той има поне едно покриващо дърво. Всяко покриващо дърво съдържа всички върхове на $G$ и точно $|V|-1$ ребра.

\begin{defn}
Нека $G=(V,E)$ е свързан тегловен граф с теглова функция $\omega$, а $T$ е покриващо дърво на $G$. Тегло на $T$ наричаме сумата
\[
\omega(T)=\sum_{e\in E(T)}\omega(e).
\]
Покриващото дърво $T$ е \emph{минимално покриващо дърво}, съкратено МПД, ако за всяко покриващо дърво $T'$ на $G$ е изпълнено
\[
\omega(T)\leq \omega(T').
\]
\end{defn}

\begin{keydefn}
Задачата за минимално покриващо дърво има следния вид.

Екземплярът е наредена двойка $\langle G,\omega\rangle$, където $G=(V,E)$ е неориентиран свързан граф, а $\omega:E\to\mathbb{R}$ е теглова функция.

Решението е минимално покриващо дърво $T$ на $G$.
\end{keydefn}

Теглата не е необходимо да са положителни и не е необходимо да са различни. При равни тегла могат да съществуват различни минимални покриващи дървета.

\section{Съгласувани множества и безопасни ребра}\label{s:7.2}

Алчните алгоритми за МПД изграждат множество от ребра $A$, което постепенно се разширява. Основният въпрос е кога към $A$ може безопасно да се добави ново ребро.

\begin{defn}
Нека $G=(V,E)$ е тегловен граф и $A\subseteq E$. Казваме, че $A$ е \emph{съгласувано} множество, ако съществува минимално покриващо дърво $T$ на $G$, за което
\[
A\subseteq E(T).
\]
\end{defn}

\begin{defn}
Нека $A\subseteq E$ е съгласувано множество и $e\in E\setminus A$. Реброто $e$ е \emph{сигурно} за $A$, ако съществува минимално покриващо дърво $T'$ на $G$, за което
\[
A\cup\{e\}\subseteq E(T').
\]
\end{defn}

\begin{defn}
Срез в графа $G=(V,E)$ е разделяне на множеството от върхове на две непразни части:
\[
S=\{V_1,V_2\},\qquad V_1\cap V_2=\varnothing,\qquad V_1\cup V_2=V.
\]
Ребро $e=\{u,v\}$ \emph{пресича} среза $S$, ако единият му край е във $V_1$, а другият --- във $V_2$.

Срезът $S$ е \emph{съобразен} с множество $A\subseteq E$, ако нито едно ребро от $A$ не пресича $S$.

Ребро, пресичащо среза, е \emph{леко} за този срез, ако теглото му е минимално измежду теглата на всички ребра, които пресичат среза.
\end{defn}

\begin{keythm}[Теорема за съгласуваното множество]
Нека $G=(V,E)$ е свързан тегловен граф с теглова функция $\omega$. Нека $A\subseteq E$ е съгласувано множество и $S$ е срез, съобразен с $A$. Ако $e$ е произволно леко ребро, което пресича $S$, то $e$ е сигурно за $A$.
\end{keythm}

\begin{proof}
Тъй като $A$ е съгласувано, съществува минимално покриващо дърво $T$, за което
\[
A\subseteq E(T).
\]
Ако $e\in E(T)$, твърдението следва веднага.

Нека сега $e\notin E(T)$. Понеже $T$ е дърво, между двата края на $e$ има точно един път в $T$. След добавяне на $e$ към $T$ се получава точно един цикъл $C$.

Реброто $e$ пресича среза $S$. Следователно, при обхождането на цикъла $C$ трябва да се премине поне още веднъж от едната част на среза в другата. Значи в $C$ съществува ребро
\[
e'\neq e,
\]
което също пресича $S$.

Тъй като срезът е съобразен с $A$, реброто $e'$ не принадлежи на $A$. Премахваме $e'$ от цикъла $C$ и получаваме графа
\[
T'=T-\{e'\}+\{e\}.
\]
Това е покриващо дърво: след премахването на едно ребро от единствения цикъл графът остава свързан и става ацикличен. Освен това
\[
A\cup\{e\}\subseteq E(T'),
\]
защото $e'\notin A$.

Тъй като $e$ е леко ребро за среза, а $e'$ също пресича този срез, имаме
\[
\omega(e)\leq \omega(e').
\]
Следователно
\[
\omega(T')
 =\omega(T)-\omega(e')+\omega(e)
 \leq \omega(T).
\]
Но $T$ е минимално покриващо дърво, следователно $\omega(T)\leq\omega(T')$. Оттук
\[
\omega(T')=\omega(T),
\]
така че $T'$ също е минимално покриващо дърво. Значи $e$ е сигурно за $A$.
\end{proof}

Теоремата обосновава общия алчен подход:

\begin{verbatim}
Generic_MST(G = (V,E), omega)
A <- empty set                 // избраните ребра; винаги гора
while |A| < |V| - 1 do             // "A още не е покриващо дърво"
    избери ребро e, сигурно за A
    A <- A union {e}
return A
\end{verbatim}

Условието на цикъла е записано чрез броя ребра, защото гора върху $|V|$ върха е покриващо дърво тогава и само тогава, когато има точно $|V|-1$ ребра.

Ако на всяка стъпка $A$ е съгласувано и се добавя сигурно ребро, то след добавянето множеството остава съгласувано. Когато то вече образува покриващо дърво, това дърво трябва да е минимално.

Схемата не казва как се намира сигурно ребро --- това е единствената разлика между двата алгоритъма по-долу. Прим държи един срез между построеното дърво и останалите върхове и взема леко ребро за него; Крускал разглежда ребрата в ненамаляващ ред и за всяко прието ребро подходящият срез се намира впоследствие.

\section{Алгоритъм на Прим}\label{s:7.3}

Алгоритъмът на Прим строи едно дърво, започвайки от избран начален връх. На всяка стъпка той добавя най-лекото налично ребро, което свързва вече построената част на дървото с връх извън нея.

\subsection{Идея и поддържани данни}\label{s:7.3.1}

Нека $r\in V$ е начален връх. Поддържа се приоритетна опашка $Q$ от върховете, които все още не са добавени в дървото. За всеки връх $v$ се пазят две полета:

\begin{itemize}
\item $v.\mathit{key}$ --- теглото на най-лекото намерено досега ребро, което свързва $v$ с вече добавен връх; това е приоритетът на $v$ в опашката;
\item $v.\pi$ --- другият край на същото това ребро, тоест предшественикът на $v$ в строящото се дърво.
\end{itemize}

Ако още не е намерено ребро от дървото към $v$, тогава $v.\mathit{key}=\infty$ и $v.\pi=\textsc{Nil}$. За началния връх се поставя ключ $0$, за да бъде избран първи.

Дървото никъде не се пази явно. То се пази разпределено, в указателите $\pi$: щом връх бъде извлечен от опашката, реброто $\{v,v.\pi\}$ е окончателно избрано и вече не се променя.

\begin{keydefn}
Нека $X=V\setminus Q$ е множеството от вече извлечените от опашката върхове. Връх от $Q$ е граничен, ако има съсед в $X$; останалите върхове от $Q$ са неизвестни. Граничните върхове имат краен ключ, а неизвестните --- ключ $\infty$.
\end{keydefn}

Върховете от $X$ са точно тези, за които реброто към дървото е избрано окончателно; $\textsc{Extract-Min}$ винаги избира граничен връх, докато $Q$ съдържа поне един такъв, тъй като графът е свързан.

\subsection{Използвани операции}\label{s:7.3.2}

Псевдокодът извиква операции на приоритетна опашка. Те не са част от алгоритъма на Прим: всяка реализация, поддържаща ги, е допустима, а изборът на реализация влияе само на сложността.

\begin{defn}[Приоритетна опашка по минимум]
Приоритетна опашка е структура от данни, която пази множество от елементи, всеки с числов ключ, и поддържа операциите:

\begin{itemize}
\item \verb|Insert(Q,x)| --- добавя елемента $x$ с текущия му ключ;
\item \verb|Extract-Min(Q)| --- премахва от $Q$ и връща елемент с минимален ключ;
\item \verb|Decrease-Key(Q,x,k)| --- заменя ключа на $x$ с новата стойност $k$, която не е по-голяма от старата, и възстановява вътрешния ред на структурата.
\end{itemize}
\end{defn}

В псевдокода се използват още две означения:

\begin{itemize}
\item $Q\leftarrow V$ --- построяване на опашка от всички върхове с текущите им ключове, тоест $|V|$ операции \verb|Insert| или еднократно построяване;
\item $\mathit{Adj}[u]$ --- списъкът на съседство на $u$, тоест списъкът на върховете, съседни на $u$.
\end{itemize}

\subsection{Псевдокод}\label{s:7.3.3}

\begin{verbatim}
MST_Prim(G = (V,E), omega, r)
// 1. никой връх още не е свързан с дървото
for each v in V do
    v.key <- infinity        // най-лекото ребро от дървото до v
    v.pi  <- NIL             // другият край на това ребро
r.key <- 0                   // за да бъде r извлечен пръв
Q <- V                       // всички върхове чакат в опашката
// 2. всяка итерация добавя по един връх към дървото
while Q is not empty do
    u <- Extract-Min(Q)      // най-евтино присъединимият връх
    for each v in Adj[u] do  // ребрата от u към съседите му
        if v in Q and omega({u,v}) < v.key then
            v.pi  <- u       // u е новият най-евтин вход към v
            v.key <- omega({u,v})
            Decrease-Key(Q, v, v.key)
// 3. дървото се чете от указателите pi
return {{v, v.pi} : v in V and v.pi != NIL}
\end{verbatim}

Означенията в псевдокода са следните.

\begin{itemize}
\item $G=(V,E)$, $\omega$ и $r$ са входът: свързан неориентиран граф, теглова функция и начален връх.
\item $Q$ е приоритетната опашка на върховете, които още не са в дървото; $V\setminus Q$ са вече добавените върхове.
\item Тялото на \verb|while| се изпълнява по веднъж за всеки връх, защото \verb|Extract-Min| премахва връх и нищо не се връща обратно в $Q$.
\item Върнатото множество съдържа по едно ребро $\{v,v.\pi\}$ за всеки връх, различен от $r$; само за $r$ остава $v.\pi=\textsc{Nil}$.
\end{itemize}

В неориентиран граф всяко ребро се среща в списъците на съседство на двата си края. Проверката $v\in Q$ може да се реализира за константно време с отделна булева или битова информация дали върхът е в опашката.

\subsection{Пример}\label{s:7.3.4}

\begin{example}
Разглеждаме графа от фигура~\ref{fig:mst-example}, започвайки от връх $a$.

\begin{lecfigure}[Примерен тегловен граф. Удебелените ребра образуват минималното покриващо дърво с тегло $13$.\label{fig:mst-example}]
\begin{tikzpicture}[scale=1.25,
  vtx/.style={circle, draw=defframe, fill=defback, inner sep=0pt,
              minimum size=6.5mm, font=\small},
  eplain/.style={gray!65, line width=0.5pt},
  etree/.style={thmframe, line width=1.6pt},
  wlab/.style={font=\footnotesize, inner sep=1.4pt, fill=white}]
  \node[vtx] (a) at (0,1)  {$a$};
  \node[vtx] (b) at (1.9,2){$b$};
  \node[vtx] (c) at (1.9,0){$c$};
  \node[vtx] (d) at (3.8,2){$d$};
  \node[vtx] (e) at (3.8,0){$e$};
  \node[vtx] (f) at (5.7,1){$f$};
  \draw[eplain] (a) -- node[wlab] {4} (b);
  \draw[etree]  (a) -- node[wlab] {3} (c);
  \draw[etree]  (b) -- node[wlab] {1} (c);
  \draw[etree]  (b) -- node[wlab] {2} (d);
  \draw[eplain] (c) -- node[wlab] {4} (d);
  \draw[etree]  (c) -- node[wlab] {5} (e);
  \draw[eplain] (d) -- node[wlab] {7} (e);
  \draw[eplain] (d) -- node[wlab] {6} (f);
  \draw[etree]  (e) -- node[wlab] {2} (f);
\end{tikzpicture}
\end{lecfigure}

След инициализацията $a.\mathit{key}=0$, а всички останали ключове са $\infty$. Всеки ред от таблицата отговаря на една итерация на цикъла \verb|while|: първо е извлеченият връх и реброто, което той внася в дървото, а после ключовете, с които останалите върхове чакат в $Q$.

\begin{center}
\begin{tabular}{llll}
\toprule
Извлечен $u$ & Ребро $\{u,u.\pi\}$ & Тегло & Ключове в $Q$ след обхождането \\
\midrule
$a$ & няма & & $b{:}4(a)$, $c{:}3(a)$, $d{:}\infty$, $e{:}\infty$, $f{:}\infty$ \\
$c$ & $\{a,c\}$ & $3$ & $b{:}1(c)$, $d{:}4(c)$, $e{:}5(c)$, $f{:}\infty$ \\
$b$ & $\{c,b\}$ & $1$ & $d{:}2(b)$, $e{:}5(c)$, $f{:}\infty$ \\
$d$ & $\{b,d\}$ & $2$ & $e{:}5(c)$, $f{:}6(d)$ \\
$e$ & $\{c,e\}$ & $5$ & $f{:}2(e)$ \\
$f$ & $\{e,f\}$ & $2$ & $Q=\varnothing$ \\
\bottomrule
\end{tabular}
\end{center}

Забележете двете места, на които \verb|Decrease-Key| променя вече намерено ребро: ключът на $b$ пада от $4(a)$ на $1(c)$ и ключът на $f$ пада от $6(d)$ на $2(e)$. Реброто $\{d,e\}$ с тегло $7$ никога не подобрява ключа на $e$, а $\{a,b\}$ и $\{c,d\}$ се губят при пресмятанията на съответните ключове.

Полученото дърво е
\[
\{a,c\},\ \{c,b\},\ \{b,d\},\ \{c,e\},\ \{e,f\},
\]
с общо тегло $3+1+2+5+2=13$.
\end{example}

\subsection{Коректност}\label{s:7.3.5}

\begin{keythm}
За всеки свързан претеглен неориентиран граф алгоритъмът на Прим връща минимално покриващо дърво.
\end{keythm}

\begin{proof}
Нека $A$ е множеството от ребрата, определени чрез указателите $\pi$ за върховете, извлечени преди текущата итерация, без началния връх. Ще покажем по индукция по итерациите на цикъла, че $A$ е съгласувано множество.

Преди първата итерация $A=\varnothing$. Празното множество е съгласувано, защото свързаният претеглен неориентиран граф има минимално покриващо дърво.

Да предположим, че преди дадена итерация $A$ е съгласувано, и нека $u$ е върхът, избран с \verb|Extract-Min|. Ако $u=r$, не се добавя ребро, $A$ остава празно и инвариантът се запазва. В останалите итерации $u\neq r$. Нека преди извличането на $u$
\[
X=V\setminus Q
\]
е непразното множество от върхове, вече добавени към построяваното дърво. Реброто
\[
e=\{u,u.\pi\}
\]
свързва $u\in V\setminus X$ с връх от $X$.

Разглеждаме среза
\[
S=\{X,V\setminus X\}.
\]
Той е съобразен с $A$, защото всички ребра от $A$ имат и двата си края в $X$. По време на обхожданията на списъците на съседство алгоритъмът поддържа за всеки връх $v\in Q$ теглото на най-лекото ребро между $v$ и $X$. Следователно изборът на $u$ с минимален ключ означава, че $\{u,u.\pi\}$ пресича среза $S$ и е леко ребро за него.

По теоремата за съгласуваното множество реброто $e$ е сигурно за $A$. Следователно
\[
A\cup\{e\}
\]
също е съгласувано. Индуктивната хипотеза се запазва.

След завършване на алгоритъма всички върхове са извлечени. Всеки връх, различен от $r$, има точно един предшественик, следователно върнатото множество съдържа $|V|-1$ ребра. Всяко ново ребро свързва нов връх с вече построената част, затова не може да образува цикъл и всички върхове са свързани. Следователно резултатът е покриващо дърво. Понеже множеството от неговите ребра е съгласувано, то се съдържа в някое минимално покриващо дърво; а и двете дървета имат по $|V|-1$ ребра. Значи върнатото дърво е минимално покриващо.
\end{proof}

\subsection{Сложност на алгоритъма на Прим}\label{s:7.3.6}

Нека $n=|V|$ и $m=|E|$. При списъци на съседство всички обхождания на редовете с \verb|for each| разглеждат общо $O(m)$ съседства. Всеки връх се извлича от приоритетната опашка точно веднъж. Операцията \verb|Decrease-Key| се извиква най-много по веднъж за всяко разгледано съседство, следователно най-много $O(m)$ пъти.

\begin{prop}
При реализация на $Q$ с двоична пирамида алгоритъмът на Прим работи за време
\[
O((n+m)\log n)=O(m\log n)
\]
за свързан граф.
\end{prop}

\begin{proof}
Построяването на двоичната пирамида от всички върхове отнема $O(n)$ време. Има $n$ операции \verb|Extract-Min|, всяка за $O(\log n)$ време, общо $O(n\log n)$. Операцията \verb|Decrease-Key| също отнема $O(\log n)$ и се извиква най-много $O(m)$ пъти, което дава $O(m\log n)$.

Следователно общото време е
\[
O(n+m+n\log n+m\log n)=O((n+m)\log n).
\]
Понеже при свързан граф $m\geq n-1$, получаваме $O(m\log n)$.
\end{proof}

\begin{prop}
При реализация на $Q$ с пирамида на Фибоначи алгоритъмът на Прим работи за амортизирано време
\[
O(m+n\log n).
\]
\end{prop}

\begin{proof}
При пирамида на Фибоначи операцията \verb|Decrease-Key| е с амортизирана цена $O(1)$, а \verb|Extract-Min| --- с амортизирана цена $O(\log n)$. Затова $O(m)$ намалявания на ключа струват $O(m)$, а $n$ извличания на минимум струват $O(n\log n)$.
\end{proof}

\begin{remark}
При двоична пирамида често се записва и оценката $O(m\log m)$, когато логаритъмът се изразява чрез броя на ребрата. За прост неориентиран граф $m< n^2$, откъдето $\log m=O(\log n)$; така двете оценки са съвместими в рамките на обичайните предположения за графа.
\end{remark}

\section{Алгоритъм на Крускал}\label{s:7.4}

Алгоритъмът на Крускал строи не едно растящо дърво, а растяща гора. Ребрата се разглеждат в ненамаляващ ред по тегло. Ребро се добавя точно когато краищата му принадлежат на различни дървета в текущата гора.

\subsection{Използвани операции}\label{s:7.4.1}

Псевдокодът извиква три операции върху разбиване на множеството от върхове на непресичащи се дялове. Всеки дял е една компонента на свързаност на текущата гора $(V,A)$.

\begin{itemize}
\item \verb|Make-Set(v)| --- създава нов дял, съдържащ само $v$;
\item \verb|Find-Set(v)| --- връща представителя (лидера) на дяла, съдържащ $v$; два върха са в един и същ дял тогава и само тогава, когато \verb|Find-Set| връща един и същ представител за тях;
\item \verb|Union(u,v)| --- слива дяловете на $u$ и $v$ в един.
\end{itemize}

Заедно те образуват структурата Union--Find, разгледана подробно след доказателството за коректност. Засега е достатъчна само тази спецификация: нейната реализация не влияе на коректността на Крускал, а само на сложността.

\subsection{Псевдокод}\label{s:7.4.2}

\begin{verbatim}
MST_Kruskal(G = (V,E), omega)
A <- empty set                   // избраните ребра
for each v in V do
    Make-Set(v)                  // всеки връх е отделен дял
sort E in nondecreasing order by omega(e)
for each edge {u,v} in E do      // в този сортиран ред
    if Find-Set(u) != Find-Set(v) then    // различни компоненти
        A <- A union {{u,v}}     // не затваря цикъл: приема се
        Union(u,v)               // компонентите се сливат
return A
\end{verbatim}

Проверката \verb|Find-Set(u) != Find-Set(v)| е точно проверката дали реброто $\{u,v\}$ би затворило цикъл: цикъл се получава тогава и само тогава, когато двата края вече са свързани чрез приетите ребра. Затова липсва отделна проверка за цикъл --- дяловете я вършат.

При свързан граф цикълът приема точно $|V|-1$ ребра, така че изпълнението може да спре веднага щом $|A|=|V|-1$; останалите ребра неизбежно биха затворили цикъл.

\subsection{Пример}\label{s:7.4.3}

\begin{example}
За графа от фигура~\ref{fig:mst-example} сортираният ред на ребрата е
\[
\{b,c\}{:}1,\ \{b,d\}{:}2,\ \{e,f\}{:}2,\ \{a,c\}{:}3,\ \{a,b\}{:}4,\ \{c,d\}{:}4,\ \{c,e\}{:}5,\ \{d,f\}{:}6,\ \{d,e\}{:}7.
\]

\begin{center}
\begin{tabular}{llll}
\toprule
Ребро & Тегло & Действие & Компоненти след стъпката \\
\midrule
$\{b,c\}$ & $1$ & приема се & $\{a\}$, $\{b,c\}$, $\{d\}$, $\{e\}$, $\{f\}$ \\
$\{b,d\}$ & $2$ & приема се & $\{a\}$, $\{b,c,d\}$, $\{e\}$, $\{f\}$ \\
$\{e,f\}$ & $2$ & приема се & $\{a\}$, $\{b,c,d\}$, $\{e,f\}$ \\
$\{a,c\}$ & $3$ & приема се & $\{a,b,c,d\}$, $\{e,f\}$ \\
$\{a,b\}$ & $4$ & отхвърля се & без промяна \\
$\{c,d\}$ & $4$ & отхвърля се & без промяна \\
$\{c,e\}$ & $5$ & приема се & $\{a,b,c,d,e,f\}$ \\
\bottomrule
\end{tabular}
\end{center}

Приети са $5=|V|-1$ ребра, така че ребрата $\{d,f\}$ и $\{d,e\}$ вече не се разглеждат. Полученото дърво съвпада с намереното от Прим и има тегло $13$, макар че двата алгоритъма добавят ребрата в различен ред.
\end{example}

\subsection{Коректност}\label{s:7.4.4}

\begin{keythm}
Алгоритъмът на Крускал връща минимално покриващо дърво.
\end{keythm}

\begin{proof}
В началото $A=\varnothing$, следователно $A$ е съгласувано. Нека разгледаме стъпка, в която алгоритъмът приема реброто
\[
e=\{u,v\}.
\]
Условието
\[
\operatorname{Find\mbox{-}Set}(u)\neq\operatorname{Find\mbox{-}Set}(v)
\]
означава, че $u$ и $v$ са в различни компоненти на гората $(V,A)$. Нека $C$ е компонентата, съдържаща $u$, и разгледаме среза
\[
S=\{C,V\setminus C\}.
\]
Този срез е съобразен с $A$, тъй като по определение няма ребро от $A$, което да свързва различни компоненти на гората.

Реброто $e$ пресича среза. То е леко за този срез: ако имаше пресичащо ребро с по-малко тегло, то би се намирало по-рано в сортирания ред. Когато то е било разгледано, краищата му са били в различни компоненти, защото компонентите могат само да се сливат, а не да се разделят. Следователно това по-леко ребро би било добавено в $A$, което противоречи на съобразеността на среза с $A$.

От теоремата за съгласуваното множество следва, че $e$ е сигурно за $A$. Следователно след всяко прието ребро множеството $A$ остава съгласувано.

Алгоритъмът никога не образува цикъл, понеже приема ребро само между различни компоненти. Тъй като графът е свързан, след разглеждане на всички ребра всички върхове принадлежат на една компонента. Тогава $A$ е свързан и ацикличен граф, тоест покриващо дърво. Понеже $A$ е съгласувано и само е покриващо дърво, то е минимално покриващо дърво.
\end{proof}

\section{Union--Find структури от данни}\label{s:7.5}

За ефикасната реализация на Крускал е необходима структура от данни, която поддържа разбиване на множество на непразни неприпокриващи се дялове.

\begin{defn}
Union--Find структура върху опорно множество $S$ поддържа следните операции:

\begin{itemize}
\item \verb|Make-Set(x)| създава едноелементния дял $\{x\}$;
\item \verb|Find-Set(x)| връща представител, наричан лидер, на дяла, съдържащ $x$;
\item \verb|Union(x,y)| слива двата дяла, съдържащи $x$ и $y$.
\end{itemize}
\end{defn}

В алгоритъма на Крускал всяка компонента на текущата гора е един дял. Така проверката дали реброто $\{u,v\}$ би образувало цикъл е проверката
\[
\operatorname{Find\mbox{-}Set}(u)
\neq
\operatorname{Find\mbox{-}Set}(v).
\]
Ако е вярна, реброто се приема и се изпълнява \verb|Union(u,v)|.

\subsection{Реализация с гора от коренови дървета}\label{s:7.5.1}

Разбиването може да се представи чрез ориентирана гора. Върховете на гората са елементите на $S$, а всяко дърво съответства на един дял. Коренът на дървото е лидерът на съответния дял. Нека $\pi[x]$ е родителят на $x$; за корен е изпълнено $\pi[x]=x$. Тази гора няма нищо общо с гората, която Крускал строи --- тя само записва кои върхове са в една компонента, а не кои ребра са приети.

Полето $\mathit{rank}[x]$ по-долу е горна граница за височината на дървото с корен $x$ и се използва само от евристиката \emph{union by rank}; при $\verb|Make-Set|$ то е $0$.

\begin{verbatim}
Make-Set(x)
    pi[x] <- x
    rank[x] <- 0

Find-Set(x)
    if x != pi[x] then
        return Find-Set(pi[x])
    else
        return x
\end{verbatim}

Без допълнителни правила операцията \verb|Union| може да образува дърво с линейна височина. Тогава \verb|Find-Set| би изисквала линейно време в най-лошия случай.

\subsection{Union by rank}\label{s:7.5.2}

\begin{defn}[Обединение по ранг]
При евристиката \emph{union by rank} коренът на дървото с по-малък ранг се закача към корена на дървото с по-голям ранг. При равни рангове единият корен се избира за нов корен и неговият ранг се увеличава с единица.
\end{defn}

\begin{verbatim}
Union(x,y)
    rx <- Find-Set(x)              // корените, а не самите x и y
    ry <- Find-Set(y)
    if rx = ry then
        return
    if rank[rx] < rank[ry] then
        pi[rx] <- ry
    else if rank[rx] > rank[ry] then
        pi[ry] <- rx
    else
        pi[ry] <- rx
        rank[rx] <- rank[rx] + 1
\end{verbatim}

\begin{lem}
Започвайки от едноелементни множества, след произволна последователност от операции \verb|Union| с union by rank височината на всяко дърво е $O(\log n)$, където $n$ е броят върхове в дървото.
\end{lem}

\begin{proof}
Поддържаме едновременно инвариантите, че височината на дървото е най-много ранга на корена и че дърво с ранг $h$ има поне $2^h$ върха. Те са верни за едноелементно дърво: височината и рангът са $0$, а броят върхове е $1$.

Ако корен с по-малък ранг се закачи към корен с по-голям ранг, новата височина е най-много по-големия ранг: височината на закаченото дърво плюс едно е най-много този ранг. Рангът на новия корен не се променя. При равни рангове $h$ новата височина е най-много $h+1$, а рангът на новия корен също става $h+1$.

За втория инвариант е достатъчно да разгледаме момента, в който рангът се увеличава от $h$ на $h+1$. Тогава се сливат две дървета с ранг $h$, всяко с поне $2^h$ върха. Новото дърво има поне
\[
2^h+2^h=2^{h+1}
\]
върха.

Следователно, ако дърво с $n$ върха има ранг $h$, то
\[
2^h\leq n,
\]
тоест
\[
h\leq \log_2 n.
\]
Понеже височината е най-много ранга, тя също е най-много $\log_2 n$. Затова \verb|Find-Set| работи за $O(\log n)$ време.
\end{proof}

\subsection{Компресия на пътя}\label{s:7.5.3}

\begin{defn}[Компресия на пътя]
Втората евристика е \emph{компресия на пътя}. При изпълнение на \verb|Find-Set(x)| всички върхове по намерения път към корена се насочват непосредствено към корена.
\end{defn}

\begin{verbatim}
Find-Set(x)
    if x != pi[x] then
        pi[x] <- Find-Set(pi[x])
    return pi[x]
\end{verbatim}

При едновременно използване на union by rank и компресия на пътя общата цена на последователност от \verb|Make-Set|, \verb|Find-Set| и \verb|Union| операции е
\[
O\bigl((n+m)\alpha(n)\bigr),
\]
където $\alpha$ е обратната функция на Акерман. Тя расте изключително бавно и за практически размери има много малка стойност.

\section{Сложност на алгоритъма на Крускал}\label{s:7.6}

Нека $n=|V|$ и $m=|E|$.

Инициализирането с $n$ операции \verb|Make-Set| отнема $O(n)$ време. Сортирането на ребрата е с време
\[
O(m\log m).
\]
При реализация с union by rank и компресия на пътя операциите на Union--Find за целия алгоритъм струват
\[
O\bigl((n+m)\alpha(n)\bigr).
\]
Следователно общата времева сложност е
\[
O\bigl(m\log m+(n+m)\alpha(n)\bigr).
\]

За прост неориентиран граф е изпълнено $m<n^2$, откъдето
\[
\log m=O(\log n).
\]
Сортирането доминира останалата работа, така че сложността на алгоритъма на Крускал обичайно се записва като
\[
O(m\log m)=O(m\log n).
\]

\begin{supp}[Бележка]
Ако се използва само union by rank, без компресия на пътя, операциите \verb|Find-Set| са $O(\log n)$, а алгоритъмът на Крускал отново има оценка $O(m\log n)$ поради сортирането. Компресията на пътя подобрява цената на операциите върху разбиването, макар че не променя доминиращата асимптотична оценка на Крускал при сравняващо сортиране.
\end{supp}

\section{Изпитен фокус}\label{s:7.7}

\begin{itemize}
\item Да се дефинират тегловен граф, покриващо дърво, тегло на дърво и минимално покриващо дърво.
\item Да се формулира задачата: неориентиран свързан граф с тегла на ребрата; търси се покриващо дърво с минимална сума от теглата.
\item Да се дефинират съгласувано множество, сигурно ребро, срез, съобразен срез и леко ребро.
\item Да се формулира и докаже теоремата за съгласуваното множество чрез добавяне на ребро, получаване на цикъл и размяна с друго ребро от същия срез.
\item Да се възпроизведе псевдокодът на Прим, значението на $key$, $\pi$ и приоритетната опашка, както и операциите \verb|Extract-Min| и \verb|Decrease-Key|, които той използва.
\item Да се обясни коректността на Прим чрез среза между вече избраните и неизбраните върхове и теоремата за съгласуваното множество.
\item Да се изведат сложностите на Прим: $O(m\log n)$ с двоична пирамида и $O(m+n\log n)$ с пирамида на Фибоначи.
\item Да се възпроизведе псевдокодът на Крускал и ролята на сортирането на ребрата.
\item Да се обясни коректността на Крускал: прието ребро свързва различни компоненти и е леко за подходящ съобразен срез.
\item Да се дефинират \verb|Make-Set|, \verb|Find-Set| и \verb|Union|, както и реализацията им чрез гора от коренови дървета.
\item Да се обяснят union by rank и компресията на пътя и приложението им в Крускал.
\item Да се изведе сложността на Крускал:
\[
O\bigl(m\log m+(n+m)\alpha(n)\bigr)=O(m\log n).
\]
\end{itemize}


\chapter{Най-къси пътища в тегловни графи.}
%% Drafted from the mapped corpus by scripts/draft_topics.py.
% Model: gpt-5.6-terra; source characters: 24927.
\begin{konspekt}
\kselect{На изпита ще се падне \textbf{или} алгоритъмът на Дийкстра \emph{и} алгоритъмът за тегловен ацикличен граф, \textbf{или} алгоритъмът на Белман--Форд \emph{и} алгоритъмът на Флойд--Уоршал. Двойките са неразделни.}
\kreq{Определение и варианти на задачата}{плюс потенциалните проблеми при отрицателни тегла --- \kref{8.1}.}
\kreq{Вариант, псевдокод, коректност, сложност при различни структури}{алгоритъм на Дийкстра --- \kref{8.3}.}
\kreq{Вариант, псевдокод, коректност, сложност, предимства пред Дийкстра}{най-къси пътища в тегловен ориентиран ацикличен граф --- \kref{8.4}.}
\kreq{Вариант, псевдокод, коректност, сложност, откриване на отрицателен цикъл}{алгоритъм на Белман--Форд --- \kref{8.5}.}
\kreq{Вариант, псевдокод, коректност, сложност по време и по памет}{алгоритъм на Флойд--Уоршал --- \kref{8.6}.}
\kreq{Литература}{[31].}
\end{konspekt}

\section{Постановка на задачата}\label{s:8.1}

\begin{defn}[Тегловен граф и тегло на маршрут]
Нека $G=(V,E)$ е тегловен граф с теглова функция
\[
\omega:E\to\mathbb{R}.
\]
Разглеждаме преди всичко ориентирани графи. Теглото на маршрут $p=(v_0,e_1,v_1,\ldots,e_k,v_k)$ е сумата от теглата на ребрата му, броени с повторенията:
\[
\omega(p)=\sum_{i=1}^{k}\omega(e_i).
\]
\end{defn}

\begin{defn}
Нека $u,v\in V$. С $\operatorname{Paths}(u,v)$ означаваме множеството от пътища от $u$ до $v$. Най-късото разстояние от $u$ до $v$ се означава с $\delta(u,v)$ и е
\[
\delta(u,v)=
\begin{cases}
+\infty, & \text{ако }\operatorname{Paths}(u,v)=\emptyset,\\
-\infty, & \text{ако е изпълнено }C(u,v),\\
\min\{\omega(p)\mid p\in\operatorname{Paths}(u,v)\}, & \text{иначе}.
\end{cases}
\]
Тук $C(u,v)$ означава, че има маршрут от $u$ до $v$, минаващ през отрицателен цикъл. Допускат се повторения на върхове и ребра; при липса на такъв цикъл минимумът се достига от прост път. Под \emph{най-къс път} ще разбираме път с най-малко тегло.
\end{defn}

\begin{defn}
В ориентиран тегловен граф \emph{отрицателен цикъл} е цикъл, чиято сума от теглата на ребрата е отрицателна.
\end{defn}

\begin{supp}[Бележка]
При неориентиран граф всяко ребро с отрицателно тегло позволява обхождане в двете посоки с отрицателно общо тегло. Следователно наличието на отрицателно ребро е съществено препятствие за стандартното понятие за най-къс път, дори ако се разглеждат само прости цикли.
\end{supp}

\subsection{Варианти на задачата}\label{s:8.1.1}

Основните варианти на задачата за най-къси пътища са следните.

\begin{itemize}
\item \emph{Най-къс път между дадена двойка върхове} (\emph{single-pair shortest paths}). Дадени са $G$, $\omega$, начален връх $s$ и краен връх $t$. Търси се $\delta(s,t)$.

\item \emph{Най-къси пътища от един източник} (\emph{single-source shortest paths}, SSSP). Дадени са $G$, $\omega$ и начален връх $s$. Търсят се всички стойности
\[
\delta(s,v),\qquad v\in V.
\]

\item \emph{Най-къси пътища до един краен връх} (\emph{single-destination shortest paths}). Дадени са $G$, $\omega$ и краен връх $t$. Търсят се всички стойности
\[
\delta(v,t),\qquad v\in V.
\]

\item \emph{Най-къси пътища между всички двойки върхове} (\emph{all-pairs shortest paths}, APSP). Дадени са $G$ и $\omega$. Търсят се всички стойности
\[
\delta(u,v),\qquad u,v\in V.
\]
\end{itemize}

Задачата с един краен връх може да се свежда до задача с един източник чрез обръщане на всички ребра. Задачата за една двойка е частен случай на задачата с един източник, но специализиран алгоритъм може да прекрати работата си по-рано, когато крайният връх е достатъчно обработен.

\subsection{Отрицателни тегла и оптимална подструктура}\label{s:8.1.2}

Отрицателните тегла на ребрата сами по себе си не правят задачата безсмислена. Проблемът възниква при отрицателен цикъл, който е достижим по път от началния до крайния връх. Тогава цикълът може да се обходи произволно много пъти и да се получават пътища с все по-малко тегло. Поради това няма път с минимално тегло и стойността е $-\infty$.

\begin{remark}
Ако от началния връх не е достижим отрицателен цикъл, за всеки достижим краен връх може да се избере прост най-къс път. Наистина, ако най-къс маршрут съдържа цикъл с неотрицателно тегло, премахването му не увеличава теглото; отрицателен цикъл по маршрута би бил достижим от началния връх. При цикъл с нулево тегло може да има и непрости най-къси маршрути, но винаги съществува прост представител със същото тегло.
\end{remark}

\begin{thm}[Оптимална подструктура на най-късите пътища]
Нека $p$ е най-къс път от $s$ до $t$ и нека $u$ предхожда $v$ по $p$. Тогава подпътят на $p$ от $u$ до $v$ е най-къс път от $u$ до $v$.
\end{thm}

\begin{proof}
Нека подпътят от $u$ до $v$ в $p$ е $q$. Ако съществува път $q'$ от $u$ до $v$ с
\[
\omega(q')<\omega(q),
\]
то можем да заменим $q$ с $q'$ в пътя $p$. Полученият път от $s$ до $t$ има тегло
\[
\omega(p)-\omega(q)+\omega(q')<\omega(p),
\]
което противоречи на избора на $p$.
\end{proof}

\section{Релаксация на ребра}\label{s:8.2}

Алгоритмите за най-къси пътища поддържат за всеки връх $v$ две полета:

\begin{itemize}
\item оценка $v.d$ за разстоянието от началния връх $s$ до $v$; тя е тегло на някакъв вече намерен маршрут от $s$ до $v$, тоест горна граница за $\delta(s,v)$;
\item предшественик $v.\pi$ --- върхът, от който този маршрут влиза в $v$. Пътят се възстановява, като се тръгне от $v$ и се следват указателите $\pi$ назад до $s$.
\end{itemize}

Използват се още две означения: $w((x,y))$ е теглото на реброто $(x,y)$, а $\mathit{adj}(x)$ е списъкът на съседство на $x$, тоест списъкът на върховете $y$, за които $(x,y)\in E$.

Инициализацията е следната; името \verb|ISS| идва от \emph{initialize single source}.

\begin{verbatim}
ISS(G, s)
    for each v in V(G)
        v.d  <- infinity       // още не е намерен маршрут до v
        v.pi <- Nil            // предшественик по този маршрут
    s.d <- 0                   // до източника води празният маршрут
\end{verbatim}

Основната операция е \emph{релаксация}.

\begin{verbatim}
Relax((x, y))
    if y.d > x.d + w((x, y))   // през x се стига по-евтино до y
        y.d  <- x.d + w((x, y))  // оценката се подобрява
        y.pi <- x                // и се запомня как е постигната
\end{verbatim}

Тя проверява дали вече известният път до $x$, последван от реброто $(x,y)$, подобрява текущата оценка за $y$. Релаксацията никога не увеличава оценка и никога не поставя стойност, която не е тегло на реален маршрут; сравнението със сума, съдържаща $x.d=+\infty$, е винаги лъжа, така че недостижимите досега върхове не се променят.

Всички алгоритми по-долу се различават само по това \emph{в какъв ред} релаксират ребрата; самата операция е една и съща.

\begin{lem}[Неравенство за ребро]
За всяко ребро $(x,y)\in E$ е изпълнено
\[
\delta(s,y)\leq \delta(s,x)+\omega((x,y)).
\]
\end{lem}

\begin{proof}
Ако $x$ не е достижим от $s$, дясната страна е $+\infty$ и неравенството е вярно. Ако до $x$ може да се достигне чрез път, съдържащ отрицателен цикъл, тогава и до $y$ може да се достигне по такъв път и и двете релевантни стойности са $-\infty$.

В останалия случай съществува най-къс път от $s$ до $x$. Добавянето на реброто $(x,y)$ дава път от $s$ до $y$ с тегло
\[
\delta(s,x)+\omega((x,y)).
\]
По определение $\delta(s,y)$ не може да е по-голямо от теглото на този път.
\end{proof}

\begin{lem}[Свойство на оценките]
След инициализация и след всяка произволна последователност от релаксации за всеки връх $v$ е изпълнено
\[
v.d\geq \delta(s,v).
\]
Освен това, щом за даден връх се достигне равенството
\[
v.d=\delta(s,v),
\]
то то се запазва при всички следващи релаксации.
\end{lem}

\begin{proof}
Твърдението се доказва с индукция по броя извършени релаксации. След инициализацията $s.d=0$, а за останалите върхове оценката е $+\infty$, следователно твърдението е вярно.

При релаксация на $(x,y)$ може да се промени само $y.d$. Ако се промени, новата му стойност е
\[
x.d+\omega((x,y))
 \geq \delta(s,x)+\omega((x,y))
 \geq \delta(s,y),
\]
където последното неравенство е от предишната лема. Релаксацията никога не увеличава оценки. Следователно достигнато точно разстояние не може по-късно да бъде развалено.
\end{proof}

\begin{lem}[Релаксация по най-къс път]
Нека $p$ е най-къс път от $s$ до $v$ и нека $(u,v)$ е последното му ребро. Ако преди релаксацията на $(u,v)$ е изпълнено
\[
u.d=\delta(s,u),
\]
то след тази релаксация е изпълнено
\[
v.d=\delta(s,v).
\]
\end{lem}

\begin{proof}
По оптималната подструктура префиксът на $p$ от $s$ до $u$ е най-къс път, следователно
\[
\delta(s,v)=\delta(s,u)+\omega((u,v)).
\]
След релаксацията
\[
v.d\leq u.d+\omega((u,v))=\delta(s,v).
\]
От свойството на оценките имаме и $v.d\geq\delta(s,v)$, откъдето следва равенство.
\end{proof}

\section{Алгоритъм на Дийкстра}\label{s:8.3}

\subsection{Вариант на задачата и идея}\label{s:8.3.1}

Алгоритъмът на Дийкстра решава задачата SSSP в ориентиран тегловен граф, когато всички тегла са неотрицателни:
\[
\omega(e)\geq 0,\qquad e\in E.
\]
Той поддържа множество $D$ от окончателно обработени върхове. На всяка стъпка се избира необработен връх с най-малка текуща оценка и той се добавя в $D$. След това се релаксират всички изходящи от него ребра.

\subsection{Използвани операции}\label{s:8.3.2}

Псевдокодът извиква операции на приоритетна опашка. Те не са част от самия алгоритъм: всяка реализация, която ги поддържа, е допустима, а изборът на реализация влияе само на сложността.

\begin{defn}[Приоритетна опашка по минимум]
Приоритетна опашка е структура от данни, която пази множество от елементи, всеки с числов ключ, и поддържа операциите:

\begin{itemize}
\item \verb|Insert(Q,x)| --- добавя елемента $x$ с текущия му ключ;
\item \verb|Extract-Min(Q)| --- премахва от $Q$ и връща елемент с минимален ключ;
\item \verb|Decrease-Key(Q,x,k)| --- заменя ключа на $x$ с новата стойност $k$, която не е по-голяма от старата, и възстановява вътрешния ред на структурата.
\end{itemize}
\end{defn}

В алгоритъма на Дийкстра ключът на връх $v$ е текущата оценка $v.d$. Записът $Q\leftarrow V(G)$ означава построяване на опашка от всички върхове с текущите им ключове, тоест $|V|$ операции \verb|Insert| или еднократно построяване.

\subsection{Псевдокод}\label{s:8.3.3}

\begin{verbatim}
Dijkstra(G, w, s)
    ISS(G, s)                  // всички оценки са infinity, s.d = 0
    D <- emptyset              // окончателно обработените върхове
    create priority queue Q    // ключ на връх v е оценката v.d
    Q <- V(G)                  // всички върхове чакат в опашката
    while Q is not empty
        x <- Extract-Min(Q)    // необработен връх с най-малък d
        D <- D union {x}       // от този момент x.d е окончателно
        for each y in adj(x)   // релаксация на изходящите ребра
            if y in Q and y.d > x.d + w((x, y))
                y.d  <- x.d + w((x, y))
                y.pi <- x
                Decrease-Key(Q, y, y.d)   // синхронизира опашката
\end{verbatim}

Означенията в псевдокода са следните.

\begin{itemize}
\item $G$, $w$ и $s$ са входът: ориентиран граф, теглова функция с неотрицателни стойности и начален връх.
\item $Q$ съдържа необработените върхове, а $D=V\setminus Q$ --- обработените; двете множества винаги се допълват, така че $D$ се пази само за удобство на изложението и може да се пропусне.
\item Тялото на \verb|while| се изпълнява по веднъж за всеки връх: \verb|Extract-Min| премахва връх и нищо не се връща обратно в $Q$.
\item Трите реда в тялото на \verb|if| са точно \verb|Relax((x,y))|, следвана от \verb|Decrease-Key|; последната само съобщава на опашката, че ключът на $y$ е намалял.
\end{itemize}

Проверката $y\in Q$ не позволява промяна на вече окончателно обработен връх; тя се реализира за константно време с отделна булева информация дали върхът е в опашката. Резултатът от алгоритъма са оценките $v.d$ и указателите $v.\pi$, които образуват дърво на най-късите пътища с корен $s$.

\subsection{Пример}\label{s:8.3.4}

\begin{example}
Разглеждаме графа от фигура~\ref{fig:dijkstra-example} с начален връх $s$.

\begin{lecfigure}[Ориентиран граф с неотрицателни тегла. Удебелените ребра са ребрата $\{(v.\pi,v)\}$ на полученото дърво на най-късите пътища.\label{fig:dijkstra-example}]
\begin{tikzpicture}[scale=1.25,
  vtx/.style={circle, draw=defframe, fill=defback, inner sep=0pt,
              minimum size=6.5mm, font=\small},
  earc/.style={gray!65, line width=0.5pt, -{Stealth[length=2mm]}},
  etree/.style={thmframe, line width=1.5pt, -{Stealth[length=2.4mm]}},
  wlab/.style={font=\footnotesize, inner sep=1.4pt, fill=white}]
  \node[vtx] (s) at (0,1)     {$s$};
  \node[vtx] (a) at (1.9,2.1) {$a$};
  \node[vtx] (b) at (1.9,-0.1){$b$};
  \node[vtx] (c) at (3.9,1.6) {$c$};
  \node[vtx] (d) at (5.8,0.7) {$d$};
  \draw[earc]  (s) -- node[wlab] {4}  (a);
  \draw[etree] (s) -- node[wlab] {2}  (b);
  \draw[etree] (b) -- node[wlab] {1}  (a);
  \draw[etree] (a) -- node[wlab] {5}  (c);
  \draw[earc]  (b) -- node[wlab] {8}  (c);
  \draw[earc]  (b) -- node[wlab] {10} (d);
  \draw[etree] (c) -- node[wlab] {2}  (d);
\end{tikzpicture}
\end{lecfigure}

След \verb|ISS| е $s.d=0$, а всички останали оценки са $\infty$. Всеки ред от таблицата отговаря на една итерация на цикъла \verb|while|: извлеченият връх, окончателната му оценка и състоянието на останалите оценки след релаксациите. Записът $3(b)$ означава $d=3$ и $\pi=b$.

\begin{center}
\begin{tabular}{llll}
\toprule
Извлечен $x$ & $x.d$ & Релаксирани ребра & Оценки в $Q$ след стъпката \\
\midrule
$s$ & $0$  & $(s,a)$, $(s,b)$          & $a{:}4(s)$, $b{:}2(s)$, $c{:}\infty$, $d{:}\infty$ \\
$b$ & $2$  & $(b,a)$, $(b,c)$, $(b,d)$ & $a{:}3(b)$, $c{:}10(b)$, $d{:}12(b)$ \\
$a$ & $3$  & $(a,c)$                   & $c{:}8(a)$, $d{:}12(b)$ \\
$c$ & $8$  & $(c,d)$                   & $d{:}10(c)$ \\
$d$ & $10$ & няма изходящи ребра       & $Q=\varnothing$ \\
\bottomrule
\end{tabular}
\end{center}

Върховете се извличат в реда $s,b,a,c,d$, тоест по нарастващи окончателни оценки $0,2,3,8,10$, а не по номер или по брой ребра. Три пъти се задейства \verb|Decrease-Key|: оценката на $a$ пада от $4(s)$ на $3(b)$, на $c$ --- от $10(b)$ на $8(a)$, и на $d$ --- от $12(b)$ на $10(c)$. Обърнете внимание, че $a$ е бил с оценка $4$, докато още е чакал в $Q$; именно затова окончателна става само оценката на извлечения връх.

Полученият резултат е
\[
\delta(s,s)=0,\quad
\delta(s,a)=3,\quad
\delta(s,b)=2,\quad
\delta(s,c)=8,\quad
\delta(s,d)=10,
\]
а указателите $\pi$ дават най-късия път до $d$:
\[
s\to b\to a\to c\to d.
\]
\end{example}

\begin{supp}[Бележка]
Алгоритъмът на Дийкстра не е коректен при отрицателни тегла. Причината е, че връх може да бъде окончателно избран, а по-късно да се открие по-лек път до него през ребро с отрицателно тегло.
\end{supp}

\begin{example}[Дийкстра греши при отрицателно тегло]
Нека графът има върхове $s,a,b$ и ребра
\[
(s,a)\text{ с тегло }2,\qquad
(s,b)\text{ с тегло }3,\qquad
(b,a)\text{ с тегло }-2.
\]
Отрицателен цикъл няма --- графът е ацикличен. Истинското разстояние е
\[
\delta(s,a)=3+(-2)=1.
\]
Алгоритъмът обаче извлича първо $s$ и получава $a{:}2$, $b{:}3$. След това минималната оценка в $Q$ е тази на $a$, затова $a$ се извлича и попада в $D$ с $a.d=2$. Когато по-късно се обработи $b$ и се релаксира $(b,a)$, проверката $y\in Q$ е лъжа, защото $a$ вече не е в опашката, и оценката остава $2\ne 1$.

Точно това не може да се случи при неотрицателни тегла: тогава удължаването на маршрут никога не го прави по-лек, което е и същината на доказателството по-долу.
\end{example}

\subsection{Коректност}\label{s:8.3.5}

\begin{keythm}[Коректност на Дийкстра]
Нека $G$ е ориентиран тегловен граф без отрицателни тегла и $s$ е негов връх. След завършване на алгоритъма на Дийкстра за всеки връх $v$ е изпълнено
\[
v.d=\delta(s,v).
\]
\end{keythm}

\begin{proof}
Ще докажем, че когато връх $v$ бъде добавен в $D$, вече е изпълнено
\[
v.d=\delta(s,v).
\]

Да допуснем противното и нека $v$ е първият връх, добавен в $D$ с неточна оценка. От свойството на оценките може да има само случая
\[
v.d>\delta(s,v).
\]
Връх $v$ не е $s$, защото $s.d=0=\delta(s,s)$.

Нека $p$ е най-къс път от $s$ до $v$. В момента непосредствено преди добавянето на $v$ в $D$ разглеждаме първия връх $u$ по $p$, който не принадлежи на $D$. Нека $a$ е предшественикът на $u$ по $p$. Тогава $a\in D$ и при обработването на $a$ реброто $(a,u)$ е било релаксирано.

Тъй като $a$ е добавен преди $v$, по избора на $v$ имаме
\[
a.d=\delta(s,a).
\]
От лемата за релаксация по най-къс път следва
\[
u.d=\delta(s,u).
\]

Върховете $u$ и $v$ са извън $D$, а $v$ е избран от приоритетната опашка с минимална оценка. Следователно
\[
v.d\leq u.d.
\]
Понеже теглата са неотрицателни и $u$ предхожда $v$ по най-късия път $p$, имаме
\[
\delta(s,u)\leq\delta(s,v).
\]
А от предположението $v.d>\delta(s,v)$ получаваме
\[
v.d\leq u.d=\delta(s,u)\leq\delta(s,v)<v.d,
\]
което е противоречие. Следователно всеки връх се добавя в $D$ с точна оценка. При завършване всички върхове са обработени.
\end{proof}

\subsection{Сложност}\label{s:8.3.6}

Нека $n=|V|$ и $m=|E|$. Времето зависи от реализацията на приоритетната опашка.

\begin{itemize}
\item При неупорядочен масив намирането и извличането на минимум струва $\Theta(n)$ на итерация. Общата сложност е
\[
\Theta(n^2+m),
\]
обикновено записвана като $\Theta(n^2)$.

\item При двоична пирамида извличането на минимум и намаляването на ключ струват $\Theta(\log n)$. Има най-много $n$ извличания и до $m$ релаксации, които могат да намалят ключ. Получава се
\[
\Theta((n+m)\log n).
\]

\item При по-сложни реализации на приоритетни опашки анализът може да бъде по-добър за операцията по намаляване на ключа, но основната идея на алгоритъма не се изменя.
\end{itemize}

\section{Най-къси пътища в тегловен DAG}\label{s:8.4}

\subsection{Алгоритъм}\label{s:8.4.1}

Ориентираният ацикличен граф, или DAG, допуска топологична наредба на върховете: за всяко ребро $(u,v)$ връх $u$ се намира преди $v$ в наредбата. Това позволява върховете да се обработят така, че преди даден връх да са обработени всички възможни предшественици по път от източника.

Алгоритъмът решава SSSP в тегловен DAG и допуска отрицателни тегла.

\begin{verbatim}
DAG-SSSP(G, w, s)
    Toposort(G)
    ISS(G, s)
    for each x in topological order
        for each y in adj(x)
            Relax((x, y))
\end{verbatim}

\subsection{Коректност}\label{s:8.4.2}

\begin{keythm}[Коректност на алгоритъма за DAG]
Нека $G$ е тегловен ориентиран ацикличен граф и $s\in V$. След завършването на \texttt{DAG-SSSP} за всеки връх $v$ е изпълнено
\[
v.d=\delta(s,v).
\]
\end{keythm}

\begin{proof}
Нека върховете са разгледани в топологична наредба. За всеки връх, стоящ преди $s$, няма път от $s$ до него, защото всеки път следва топологичната наредба. Следователно за него
\[
v.d=+\infty=\delta(s,v).
\]
За $s$ след инициализацията е изпълнено
\[
s.d=0=\delta(s,s),
\]
тъй като в DAG няма цикли, а следователно няма и отрицателни цикли.

Нека $v$ е произволен връх след $s$ в топологичната наредба. Всеки вътрешен връх на път от $s$ до $v$ стои преди $v$. Нека
\[
U=\{u\in V\mid (u,v)\in E \text{ и има път от }s\text{ до }u\}.
\]
До момента, в който започне обработването на $v$, всички върхове от $U$ вече са обработени. По индуктивното предположение за тях е вярно
\[
u.d=\delta(s,u).
\]
При релаксиране на всички входящи към $v$ възможности, реализирани чрез обработката на техните начала, оценката става
\[
v.d=\min_{u\in U}\{\delta(s,u)+\omega((u,v))\}.
\]
По оптималната подструктура това е точно $\delta(s,v)$. Ако $U=\emptyset$, връх $v$ не е достижим от $s$; тогава и двете стойности са $+\infty$.
\end{proof}

\subsection{Сложност и предимства}\label{s:8.4.3}

Топологичното сортиране има сложност $\Theta(n+m)$. Вложеният цикъл обхожда всеки връх и всяко ребро най-много веднъж, затова общата сложност е
\[
\Theta(n+m).
\]

Алгоритъмът има две важни предимства пред алгоритъма на Дийкстра върху DAG:

\begin{itemize}
\item работи с отрицателни тегла, защото в DAG няма цикли и следователно няма отрицателни цикли;
\item има линейна сложност $\Theta(n+m)$, която е асимптотично по-добра от сложността на Дийкстра с приоритетна опашка.
\end{itemize}

\begin{remark}
Чрез обръщане на неравенството в операцията за релаксация алгоритъмът за DAG може аналогично да намира най-дълги пътища. Това е възможно именно поради липсата на цикли.
\end{remark}

\section{Алгоритъм на Белман--Форд}\label{s:8.5}

\subsection{Вариант на задачата и псевдокод}\label{s:8.5.1}

Алгоритъмът на Белман--Форд решава SSSP в ориентиран тегловен граф, който може да съдържа отрицателни ребра. Освен това открива наличие на отрицателен цикъл, достижим от началния връх.

\begin{verbatim}
Bellman-Ford(G, w, s)
    ISS(G, s)
    for i <- 1 to n - 1
        for each (x, y) in E(G)
            Relax((x, y))
    for each (x, y) in E(G)
        if y.d > x.d + w((x, y))
            return False
    return True
\end{verbatim}

Ако алгоритъмът върне \texttt{True}, отрицателен цикъл, достижим от $s$, не е открит и оценките представляват най-късите разстояния. Ако върне \texttt{False}, съществува достижим отрицателен цикъл.

\subsection{Коректност при липса на достижими отрицателни цикли}\label{s:8.5.2}

\begin{lem}
Нека в графа няма отрицателен цикъл, достижим от $s$. След изпълнение на първите $n-1$ обхождания на всички ребра от алгоритъма на Белман--Форд за всеки връх $v$ е изпълнено
\[
v.d=\delta(s,v).
\]
\end{lem}

\begin{proof}
При липса на отрицателни цикли, достижими от $s$, могат да се изберат прости най-къси пътища от $s$ до достижимите върхове и техните предшественици да се разглеждат като дърво на най-късите пътища с корен $s$. Недостижимите върхове запазват оценка $+\infty$: ако към такъв връх имаше ребро от достижим връх, той също би бил достижим.

Ще използваме индукция по нивата на това дърво. Преди първото обхождане е вярно
\[
s.d=0=\delta(s,s).
\]
Да допуснем, че след някакъв брой обхождания всички върхове до ниво $i-1$ имат точни оценки. Нека $t$ е връх на ниво $i$, а $v$ е неговият родител в дървото на най-късите пътища. При следващото обхождане реброто $(v,t)$ се релаксира. Тъй като
\[
v.d=\delta(s,v),
\]
от лемата за релаксация по най-къс път следва
\[
t.d=\delta(s,t).
\]
Това равенство се запазва.

Всеки прост път съдържа най-много $n-1$ ребра. Следователно след $n-1$ пълни обхождания са достигнати всички нива на дървото на най-късите пътища и всички оценки са точни.
\end{proof}

\begin{lem}
Ако в графа няма отрицателен цикъл, достижим от $s$, то след първите $n-1$ обхождания за всяко ребро $(x,y)$ е изпълнено
\[
y.d\leq x.d+\omega((x,y)).
\]
\end{lem}

\begin{proof}
Ако $x$ е достижим от $s$, тогава и $y$ е достижим и след първите $n-1$ обхождания имаме $x.d=\delta(s,x)$ и $y.d=\delta(s,y)$. Твърдението следва от неравенството за ребро. Ако $x$ не е достижим, тогава $x.d=+\infty$ и неравенството е тривиално.
\end{proof}

\begin{thm}
Ако в графа няма отрицателен цикъл, достижим от $s$, алгоритъмът на Белман--Форд връща \texttt{True} и за всеки връх $v$ е изпълнено
\[
v.d=\delta(s,v).
\]
\end{thm}

\begin{proof}
Точността на оценките след първия двоен цикъл следва от първата лема в този подраздел. От следващата лема тогава за никое ребро условието в последния цикъл не е изпълнено, поради което алгоритъмът достига \texttt{return True}.
\end{proof}

\subsection{Откриване на отрицателен цикъл}\label{s:8.5.3}

\begin{lem}
Ако в графа има отрицателен цикъл, достижим от $s$, то при проверката след $n-1$ обхождания съществува ребро $(x,y)$, за което
\[
y.d>x.d+\omega((x,y)).
\]
\end{lem}

\begin{proof}
Да допуснем противното: за всяко ребро $(x,y)$ е изпълнено
\[
y.d\leq x.d+\omega((x,y)).
\]
Нека отрицателният цикъл е
\[
v_1,e_1,v_2,e_2,\ldots,v_k,e_k,v_1.
\]
Тъй като цикълът е достижим от $s$, след извършените обхождания всички негови върхове имат крайни оценки. Прилагаме предположеното неравенство за всички ребра на цикъла:
\[
d(v_2)\leq d(v_1)+\omega(e_1),
\]
\[
d(v_3)\leq d(v_2)+\omega(e_2),
\]
\[
\cdots
\]
\[
d(v_1)\leq d(v_k)+\omega(e_k).
\]
Като съберем неравенствата, всички оценки се съкращават и получаваме
\[
0\leq \omega(e_1)+\omega(e_2)+\cdots+\omega(e_k).
\]
Това противоречи на отрицателното тегло на цикъла.
\end{proof}

\begin{keythm}[Коректност на Белман--Форд]
Алгоритъмът на Белман--Форд връща \texttt{True} точно когато в графа няма отрицателен цикъл, достижим от $s$; в този случай за всеки връх $v$ е изпълнено
\[
v.d=\delta(s,v).
\]
Ако върне \texttt{False}, има достижим отрицателен цикъл.
\end{keythm}

\begin{proof}
Ако няма достижим отрицателен цикъл, всеки достижим връх има най-къс прост път с най-много $|V|-1$ ребра. По индукция по дължината на такъв път след $k$ пълни обхода на ребрата е получена правилната стойност за всеки връх с най-къс път от най-много $k$ ребра: поредният обход релаксира последното ребро след вече получения префикс. Оценките не могат да са по-малки от истинското разстояние, защото всяка крайна оценка е тегло на реален маршрут. След $|V|-1$ обхода всички оценки са точни и никое ребро не допуска подобрение; алгоритъмът връща \texttt{True}.

Ако има достижим отрицателен цикъл, всички негови върхове имат крайни оценки след тези обходи. Ако никое негово ребро не допуска подобрение, за всяко $(u,v)$ от цикъла е $v.d\le u.d+w(u,v)$. Сумирането съкращава оценките и дава $0\le w(C)$, противоречие. Значи проверката открива подобрение и връща \texttt{False}. Недостижимите върхове остават с $+\infty$ и не предизвикват такова подобрение.
\end{proof}

\subsection{Сложност}\label{s:8.5.4}

Външният цикъл се изпълнява $n-1$ пъти, а при всяко изпълнение се обхождат всички $m$ ребра. Последната проверка има сложност $\Theta(m)$. Следователно времевата сложност е
\[
\Theta(nm).
\]
В най-лошия случай, когато $m=\Theta(n^2)$, тя е кубична:
\[
\Theta(n^3).
\]

\section{Алгоритъм на Флойд--Уоршал}\label{s:8.6}

\subsection{Вариант на задачата}\label{s:8.6.1}

Алгоритъмът на Флойд--Уоршал решава задачата APSP: намира най-късите разстояния между всички двойки върхове. Графът е ориентиран и е номериран с върхове
\[
V=\{1,\ldots,n\}.
\]

Входът се представя с матрица на теглата $\omega$, където
\[
\omega[i,j]=
\begin{cases}
0, & i=j,\\
\omega((i,j)), & i\ne j \text{ и }(i,j)\in E,\\
+\infty, & i\ne j \text{ и }(i,j)\notin E.
\end{cases}
\]

\subsection{Динамично програмиране}\label{s:8.6.2}

За път $p$ нека $\operatorname{intern}(p)$ е множеството от неговите вътрешни върхове. Нека $D^{(k)}[i,j]$ означава теглото на най-къс път от $i$ до $j$, чиито вътрешни върхове са само измежду
\[
\{1,\ldots,k\}.
\]

При $k=0$ не се разрешават вътрешни върхове, затова
\[
D^{(0)}=\omega.
\]

За $k>0$ има две възможности за най-къс път от $i$ до $j$ с разрешени вътрешни върхове от $\{1,\ldots,k\}$:

\begin{itemize}
\item връх $k$ не участва като вътрешен връх; тогава стойността е $D^{(k-1)}[i,j]$;
\item връх $k$ участва като вътрешен връх; тогава пътят се разделя на път от $i$ до $k$ и път от $k$ до $j$, чиито вътрешни върхове са от $\{1,\ldots,k-1\}$.
\end{itemize}

Следователно рекурентната формула е
\[
D^{(k)}[i,j]
=
\min\left\{
D^{(k-1)}[i,j],
D^{(k-1)}[i,k]+D^{(k-1)}[k,j]
\right\}.
\]

\subsection{Псевдокод}\label{s:8.6.3}

\begin{verbatim}
Floyd-Warshall(w)
    D <- w
    for k <- 1 to n
        for i <- 1 to n
            for j <- 1 to n
                D[i, j] <- min(D[i, j], D[i, k] + D[k, j])
    return D
\end{verbatim}

\subsection{Коректност}\label{s:8.6.4}

\begin{keythm}[Коректност на Флойд--Уоршал]
При отсъствие на отрицателни цикли алгоритъмът на Флойд--Уоршал връща матрица $D$, за която
\[
D[i,j]=\delta(i,j)
\]
за всички $i,j\in\{1,\ldots,n\}$.
\end{keythm}

\begin{proof}
Ще докажем с индукция по $k$, че след завършване на итерацията за $k$ матрицата съдържа точно $D^{(k)}$.

За $k=0$ това е вярно по дефиницията на входната матрица $\omega$: позволени са само пътища без вътрешни върхове.

Нека твърдението е вярно за $k-1$. За всяка двойка $i,j$ рекурентната формула разглежда точно двата възможни случая: връх $k$ не помага или $k$ е вътрешен връх на пътя. Във втория случай оптималният път се разделя на два подпътя с вътрешни върхове, номерирани най-много с $k-1$. Следователно присвояването изчислява точно $D^{(k)}[i,j]$.

При $k=n$ всички върхове са разрешени като вътрешни. При липса на отрицателни цикли всеки най-къс път може да бъде избран прост, така че получената стойност е $\delta(i,j)$.
\end{proof}

\subsection{Сложност по време и памет}\label{s:8.6.5}

Трите вложени цикъла извършват $\Theta(n^3)$ операции. Следователно времевата сложност е
\[
\Theta(n^3).
\]

Буквалното съхраняване на всички матрици
\[
D^{(0)},D^{(1)},\ldots,D^{(n)}
\]
изисква
\[
\Theta(n^3)
\]
памет. Достатъчно е обаче да се пазят само две работни матрици: една за $D^{(k-1)}$ и една за $D^{(k)}$. Така използваната памет става
\[
\Theta(n^2).
\]

Възможна е и реализация на място, при която се използва една матрица $D$ и се изпълнява
\[
D[i,j]\leftarrow
\min\{D[i,j],D[i,k]+D[k,j]\}.
\]
Тогава извън входната матрица е нужна $\Theta(1)$ допълнителна памет, но входната матрица се презаписва. При липса на отрицателни цикли случаите $i=k$ или $j=k$ не пречат на коректността, защото диагоналните стойности са нула и връх $k$ не може да бъде вътрешен връх на прост път, който започва или завършва в $k$.

\section{Изпитен фокус}\label{s:8.7}

\begin{itemize}
\item Да се дефинират тегло на път, $\delta(u,v)$, отрицателен цикъл и четирите варианта на задачата за най-къси пътища.

\item Да се обясни защо отрицателен цикъл, разположен по път от $u$ до $v$, прави стойността $\delta(u,v)$ равна на $-\infty$, както и защо при липса на достижим отрицателен цикъл може да се избере прост най-къс път.

\item Да се възпроизведат и използват оптималната подструктура на най-късите пътища, инициализацията \texttt{ISS}, операцията \texttt{Relax} и инвариантът $v.d\geq\delta(s,v)$.

\item За Дийкстра: да се посочат условието за неотрицателни тегла, псевдокодът, множеството $D$, операциите \verb|Extract-Min| и \verb|Decrease-Key| на приоритетната опашка и ролята на проверката $y\in Q$, доказателствената идея с първия неправилно избран връх и сложностите при различни реализации на опашката.

\item Да се посочи контрапример, показващ, че при отрицателно тегло Дийкстра дава грешен резултат дори без отрицателен цикъл.

\item За DAG-SSSP: да се посочат топологичната наредба, релаксирането на всички ребра в тази наредба, доказателството по реда на върховете, сложността $\Theta(n+m)$ и предимствата пред Дийкстра.

\item За Белман--Форд: да се възпроизведат $n-1$ пълни обхождания на всички ребра, финалната проверка, доказателството чрез нивата в дърво на най-късите пътища и доказателството за откриване на отрицателен цикъл чрез сумиране на неравенства по цикъла.

\item За Флойд--Уоршал: да се дефинира $D^{(k)}[i,j]$, да се изведе рекурентната формула според това дали връх $k$ участва, да се даде псевдокодът, индуктивното доказателство и сложностите $\Theta(n^3)$ по време и $\Theta(n^2)$ памет при реализация с една или две матрици.
\end{itemize}


\part{Ядро на компютърните науки}

\chapter{Компютърни архитектури. Формати на данните. Вътрешна структура на централен процесор.}
%% Drafted from the mapped corpus by scripts/draft_topics.py.
% Model: gpt-5.6-luna; source characters: 10195.
\begin{konspekt}
\kreq{Официално заглавие}{„Компютърни архитектури --- Формати на данните. Вътрешна структура на централен процесор --- \emph{блокове и конвейерна обработка, инструкции}.“ Заглавието на главата е съкратено; конвейерът е изискване, не допълнение.}
\kreq{Описание}{обща структура на компютрите, запомнена програма --- \kref{9.1}.}
\kreq{Описание}{концептуално изпълнение на инструкциите --- \kref{9.4}.}
\kreq{Формати на данните}{цели двоични числа, двоично-десетични числа, числа с плаваща запетая, знакови данни и кодови таблици --- \kref{9.2}.}
\kreq{Централен процесор}{регистри, АЛУ, регистри на състоянията и флаговете, блокове за управление, връзка с паметта, дешифрация на инструкциите, преходи --- \kref{9.3}.}
\kreq{Конвейерна обработка}{\kref{9.5}.}
\kreq{Литература}{[5], [36], [45].}
\end{konspekt}

\section{Обща структура на компютрите и компютърни архитектури}\label{s:9.1}

Компютърната архитектура описва основните компоненти на компютъра, начина на свързването им и принципите, по които се изпълняват програмите. В общия случай компютърът се състои от:

\begin{itemize}
    \item централен процесор;
    \item основна памет;
    \item входно-изходни устройства.
\end{itemize}

Тези компоненти са взаимносвързани и работят съвместно, за да изпълняват програми. Централният процесор обработва данните и управлява изпълнението на инструкциите. Основната памет съхранява инструкции и данни. Входно-изходните устройства осигуряват обмен на информация между компютъра и външната среда.

Връзката между компонентите се осъществява чрез шини.

\begin{defn}[Шина]
Шината е съвкупност от проводници, по които се предава информация между компонентите на компютъра. Обичайно се разграничават три шини:

\begin{itemize}
\item \emph{адресна шина} --- по нея процесорът подава адреса на клетката от паметта или на устройството, към което се обръща;
\item \emph{шина за данни} --- по нея се пренася самата стойност, която се чете или записва;
\item \emph{контролна шина} --- по нея се предават управляващите сигнали, например дали текущото обръщение е четене или запис, и сигналите за прекъсване.
\end{itemize}
\end{defn}

\subsection{Архитектура на фон Нойман}\label{s:9.1.1}

Повечето съвременни компютри се основават на архитектурата на Джон фон Нойман. Тя се характеризира със следните основни идеи:

\begin{enumerate}
    \item Инструкциите и данните се съхраняват в памет, която позволява четене и запис.
    \item Съдържанието на паметта се адресира чрез адрес, независимо от вида на данните, които се намират на съответното място.
    \item Инструкциите се изпълняват последователно, от една инструкция към следващата, освен ако текущата инструкция не промени реда на изпълнение.
\end{enumerate}

Последната особеност е съществена за разбирането на програмния брояч. Той съдържа адреса на инструкцията, която трябва да бъде извлечена и изпълнена. При обикновено последователно изпълнение програмният брояч се насочва към следващата инструкция. При преход или друга инструкция, която променя управлението, в него се записва различен адрес.

\begin{defn}[Запомнена програма]
Запомнена програма е програма, чиито инструкции са представени в двоичен вид и са съхранени в паметта заедно с данните, върху които работи програмата. Процесорът извлича инструкциите от паметта и ги изпълнява последователно, като редът може да бъде променен чрез инструкции за преход.
\end{defn}

Така програмата не е външна поредица от действия, която процесорът следва по специален начин, а е информация, записана в паметта. Процесорът интерпретира тази информация като инструкции и данни според съответния формат.

\section{Формати на данните}\label{s:9.2}

\begin{defn}[Формат на данни]
Форматът на данните определя каква битова последователност представя дадена стойност. Една и съща последователност от битове може да има различно значение в зависимост от това дали се разглежда като цяло число, двоично-десетично число, число с плаваща запетая или код на символ.
\end{defn}

\subsection{Цели двоични числа}\label{s:9.2.1}

Целите двоични числа се представят в двоична бройна система. При знаковите числа един от битовете обикновено се използва за представяне на знака. Разглеждат се различни кодове за представяне на отрицателните числа.

\subsubsection{Прав код}

\begin{defn}[Прав код]
При правия код най-старшият бит в $n$-битовото поле е знаков. Стойност $0$ означава положително число, а стойност $1$ означава отрицателно число. Останалите битове представят абсолютната стойност на числото в двоична бройна система.
\end{defn}

Следователно в прав код има отделно поле за знака и отделно поле за абсолютната стойност. Основните недостатъци на този начин на представяне са:

\begin{itemize}
    \item нулата има две представяния -- положителна и отрицателна;
    \item събирането и изваждането изискват по-сложна апаратна реализация.
\end{itemize}

Диапазонът на представимите стойности е от $-(2^{n-1}-1)$ до $2^{n-1}-1$.

\subsubsection{Двоично-допълнителен код}

\begin{defn}[Двоично-допълнителен код]
При двоично-допълнителния код $n$-битовата последователност $b_{n-1}b_{n-2}\ldots b_0$ представя числото
\[
-b_{n-1}2^{n-1}+\sum_{i=0}^{n-2}b_i2^i .
\]
Тоест всички битове имат обичайните си тегла, но най-старшият участва с отрицателно тегло. Затова той отново определя знака: стойност $1$ дава отрицателно число.

Еквивалентно, отрицателното число $x<0$ се записва като двоичния запис на $2^n+x$; оттук идва и името на кода.
\end{defn}

От определението следва диапазонът на представимите стойности:
\[
[-2^{n-1},\,2^{n-1}-1],
\]
защото най-малката стойност се получава при $b_{n-1}=1$ и всички останали битове нула, а най-голямата --- при $b_{n-1}=0$ и всички останали единици.

Този код използва едно представяне на нулата и позволява операциите със знакови числа да се реализират чрез обикновено двоично събиране, като се работи с фиксиран брой битове.

Смяната на знака на число в двоично-допълнителен код се извършва в две стъпки:

\begin{enumerate}
    \item всички битове се инвертират побитово;
    \item към получения резултат се прибавя $1$.
\end{enumerate}

\begin{prop}
Нека $x$ е представено в $n$-битов двоично-допълнителен код и нека $\overline{x}$ е побитовото му инвертиране. Тогава $\overline{x}+1$ е представянето на $-x$ при пресмятане с $n$ бита.
\end{prop}

\begin{proof}
Инвертирането на бит го заменя с $1$ минус стойността му, затова при разглеждане на битовете като безнакови тегла
\[
x+\overline{x}=\underbrace{11\ldots1}_{n}=2^n-1 .
\]
Следователно
\[
\overline{x}+1=2^n-x .
\]
При работа с $n$ бита преносът извън най-старшия бит се отхвърля, тоест пресмята се по модул $2^n$, а
\[
2^n-x\equiv -x \pmod{2^n}.
\]
\end{proof}

\begin{example}
Нека числото $3_{10}$ бъде представено с четири бита:

\[
3_{10}=0011_2.
\]

Инвертирането на битовете дава

\[
0011 \longrightarrow 1100.
\]

След прибавяне на $1$ се получава

\[
1100+1=1101.
\]

Следователно $1101_2$ е представянето на $-3_{10}$ в четирибитов двоично-допълнителен код.
\end{example}

При събиране на число с противоположното му число резултатът е нула в рамките на използвания брой битове. Преносът извън най-старшия бит не се включва в резултата.

\begin{supp}[Бележка]
При работа с двоично-допълнителен код броят на битовете трябва да бъде фиксиран предварително. Например в четирибитово поле резултатът от събиране може да съдържа пренос извън полето, който не е част от представената стойност.
\end{supp}

\subsection{Двоично-десетични числа}\label{s:9.2.2}

\begin{defn}[Двоично-десетичен код]
При двоично-десетичното представяне всяка десетична цифра се заменя с уникална двоична последователност. За една десетична цифра се използват четири бита. Понеже с четири бита могат да се образуват $16$ комбинации, а десетичните цифри са само десет, шест комбинации остават неизползвани.
\end{defn}

Предимството на двоично-десетичното представяне е, че числата са удобни за извеждане и за работа с десетични цифри. Недостатък е усложненото изпълнение на аритметичните операции.

Различават се две основни форми:

\begin{description}
    \item[Непакетирана форма.] Всяка десетична цифра се съхранява в отделен байт.
    \item[Пакетирана форма.] В един байт се съхраняват две десетични цифри, като всяка цифра заема по четири бита.
\end{description}

При непакетираната форма се използва повече памет, но обработката на отделните цифри е по-пряка. При пакетираното представяне паметта се използва по-икономично, защото два четирибитови кода се съхраняват в един байт.

\subsection{Двоични числа с плаваща запетая}\label{s:9.2.3}

Числата с плаваща запетая се използват за представяне на нецели числа. Общата идея е стойността да се представи чрез знак, мантиса и характеристика, тоест като произведение на мантиса и степен на основата:

\[
\pm S \times B^{\pm E}.
\]

\begin{defn}[Мантиса, характеристика, отклонение]
В представянето $\pm S\times B^{\pm E}$:

\begin{itemize}
\item $S$ е \emph{мантисата} --- значещите цифри на числото;
\item $B$ е \emph{основата} на представянето, при двоичните формати $B=2$;
\item $E$ е \emph{характеристиката}, наричана още \emph{експонента} --- степента, на която се повдига основата, тоест с колко позиции се измества запетаята.
\end{itemize}

Полето за характеристиката съдържа не самата стойност $E$, а изместена стойност, за да се представят и отрицателни експоненти без отделен знаков бит. Фиксираната добавка се нарича \emph{отклонение}: истинската експонента се получава, като отклонението се извади от записаното в полето.
\end{defn}

В един двоичен формат числото може да бъде разделено на три части:

\begin{itemize}
    \item един бит за знака;
    \item поле за характеристиката, или експонентата;
    \item поле за мантисата.
\end{itemize}

Посоченият пример използва формат с общо $32$ бита:

\begin{itemize}
    \item $1$ бит за знака;
    \item $8$ бита за характеристиката;
    \item $23$ бита за мантисата.
\end{itemize}

За характеристиката обикновено се използва отклонение. Отклонението е фиксирана стойност, която се изважда от записаната характеристика, за да се получи действителната стойност на експонентата. При $k$ бита за характеристиката обичайно се използва отклонение

\[
2^{k-1}-1.
\]

При $k=8$ това отклонение е $127$. Източникът посочва диапазон на действителните стойности на характеристиката от $-127$ до $128$.

За да се избегнат различни представяния на една и съща стойност, числата се нормализират. Без такова условие едно и също число има много записи: например
\[
1{,}5=1{,}5\times2^0=0{,}75\times2^1=0{,}375\times2^2=\ldots
\]

\begin{defn}[Нормализирано число]
Ненулево число с плаваща запетая е \emph{нормализирано}, ако характеристиката му е избрана така, че мантисата да попада в предварително фиксиран интервал. При основа $B=2$ това означава
\[
1\leq S<2,
\]
тоест числото се записва във вида
\[
\pm 1{,}f\times 2^{E}.
\]
При това условие всяко ненулево число има точно едно представяне.
\end{defn}

Тъй като при двоична основа старшата цифра на нормализирана мантиса е винаги $1$, тя не се записва в полето за мантисата; говори се за \emph{скрит бит}. Така при $23$ бита се получава точност от $24$ значещи двоични цифри. Нулата не може да се нормализира и за нея се запазва отделно представяне.

\begin{supp}[Бележка]
Изречението, което определя нормализираното число, е останало незавършено в източника: формулата след него липсва. Дадената по-горе форма е стандартната за двоичните формати и е добавена, за да бъде определението използваемо.
\end{supp}

Аритметиката на числата с плаваща запетая е по-сложна от аритметиката на целите числа. Стандартното аритметико-логическо устройство не е достатъчно за непосредственото изпълнение на всички операции върху тях. Затова се използва отделно изчислително звено -- блок за плаваща запетая, наричан още \emph{Floating Point Unit}.

Инструкциите за работа с числа с плаваща запетая имат специфични особености. Поради това те се разглеждат отделно от инструкциите за обикновени цели двоични числа.

\begin{supp}[Бележка]
При числата с плаваща запетая точният диапазон и конкретните специални представяния зависят от избрания формат. Тук се разглежда само общата структура, дадена в източника: знак, характеристика и мантиса.
\end{supp}

\subsection{Знакови данни и кодови таблици}\label{s:9.2.4}

\begin{defn}[Кодова таблица]
Кодова таблица е съответствие между символите от определена азбука и битови последователности с определена дължина, чрез които тези символи се съхраняват и обработват в компютъра.
\end{defn}

Кодовите таблици позволяват на процесора и на програмите да обработват букви, цифри, препинателни знаци и други символи като цифрови данни.

Източникът посочва два стандарта:

\begin{description}
    \item[ASCII.] Използва седем бита за кодиране на английската азбука, цифрите, пунктуацията и някои други символи.
    \item[Unicode.] Стандарт за представяне на символи чрез кодови стойности, предназначен за работа с различни азбуки и знакови системи.
\end{description}

\section{Централен процесор}\label{s:9.3}

Централният процесор изпълнява инструкциите на програмата и управлява обработката на данните. За тази цел той трябва да извършва няколко основни действия:

\begin{enumerate}
    \item да извлича инструкции от паметта;
    \item да интерпретира инструкциите;
    \item да извлича необходимите операнди;
    \item да обработва данните;
    \item да записва резултатите.
\end{enumerate}

Извличането на инструкция означава прочитане на инструкция от достъпно за процесора място -- регистри, кеш или основна памет. Интерпретирането включва дешифрация на инструкцията, за да се установи каква операция трябва да бъде изпълнена и какви са нейните операнди.

Изпълнението може да изисква четене на данни от паметта или от входно-изходен модул. Самата обработка може да представлява аритметична или логическа операция. След приключването й резултатът се записва в регистър, памет или входно-изходно устройство.

Основните компоненти на процесора са:

\begin{itemize}
    \item аритметико-логическо устройство;
    \item блок за управление;
    \item регистри.
\end{itemize}

\subsection{Регистри}\label{s:9.3.1}

\begin{defn}[Регистър]
Регистрите са вътрешни места за временно съхраняване на данни, адреси, инструкции и информация за състоянието на процесора. Те позволяват на процесора да работи с междинни резултати и да запомня позицията си в програмата.
\end{defn}

Регистрите могат да се разделят на две основни групи:

\begin{enumerate}
    \item потребителски видими регистри;
    \item регистри за управление и състояние.
\end{enumerate}

\subsubsection{Потребителски видими регистри}

\begin{defn}[Потребителски видим регистър]
Потребителски видими са регистрите, които могат да бъдат достъпвани чрез машинни инструкции. Те позволяват на програмиста или на компилатора да съхранява междинни резултати и да намалява необходимостта от обръщения към основната памет.
\end{defn}

Основните подкатегории са:

\begin{description}
    \item[Регистри с общо предназначение.] Съхраняват междинни резултати и могат да се използват при адресиране.
    \item[Регистри за данни.] Предназначени са за съхраняване на данни, но не и за изчисляване на адрес на операнд.
    \item[Регистри за адреси.] Използват се за адресиране. Те могат да бъдат регистри с общо предназначение или да са предназначени за конкретен режим на адресиране, например като индексни регистри, указатели на стека или сегментни указатели.
\end{description}

Границата между потребителски видимите регистри и регистрите за управление и състояние не е еднаква при всички процесори. Някои регистри, използвани основно от операционната система, могат да бъдат достъпни чрез машинни инструкции, макар да не се използват пряко от обикновените програми.

\subsubsection{Регистри за управление и състояние}

Четири регистъра са особено важни за изпълнението на инструкциите:

\begin{description}
    \item[Програмен брояч.] Съдържа адреса на инструкцията, която трябва да бъде прочетена. При последователно изпълнение този адрес се променя така, че да сочи към следващата инструкция.
    \item[Регистър на инструкциите.] Съдържа последно извлечената инструкция, която се декодира и изпълнява.
    \item[Регистър на адреса на паметта (MAR).] Съдържа адреса на място в паметта, към което ще се извърши обръщение.
    \item[Буферен регистър на паметта (MBR).] Съдържа данните, които трябва да бъдат записани в паметта, или последно прочетените от паметта данни.
\end{description}

Много процесори имат и регистър на флаговете.

\begin{defn}[Флаг и регистър на състоянието]
\emph{Флаг} е отделен бит, който записва едно двоично условие --- най-често свойство на резултата от последната операция на аритметико-логическото устройство или режим на работа на процесора. Регистърът, в който са събрани тези битове, се нарича \emph{регистър на състоянието} или \emph{регистър на флаговете}.

Флаговете не се четат като обикновени данни, а се използват от инструкциите за условен преход: те решават коя да бъде следващата инструкция.
\end{defn}

Сред често срещаните флагове са:

\begin{itemize}
    \item флаг за знак;
    \item флаг за нулев резултат;
    \item флаг за пренос;
    \item флаг за препълване;
    \item флаг за разрешени или забранени прекъсвания.
\end{itemize}

Флагът за знак показва знака на резултата, флагът за нула показва дали резултатът е нулев, а флаговете за пренос и препълване описват определени условия при аритметични операции. Флагът за прекъсване участва в управлението на възможността процесорът да приема прекъсвания.

\subsection{Аритметико-логическо устройство}\label{s:9.3.2}

\begin{defn}[Аритметико-логическо устройство]
Аритметико-логическото устройство, съкратено АЛУ, извършва аритметични и логически операции върху операнди. То получава данни от регистрите или от пътя за данни, изпълнява избраната операция и връща резултата към регистър или към следващ етап на обработката.
\end{defn}

Работата на АЛУ се определя от управляващи сигнали. Те задават каква операция трябва да бъде извършена и как резултатът да бъде прехвърлен. В зависимост от резултата АЛУ може да промени и флаговете в регистъра на състоянието.

\subsection{Блок за управление}\label{s:9.3.3}

\begin{defn}[Микрооперация]
Микрооперацията е най-малкото неделимо действие, което процесорът извършва за един такт --- например прехвърляне на съдържанието на един регистър в друг, подаване на адрес към паметта или активиране на конкретна функция на аритметико-логическото устройство. Изпълнението на една машинна инструкция се разлага на последователност от микрооперации.
\end{defn}

Блокът за управление координира работата на процесора. Той изпълнява две основни задачи:

\begin{enumerate}
    \item задвижва процесора през поредица от микрооперации в правилната последователност;
    \item генерира управляващи сигнали, които предизвикват изпълнението на микрооперациите.
\end{enumerate}

Управляващите сигнали контролират прехвърлянето на данни между регистрите, избора на операция на АЛУ и обмена с паметта и входно-изходните устройства. Те управляват отварянето и затварянето на логически пътища, чрез които данните достигат до съответните регистри и функционални блокове.

Входове за блока за управление са:

\begin{itemize}
    \item тактовият сигнал;
    \item съдържанието на регистъра на инструкциите;
    \item флаговете на процесора;
    \item управляващи сигнали от контролната шина.
\end{itemize}

Тактовият сигнал задава последователността на микрооперациите. Регистърът на инструкциите предоставя кода на операцията и информацията за режима на адресиране. Флаговете показват състоянието на процесора и резултатите от предишни операции. Управляващите сигнали от контролната шина осигуряват информация за състоянието на паметта и входно-изходните устройства.

Изходите на блока за управление са:

\begin{itemize}
    \item управляващи сигнали вътре в процесора;
    \item управляващи сигнали към контролната шина.
\end{itemize}

Сигналите вътре в процесора осигуряват прехвърляне на данни от един регистър в друг и активират конкретни функции на АЛУ. Сигналите към контролната шина управляват обмена с паметта и входно-изходните устройства.

\subsection{Инструкции}\label{s:9.3.4}

Машинната инструкция е двоично кодирана команда, която определя операцията и необходимите за нея данни. Обикновено една инструкция може да съдържа:

\begin{itemize}
    \item код на операцията;
    \item указател към един или повече входни операнди;
    \item указател към изходен операнд;
    \item указател към следващата инструкция.
\end{itemize}

\begin{defn}[Код на операцията]
Кодът на операцията, или \emph{opcode}, е двоична част от инструкцията, която определя какво действие трябва да извърши процесорът.
\end{defn}

\begin{defn}[Режим на адресиране]
Режимът на адресиране определя как от полето на инструкцията се получава действителният адрес, или самата стойност, на операнда. Сред основните режими са:

\begin{itemize}
\item \emph{непосредствен} --- операндът е записан в самата инструкция;
\item \emph{регистров} --- операндът е в посочения регистър;
\item \emph{директен} --- инструкцията съдържа адреса на операнда в паметта;
\item \emph{недиректен} --- инструкцията съдържа адрес, на който е записан адресът на операнда, така че е необходимо още едно обръщение към паметта.
\end{itemize}
\end{defn}

Операндите могат да се намират в регистри, в паметта или във входно-изходно устройство. За да ги използва, процесорът трябва да определи местоположението им според зададения режим на адресиране.

\subsection{Връзка с паметта}\label{s:9.3.5}

Връзката между процесора и паметта се реализира посредством шини, сред които особено значение има контролната шина. Процесорът подава управляващи сигнали, които указват дали се извършва четене или запис, а чрез регистрите MAR и MBR се задават адресът и обменяните данни.

При четене адресът на необходимото място се поставя в MAR. След осъществяване на обръщението прочетените данни се приемат в MBR. При запис адресът се поставя в MAR, а данните за запис -- в MBR. Блокът за управление подава съответния управляващ сигнал към паметта.

\section{Концептуално изпълнение на инструкциите}\label{s:9.4}

Състоянието на процесора образува контекста, необходим за изпълнението на следващата инструкция. Важна част от този контекст е програмният брояч, който съдържа адреса на инструкцията, която трябва да бъде извлечена.

На концептуално ниво изпълнението на инструкцията включва следните фази:

\begin{enumerate}
    \item извличане;
    \item декодиране и интерпретиране;
    \item изчисляване на адресите на операндите;
    \item извличане на операндите;
    \item изпълнение;
    \item записване на резултата;
    \item проверка и обслужване на прекъсване.
\end{enumerate}

\subsection{Извличане на инструкция}\label{s:9.4.1}

При извличането процесорът използва адреса от програмния брояч. Този адрес се подава към паметта, а инструкцията се прочита и се поставя в регистъра на инструкциите. След извличането програмният брояч обикновено се насочва към следващата инструкция.

\subsection{Декодиране}\label{s:9.4.2}

Блокът за управление анализира съдържанието на регистъра на инструкциите. От кода на операцията се определя какво действие трябва да бъде извършено. Разпознават се също операндите и режимът на адресиране.

Дешифрацията на инструкциите е необходима, защото една и съща обща процедура трябва да обслужва различни операции -- например аритметични операции, логически операции, прехвърляния на данни и преходи.

\subsection{Изчисляване и извличане на операнди}\label{s:9.4.3}

Ако инструкцията използва операнд от паметта или от входно-изходно устройство, процесорът изчислява адреса на операнда. След това извършва необходимото обръщение и извлича данните.

Ако операндът вече се намира в регистър, допълнително обръщение към основната памет не е необходимо. Във всички случаи блокът за управление трябва да осигури подаването на правилните данни към АЛУ или към друг функционален блок.

\subsection{Изпълнение и записване на резултата}\label{s:9.4.4}

След извличането на операндите се изпълнява операцията, зададена от инструкцията. При аритметична или логическа операция АЛУ обработва данните. Резултатът може да бъде записан в регистър, в паметта или към входно-изходно устройство.

Ако операцията променя състоянието на процесора, се променят и съответните флагове. След това процесорът продължава със следващата инструкция или изпълнява преход, ако инструкцията е променила програмния брояч.

\subsection{Недиректна стъпка}\label{s:9.4.5}

След извличането на инструкцията може да бъде необходимо да се извърши недиректна стъпка. Тя се прилага, когато инструкцията използва недиректна адресация. В такъв случай извлечената от инструкцията информация не е непосредствено крайният адрес на операнда. Процесорът трябва да извърши допълнително обръщение, за да получи необходимия адрес.

\subsection{Прекъсвания}\label{s:9.4.6}

\begin{defn}[Прекъсване]
Прекъсването е сигнал, който кара процесора да спре временно изпълнението на текущата програма и да премине към предварително определена обслужваща процедура. Сигналът може да идва от входно-изходно устройство, от таймер или от самото изпълнение на инструкция, например при опит за деление на нула.

Прекъсването не се проверява по средата на инструкция: то се обслужва след завършването на текущата инструкция, което е и последната фаза от концептуалното изпълнение.
\end{defn}

Ако прекъсванията са разрешени и възникне прекъсване, процесорът запазва текущото състояние на програмата и преминава към обслужване на прекъсването.

Запазеното състояние включва:

\begin{itemize}
    \item съдържанието на програмния брояч, тоест адреса на следващата инструкция;
    \item релевантна информация за изпълняваната програма;
    \item състоянието на регистрите;
    \item състоянието на входно-изходните устройства.
\end{itemize}

След обслужването на прекъсването запазеното състояние позволява изпълнението на прекъснатата програма да бъде продължено.

\subsection{Преходи}\label{s:9.4.7}

При обикновено изпълнение програмният брояч сочи към следващата инструкция. Инструкцията за преход променя това поведение, като записва нов адрес в програмния брояч. Новият адрес може да зависи от флаговете на процесора.

Например при инструкция от вида \verb|increment-and-skip-if-zero| блокът за управление увеличава програмния брояч, ако флагът за нулев резултат е установен. Така флаговете могат да участват пряко в избора на следващата инструкция.

\begin{supp}[Бележка]
Формулировката е тази от източника. Става дума за \emph{допълнително} увеличаване: програмният брояч се увеличава при всяко извличане на инструкция, а при изпълнена проверка се увеличава още веднъж, с което следващата инструкция се прескача. Оттам идва и \emph{skip} в името на инструкцията.
\end{supp}

\section{Конвейерна обработка}\label{s:9.5}

Изпълнението на една инструкция може да бъде представено като последователност от фази: извличане, декодиране, извличане на операнди, изпълнение и записване на резултата. В процесора тези фази се реализират чрез поредица от микрооперации, управлявани от блока за управление.

\begin{defn}[Конвейерна обработка]
При конвейерната обработка изпълнението на инструкциите се разделя на последователни етапи, всеки от които се обслужва от отделна част на процесора. Докато една инструкция преминава през даден етап, следващата вече може да се обработва в предходния. По този начин отделните части на процесора работят едновременно върху различни инструкции.

Конвейерът не съкращава времето за изпълнение на една инструкция, а увеличава броя инструкции, завършвани за единица време.
\end{defn}

Основата на конвейерната обработка е съгласуваното управление на етапите чрез тактовия сигнал. Блокът за управление генерира управляващите сигнали в правилната последователност, така че данните да преминават между регистрите, АЛУ, паметта и останалите блокове.

При инструкции за преход нормалната последователност може да бъде променена, тъй като програмният брояч получава нов адрес. Затова преходите зависят от състоянието на процесора и от резултатите на предходни операции.

\section{Изпитен фокус}\label{s:9.6}

При подготовка за устен изпит трябва да могат да бъдат възпроизведени:

\begin{itemize}
    \item трите основни компонента на компютъра -- централен процесор, основна памет и входно-изходни устройства;
    \item трите основни идеи на архитектурата на фон Нойман;
    \item понятието запомнена програма;
    \item представянето на цели числа в прав и двоично-допълнителен код;
    \item алгоритъмът за смяна на знака в двоично-допълнителен код и защо той работи: $x+\overline{x}=2^n-1$, откъдето $\overline{x}+1=2^n-x\equiv-x \pmod{2^n}$;
    \item пакетираното и непакетираното двоично-десетично представяне;
    \item частите на числото с плаваща запетая -- знак, характеристика и мантиса --- и ролята на отклонението при характеристиката;

    \item какво означава нормализирано число и защо старшият бит на мантисата може да не се записва;
    \item ролята на кодовите таблици, ASCII и Unicode;
    \item основните блокове на процесора -- регистри, АЛУ и блок за управление;
    \item предназначението на програмния брояч, регистъра на инструкциите, MAR и MBR;
    \item флаговете за знак, нула, пренос, препълване и прекъсване;
    \item задачите и входовете и изходите на блока за управление;
    \item съставът на машинната инструкция -- opcode и указатели към операнди и следваща инструкция;
    \item трите шини --- адресна, за данни и контролна --- и ролята на контролната шина при връзката с паметта;

    \item понятието режим на адресиране и разликата между директно и недиректно адресиране;

    \item понятието микрооперация и връзката му с работата на блока за управление;
    \item фазите на изпълнение на инструкцията;
    \item изчисляването на адреси и извличането на операнди;
    \item какво е прекъсване и защо се обслужва след завършването на текущата инструкция;

    \item запазването на състоянието при прекъсване;
    \item ролята на програмния брояч и флаговете при изпълнение на преходи;
    \item идеята на конвейерната обработка и последователността от микрооперации.
\end{itemize}


\chapter{Структура и йерархия на паметта. Сегментна и странична преадресация. Прекъсвания.}
% AUTO-GENERATED by scripts/build_topics.py — do not edit.
% Edit topics/bodies/topic_10.tex or topics/preamble.tex instead.
%% Placeholder. Replace with the written chapter.
\textit{Източници:} PDF \texttt{10.pdf}

\subsection*{Анотация от конспекта}
Структура и йерархия на паметта. Сегментна и странична преадресация. Система за прекъсване – приоритети и обслужване. Структура на основната памет. Йерархия – кеш, основна и виртуална памет. Сегметна и странична преадресация – селектор, дескриптор, таблици и регистри при сегметна преадресация; каталог на страниците, описател, стратегии на подмяна на страниците при странична преадресация. Система за прекъсване – видове прекъсвания, структура и обработка, конкурентност и приоритети, контролери на прекъсванията. Литература: [5], [45], [50].


\chapter{Файлова система. Функции, структура и реализация.}
% AUTO-GENERATED by scripts/build_topics.py — do not edit.
% Edit topics/bodies/topic_11.tex or topics/preamble.tex instead.

\begin{konspekt}
\kreq{Описание}{файлове и основни атрибути; единно пространство от именувани обекти; структура, физическо представяне, точки на монтиране; видове обекти --- файлове, директории, устройства, информационни канали; основни функции --- достъп, изолация и защита, управление на правата --- \kref{11.1}.}
\kreq{Реализация}{кратък обзор на ext2/ext3 --- \kref{11.2}.}
\kreq{Ускоряване и надеждност}{кеширане на вход/изхода, отлагане на записа, алгоритъм на асансьора; журнална файлова система, RAID1 и RAID5 --- \kref{11.3}.}
\kreq{POSIX функции}{\texttt{open}, \texttt{close}, \texttt{read}, \texttt{write}, \texttt{lseek}, включително частичен вход/изход --- \kref{11.4}.}
\kreq{Задача 1}{извличане на интервали от двоичен файл --- схема в \kr{11.5.1}, пълна програма в \kref{11.6}.}
\kreq{Задача 2}{сортиране на файл от 8-битови числа --- \kref{11.5.2}.}
\kreq{Задача 3}{прилагане на двоична кръпка --- \kref{11.5.3}.}
\kstatus{изрично се оценява скорост, пропорционална на реално обработваните данни. Задача 1 е дадена като пълна програма; задачи 2 и 3 са схеми, не програми. Първичните източници за този въпрос липсват в корпуса --- главата следва официалната анотация и подлежи на проверка по литературата.}
\kreq{Литература}{[42, четвърта глава], [48], [49, глави 10--12].}
\end{konspekt}

\begin{supp}[Статус на източниците]
В наличния корпус липсват първичните материали за този въпрос. Главата следва подробната официална анотация и е сверена с предоставения сравнителен конспект. Преди окончателно академично издание тя трябва да се провери по посочената литература: [42, четвърта глава], [48], [49, глави 10--12].
\end{supp}

\section{Файлове и файлова система}\label{s:11.1}

\begin{defn}[Файл]
Файлът е именуван информационен обект, чието съдържание се пази устойчиво и се достъпва като последователност от байтове или записи.
\end{defn}

Операционната система свързва името на файла с метаданни и с физическите блокове, в които е разположено съдържанието. Типични атрибути са тип, размер, собственик, група, права за достъп и времена за последна промяна на съдържанието, метаданните и достъпа. Наличието на отделно време на създаване зависи от конкретната файлова система.

\begin{defn}[Файлова система]
Файловата система е абстракция и реализация за именуване, организиране, съхраняване, намиране, защита и споделяне на информационни обекти.
\end{defn}

\begin{defn}[Видове обекти и директория]
В Unix-подобните системи обектите образуват единно дърво с корен \texttt{/}. В него участват обикновени файлове, директории, символни връзки, специални файлове за устройства, именувани програмни канали (FIFO) и сокети. Директорията съпоставя име с обект на файловата система; самото съдържание на обикновения файл не съдържа името му.
\end{defn}

\begin{defn}[Монтиране]
Отделна файлова система се включва в дървото чрез \emph{монтиране} върху съществуваща директория --- точка на монтиране. Така различни устройства и реализации се виждат през едно пространство от пътища.
\end{defn}

Изолацията и защитата се основават на идентичността на процеса и метаданните на обекта. Класическият Unix модел различава права за четене, запис и изпълнение за собственик, група и останали. За директорията правото за изпълнение означава преминаване и търсене по имена, а правото за запис --- промяна на записите в нея.

\section{Физическо представяне и ext2/ext3}\label{s:11.2}

Дискът се дели на блокове. Файловата система пази служебни структури за свободните и заетите блокове, метаданните на файловете и записите в директориите. При ext2 основни елементи са:

\begin{itemize}
\item \emph{superblock} с общи параметри и състояние на файловата система;
\item групи от блокове с битови карти за свободни блокове и inode-и;
\item таблица от inode структури;
\item блокове с данни и директории.
\end{itemize}

\begin{defn}[Inode]
Inode е структура с метаданните на обекта и указатели към блоковете на съдържанието му. Името се пази в запис на директория, който сочи към номер на inode.
\end{defn}

Малките файлове се адресират чрез преки указатели, а по-големите --- и чрез единично, двойно и тройно косвени блокове с указатели. Твърдата връзка е допълнителен запис на директория към същия inode; символичната връзка е отделен обект, който съдържа път.

\begin{defn}[Журнална файлова система]
ext3 развива ext2 чрез журнал. Преди критични промени те се записват като транзакция в журнала, което позволява след срив файловата система да възстанови съгласувано състояние без пълно сканиране на всички блокове. Режимът на журнализиране определя дали се журнализират само метаданните или и потребителските данни.
\end{defn}

\section{Ускоряване и надеждност}\label{s:11.3}

\begin{defn}[Кеш, отложен запис и алгоритъм на асансьора]
\begin{description}
\item[Кеширане.] Често използвани блокове и файлови страници се пазят в основната памет. Четенето напред (read-ahead) зарежда вероятно нужни следващи блокове.
\item[Отложен запис.] Промените първо се натрупват като замърсени кеширани страници и по-късно се записват групово. Това ускорява работата, но при срив незаписаните данни могат да се загубят; \texttt{fsync} изисква устойчиво записване на съответните промени.
\item[Алгоритъм на асансьора.] При въртящ диск заявките се подреждат подобно на SCAN: главата обслужва позиции в едната посока, после сменя посоката. Целта е по-малко механично преместване; ползата не е същата при SSD.
\end{description}
\end{defn}

\begin{defn}[RAID1 и RAID5]
RAID използва няколко диска като едно блоково устройство. RAID1 огледално дублира данните и понася отказ на един диск във всяка огледална двойка. RAID5 разпределя данните и паритета между поне три диска и обичайно понася отказ на един диск, но записът изисква обновяване на данни и паритет. RAID подобрява наличността, но не е резервно копие: не пази от изтриване, повреда от софтуер или загуба на целия масив.
\end{defn}

\section{POSIX функции за файлове}\label{s:11.4}

\begin{verbatim}
int     open(const char *path, int flags, ...);
int     close(int fd);
ssize_t read(int fd, void *buf, size_t count);
ssize_t write(int fd, const void *buf, size_t count);
off_t   lseek(int fd, off_t offset, int whence);
\end{verbatim}

\texttt{open} връща файлов дескриптор или $-1$. Флаговете задават например четене, запис, създаване и съкращаване. \texttt{read} връща броя реално прочетени байтове, $0$ при край на файл и $-1$ при грешка. \texttt{write} също може да запише по-малко от поисканото, затова надеждната програма повтаря операцията до обработване на целия буфер. Прекъсване със сигнал може да даде \texttt{EINTR} и да изисква повторение.

\texttt{lseek} променя текущото байтово отместване спрямо началото, текущата позиция или края. Тя не извършва вход/изход и не е приложима към всеки тип дескриптор, например към обикновен програмен канал. Всички резултати трябва да се проверяват, а отворените дескриптори да се затварят и при грешка.

\section{Типови задачи}\label{s:11.5}

\subsection{Извличане на интервали от двоичен файл}\label{s:11.5.1}

Файлът \texttt{f1} съдържа двойки $\langle x_i,y_i\rangle$ от \texttt{uint32\_t}, а \texttt{f2} --- масив от такива числа. За всяка двойка се проверява, че интервалът е в границите, позицията се задава на
\[
x_i\cdot\operatorname{sizeof}(\texttt{uint32\_t}),
\]
и се прехвърлят точно $y_i$ числа към \texttt{f3} на блокове. Може да се използва \texttt{lseek} и \texttt{read} или позиционно четене. Не трябва да се обхождат пропуснатите части на \texttt{f2}; времето е пропорционално на броя двойки и размера на изхода. Проверяват се препълване при умножението, частично четене и преждевременен край на файла.

\subsection{Сортиране на файл от 8-битови числа}\label{s:11.5.2}

Понеже има само 256 възможни стойности, се използва броене:

\begin{enumerate}
\item входът се чете последователно на блокове;
\item за всеки байт се увеличава един от 256 брояча с достатъчно широк тип;
\item в изхода се записват последователно съответният брой нули, единици и т.н. до 255.
\end{enumerate}

Така времето е $O(n+256)=O(n)$, а допълнителната памет е $O(256)$ независимо от размер до 2 GB. Изходът също се записва на блокове, а не байт по байт.

\subsection{Прилагане на двоична кръпка}\label{s:11.5.3}

Първо \texttt{f1} се копира в \texttt{f2}. Всеки запис в \texttt{patch} е тройка
\[
(\texttt{uint16\_t offset},\texttt{uint8\_t old},\texttt{uint8\_t new}).
\]
Тройките се обработват последователно. За всяка се проверява, че отместването
съществува и текущият байт в работното копие \texttt{f2} е \texttt{old}; тогава
той се заменя с \texttt{new}. При несъответствие, непълна тройка или I/O грешка
програмата прекратява по контролиран начин.

\begin{supp}[Приета интерпретация: повторено отместване]
Официалната формулировка не казва какво става, когато две тройки сочат едно и
също отместване. Тук е приета следната интерпретация и тя трябва да се \emph{каже
на глас} при отговор: сравнението е винаги с \emph{текущото} състояние на
работното копие \texttt{f2}, а не с неизменния \texttt{f1}. Следствие: при
тройки $(o,a,b)$ и $(o,b,c)$ и двете се прилагат и байтът става $c$; при $(o,a,b)$
и $(o,a,c)$ втората се проваля, защото байтът вече е $b\ne a$. Другата
интерпретация --- сравнение винаги с \texttt{f1} --- е също защитима, но дава
различно поведение и трябва да бъде обявена преди решаването, а не след това.
\end{supp}

Двоичният формат поставя и въпроси за размера и подредбата на числата по байтове. Ако файловете се обменят между различни платформи, форматът трябва изрично да зададе byte order; не бива сляпо да се разчита на представянето на C структура в паметта.

\section{Пълна програма: извличане на интервали}\label{s:11.6}

Разделът по-горе описва решенията в схема. Схема и работеща програма са различни
неща: разликата е в обработката на грешки, частичния вход/изход и затварянето на
дескрипторите. По-долу задача 1 е написана докрай; другите две следват същата
рамка --- сменя се само тялото на цикъла.

\begin{verbatim}
/* Задача 1. f1 съдържа двойки <x, y> от uint32_t;
   f2 е масив от uint32_t. За всяка двойка в f3 се
   записват y числа от f2, започвайки от индекс x.
   Компилиране: cc -std=c11 -O2 -o extract extract.c   */
#include <errno.h>
#include <fcntl.h>
#include <inttypes.h>
#include <stdio.h>
#include <unistd.h>

#define CHUNK 1024

/* read може да върне по-малко от поисканото;
   EINTR не е грешка, а прекъсване.               */
static ssize_t read_full(int fd, void* buf, size_t n) {
  size_t done = 0;
  while (done < n) {
    ssize_t r = read(fd, (char*)buf + done, n - done);
    if (r < 0) {
      if (errno == EINTR) continue;
      return -1;
    }
    if (r == 0) break;                    /* край на файла */
    done += (size_t)r;
  }
  return (ssize_t)done;
}

/* write също може да запише частично --- това е
   най-често пропусканото място.                 */
static int write_full(int fd, const void* buf, size_t n) {
  size_t done = 0;
  while (done < n) {
    ssize_t w = write(fd, (const char*)buf + done, n - done);
    if (w < 0) {
      if (errno == EINTR) continue;
      return -1;
    }
    done += (size_t)w;
  }
  return 0;
}

int main(int argc, char** argv) {
  if (argc != 4) {
    fprintf(stderr, "употреба: %s f1 f2 f3\n", argv[0]);
    return 2;
  }

  int f1 = open(argv[1], O_RDONLY);
  int f2 = open(argv[2], O_RDONLY);
  int f3 = open(argv[3], O_WRONLY | O_CREAT | O_TRUNC, 0644);
  if (f1 < 0 || f2 < 0 || f3 < 0) {
    perror("open");
    return 1;
  }

  /* колко числа има в f2 */
  off_t size2 = lseek(f2, 0, SEEK_END);
  if (size2 < 0) {
    perror("lseek");
    return 1;
  }
  uint64_t count2 = (uint64_t)size2 / sizeof(uint32_t);

  int status = 0;
  uint32_t pair[2];
  ssize_t got;

  const ssize_t PAIR = (ssize_t)sizeof pair;

  while ((got = read_full(f1, pair, sizeof pair)) == PAIR) {
    uint64_t from = pair[0];
    uint64_t len  = pair[1];

    /* Проверката е написана така, че from + len да не прелее. */
    if (from > count2 || len > count2 - from) {
      fprintf(stderr, "интервал извън f2: [%" PRIu64
              ", %" PRIu64 ")\n", from, from + len);
      status = 1;
      break;
    }

    /* Позиционираме се; пропуснатите части на f2
       изобщо не се четат. */
    if (lseek(f2, (off_t)(from * sizeof(uint32_t)), SEEK_SET) < 0) {
      perror("lseek");
      status = 1;
      break;
    }

    uint32_t buf[CHUNK];
    while (len > 0) {
      size_t take  = len < CHUNK ? (size_t)len : CHUNK;
      size_t bytes = take * sizeof(uint32_t);
      if (read_full(f2, buf, bytes) != (ssize_t)bytes) {
        fprintf(stderr, "непълно четене от f2\n");
        status = 1;
        break;
      }
      if (write_full(f3, buf, bytes) != 0) {
        perror("write");
        status = 1;
        break;
      }
      len -= take;
    }
    if (status != 0) break;
  }

  if (got < 0) {
    perror("read");
    status = 1;
  } else if (status == 0 && got != 0) {
    fprintf(stderr, "непълна двойка в края на f1\n");
    status = 1;
  }

  if (close(f1) < 0 || close(f2) < 0 || close(f3) < 0) {
    /* грешка при close е загуба на данни */
    perror("close");
    status = 1;
  }
  return status;
}
\end{verbatim}

Четирите решения, които отличават работеща програма от схема:

\begin{itemize}
\item \texttt{read\_full} и \texttt{write\_full} --- частичен вход/изход и \texttt{EINTR}
се обработват на едно място, вместо да се повтарят във всеки цикъл;
\item проверката \verb|len > count2 - from| вместо \verb|from + len > count2| ---
втората прелива при големи стойности и пропуска невалиден интервал;
\item \texttt{lseek} вместо четене на пропуснатите части --- времето остава
пропорционално на броя двойки плюс размера на изхода, а не на размера на \texttt{f2};
\item резултатът на \texttt{close} се проверява: при отложен запис грешката се
проявява точно там.
\end{itemize}

\begin{supp}[Обхват]
Другите две задачи (\S11.5.2 и \S11.5.3) са дадени като схеми, не като пълни
програми. За сортирането на байтове това е достатъчно --- схемата от три реда е
целият алгоритъм и рамката по-горе се пренася дословно. За кръпката пълната
програма изисква и разрешаването на въпроса с повторените отмествания, обявен
по-горе.
\end{supp}

\section{Изпитен фокус}\label{s:11.7}

\begin{itemize}
\item Файловата система като единно именувано дърво; обекти, метаданни, права и точки на монтиране.
\item Ролята на superblock, inode, записите в директориите и преките/косвените блокове при ext2; журналът при ext3.
\item Кеш, read-ahead, отложен запис, алгоритъм на асансьора, RAID1 и RAID5; защо RAID не е backup.
\item Семантика и проверка на резултатите на \texttt{open}, \texttt{close}, \texttt{read}, \texttt{write} и \texttt{lseek}, включително частичен вход/изход.
\item Решение на трите типови задачи с време, пропорционално на реално обработваните данни.
\item Пълна програма за задача 1 (\S11.6): цикли за частичен \texttt{read}/\texttt{write}, обработка на \texttt{EINTR}, проверка без препълване, \texttt{lseek} вместо обхождане и проверка на \texttt{close}.
\item Обявяване на приетата интерпретация при повторено отместване в задачата за кръпката, преди решаването.
\end{itemize}


\chapter{Управление на процеси и междупроцесни комуникации.}
% AUTO-GENERATED by scripts/build_topics.py — do not edit.
% Edit topics/bodies/topic_12.tex or topics/preamble.tex instead.

\begin{konspekt}
\kreq{Примитиви}{създаване на процес, изпълнение на програма, завършване, синхронизация със завършването на процеса-син --- \kref{12.1}.}
\kreq{Права и идентичности}{потребителски идентификатори на процес, групи процеси и сесия --- \kref{12.2}.}
\kreq{Междупроцесна комуникация}{сигнали и програмни канали --- \kref{12.3}.}
\kreq{Междупроцесна комуникация}{IPC пакетът на UNIX System V --- обща памет, семафори, съобщения --- \kref{12.4}.}
\kreq{Схема}{родител--дете --- \kr{12.5.1}, като пълна програма в \kref{12.6}.}
\kreq{Схема}{производител--потребител --- \kref{12.5.2}.}
\kreq{Типове задачи}{задачи, съответни на съдържанието на въпроса.}
\kstatus{схемата родител--дете е дадена като пълна програма; производител--потребител остава схема. Първичните източници липсват в корпуса --- главата подлежи на проверка по литературата.}
\kreq{Литература}{[21], [46].}
\end{konspekt}

\begin{supp}[Статус на източниците]
В наличния корпус липсват първичните материали за този въпрос. Главата следва подробната официална анотация и е сверена с предоставения сравнителен конспект. Преди окончателно академично издание тя трябва да се провери по посочената литература [21] и [46].
\end{supp}

\section{Процес и основни примитиви}\label{s:12.1}

\begin{defn}[Процес]
Процесът е изпълняващ се екземпляр на програма заедно с адресното му пространство, регистрите, отворените файлови дескриптори, идентичността и останалото състояние, което ядрото управлява.
\end{defn}

В Unix създаването и зареждането на програма са отделни операции.

\begin{defn}[Системни примитиви за процеси]
\begin{description}
\item[\texttt{fork}.] Създава процес-син като логическо копие на родителя. И двата процеса продължават след извикването: в сина резултатът е $0$, а в родителя --- PID на сина. При грешка родителят получава $-1$. Реализацията обичайно използва copy-on-write.
\item[\texttt{exec}.] Семейството \texttt{exec} заменя програмния образ на текущия процес с нова програма. При успех не се връща; PID остава същият, а запазването на дескриптори зависи от флага close-on-exec.
\item[\texttt{exit} и \texttt{\_exit}.] Завършват процеса със статус. \texttt{exit} изпълнява регистрираните потребителски обработчици и flush на C потоците, докато \texttt{\_exit} пропуска тази потребителска финализация и предава завършването на ядрото. Ядрото все пак трябва да освободи ресурсите и да затвори дескрипторите, което може да изисква време. След неуспешен \texttt{exec} в дете често се използва \texttt{\_exit}.
\item[\texttt{wait}/\texttt{waitpid}.] Родителят изчаква промяна на състоянието на дете и получава статуса му. Завършило, но неприбрано дете е zombie. Ако родителят не иска да блокира, \texttt{waitpid} позволява подходящи опции.
\end{description}
\end{defn}

Класическият шаблон е: \texttt{fork}; детето подготвя дескрипторите и извиква \texttt{exec}; родителят продължава паралелно или извиква \texttt{waitpid}. Кодът винаги различава трите резултата на \texttt{fork}.

\section{Права, групи и сесии}\label{s:12.2}

\begin{defn}[Идентичност на процес]
Всеки процес има реални и ефективни потребителски и групови идентификатори. Реалните UID/GID описват произхода, а ефективните обичайно участват в проверките за достъп. Допълнителни групи и механизми като set-user-ID променят конкретните правила, затова не е достатъчно да се гледа само PID.
\end{defn}

\begin{defn}[Процесна група]
Процесната група е множество процеси с общ PGID, към което терминалът или друг процес може да адресира сигнал като към едно цяло.
\end{defn}

\begin{defn}[Сесия]
Сесията е множество процесни групи. Обичайно обвивката създава отделна група за всяка задача и използва сесията и управляващия терминал за foreground/background управление.
\end{defn}

Всяка процесна група принадлежи на точно една сесия. \texttt{setpgid} създава или променя принадлежността към процесна група при допустимите ограничения, а \texttt{getpgrp} връща текущата група. \texttt{setsid} създава нова сесия, нова процесна група и прави извикващия процес техен лидер; \texttt{getsid} връща идентификатора на сесията. Тези операции стоят в основата на управлението на задачи от командната обвивка.

\section{Сигнали и програмни канали}\label{s:12.3}

\begin{defn}[Сигнал]
Сигналът е асинхронно известие за събитие. Процесът може да установи обработчик, да остави стандартното действие или за някои сигнали да ги игнорира. Сигналите \texttt{SIGKILL} и \texttt{SIGSTOP} не могат да бъдат прихващани, блокирани или игнорирани. Маска блокира временно доставянето на други сигнали; обработчикът трябва да използва само операции, безопасни в сигнален контекст. \texttt{SIGCHLD} уведомява родителя за промяна в състоянието на дете. При обичайното поведение родителят прибира статуса чрез \texttt{wait}/\texttt{waitpid}; при изрично игнориране на \texttt{SIGCHLD} или флаг \texttt{SA\_NOCLDWAIT} завършилото дете не се оставя като чакащ zombie процес.
\end{defn}

\begin{defn}[Програмен канал]
Програмният канал (pipe) е еднопосочен байтов поток с край за четене и край за запис. \texttt{pipe} връща два файлови дескриптора, които могат да се наследят през \texttt{fork}.
\end{defn}

След \texttt{fork} всеки процес затваря неизползваните краища. Четенето връща край на файл едва когато всички дескриптори за запис са затворени; забравен записващ край може да причини безкрайно блокиране. Записите до определена малка граница са атомарни, но pipe не пази структура на произволно големи съобщения. Именуваният канал FIFO позволява несвързани процеси да отворят общ обект от файловата система.

\section{UNIX System V IPC}\label{s:12.4}

System V IPC идентифицира обекти чрез ключове и системни идентификатори; правата им са отделни от обикновените файлови дескриптори.

\subsection{Обща памет}\label{s:12.4.1}

\begin{defn}[Обща памет]
Общата памет позволява няколко процеса да прикачат един сегмент към адресните си пространства. Типичните операции са създаване/намиране, прикачване, откачване и управление чрез \texttt{shmget}, \texttt{shmat}, \texttt{shmdt} и \texttt{shmctl}. След прикачването обменът е бърз, защото няма копиране през ядрото при всяко обръщение, но общата памет сама по себе си не осигурява взаимно изключване или ред на достъп.
\end{defn}

\subsection{Семафори}\label{s:12.4.2}

\begin{defn}[Семафор]
Семафорът е целочислен синхронизационен обект, променян атомарно. Операцията $P$ (wait/down) изчаква положителна стойност и я намалява, а $V$ (signal/up) я увеличава и може да събуди чакащ процес. Двоичен семафор може да пази критична секция; броящ семафор представя наличен брой еднотипни ресурси. При System V \texttt{semget} създава или намира множество от семафори, \texttt{semop} прилага атомарно вектор от операции, а \texttt{semctl} управлява множеството. Неправилен ред на придобиване може да доведе до взаимно блокиране.
\end{defn}

\subsection{Опашки от съобщения}\label{s:12.4.3}

\begin{defn}[Опашка от съобщения]
Опашката пази отделни съобщения с тип и съдържание. \texttt{msgget} създава или намира опашка, \texttt{msgsnd} изпраща, \texttt{msgrcv} получава, а \texttt{msgctl} управлява жизнения цикъл и параметрите й. Получателят може да избере следващото съобщение или съобщение от определен тип. За разлика от pipe границите между съобщенията се запазват. Капацитетът е ограничен и операциите могат да блокират или да се изпълнят неблокиращо.
\end{defn}

System V обектите могат да останат в ядрото след приключване на процеса, докато не бъдат изрично премахнати. Затова програмата трябва да има ясен собственик и път за почистване.

\section{Синхронизационни схеми}\label{s:12.5}

\subsection{Родител и дете}\label{s:12.5.1}

За изпълнение на програма и получаване на резултата родителят създава дете с \texttt{fork}; детето извиква \texttt{exec}; родителят извиква \texttt{waitpid} и проверява дали детето е завършило нормално и какъв е статусът му. Ако между тях има pipe, краищата се създават преди \texttt{fork}, пренасочват се при нужда с \texttt{dup2}, а всички излишни копия се затварят преди чакането.

\subsection{Производител--потребител}\label{s:12.5.2}

Нека общ кръгов буфер има $N$ места. Използват се броящи семафори
\[
empty=N,\qquad full=0
\]
и двоичен семафор
\[
mutex=1.
\]

Производителят изпълнява $P(empty)$, $P(mutex)$, добавя елемент, после $V(mutex)$ и $V(full)$. Потребителят изпълнява $P(full)$, $P(mutex)$, премахва елемент, после $V(mutex)$ и $V(empty)$. Броящите семафори предотвратяват препълване и четене от празен буфер, а \texttt{mutex} пази индексите и съдържанието. Изчакването за свободно/заето място трябва да е преди заключването на \texttt{mutex}; обратният ред може да задържи критичната секция и да блокира другата страна.

\section{Пълна програма: родител, дете и програмен канал}\label{s:12.6}

Схемата в \S12.5.1 казва \emph{кои} примитиви се използват. Работещата програма
трябва да реши и четирите въпроса, които схемата премълчава: кой край на канала
кой затваря, защо детето излиза с \texttt{\_exit}, какво се прави при
\texttt{EINTR} и как се чете статусът.

\begin{verbatim}
/* Родителят стартира външна програма в дете и чете
   изхода й през канал.
   Компилиране: cc -std=c11 -O2 -o run run.c        */
#include <errno.h>
#include <stdio.h>
#include <sys/wait.h>
#include <unistd.h>

int main(void) {
  int fd[2];                 /* [0] за четене, [1] за писане */
  if (pipe(fd) < 0) {
    perror("pipe");
    return 1;
  }

  pid_t pid = fork();
  if (pid < 0) {
    perror("fork");
    return 1;
  }

  if (pid == 0) {              /* ------- дете ------- */
    /* четящият край не му трябва */
    if (close(fd[0]) < 0) _exit(127);
    if (dup2(fd[1], STDOUT_FILENO) < 0) _exit(127);
    /* дубликатът вече е stdout */
    if (close(fd[1]) < 0) _exit(127);
    execlp("ls", "ls", "-l", (char*)NULL);
    _exit(127);              /* стига се само при неуспех */
  }

  /* ---------- родител ---------- */
  /* без това read никога не вижда край на файл */
  if (close(fd[1]) < 0) {
    perror("close");
    return 1;
  }

  int status = 0;
  char buf[4096];
  for (;;) {
    ssize_t n = read(fd[0], buf, sizeof buf);
    /* детето е затворило канала */
    if (n == 0) break;
    if (n < 0) {
      if (errno == EINTR) continue;
      perror("read");
      status = 1;
      break;
    }
    /* изходът също може да се провали: пълен диск, счупена
       тръба, дескриптор само за четене */
    if (fwrite(buf, 1, (size_t)n, stdout) != (size_t)n) {
      perror("fwrite");
      status = 1;
      break;
    }
  }
  close(fd[0]);

  /* буферът се изпразва тук, докато има кой да съобщи за
     грешката; при изход от main е вече късно */
  if (fflush(stdout) != 0 || ferror(stdout)) {
    perror("stdout");
    status = 1;
  }

  int wstatus;
  while (waitpid(pid, &wstatus, 0) < 0) {
    if (errno != EINTR) {
      perror("waitpid");
      return 1;
    }
  }

  if (WIFSIGNALED(wstatus)) {
    fprintf(stderr, "детето е прекратено от сигнал %d\n",
            WTERMSIG(wstatus));
    return 1;
  }
  if (!WIFEXITED(wstatus)) return 1;

  fprintf(stderr, "детето завърши с код %d\n",
          WEXITSTATUS(wstatus));
  /* успех само ако и детето, и изходът са наред */
  return (WEXITSTATUS(wstatus) == 0 && status == 0) ? 0 : 1;
}
\end{verbatim}

Решенията, които се питат на изпита:

\begin{itemize}
\item \textbf{Затварянето на неизползваните краища не е чистота, а коректност.}
Каналът дава край на файл при четене чак когато \emph{всички} записващи
дескриптори са затворени. Ако родителят задържи \verb|fd[1]|, неговият
\texttt{read} блокира вечно, макар детето отдавна да е завършило.
\item \textbf{Детето излиза с \texttt{\_exit}, не с \texttt{exit}.} След
\texttt{fork} детето има копие на буферите на стандартната библиотека;
\texttt{exit} би ги изчистило втори път и същият текст би се отпечатал два пъти.
\texttt{\_exit} завършва процеса без тези действия.
\item \textbf{\texttt{waitpid} се повтаря при \texttt{EINTR}.} Пристигнал сигнал
прекъсва чакането; това не е грешка и не означава, че детето е завършило.
\item \textbf{Без \texttt{waitpid} детето остава \emph{zombie}}: записът в
таблицата на процесите се пази, докато родителят не прочете статуса му.
\item \texttt{execlp} се връща само при неуспех --- затова редът след него е
изход с код, а не продължение на програмата.
\item \textbf{Записът към изхода се проверява точно както четенето.}
\texttt{fwrite} връща броя записани елементи, а не признак за успех; при пълен
диск, затворена тръба надолу по конвейера или дескриптор, отворен само за
четене, тя записва по-малко и програмата трябва да го забележи. Пропуснато ли
се, програмата връща $0$, след като е загубила част от изхода --- най-лошият
възможен изход, защото извикващият скрипт продължава напред.
\item \textbf{\texttt{fflush} се извиква изрично, преди да се докладва успех.}
Стандартната библиотека буферира; ако изпразването се остави на изхода от
\texttt{main}, грешката настъпва след последната възможност да бъде съобщена и
кодът за връщане е вече определен. \texttt{ferror} улавя и грешка, отбелязана в
по-ранно частично записване.
\item \textbf{Кодът за връщане обединява двете условия.} Успех се докладва само
когато детето е завършило с $0$ \emph{и} целият изход е стигнал до
местоназначението си.
\end{itemize}

\begin{supp}[Обхват]
Схемата производител--потребител (\S12.5.2) е дадена със семафорни операции, но
не като пълна програма: тя изисква и създаване на System V обектите, права,
обработка на \texttt{EINTR} при \texttt{semop} и изрично премахване на обектите
при изход, защото те надживяват процеса. Това остава да се напише.
\end{supp}

\section{Изпитен фокус}\label{s:12.7}

\begin{itemize}
\item Разлика между процес и програма; семантика и резултати на \texttt{fork}, \texttt{exec}, \texttt{exit}/\texttt{\_exit}, \texttt{wait}/\texttt{waitpid}; zombie процес.
\item Реални и ефективни UID/GID, процесни групи, сесии и управляващ терминал.
\item Сигнали и маски; специалният статус на \texttt{SIGKILL} и \texttt{SIGSTOP}; \texttt{SIGCHLD}.
\item Pipe/FIFO, наследяване и правилно затваряне на краищата.
\item Обща памет, семафори и опашки от съобщения в System V IPC, включително основните семейства \texttt{shm*}, \texttt{sem*} и \texttt{msg*}, жизнен цикъл и права.
\item Синхронизация родител--дете и схема производител--потребител с \texttt{empty}, \texttt{full} и \texttt{mutex}.
\item Пълна програма родител--дете през канал (\S12.6): ред на затваряне на краищата, \texttt{dup2}, \texttt{\_exit} вместо \texttt{exit} в детето, повтаряне при \texttt{EINTR} и четене на статуса с \texttt{WIFEXITED}/\texttt{WEXITSTATUS}.
\item Защо незатворен записващ край блокира \texttt{read} завинаги и защо липсващият \texttt{waitpid} оставя zombie процес.
\end{itemize}


\chapter{Компютърни мрежи и протоколи. OSI модел. Маршрутизация. IPv4, IPv6, TCP, DNS.}
% AUTO-GENERATED by scripts/build_topics.py — do not edit.
% Edit topics/bodies/topic_13.tex or topics/preamble.tex instead.
%% Placeholder. Replace with the written chapter.
\textit{Източници:} PDF \texttt{13.pdf}

\subsection*{Анотация от конспекта}
Компютърни мрежи и протоколи – OSI модел. Маршрутизация. Протоколи IPv4, IPv6, TCP, DNS. OSI модел – най-обща характеристика на нивата, съпоставяне с модела TCP/IP. Разпределена маршрутизация – алгоритми с дистантен вектор и следене на състоянието на линията. IPv4 адресация – класова и безкласова. Основни характеристики на протокол IPv6. TCP – процедура на трикратно договаряне. Основни характеристики на протоколи DNS (резолвинг на имената по IPv4 и IPv6). Литература:[3], [35], [51].


\chapter{Процедурно програмиране – основни конструкции.}
%% Drafted from the mapped corpus by scripts/draft_topics.py.
% Model: gpt-5.6-luna; source characters: 24250.
\section{Процедурно програмиране}

\begin{defn}[Процедурно програмиране]
Процедурното програмиране е програмна парадигма, при която поведението на програмата се реализира чрез процедури, наричани още функции. Процедурите представляват самостоятелни части от програмата, които изпълняват определена задача и могат да бъдат извиквани многократно. Една процедура може да извиква друга процедура, поради което цялата програма се представя като поредица от стъпки, организирани в йерархия от извиквания.
\end{defn}

Процедурното програмиране се класифицира като императивно програмиране. При него програмата съдържа директни команди, които определят какво да бъде изпълнено и в какъв ред. Основни понятия са променливата, операторът, блокът, функцията и областта на действие.

\begin{defn}[Блок]
Блок е последователност от оператори, оградена с фигурни скоби $\{\}$.
Блокът се разглежда като една съставна конструкция и обикновено определя областта на действие на локално дефинираните в него променливи.
\end{defn}

\begin{defn}[Област на действие]
Областта на действие на даден идентификатор е частта от програмния текст, в която този идентификатор може да бъде използван.
\end{defn}

Понятията блок и област на действие са съществени за структурирането на програмата. Чрез тях се ограничава достъпът до данните и се отделят отделните части на алгоритъма.

\subsection{Принципи на структурното програмиране}

Основната идея на структурното програмиране е изграждането на ясна и разбираема структура на програмата. Програмата се разделя на отделни същности, всяка от които има ясно определена задача и граници. Пример за такава същност е функцията.

Основен метод за постигане на ясна структура е декомпозицията.

\begin{defn}[Декомпозиция]
Декомпозицията е разделяне на общата задача на по-малки подзадачи, всяка от които може да бъде реализирана чрез отделна процедура или функция.
\end{defn}

При декомпозиция сложната задача се представя чрез йерархия от по-прости задачи. Функциите, които реализират тези задачи, могат да бъдат извиквани от други функции. Така програмата се организира като множество от взаимосвързани, но логически обособени компоненти.

\begin{keydefn}[Модулност]
Модулността е принцип, според който функционалността на програмата се разпределя между отделни, взаимозаменяеми модули. Всеки модул трябва да съдържа всичко необходимо за изпълнението на един ясно определен аспект от програмата.
\end{keydefn}

Модулността има две основни цели:

\begin{itemize}
\item всяка част от програмата да има ясно определена отговорност;
\item отделните части да могат да бъдат използвани, проверявани и при необходимост заменяни независимо една от друга.
\end{itemize}

Добре структурираната процедурна програма се характеризира с последователност от ясно разграничени действия, функции с определено предназначение и ограничена област на действие на данните.

\section{Управление на изчислителния процес}

\begin{defn}[Изчислителен процес]
Изчислителният процес представлява изпълнение на поредица от стъпки в определен ред. Той има начална точка, последователност от междинни действия и крайна точка.
\end{defn}

В C++ изпълнението на програмата започва от функцията \texttt{main}. Функцията \texttt{main} може да извиква други функции. При нормално изпълнение програмата завършва, когато приключи изпълнението на \texttt{main}.

Една от специалните форми на дефиниция на \texttt{main} е:

\begin{verbatim}
int main(int argc, char* argv[])
\end{verbatim}

Параметърът \texttt{argc} указва броя на елементите в масива \texttt{argv}, а \texttt{argv} съдържа параметрите, подадени при стартиране на програмата от терминала. Възможна е и форма без параметри:

\begin{verbatim}
int main()
\end{verbatim}

Управлението на изчислителния процес се осъществява чрез управляващи конструкции. Те променят нормалния линеен ред на изпълнение, като позволяват:

\begin{itemize}
\item избор между различни действия;
\item многократно изпълнение на дадена последователност от оператори;
\item прекратяване на част от текущото изпълнение.
\end{itemize}

Основните управляващи конструкции са условните оператори и операторите за цикъл.

\section{Условни оператори}

Условният оператор позволява даден оператор или блок да бъде изпълнен само когато определено условие е вярно.

\subsection{Оператор \texttt{if}}

Общият синтаксис е:

\begin{verbatim}
if (условие_0)
    оператор_0
else if (условие_1)
    оператор_1
else
    оператор_else
\end{verbatim}

Частите \texttt{else if} могат да се повторят произволен брой пъти, а частта \texttt{else} е незадължителна и може да се срещне най-много веднъж.

Когато операторът се състои от повече от една команда, командите се обединяват в блок:

\begin{verbatim}
if (условие) {
    оператор_1
    оператор_2
}
\end{verbatim}

Изпълнението протича по следния начин:

\begin{enumerate}
\item Оценява се първото условие.
\item Ако то има стойност \texttt{true}, изпълнява се съответният оператор или блок.
\item След това останалите части на условния оператор се прескачат.
\item Ако условието е \texttt{false}, се проверява следващото условие в \texttt{else if}.
\item При първото условие със стойност \texttt{true} се изпълнява съответният оператор и проверката приключва.
\item Ако всички условия са \texttt{false} и съществува \texttt{else}, се изпълнява операторът след \texttt{else}.
\end{enumerate}

\begin{example}
Следната програма проверява знака на число:

\begin{verbatim}
int x;
cin >> x;

if (x > 0) {
    cout << "positive";
} else if (x < 0) {
    cout << "negative";
} else {
    cout << "zero";
}
\end{verbatim}

При изпълнението се извършва само един от трите блока.
\end{example}

Ако няма част \texttt{else} и всички условия са неверни, не се изпълнява нито един от операторите в конструкцията.

\subsection{Условна, или тернарна, операция}

Условната операция е триместна операция със синтаксис:

\[
\langle\text{булев израз}\rangle
\mathrel{?}
\langle\text{израз}_1\rangle
:
\langle\text{израз}_2\rangle
\]

Първо се оценява булевият израз. Ако резултатът е \texttt{true}, се оценява и връща резултатът на $\langle\text{израз}_1\rangle$. Ако резултатът е \texttt{false}, се оценява и връща резултатът на $\langle\text{израз}_2\rangle$.

Пример:

\begin{verbatim}
int x = (y < 2) ? y + 1 : y - 2;
\end{verbatim}

Ако $y=1$, на \texttt{x} се присвоява стойност $2$. Ако $y=2$, на \texttt{x} се присвоява стойност $0$.

Тернарната операция е подходяща, когато трябва да се избере една от две стойности. Тя връща резултат и затова може да се използва като част от по-голям израз.

\subsection{Оператор за многозначен избор \texttt{switch}}

Операторът \texttt{switch} служи за проверка на стойността на израз спрямо множество зададени стойности.

Синтаксисът има вида:

\begin{verbatim}
switch (израз) {
    case константен_израз_1:
        оператори_1
        break;

    case константен_израз_2:
        оператори_2
        break;

    default:
        оператори_default
}
\end{verbatim}

Оценяваната стойност на израза се сравнява със стойностите след етикетите \texttt{case}. Когато бъде намерено съвпадение, изпълнението започва от съответния \texttt{case}.

\begin{example}
\begin{verbatim}
int day;
cin >> day;

switch (day) {
    case 1:
        cout << "Monday";
        break;
    case 2:
        cout << "Tuesday";
        break;
    case 3:
        cout << "Wednesday";
        break;
    default:
        cout << "Unknown day";
}
\end{verbatim}
\end{example}

Ако не е намерен съответстващ \texttt{case} и съществува етикет \texttt{default}, се изпълняват операторите след \texttt{default}. Ако няма съвпадение и няма \texttt{default}, операторът приключва без изпълнение на case-блок.

След намиране на съответстващ \texttt{case} изпълнението продължава последователно надолу в блока на \texttt{switch}. Самото достигане на следващ етикет \texttt{case} не прекратява изпълнението. За прекратяване се използва операторът \texttt{break}.

\begin{supp}[Бележка]
Ако в съответстващия \texttt{case} липсва \texttt{break}, се осъществява последователно преминаване към следващите case-блокове. Това поведение е част от описаната семантика на оператора \texttt{switch}, а не грешка в синтаксиса.
\end{supp}

Изпълнението на \texttt{switch} приключва при едно от следните събития:

\begin{itemize}
\item достигнат е краят на блока на \texttt{switch};
\item изпълнен е оператор за управление на потока като \texttt{break} или \texttt{return};
\item нормалното изпълнение е прекъснато по друга причина.
\end{itemize}

Ползата от \texttt{switch} спрямо множество \texttt{if--else if--else} конструкции е, че управляващият израз се оценява веднъж и след това се сравнява с различните константни стойности.

\section{Оператори за цикъл}

Цикълът е конструкция за управление на изчислителния процес, чрез която даден оператор или блок се изпълнява многократно. Повторението продължава, докато условието има стойност \texttt{true}. Когато условието стане \texttt{false}, цикълът приключва.

В C++ основните оператори за цикъл са \texttt{for}, \texttt{while} и \texttt{do while}.

\subsection{Оператор \texttt{for}}

Синтаксисът е:

\begin{verbatim}
for (инициализация; условие; корекция)
    тяло
\end{verbatim}

При използване на блок:

\begin{verbatim}
for (инициализация; условие; корекция) {
    оператори
}
\end{verbatim}

Изпълнението протича в следния ред:

\begin{enumerate}
\item Изпълнява се инициализацията. Тя се изпълнява само веднъж при влизане в цикъла.
\item Оценява се условието.
\item Ако условието е \texttt{true}, изпълнява се тялото на цикъла.
\item След тялото се изпълнява корекцията.
\item Отново се оценява условието.
\item Процесът се повтаря, докато условието стане \texttt{false}.
\end{enumerate}

Най-често инициализацията служи за дефиниране и начално задаване на стойност на променлива, а корекцията увеличава или намалява тази променлива.

\begin{example}
\begin{verbatim}
for (int i = 0; i < 5; i++) {
    cout << i << " ";
}
\end{verbatim}

Изпълняват се стойностите $0$, $1$, $2$, $3$ и $4$. След като \texttt{i} стане $5$, условието \texttt{i < 5} е невярно и цикълът приключва.
\end{example}

Променлива, дефинирана в инициализационната част на \texttt{for}, има област на действие в цикъла, тоест може да бъде използвана до края на блока на цикъла.

\subsection{Оператор \texttt{while}}

Синтаксисът е:

\begin{verbatim}
while (условие)
    тяло
\end{verbatim}

или:

\begin{verbatim}
while (условие) {
    оператори
}
\end{verbatim}

При изпълнение първо се оценява условието. Ако то е \texttt{true}, се изпълнява тялото и след това отново се проверява условието. Ако условието е \texttt{false}, тялото не се изпълнява и цикълът приключва.

\begin{example}
\begin{verbatim}
int i = 0;

while (i < 5) {
    cout << i << " ";
    i++;
}
\end{verbatim}
\end{example}

При \texttt{while} е необходимо тялото на цикъла или друга част от изпълнението да води до промяна, която в даден момент да направи условието невярно. В противен случай повторението може да продължи без край.

\subsection{Оператор \texttt{do while}}

Операторът \texttt{do while} е сходен с \texttt{while}, но проверката на условието се извършва след изпълнението на тялото.

Синтаксисът е:

\begin{verbatim}
do {
    оператори
} while (условие);
\end{verbatim}

Редът на изпълнение е:

\begin{enumerate}
\item изпълнява се тялото на цикъла;
\item оценява се условието;
\item ако условието е \texttt{true}, тялото се изпълнява отново;
\item ако условието е \texttt{false}, цикълът приключва.
\end{enumerate}

Следователно тялото на \texttt{do while} се изпълнява поне веднъж, преди да бъде направена първата проверка.

\begin{example}
\begin{verbatim}
int x;

do {
    cin >> x;
} while (x < 0);
\end{verbatim}

Тялото се изпълнява, след което се проверява дали въведеното число е отрицателно.
\end{example}

\section{Променливи и присвояване}

\begin{defn}[Променлива]
Променливата е именувана област в паметта, свързана с име, тип, адрес и стойност.
\end{defn}

Името на променливата е идентификатор, чрез който тя се използва в програмния текст. Адресът указва мястото в паметта, отредено за променливата. Типът определя как компилаторът интерпретира съдържанието на това място и какви операции могат да се извършват със стойността. Стойността е текущото съдържание на променливата.

\subsection{Дефиниране и инициализиране}

Общият синтаксис за дефиниране на променливи е:

\[
\langle\text{тип}\rangle\ \langle\text{идентификатор}\rangle
[\mathrel{=}\langle\text{израз}\rangle]
\{
,\ \langle\text{идентификатор}\rangle
[\mathrel{=}\langle\text{израз}\rangle]
\}\ ;
\]

Например:

\begin{verbatim}
int a;
double x = 3.5;
int p = 1, q = 2;
\end{verbatim}

Дефинирането създава променлива от даден тип. Променливата е инициализирана, когато при създаването й е зададена начална стойност.

В C++ могат да се използват например следните форми:

\begin{verbatim}
int x = 5;
int y(5);
int z{5};
\end{verbatim}

При използване на присвояване с \texttt{=} се използва копираща инициализация. При формите с кръгли или фигурни скоби началната стойност се задава непосредствено при инициализацията.

\begin{supp}[Бележка]
Източниците описват неинициализираната локална променлива като съдържаща непредвидима стойност. Практическото следствие е, че такава променлива не трябва да се използва, преди да й бъде зададена определена стойност.
\end{supp}

\subsection{Оператор за присвояване}

Синтаксисът на присвояването е:

\[
\langle\text{идентификатор}\rangle
=
\langle\text{израз}\rangle;
\]

Операторът оценява израза отдясно и записва получената стойност в променливата отляво.

\begin{verbatim}
int x;
x = 7;
x = x + 1;
\end{verbatim}

Резултатът от присвояването е стойността, записана в променливата. Поради това присвояванията могат да се влагат едно в друго.

\begin{defn}[lvalue и rvalue]
$lvalue$ е място в паметта, което има стойност и може да бъде променяно. Обикновена променлива е пример за $lvalue$.

$rvalue$ е стойност, която не представлява самостоятелно променяемо място в паметта. Константа, литералът и резултатът от пресмятане са примери за $rvalue$.
\end{defn}

При израза

\[
a=b=c=2
\]

присвояването се извършва отдясно наляво:

\[
a=(b=(c=2)).
\]

Следователно първо на \texttt{c} се присвоява $2$, полученият резултат се присвоява на \texttt{b}, а след това и на \texttt{a}.

\subsection{Типове на променливите}

Типът указва на компилатора как да интерпретира съдържанието на паметта и носи семантична информация. Променливите от един и същи тип имат сходни характеристики, включително размер и допустими операции.

Типовете могат да бъдат разглеждани като:

\begin{itemize}
\item скаларни типове;
\item съставни типове.
\end{itemize}

Към скаларните типове спадат булевият, целочисленият, символният, изброимият тип, типовете за числа с плаваща запетая, указателите и референциите.

Към съставните типове спадат масивите, структурите и класовете.

\subsection{Локални и глобални променливи}

Областта на действие определя къде в програмния текст променливата може да бъде достъпвана.

\begin{defn}[Локална променлива]
Локална е променлива, дефинирана в блок или във функция. Областта й на действие започва от мястото на дефинирането и продължава до края на блока, в който е дефинирана.
\end{defn}

\begin{example}
\begin{verbatim}
int main() {
    int x = 1;

    if (x > 0) {
        int y = 2;
        cout << y;
    }

    // y вече не е достъпна тук
}
\end{verbatim}
\end{example}

\begin{defn}[Глобална променлива]
Глобална е променлива, дефинирана извън тялото на функция. Областта й на действие е от мястото на дефинирането до края на файла.
\end{defn}

Локалната област на действие ограничава използването на променливата до частта от програмата, в която тя е необходима. Глобалната променлива може да бъде достъпвана от по-голяма част от програмния файл.

\section{Функции и процедури}

\begin{defn}[Функция]
Функцията е самостоятелен фрагмент от програмата, съдържащ поредица от команди, предназначени за изпълнение на определена задача и позволяващи многократно използване.
\end{defn}

Функцията извършва пресмятане и връща резултат. Процедурата изпълнява поредица от оператори, без задължително да връща резултат. В C++ и двете понятия се реализират чрез функции. Функция, която не връща резултат, има резултатен тип \texttt{void}.

\subsection{Дефиниране на функция}

Общият синтаксис е:

\begin{verbatim}
тип_резултат име(формални_параметри) {
    тяло
}
\end{verbatim}

Например:

\begin{verbatim}
int square(int x) {
    return x * x;
}
\end{verbatim}

В дефиницията се съдържат:

\begin{itemize}
\item резултатен тип;
\item име на функцията;
\item списък от формални параметри;
\item тяло на функцията.
\end{itemize}

Резултатният тип може да бъде \texttt{void} или друг допустим тип, но в съвременния C++ трябва да бъде зададен изрично. Пропуснат резултатен тип не означава подразбиращ се \texttt{int}; такава стара конвенция е съществувала в ранни версии на C, но не е правило на C++.

Функция без параметри може да бъде представена с празен списък или с \texttt{void}:

\begin{verbatim}
void printMessage() {
    cout << "Hello";
}
\end{verbatim}

Формалните параметри имат вида:

\[
\langle\text{тип}\rangle\ \langle\text{идентификатор}\rangle
\{
,\ \langle\text{тип}\rangle\ \langle\text{идентификатор}\rangle
\}.
\]

Например:

\begin{verbatim}
int add(int a, int b) {
    return a + b;
}
\end{verbatim}

\subsection{Декларация и извикване}

Функцията може да бъде само декларирана, без да се задава тялото й. Декларацията представлява обещание, че дефиниция на функцията ще бъде предоставена.

Примерна декларация е:

\begin{verbatim}
int add(int a, int b);
\end{verbatim}

Извикването на функция се записва чрез нейното име и фактическите параметри:

\begin{verbatim}
име(фактически_параметри)
\end{verbatim}

Например:

\begin{verbatim}
int result = add(2, 3);
\end{verbatim}

Фактическите параметри, наричани още аргументи, са изрази:

\[
\langle\text{израз}\rangle
\{
,\ \langle\text{израз}\rangle
\}.
\]

При извикването типът на всеки фактически параметър се съпоставя с типа на съответния формален параметър. Ако е необходимо, се извършва преобразуване на типовете.

\subsection{Връщане на резултат}

Операторът за връщане има вида:

\begin{verbatim}
return израз;
\end{verbatim}

Той връща стойността на израза като резултат от извикването на функцията и незабавно прекратява изпълнението на функцията.

\begin{example}
\begin{verbatim}
int maximum(int a, int b) {
    if (a > b) {
        return a;
    }
    return b;
}
\end{verbatim}
\end{example}

Типът на израза след \texttt{return} се съпоставя с резултатния тип на функцията. Ако е необходимо, се извършва преобразуване.

Функция с резултатен тип \texttt{void} може да изпълнява оператори, без да връща стойност. В този случай работата й приключва при достигане на края на тялото или чрез \texttt{return} без резултат.

\subsection{Стекова рамка при извикване}

При извикване на функция се заделя нова стекова рамка. В нея се пазят фактическите параметри, адресът за връщане и локалните променливи на функцията. Рамката се намира на върха на програмния стек.

Следователно всяко извикване има собствен контекст. Локалните променливи и параметрите на едно извикване са свързани с неговата стекова рамка. При приключване на функцията изпълнението се връща към адреса, от който тя е била извикана.

\subsection{Предаване на параметри}

\subsubsection{Предаване по стойност}

При предаване по стойност първо се пресмята стойността на фактическия параметър. В стековата рамка на функцията се създава копие на тази стойност.

\begin{verbatim}
void change(int x) {
    x = 10;
}

int a = 3;
change(a);
\end{verbatim}

Промяната на \texttt{x} остава локална за функцията. Стойността на \texttt{a} не се променя. При завършване на функцията копието и направените върху него промени изчезват.

\subsubsection{Предаване чрез референция}

Предаването чрез референция се използва, когато промените във формалния параметър трябва да се отразят върху фактическия параметър.

Синтаксисът на параметър-референция е:

\begin{verbatim}
тип& идентификатор
\end{verbatim}

Пример:

\begin{verbatim}
int add5(int& x) {
    x += 5;
    return x;
}

int a = 3;
cout << add5(a) << " " << a;
\end{verbatim}

Резултатът е:

\begin{verbatim}
8 8
\end{verbatim}

В този случай параметърът \texttt{x} се свързва с променливата \texttt{a}. Затова промяната на \texttt{x} променя и \texttt{a}. При предаване чрез референция фактическият параметър трябва да бъде $lvalue$, тоест променяемо място в паметта.

\subsubsection{Съответствие на типовете}

По време на компилация се проверява съответствието между типовете на формалните и фактическите параметри. Проверява се и съответствието между типа на връщания израз и резултатния тип на функцията. Когато е необходимо, се извършва преобразуване на типовете.

\begin{keydefn}[Модел на нормално извикване и връщане]
При извикване на функция се оценяват фактическите параметри, съпоставят се със съответните формални параметри и се създава контекст за изпълнение на функцията. При нормално завършване функцията връща управлението към мястото на извикването. При функция, връщаща стойност, резултатът се задава чрез \texttt{return}; достигането на края е допустимо за \texttt{void} функции и за \texttt{main}. Описанието не определя реда на оценяване на отделните аргументи.
\end{keydefn}

Това твърдение обобщава връзката между процедурната декомпозиция, предаването на данни и управлението на изчислителния процес.

\section{Символни низове}

\begin{defn}[Символен низ]
Символният низ е последователност от символи. В C++ той се представя като масив от елементи от тип \texttt{char}, след последния символ на който се записва терминиращият символ \texttt{'\textbackslash 0'}.
\end{defn}

Терминиращият символ има код $0$ в ASCII таблицата. Той показва къде завършва низът, независимо от размера на масива, в който е записан.

Пример:

\begin{verbatim}
char word[] = {'H', 'e', 'l', 'l', 'o', '\0'};
\end{verbatim}

Същият низ може да бъде зададен чрез литерал:

\begin{verbatim}
char word[] = "Hello";
\end{verbatim}

При дефиниране на по-голям буфер:

\begin{verbatim}
char word1[100] = "Hello";
\end{verbatim}

се заделя място за $100$ символа. Низът заема символите на \texttt{Hello} и терминиращия символ, а останалото място остава част от масива. Максималната дължина на низа в такъв буфер е $99$ символа, тъй като един елемент трябва да бъде запазен за \texttt{'\textbackslash 0'}.

\subsection{Основни операции със символни низове}

Въвеждането чрез \texttt{cin >>} чете до разделител, например интервал, табулация или нов ред.

\begin{verbatim}
char word[100];
cin >> word;
\end{verbatim}

Функцията \texttt{getline} въвежда символи до нов ред, но трябва да й бъде зададен размер на буфера:

\begin{verbatim}
cin.getline(word, 100);
\end{verbatim}

В този случай се въвеждат най-много $99$ символа, така че да остане място за терминиращия символ.

Извеждането се извършва чрез \texttt{cout}:

\begin{verbatim}
cout << word;
\end{verbatim}

До отделен символ се достига чрез индексиране:

\begin{verbatim}
char word[] = "Hello";
cout << word[1];
\end{verbatim}

Резултатът е символът \texttt{'e'}, защото индексирането започва от $0$.

\subsection{Функции от библиотеката \texttt{cstring}}

За работа със символни низове се използват функции от заглавния файл \texttt{cstring}.

\begin{itemize}
\item \texttt{strlen(низ)} връща дължината на низа, тоест броя символи до първото срещане на \texttt{'\textbackslash 0'};
\item \texttt{strcpy(буфер, низ)} прехвърля символите на низа в буфера и връща буфера;
\item \texttt{strcmp(низ1, низ2)} сравнява два низа лексикографски;
\item \texttt{strcat(низ1, низ2)} добавя символите на \texttt{низ2} в края на \texttt{низ1};
\item \texttt{strchr(низ, символ)} търси първото срещане на символ и връща адрес към него или нулев указател;
\item \texttt{strstr(низ, подниз)} търси първото срещане на подниз и връща адрес към него или нулев указател.
\end{itemize}

Функцията \texttt{strcmp} връща:

\begin{itemize}
\item $0$, ако низовете съвпадат;
\item число, по-малко от $0$, ако \texttt{низ1} е преди \texttt{низ2};
\item число, по-голямо от $0$, ако \texttt{низ1} е след \texttt{низ2}.
\end{itemize}

При \texttt{strcat} програмистът трябва да осигури достатъчно място в първия буфер за символите на двата низа и за новия терминиращ символ.

Примерна реализация на \texttt{strlen} е:

\begin{verbatim}
int my_strlen(const char* str) {
    int length = 0;

    while (str[length] != '\0') {
        ++length;
    }

    return length;
}
\end{verbatim}

Алгоритъмът започва с дължина $0$ и последователно проверява символите на низа. При всеки символ, различен от \texttt{'\textbackslash 0'}, дължината се увеличава с единица. Когато бъде достигнат терминиращият символ, текущата дължина се връща.

\section{Изпитен фокус}

При устен отговор трябва да могат да бъдат възпроизведени:

\begin{itemize}
\item определението за процедурно програмиране и връзката му с императивното програмиране;
\item ролята на блоковете и областта на действие;
\item принципите на декомпозицията и модулността;
\item началото и краят на изпълнението чрез функцията \texttt{main};
\item синтаксисът и семантиката на \texttt{if}, \texttt{else if}, \texttt{else}, тернарния оператор и \texttt{switch};
\item поведението на \texttt{switch} при липса на \texttt{break};
\item синтаксисът и редът на изпълнение при \texttt{for}, \texttt{while} и \texttt{do while};
\item понятието променлива, нейните име, тип, адрес и стойност;
\item дефиниране, инициализиране и присвояване;
\item разликата между $lvalue$ и $rvalue$;
\item разликата между локална и глобална променлива;
\item видовете скаларни и съставни типове;
\item структурата на функцията, декларацията, дефиницията, извикването и \texttt{return};
\item предаването на параметри по стойност и чрез референция;
\item съпоставянето на типовете на формалните и фактическите параметри;
\item стековата рамка при извикване на функция;
\item представянето на символен низ чрез масив от \texttt{char} и терминиращ символ;
\item основните операции и функциите от \texttt{cstring};
\item алгоритъма за реализация на \texttt{strlen}.
\end{itemize}


\chapter{Процедурно програмиране – указатели, масиви и рекурсия.}
%% Drafted from the mapped corpus by scripts/draft_topics.py.
% Model: gpt-5.6-luna; source characters: 19891.
\section{Указатели}

\begin{defn}
Указателят е обект, чиято стойност е адрес в паметта. Адресът обикновено е адресът на началото на стойност, а типът на указателя указва какъв е типът на стойността, намираща се на този адрес.
\end{defn}

Указател се декларира чрез оператора $*$:

\[
\texttt{<тип>* <име> [= <израз>];}
\]

Например:

\begin{verbatim}
int x{53};
int* ptr{&x};
\end{verbatim}

В този пример променливата \texttt{x} има стойност $53$, а \texttt{ptr} съдържа адреса на \texttt{x}. На 32-битова архитектура указателят обикновено заема $4$ байта, а на 64-битова архитектура --- $8$ байта. Физически адресът може да се представи като цяло число, но указателят не трябва да се разглежда просто като обикновено цяло число: неговият тип определя как се интерпретира адресът и как се извършват операциите с него.

\subsection{Оператори за работа с указатели}

Операторът за рефериране (достъпване до адрес) \texttt{\&} е унарен оператор. Той връща адреса на своя операнд:

\begin{verbatim}
int x{53};
cout << &x << '\n';
\end{verbatim}

Извежданата стойност е адресът на променливата \texttt{x}.

Операторът за дереференция \texttt{*} също е унарен. Той приема указател и връща стойността, съхранявана на адреса, към който указателят сочи:

\begin{verbatim}
int x{53};
cout << *(&x) << '\n';
\end{verbatim}

Резултатът от програмния фрагмент е $53$. Ако \texttt{ptr} сочи към \texttt{x}, изразът \texttt{*ptr} обозначава самата стойност на \texttt{x} и може да участва както в четене, така и в промяна на стойността, ако това е допустимо.

Указател, който не е инициализиран, съдържа неопределена стойност. Използването му като валиден адрес може да доведе до неправилно поведение. Специалната стойност \texttt{nullptr} означава, че указателят не сочи към обект:

\begin{verbatim}
int* pi;
double* pd = nullptr;
\end{verbatim}

Указателят \texttt{pd} не сочи към никакъв адрес. Дереферирането на \texttt{nullptr} е недопустимо.

Основните операции с указатели са:

\begin{itemize}
\item получаване на адрес чрез \texttt{\&};
\item дерефериране чрез \texttt{*};
\item сравнение чрез операторите \texttt{==}, \texttt{!=}, \texttt{<}, \texttt{>}, \texttt{<=}, \texttt{>=};
\item извеждане чрез \texttt{<<}.
\end{itemize}

Указател не се въвежда директно чрез оператора \texttt{>>}. Обикновено се въвежда стойност в обект, след което указателят се насочва към този обект.

\subsection{Указатели към указатели}

Тъй като указателят е обект, той също има адрес. Следователно може да се дефинира указател към указател:

\begin{verbatim}
double d = 1.23;
double* dptr = &d;
double** dptrptr = &dptr;
\end{verbatim}

Тук \texttt{dptr} сочи към \texttt{d}, а \texttt{dptrptr} сочи към самия указател \texttt{dptr}. Изразът \texttt{*dptrptr} дава указателя \texttt{dptr}, а \texttt{**dptrptr} дава стойността $1.23$.

\subsection{Указатели и константи}

Различават се две важни конструкции.

Константният указател е указател, чиято посока не може да се променя:

\begin{verbatim}
int x = 5;
int* p = &x;
int* const q = p;

*q = 5;
q = p + 2;       // грешка
\end{verbatim}

Посочваната стойност може да бъде променяна чрез \texttt{q}, но на \texttt{q} не може да се присвои друг адрес.

Указателят към константа не позволява промяна на стойността чрез този указател:

\begin{verbatim}
int x = 5;
const int* q = &x;

q++;             // допустимо: променя се адресът
*q = 5;          // грешка
\end{verbatim}

Записът \texttt{const int*} е еквивалентен на \texttt{int const*}. Ако \texttt{x} е константа, то \texttt{\&x} е указател към константа. Не може безусловно да се присвои указател към константа на обикновен указател, защото така би се премахнала гаранцията, че стойността няма да бъде променяна.

\section{Масиви}

\begin{defn}
Масивът е съставен тип данни, представляващ крайна редица от елементи от един и същи тип с фиксирана дължина. Всеки елемент се достига по индекс. Индексите са последователни цели числа, като първият елемент има индекс $0$.
\end{defn}

Едномерни масиви се дефинират например така:

\begin{verbatim}
bool b[10];
double x[3] = {1.2, 0.5, 3.14};
int a[] = {2, 3};
\end{verbatim}

Масивът \texttt{b} има десет елемента. Ако не е инициализиран, елементите му съдържат неопределени стойности. При \texttt{x} дължината и инициализиращият списък са зададени явно. При \texttt{a} размерът е пропуснат и се извежда от броя на елементите в списъка.

Елементите на масива са разположени последователно в паметта. Името на масива се използва като константен указател към първия му елемент. Това обяснява тясната връзка между масивите и указателите.

\subsection{Достъп до елементи}

Достъпът до елемент се извършва чрез оператора \texttt{[]}:

\[
\texttt{<масив>[<цяло число>]}
\]

Например:

\begin{verbatim}
int a[] = {8, 2};
int y = a[1];
\end{verbatim}

След изпълнението \texttt{y} има стойност $2$. За масив \texttt{x[3]} допустимите индекси са $0$, $1$ и $2$. В C++ не се извършва автоматична проверка дали индексът е в допустимите граници. Отрицателен индекс или индекс, по-голям от последния валиден индекс, може да доведе до достъп до чужда памет.

За вградените масиви няма присвояване на цели масиви и няма автоматично поелементно сравнение. Например два масива не могат да бъдат сравнени с един оператор така, че да се сравнят всичките им елементи.

\subsection{Многомерни масиви}

Масив, чиито елементи са едномерни масиви, се нарича двумерен масив. По аналогичен начин се дефинират тримерни и по-общо $N$-мерни масиви:

\[
\texttt{<тип> <име>[<размер1>][<размер2>] \ldots [<размерN>];}
\]

Пример за двумерен масив е:

\begin{verbatim}
int a[2][3] = {{1, 2, 3}, {4, 5, 6}};
\end{verbatim}

Той има две редици с по три елемента. Достъпът до елемент се извършва с два индекса, например \texttt{a[1][2]}.

При инициализиране на многомерен масив първата размерност може да бъде пропусната, ако тя може да се изведе от инициализиращия списък. Останалите размерности трябва да бъдат известни, за да може компилаторът да изчислява адреса на елементите.

\section{Указателна аритметика}

Указателната аритметика позволява чрез указател да се достигат съседни елементи в паметта. Ако \texttt{p} е указател към тип \texttt{T}, тогава:

\[
\texttt{p+i}
\]

обозначава адрес, отместен с $i$ елемента от тип \texttt{T}, а не просто с $i$ байта. Реалното отместване в байтове е:

\[
i\cdot\operatorname{sizeof}(T).
\]

Съответно \texttt{p-i} се отмества назад с $i$ елемента. Например, ако \texttt{int} заема $4$ байта, то \texttt{p+2} е адрес, отстоящ на $8$ байта от \texttt{p}.

При масив \texttt{a} имаме равенствата:

\[
\texttt{a[i]}\Longleftrightarrow\texttt{*(a+i)}
\Longleftrightarrow\texttt{*(i+a)}
\Longleftrightarrow\texttt{i[a]}.
\]

Следователно следният код е еквивалентен по смисъл:

\begin{verbatim}
int a[5] = {1, 2, 3, 4, 5};

cout << *a << '\n';
*(a + 1) = 7;
*(a + 4) = 4;
\end{verbatim}

След присвояването вторият елемент е $7$, а последният елемент е $4$. Изразът \texttt{a} се интерпретира като адреса на първия елемент, а \texttt{a+i} като адреса на елемента с индекс $i$.

\section{Указатели и низови константи}

Символният низ в C++ може да се представи като масив от символи, завършващ с терминиращия символ \texttt{'\textbackslash 0'}. Името на масива от символи е константен указател към първия символ. Низовата константа се разглежда като указател към константни символи.

\begin{verbatim}
char const* p = "Hello";
char q[] = "Hello1";

p = q;
q[1] = 'o';
\end{verbatim}

Чрез \texttt{p} символите на низовата константа не могат да се променят. За разлика от това, \texttt{q} е масив от символи и неговите елементи могат да бъдат променяни, ако масивът не е константен.

\section{Сортиране и търсене в едномерен масив}

\subsection{Selection sort}

Selection sort разделя масива на сортирана и несортирана част. На всяка стъпка се намира най-малкият елемент в несортираната част и се разменя с нейния най-ляв елемент. Така сортираната част нараства с един елемент.

\begin{verbatim}
void selectionSort(int* a, int n) {
    for (int i = 0; i < n - 1; ++i) {
        int min = i;
        for (int j = i + 1; j < n; ++j) {
            if (a[j] < a[min]) {
                min = j;
            }
        }
        if (min != i) {
            swap(a[i], a[min]);
        }
    }
}
\end{verbatim}

Алгоритъмът има времева сложност $O(n^2)$.

\subsection{Bubble sort}

Bubble sort многократно преминава през несортираната част. Сравнява съседни елементи и ги разменя, ако левият е по-голям от десния. След едно преминаване най-големият елемент на разглежданата част се оказва най-вдясно. Процесът се повтаря до сортиране или до преминаване без нито една размяна.

\begin{verbatim}
void bubbleSort(int* a, int n) {
    for (int i = 0; i < n - 1; ++i) {
        bool swapped = false;

        for (int j = 0; j < n - 1 - i; ++j) {
            if (a[j] > a[j + 1]) {
                swap(a[j], a[j + 1]);
                swapped = true;
            }
        }

        if (!swapped) {
            break;
        }
    }
}
\end{verbatim}

Времевата сложност е $O(n^2)$.

\subsection{Insertion sort}

Insertion sort поддържа вече сортирана част от масива. На всяка итерация избира следващ елемент, наречен ключ, и го вмъква на подходящото място, като по-големите елементи се изместват надясно.

\begin{verbatim}
void insertionSort(int* a, int len) {
    int key;

    for (int i = 1; i < len; ++i) {
        key = a[i];
        int j = i - 1;

        while (j >= 0 && key < a[j]) {
            a[j + 1] = a[j];
            --j;
        }

        a[j + 1] = key;
    }
}
\end{verbatim}

Времевата сложност, посочена за алгоритъма, е $O(n^2)$.

\subsection{Quick sort}

Quick sort избира разделящ елемент, наречен pivot, и разделя останалите елементи на подмасиви според това дали са по-малки или по-големи от pivot-а. След това двата подмасива се сортират рекурсивно.

Една възможна схема за разделяне е:

\begin{verbatim}
int partition(int* a, int l, int r) {
    int pivot = a[r];
    int i = l - 1;

    for (int j = l; j < r; ++j) {
        if (a[j] <= pivot) {
            ++i;
            swap(a[i], a[j]);
        }
    }

    swap(a[i + 1], a[r]);
    return i + 1;
}

void quickSort(int* a, int l, int r) {
    if (l >= r) {
        return;
    }

    int p = partition(a, l, r);
    quickSort(a, l, p - 1);
    quickSort(a, p + 1, r);
}
\end{verbatim}

\begin{supp}[Бележка]
В предоставените материали е посочено както $O(n^2)$, така и средна сложност $O(n\log n)$ за quick sort. Тези твърдения се отнасят до различни случаи: $O(n^2)$ е посочено за неблагоприятния случай, а $O(n\log n)$ --- като средна сложност. Изборът на pivot влияе върху полученото разделяне.
\end{supp}

\subsection{Merge sort}

Merge sort разделя масива на две части, след това всяка част отново се разделя, докато се получат подмасиви с по един елемент. После подмасивите се сливат два по два, като при сливането се запазва сортираната наредба.

Сливането на два вече сортирани интервала използва помощен масив:

\begin{verbatim}
void merge(int* a, int l, int m, int r, int* tmp) {
    int i = l;
    int j = m + 1;
    int k = 0;

    while (i <= m && j <= r) {
        if (a[i] <= a[j]) {
            tmp[k++] = a[i++];
        } else {
            tmp[k++] = a[j++];
        }
    }

    while (i <= m) {
        tmp[k++] = a[i++];
    }

    while (j <= r) {
        tmp[k++] = a[j++];
    }

    memcpy(a + l, tmp, k * sizeof(int));
}

void mergeSort(int* a, int l, int r, int* tmp) {
    if (l < r) {
        int m = (l + r) / 2;
        mergeSort(a, l, m, tmp);
        mergeSort(a, m + 1, r, tmp);
        merge(a, l, m, r, tmp);
    }
}
\end{verbatim}

Посочената времева сложност е $O(n\log n)$.

\subsection{Линейно търсене}

При линейното търсене елементите се разглеждат последователно от началото на масива. Търсенето приключва при намиране на търсения елемент или при достигане на края на масива. Времевата сложност е $O(n)$.

\subsection{Двоично търсене}

Двоичното търсене се прилага върху сортиран масив. На всяка стъпка се разглежда средният елемент. Ако той е търсеният, алгоритъмът приключва. Ако средният елемент е по-малък от търсения, се продължава в дясната половина; в противен случай --- в лявата половина.

\begin{verbatim}
int binarySearch(int* a, int n, int x) {
    int l = 0;
    int r = n - 1;

    while (l <= r) {
        int m = l + (r - l) / 2;

        if (a[m] == x) {
            return m;
        }

        if (a[m] < x) {
            l = m + 1;
        } else {
            r = m - 1;
        }
    }

    return -1;
}
\end{verbatim}

Ако елементът бъде намерен, функцията връща индекса му. При липса на елемента връща $-1$.

\section{Рекурсия}

\begin{defn}
Рекурсивна е функция, която извиква себе си пряко или косвено.
\end{defn}

При пряка рекурсия функцията извиква сама себе си в собственото си тяло. При косвена рекурсия се образува цикъл от извиквания:

\[
F_1\to F_2\to\cdots\to F_n\to F_1.
\]

Коректната рекурсивна функция трябва да има условие за край, наричано базов случай. Освен това всяко рекурсивно извикване трябва да води към този случай.

\subsection{Линейна рекурсия}

При линейната рекурсия функцията извършва едно рекурсивно извикване за решаването на даден подпроблем. Пример е факториелът:

\begin{verbatim}
int factorial(int n) {
    if (n == 0) {
        return 1;
    }

    return n * factorial(n - 1);
}
\end{verbatim}

Базовият случай е \texttt{n == 0}. При всяко следващо извикване аргументът намалява с единица.

\begin{supp}[Бележка]
В предоставения материал до примера за факториел е записана сложност $n\log n$. Самата рекурсивна схема извършва по едно извикване за всяка стойност от $n$ до $0$, поради което посоченото твърдение не съответства на показания код. В изложението се използва структурата на алгоритъма: линейна рекурсия с линейно число извиквания.
\end{supp}

\subsection{Разклонена рекурсия}

При разклонената рекурсия функцията извиква себе си два или повече пъти за един подпроблем. Типичен пример е рекурсивната реализация на числата на Фибоначи:

\begin{verbatim}
int fib(int n) {
    if (n == 0 || n == 1) {
        return 1;
    }

    return fib(n - 2) + fib(n - 1);
}
\end{verbatim}

Тук всяко нетривиално извикване поражда две нови извиквания. Получава се дърво на извикванията, а не една линейна верига.

Рекурсията е особено удобна при рекурсивни структури, обхождане на графи, търсене с връщане назад и алгоритми от типа ``разделяй и владей''. Quick sort и merge sort също естествено се описват рекурсивно, тъй като решават същия вид задача върху по-малки подмасиви.

Основен недостатък на рекурсията е скритото използване на памет за стекови рамки. При всяко извикване се съхраняват параметри, локални променливи и информация за връщането. При прекалено дълбока рекурсия стекът може да се изчерпи.

\section{Символни низове}

\begin{defn}
Символният низ е поредица от символи. В C++ той може да се представи като масив от тип \texttt{char}, след последния значим символ на който се записва терминиращият символ \texttt{'\textbackslash 0'}.
\end{defn}

Например:

\begin{verbatim}
char word[] = {'H', 'i', '\0'};
\end{verbatim}

Терминиращият символ има код $0$ и показва края на низа.

Основните операции са:

\begin{itemize}
\item операторът \texttt{>>} въвежда символи до разделител, например интервал, нов ред или табулация;
\item \texttt{cin.getline} въвежда до нов ред, но не повече от зададения брой символи;
\item операторът \texttt{<<} извежда низа;
\item операторът \texttt{[]} осигурява достъп до отделен символ.
\end{itemize}

Функциите за работа с низове се намират в заглавния файл \texttt{<cstring>}:

\begin{itemize}
\item \texttt{strlen(str)} връща броя на символите преди \texttt{'\textbackslash 0'};
\item \texttt{strcpy(buffer, str)} копира низа в буфера и връща буфера;
\item \texttt{strcmp(str1, str2)} извършва лексикографско сравнение;
\item \texttt{strcat(str1, str2)} добавя втория низ към края на първия;
\item \texttt{strchr(str, symbol)} връща адреса на първото срещане на символа или нулев указател;
\item \texttt{strstr(str, sub)} връща адреса на първото срещане на подниза или нулев указател.
\end{itemize}

При \texttt{strcmp} резултатът е $0$, ако низовете съвпадат, отрицателен, ако първият низ предхожда втория, и положителен, ако първият следва втория.

При \texttt{strcat} програмистът трябва да осигури достатъчно място в първия масив за получения низ и неговия терминиращ символ.

Една възможна реализация на \texttt{strlen} е:

\begin{verbatim}
size_t strlen(const char* str) {
    int length = 0;

    while (str[length] != '\0') {
        ++length;
    }

    return length;
}
\end{verbatim}

Функцията обхожда символите последователно, докато достигне \texttt{'\textbackslash 0'}, и брои видимите символи.

\section{Изпитен фокус}

\begin{itemize}
\item Определение на указател и предназначение на операторите \texttt{\&} и \texttt{*}.
\item Разлика между \texttt{nullptr}, неинициализиран указател, константен указател и указател към константа.
\item Представяне на масив в паметта, индексиране и липса на автоматична проверка на границите.
\item Еквивалентността $\texttt{a[i] = *(a+i) = *(i+a) = i[a]}$.
\item Формулата за указателно отместване: $i\cdot\operatorname{sizeof}(T)$.
\item Едномерни и многомерни масиви.
\item Идеята и сложността на selection sort, bubble sort, insertion sort, quick sort и merge sort.
\item Разлика между линейно и двоично търсене и условието двоичното търсене да се извършва върху сортиран масив.
\item Определение за пряка и косвена рекурсия.
\item Разлика между линейна и разклонена рекурсия, включително ролята на базовия случай.
\item Предимства и недостатъци на рекурсивния подход, по-специално използването на стекови рамки.
\item Представяне на символен низ чрез масив от \texttt{char} с терминиращ символ \texttt{'\textbackslash 0'} и предназначение на основните функции от \texttt{<cstring>}.
\end{itemize}


\chapter{ООП. Основни принципи. Класове и обекти. Наследяване и капсулация.}
% AUTO-GENERATED by scripts/build_topics.py — do not edit.
% Edit topics/bodies/topic_16.tex or topics/preamble.tex instead.
%% Placeholder. Replace with the written chapter.
\textit{Източници:} PDF \texttt{16.pdf}, снимки \texttt{16та тема}, \texttt{oop2/oop2.pdf}

\subsection*{Анотация от конспекта}
Обектно-ориентирано програмиране. Основни принципи. Класове и обекти. Наследяване и капсулация.


\chapter{ООП. Подтипов и параметричен полиморфизъм. Множествено наследяване.}
% AUTO-GENERATED by scripts/build_topics.py — do not edit.
% Edit topics/bodies/topic_17.tex or topics/preamble.tex instead.
%% Placeholder. Replace with the written chapter.
\textit{Източници:} снимки \texttt{17та тема}, \texttt{oop2/oop2.pdf}

\subsection*{Анотация от конспекта}
Обектно-ориентирано програмиране. Подтипов и параметричен полиморфизъм. Множествено наследяване.


\chapter{Структури от данни. Стек, опашка, списък, дърво.}
%% Drafted from the mapped corpus by scripts/draft_topics.py.
% Model: gpt-5.6-luna; source characters: 9356.
\begin{konspekt}
\kreq{Официално заглавие}{„Структури от данни. Стек, опашка, списък, дърво. \emph{Основни операции върху тях. Реализация}.“}
\kselect{За изпита ще бъдат избрани \textbf{две} от посочените структури, които да бъдат описани. Възможни са и задачи за \emph{практическата} част на изпита.}
\kreq{Дефиниция}{на понятието структура от данни --- \kref{18.1}.}
\kreq{За всяка структура се иска еднакъв набор}{логическо описание; характеристики на реализациите; сложност на операциите; \textbf{дефиниране на клас}, използващ една от реализациите.}
\kreq{Реализации}{статична, динамична и свързана --- назованите от конспекта три категории --- \kref{18.2}.}
\kreq{Списък}{с една и две връзки; сложност на добавяне, премахване и намиране --- \kr{18.5}; клас --- \kref{18.5.3}.}
\kreq{Стек}{логическо описание и представяния --- \kr{18.3}; клас --- \kref{18.3.3}.}
\kreq{Опашка}{логическо описание и представяния --- \kr{18.4}; клас --- \kref{18.4.3}.}
\kreq{Дървовидни структури}{кореново и двоично кореново дърво; начини за представяне в паметта --- \kr{18.6}; клас --- \kref{18.6.4}.}
\kreq{Двоично кореново дърво за търсене}{логическо описание, представяне, сложност --- \kr{18.7}; клас --- \kref{18.7.5}.}
\kreq{Обобщение}{таблица със сложностите и допусканията --- \kref{18.8.1}.}
\kreq{Литература}{[11], [15], [25], [28].}
\end{konspekt}

\section{Структури от данни}\label{s:18.1}

\begin{defn}[Структура от данни]
Структура от данни е организирана информация, която може да бъде описана, създадена и обработвана с помощта на програма.
\end{defn}

При определяне на една структура от данни се разграничават два аспекта:

\begin{itemize}
    \item логическо описание на структурата -- описание чрез нейната декомпозиция, чрез по-прости структури и чрез операциите над нея, сведени до по-прости операции;
    \item физическо представяне -- методът, по който структурата се разполага в паметта на компютъра.
\end{itemize}

Логическата структура определя какви са елементите, как са свързани и какви операции са допустими. Физическото представяне определя как тази информация се реализира чрез клетки памет, масиви, указатели и други програмни средства.

В настоящата глава се разглеждат четири основни структури от данни: стек, опашка, списък и дърво.

\section{Трите категории реализации}\label{s:18.2}

Конспектът различава \emph{статична}, \emph{динамична} и \emph{свързана}
реализация за стека и за опашката. Разликата е в това как се управлява паметта,
а не в логическото описание, което остава едно и също.

\begin{defn}[Статична (последователна) реализация]
Елементите се пазят в масив с \emph{фиксиран} капацитет $N$, заделен предварително.
Достъпът по индекс е константен, но при $N$ елемента следващото включване е
невъзможно: настъпва \emph{препълване}. Паметта е заета изцяло независимо от
броя на реално съхранените елементи.
\end{defn}

\begin{defn}[Динамична реализация (масив с преоразмеряване)]
Елементите се пазят в масив, заделен в динамичната памет, чийто капацитет се
\emph{удвоява}, когато се запълни: заделя се нов блок с двоен размер, елементите
се копират и старият блок се освобождава. Няма предварителна граница за броя
елементи; достъпът по индекс остава константен.
\end{defn}

\begin{prop}[Амортизирана цена на включването]
При удвояване на капацитета всяко включване струва $O(1)$ \emph{амортизирано},
макар отделно включване да струва $O(n)$, когато предизвика преоразмеряване.
\end{prop}

\begin{proof}
Разглеждаме редица от $n$ включвания в стек, започнал с капацитет $c_0\ge 1$.
Преоразмеряванията стават при размери $c_0,2c_0,4c_0,\ldots$, а цената на всяко е
пропорционална на броя копирани елементи. Общата цена на копиранията е
\[
c_0+2c_0+4c_0+\cdots+2^kc_0
<2\cdot 2^kc_0\le 2n\cdot\frac{2^kc_0}{n},
\]
където $2^kc_0\le n$ е последният капацитет преди спиране; следователно сумата е
по-малка от $2n$. Заедно с $n$-те константни записа общата цена е $O(n)$, тоест
$O(1)$ на включване средно.
\end{proof}

\begin{defn}[Свързана реализация]
Всеки елемент се пази в отделен възел с указател към следващия. Няма
преоразмеряване и няма препълване, но за всеки елемент се плаща допълнителна
памет за указателя, а достъпът до $i$-тия елемент е $O(i)$, защото се минава
през веригата.
\end{defn}

\begin{remark}
Изборът не е въпрос на вкус. Статичната реализация се използва при известна
горна граница и при забрана за динамична памет; динамичната --- когато е нужен
пряк достъп по индекс без такава граница; свързаната --- когато преобладават
включвания и изключвания в известна позиция.
\end{remark}

\section{Стек}\label{s:18.3}

\begin{defn}[Стек]
Стекът е крайна редица от елементи от един и същи тип, при която включването и изключването на елементи са допустими само в единия край на редицата. Този край се нарича връх на стека.
\end{defn}

Стекът е линейна динамична структура от данни. При него е възможен пряк достъп само до елемента на върха. За останалите елементи достъпът се осъществява последователно, започвайки от върха.

\begin{keydefn}[LIFO]
Правилото LIFO (\emph{Last In, First Out}) означава, че последно включеният в стека елемент е първият изключен от него.
\end{keydefn}

Основните операции върху стек са:

\begin{itemize}
    \item \emph{включване} (\emph{push}) -- добавяне на елемент на върха;
    \item \emph{изключване} (\emph{pop}) -- премахване на елемента от върха;
    \item \emph{достъп до върха} -- прочитане на елемента на върха без премахването му;
    \item \emph{проверка за празен стек};
    \item \emph{търсене} на елемент.
\end{itemize}

Включването и изключването на елемент са константни операции, тоест имат сложност $O(1)$. Търсенето на елемент има линейна сложност $O(n)$, тъй като в най-общия случай трябва да се прегледат всички елементи.

\subsection{Последователно представяне на стек}\label{s:18.3.1}

При последователното, или статичното, представяне предварително се заделя блок от памет, обикновено реализиран чрез масив. В този блок стекът нараства и се свива. Поддържа се индекс или указател към върха.

Ако масивът е с размер $N$, стекът може да бъде описан чрез масива \verb|data| и целочислен индекс \verb|top|. При празен стек \verb|top| има стойност, която показва, че няма валиден елемент. При включване стойността на \verb|top| се увеличава и новият елемент се записва на съответната позиция. При изключване елементът на \verb|top| се извлича, след което \verb|top| се намалява.

Примерна реализация на стек от цели числа чрез масив е:

\begin{verbatim}
const int Capacity = 100;

struct Stack {
  int data[Capacity];
  int top;
};

void init(Stack& s) {
  s.top = -1;
}

bool empty(const Stack& s) {
  return s.top == -1;
}

bool full(const Stack& s) {
  return s.top == Capacity - 1;
}

bool push(Stack& s, int value) {
  if (full(s)) {
    return false;
  }

  ++s.top;
  s.data[s.top] = value;
  return true;
}

bool pop(Stack& s, int& value) {
  if (empty(s)) {
    return false;
  }

  value = s.data[s.top];
  --s.top;
  return true;
}
\end{verbatim}

При тази реализация възниква състояние \emph{препълване}, ако се направи опит за включване в стек, който е достигнал капацитета си. При изключване от празен стек възниква състояние \emph{празен стек} или неуспешна операция. Размерът на последователно представения стек е ограничен от предварително заделения блок.

\subsection{Свързано представяне на стек}\label{s:18.3.2}

При свързаното представяне последователните елементи на стека могат да се намират на различни места в паметта. Връзката между два съседни елемента се осъществява чрез указател към следващия елемент. Указателят на последния елемент обикновено има стойност \verb|nullptr|.

За задаване на стека е достатъчен указател към неговия връх. Един елемент на стек от цели числа може да има вида:

\begin{verbatim}
struct StackItem {
  int data;
  StackItem* link;
};
\end{verbatim}

Полето \verb|data| съдържа информацията, а полето \verb|link| съдържа адреса на следващия елемент към вътрешността на стека.

\begin{verbatim}
struct LinkedStack {
  StackItem* top;
};

void init(LinkedStack& s) {
  s.top = nullptr;
}

bool empty(const LinkedStack& s) {
  return s.top == nullptr;
}

void push(LinkedStack& s, int value) {
  StackItem* item = new StackItem;
  item->data = value;
  item->link = s.top;
  s.top = item;
}

bool pop(LinkedStack& s, int& value) {
  if (empty(s)) {
    return false;
  }

  StackItem* item = s.top;
  value = item->data;
  s.top = item->link;
  delete item;
  return true;
}
\end{verbatim}

При включване новият елемент се свързва с досегашния връх и става нов връх. При изключване се запазва стойността на върха, указателят се премества към следващия елемент, а старият елемент се освобождава. И двете операции са с константна сложност.

Свързаното представяне не изисква предварително зададен максимален брой елементи, но за всеки елемент се съхранява и указател. То е подходящо, когато размерът на стека се изменя динамично.


\subsection{Клас за стек (динамична реализация)}\label{s:18.3.3}

Конспектът иска \emph{дефиниране на клас} за стек, използващ една от
реализациите. Функциите по-горе работят върху отворена структура: всеки може да
пипне \verb|top| и да наруши съгласуваността. Класът затваря представянето и
поема отговорност за паметта.

Инвариант на класа: \verb|data| сочи блок от точно \verb|capacity| клетки в
динамичната памет, $0\le$ \verb|size| $\le$ \verb|capacity|, валидни са
елементите \verb|data[0..size-1]|, а върхът е \verb|data[size-1]|.

\begin{verbatim}
#include <stdexcept>

class Stack {
  int* data;       // блок с capacity клетки
  int size;        // брой съхранени елементи
  int capacity;    // размер на блока

  void resize(int newCapacity) {
    int* moved = new int[newCapacity];
    for (int i = 0; i < size; ++i) moved[i] = data[i];
    delete[] data;
    data = moved;
    capacity = newCapacity;
  }

public:
  Stack() : data(new int[4]), size(0), capacity(4) {}

  Stack(const Stack& other)
      : data(new int[other.capacity]),
        size(other.size), capacity(other.capacity) {
    for (int i = 0; i < size; ++i) data[i] = other.data[i];
  }

  Stack& operator=(const Stack& other) {
    if (this != &other) {
      int* moved = new int[other.capacity];
      for (int i = 0; i < other.size; ++i)
        moved[i] = other.data[i];
      delete[] data;
      data = moved;
      size = other.size;
      capacity = other.capacity;
    }
    return *this;
  }

  ~Stack() { delete[] data; }

  bool empty() const { return size == 0; }

  void push(int value) {
    if (size == capacity) resize(2 * capacity);
    data[size++] = value;
  }

  int pop() {
    if (empty())
      throw std::underflow_error("pop от празен стек");
    return data[--size];
  }

  int top() const {
    if (empty())
      throw std::underflow_error("top от празен стек");
    return data[size - 1];
  }
};
\end{verbatim}

Трите члена --- копиращ конструктор, операция за присвояване и деструктор ---
вървят заедно. Ако липсват, копирането на \verb|Stack| копира указателя, двата
обекта освобождават един и същи блок и програмата се разрушава. Това е
управлението на ресурси, което Въпрос 16 нарича RAII: заделянето става в
конструктора, освобождаването --- в деструктора, и никой път през кода не може
да ги пропусне.

Същият клас в \emph{статична} реализация се получава така: \verb|data| става
масив с фиксиран капацитет, \verb|resize| отпада, а \verb|push| хвърля
\verb|std::overflow_error| при пълен стек. В \emph{свързана} реализация трите
полета се заменят с единствен указател \verb|top| към възел, както в \S18.3.2.

\section{Опашка}\label{s:18.4}

\begin{defn}[Опашка]
Опашката е крайна редица от елементи от един и същи тип, при която включването е допустимо в единия край, а изключването -- в другия край.
\end{defn}

Краят, в който се добавят елементи, се нарича \emph{край на опашката}, а краят, от който се премахват елементи, се нарича \emph{начало на опашката}. Пряк достъп е възможен само до елемента в началото на опашката.

\begin{keydefn}[FIFO]
Правилото FIFO (\emph{First In, First Out}) означава, че първият включен в опашката елемент е първият изключен от нея.
\end{keydefn}

Основните операции върху опашка са:

\begin{itemize}
    \item включване на елемент в края;
    \item изключване на елемент от началото;
    \item достъп до елемента в началото;
    \item проверка за празна опашка;
    \item търсене на елемент.
\end{itemize}

Добавянето и премахването на елемент са операции с константна сложност $O(1)$. Търсенето, когато е разрешено, е линейно по броя на елементите.

\subsection{Последователно представяне на опашка}\label{s:18.4.1}

При последователното представяне предварително се заделя блок памет, реализиран чрез масив. Обикновено блокът се разглежда като цикличен. Това означава, че след достигане на последната позиция следва първата позиция, ако в началото има освободено място.

Необходимо е да се пазят поне два индекса:

\begin{itemize}
    \item индекс на началото на опашката;
    \item индекс на следващата свободна позиция в края на опашката.
\end{itemize}

Необходимо е също да се знае броят на елементите или да се използва допълнителен признак, който различава празна от пълна опашка.

\begin{verbatim}
const int Capacity = 100;

struct Queue {
  int data[Capacity];
  int first;
  int last;
  int count;
};

void init(Queue& q) {
  q.first = 0;
  q.last = 0;
  q.count = 0;
}

bool empty(const Queue& q) {
  return q.count == 0;
}

bool full(const Queue& q) {
  return q.count == Capacity;
}

bool enqueue(Queue& q, int value) {
  if (full(q)) {
    return false;
  }

  q.data[q.last] = value;
  q.last = (q.last + 1) % Capacity;
  ++q.count;
  return true;
}

bool dequeue(Queue& q, int& value) {
  if (empty(q)) {
    return false;
  }

  value = q.data[q.first];
  q.first = (q.first + 1) % Capacity;
  --q.count;
  return true;
}
\end{verbatim}

При включване елементът се записва на позицията, определена от \verb|last|, след което индексът се придвижва циклично. При изключване се извлича елементът на позиция \verb|first| и началният индекс също се придвижва циклично.

\subsection{Свързано представяне на опашка}\label{s:18.4.2}

Свързаното представяне на опашка е подобно на свързаното представяне на стек. Елементите се намират на различни места в паметта и са свързани чрез указатели. За разлика от стека, при опашката са нужни два указателя:

\begin{itemize}
    \item указател към началото на опашката;
    \item указател към края на опашката.
\end{itemize}

Елементът на опашката може да бъде описан така:

\begin{verbatim}
struct QueueItem {
  int data;
  QueueItem* link;
};
\end{verbatim}

При празна опашка и двата указателя имат стойност \verb|nullptr|. При първото включване и двата трябва да бъдат насочени към новия елемент. При следващи включвания новият елемент се свързва след досегашния край, а указателят към края се премества към него. При изключване се премахва елементът от началото и указателят към началото се премества към следващия елемент. Ако след премахването опашката стане празна, указателят към края също се установява на \verb|nullptr|.

\begin{verbatim}
struct LinkedQueue {
  QueueItem* first;
  QueueItem* last;
};

void init(LinkedQueue& q) {
  q.first = nullptr;
  q.last = nullptr;
}

bool empty(const LinkedQueue& q) {
  return q.first == nullptr;
}

void enqueue(LinkedQueue& q, int value) {
  QueueItem* item = new QueueItem;
  item->data = value;
  item->link = nullptr;

  if (q.last == nullptr) {
    q.first = item;
    q.last = item;
  } else {
    q.last->link = item;
    q.last = item;
  }
}

bool dequeue(LinkedQueue& q, int& value) {
  if (empty(q)) {
    return false;
  }

  QueueItem* item = q.first;
  value = item->data;
  q.first = item->link;

  if (q.first == nullptr) {
    q.last = nullptr;
  }

  delete item;
  return true;
}
\end{verbatim}

При наличие на указател към края включването е константна операция. Без такъв указател намирането на последния елемент би изисквало преминаване през опашката.


\subsection{Клас за опашка (свързана реализация)}\label{s:18.4.3}

Инвариант: или \verb|head == nullptr| и \verb|tail == nullptr| (празна опашка),
или \verb|head| сочи първия възел, \verb|tail| --- последния,
\verb|tail->next == nullptr|, и всеки възел е достижим от \verb|head|.

\begin{verbatim}
#include <stdexcept>

class Queue {
  struct Node {
    int data;
    Node* next;
  };

  Node* head;   // оттук се изключва
  Node* tail;   // тук се включва

  void clear() {
    while (head != nullptr) {
      Node* item = head;
      head = head->next;
      delete item;
    }
    tail = nullptr;
  }

  void copyFrom(const Queue& other) {
    head = tail = nullptr;
    for (Node* it = other.head; it != nullptr; it = it->next)
      enqueue(it->data);
  }

public:
  Queue() : head(nullptr), tail(nullptr) {}
  Queue(const Queue& other) { copyFrom(other); }

  Queue& operator=(const Queue& other) {
    if (this != &other) {
      clear();
      copyFrom(other);
    }
    return *this;
  }

  ~Queue() { clear(); }

  bool empty() const { return head == nullptr; }

  void enqueue(int value) {
    Node* item = new Node{value, nullptr};
    if (tail != nullptr) tail->next = item;
    else head = item;
    tail = item;
  }

  int dequeue() {
    if (empty())
      throw std::underflow_error("dequeue от празна опашка");
    Node* item = head;
    int value = item->data;
    head = head->next;
    if (head == nullptr) tail = nullptr;
    delete item;
    return value;
  }

  int front() const {
    if (empty())
      throw std::underflow_error("front от празна опашка");
    return head->data;
  }
};
\end{verbatim}

Двата указателя са това, което прави и включването, и изключването $O(1)$: без
\verb|tail| включването в края би изисквало обхождане на цялата верига.
Присвояването на \verb|tail = nullptr| при изпразване не е козметика --- без
него \verb|tail| остава да сочи освободена памет.

\section{Списък}\label{s:18.5}

\begin{defn}[Списък]
Списъкът е крайна редица от елементи от един и същи тип, при която включването и изключването са допустими от произволно място в редицата.
\end{defn}

За разлика от стека и опашката списъкът позволява работа с произволна позиция. Възможен е достъп до всички елементи, който може да бъде пряк или непряк в зависимост от физическото представяне.

Разглеждат се три основни начина за представяне на списък:

\begin{itemize}
    \item последователно представяне;
    \item едносвързано представяне;
    \item двусвързано представяне.
\end{itemize}

Последователното представяне не се използва като основен начин за списък, тъй като включването и изключването в произволно място обикновено изискват преместване на елементи.

\subsection{Едносвързан списък}\label{s:18.5.1}

При едносвързаното представяне елементите се съхраняват на различни места в паметта. Всеки елемент съдържа информация и указател към следващия елемент. Указателят на последния елемент е \verb|nullptr|.

За задаване на списъка е достатъчен указател към началото. Ако се добави и указател към края, включването в края може да се извършва за константно време.

\begin{verbatim}
struct ListItem {
  int data;
  ListItem* next;
};

struct List {
  ListItem* first;
  ListItem* last;
};
\end{verbatim}

Основните операции са:

\begin{itemize}
    \item включване в началото;
    \item включване в края;
    \item включване след даден елемент;
    \item изключване от началото;
    \item изключване след даден елемент;
    \item търсене;
    \item обхождане на списъка.
\end{itemize}

Включването и изключването в началото са константни операции. Включването в края също е константно, ако се поддържа указател към последния елемент. Включването или изключването в произволна позиция, когато позицията трябва първо да бъде намерена, е линейно. Търсенето също е линейна операция.

\begin{verbatim}
void insertFirst(List& list, int value) {
  ListItem* item = new ListItem;
  item->data = value;
  item->next = list.first;
  list.first = item;

  if (list.last == nullptr) {
    list.last = item;
  }
}

bool removeFirst(List& list, int& value) {
  if (list.first == nullptr) {
    return false;
  }

  ListItem* item = list.first;
  value = item->data;
  list.first = item->next;

  if (list.first == nullptr) {
    list.last = nullptr;
  }

  delete item;
  return true;
}

void insertLast(List& list, int value) {
  ListItem* item = new ListItem;
  item->data = value;
  item->next = nullptr;

  if (list.last == nullptr) {
    list.first = item;
    list.last = item;
  } else {
    list.last->next = item;
    list.last = item;
  }
}
\end{verbatim}

За изключване на елемент, който не е първи, е необходим указател към предходния елемент. След промяната указателят на предходния елемент трябва да прескочи изключвания елемент и да сочи към следващия.

\subsection{Двусвързан списък}\label{s:18.5.2}

При двусвързания списък всеки елемент има два указателя: към следващия и към предходния елемент. Елементът може да бъде описан чрез:

\begin{verbatim}
struct DLListItem {
  int data;
  DLListItem* next;
  DLListItem* prev;
};

struct DLList {
  DLListItem* first;
  DLListItem* last;
};
\end{verbatim}

За списъка обикновено се пазят указатели към началото и края. Допълнително може да се пази указател към текущия елемент, което улеснява някои операции по вмъкване, изключване и търсене.

При едносвързан списък движението по връзките е само напред. При двусвързан списък може да се преминава и в двете посоки. Изключването на текущ елемент се извършва чрез свързване на неговия предходен и следващ елемент един с друг. При крайни елементи единият от съответните указатели е \verb|nullptr|.

\begin{verbatim}
void insertAfter(DLList& list,
                 DLListItem* current,
                 int value) {
  if (current == nullptr) {
    return;
  }

  DLListItem* item = new DLListItem;
  item->data = value;
  item->prev = current;
  item->next = current->next;

  if (current->next != nullptr) {
    current->next->prev = item;
  } else {
    list.last = item;
  }

  current->next = item;
}
\end{verbatim}

Двусвързаното представяне изисква повече памет заради втория указател, но улеснява изключването и движението в обратна посока.


\subsection{Клас за списък (двусвързана реализация)}\label{s:18.5.3}

Инвариант: \verb|count| е броят на възлите; предшественикът на \verb|first| и
наследникът на \verb|last| са \verb|nullptr|; за всеки възел $v$ с наследник е
изпълнено \verb|v->next->prev == v|; при празен списък и двата указателя са
\verb|nullptr|.

\begin{verbatim}
class List {
  struct Node {
    int data;
    Node* prev;
    Node* next;
  };

  Node* first;
  Node* last;
  int count;

  Node* nodeAt(int pos) const {      // изисква 0 <= pos < count
    Node* it = first;
    for (int i = 0; i < pos; ++i) it = it->next;
    return it;
  }

  void clear() {
    while (first != nullptr) {
      Node* item = first;
      first = first->next;
      delete item;
    }
    last = nullptr;
    count = 0;
  }

  void copyFrom(const List& other) {
    first = last = nullptr;
    count = 0;
    for (Node* it = other.first; it != nullptr; it = it->next)
      pushBack(it->data);
  }

public:
  List() : first(nullptr), last(nullptr), count(0) {}
  List(const List& other) { copyFrom(other); }

  List& operator=(const List& other) {
    if (this != &other) {
      clear();
      copyFrom(other);
    }
    return *this;
  }

  ~List() { clear(); }

  int size() const { return count; }
  bool empty() const { return count == 0; }

  void pushBack(int value) {
    Node* item = new Node{value, last, nullptr};
    if (last != nullptr) last->next = item;
    else first = item;
    last = item;
    ++count;
  }

  void insertAt(int pos, int value) {  // изисква 0 <= pos <= count
    if (pos == count) {
      pushBack(value);
      return;
    }
    Node* next = nodeAt(pos);
    Node* item = new Node{value, next->prev, next};
    if (next->prev != nullptr) next->prev->next = item;
    else first = item;
    next->prev = item;
    ++count;
  }

  void removeAt(int pos) {           // изисква 0 <= pos < count
    Node* item = nodeAt(pos);
    if (item->prev != nullptr) item->prev->next = item->next;
    else first = item->next;
    if (item->next != nullptr) item->next->prev = item->prev;
    else last = item->prev;
    delete item;
    --count;
  }

  // индекс на първото срещане или -1
  int find(int value) const {
    int i = 0;
    for (Node* it = first; it != nullptr; it = it->next, ++i)
      if (it->data == value) return i;
    return -1;
  }
};
\end{verbatim}

Двете връзки правят премахването на \emph{известен} възел константно: при
едносвързан списък трябва да се търси предшественикът, тоест $O(n)$. Цената на
\verb|insertAt| и \verb|removeAt| тук е $O(\mathit{pos})$ --- разходът е в
намирането на позицията, не в пренавързването, което е $O(1)$.

\section{Дърво}\label{s:18.6}

\begin{defn}[Дърво]
Дърво е свързан ацикличен граф.
\end{defn}

Дървото е разклонена структура от данни. Елементите му се наричат \emph{върхове} или \emph{възли}. Един от върховете се разглежда като корен. Всеки връх може да има произволен брой наследници. Връзките между върховете определят йерархична структура.

Рекурсивното определение на дърво е:

\begin{itemize}
    \item празното дърво е дърво;
    \item дървото с единствен връх, наречен корен, е дърво;
    \item дърво с корен и произволен брой наследници, всеки от които е корен на дърво, също е дърво.
\end{itemize}

Предоставя се пряк достъп до корена, а до останалите върхове достъпът се осъществява чрез връзките в структурата.

\subsection{Двойно дърво}\label{s:18.6.1}

\begin{defn}[Двойно дърво]
Двойното дърво от тип $T$ е рекурсивна структура от данни, която е или празна, или е образувана от данна от тип $T$, наречена корен, двойно дърво от тип $T$, наречено ляво поддърво, и двойно дърво от тип $T$, наречено дясно поддърво.
\end{defn}

Множеството от върховете на двойното дърво се определя рекурсивно. Празното двойно дърво няма върхове. Непразното двойно дърво има корен и всички върхове на лявото и дясното поддърво.

Левият наследник на връх е коренът на лявото му поддърво, ако такова поддърво съществува. Аналогично се определя десният наследник.

\begin{defn}[Листа и вътрешни върхове]
Върхове без наследници се наричат \emph{листа}. Вътрешни върхове са върховете с поне едно дете, включително корена, когато той има дете.
\end{defn}

Свързаното представяне на двойно дърво се реализира чрез указател към структура с три полета: информационно поле и два адреса за лявото и дясното поддърво.

\begin{verbatim}
struct Node {
  int data;
  Node* left;
  Node* right;
};
\end{verbatim}

Празното поддърво се представя с указател \verb|nullptr|.

\subsection{Операции върху двойно дърво}\label{s:18.6.2}

Основните операции върху двойно дърво са:

\begin{itemize}
    \item достъп до връх;
    \item включване на връх;
    \item изключване на връх;
    \item обхождане.
\end{itemize}

До корена се осъществява пряк достъп, а до останалите върхове -- непряк чрез лявата и дясната връзка. Включването и изключването могат да се извършват на произволно място, като резултатът трябва да бъде двойно дърво от същия тип.

\begin{defn}[Обхождане]
Обхождането е метод, чрез който се осъществява достъп до всеки връх на дървото точно един път.
\end{defn}

При рекурсивното обхождане се изпълняват три действия в определен ред:

\begin{enumerate}
    \item обхождане на корена;
    \item обхождане на лявото поддърво;
    \item обхождане на дясното поддърво.
\end{enumerate}

В зависимост от реда на тези действия се получават различни обходи. Когато редът е корен, ляво поддърво, дясно поддърво, се получава предварително обхождане. Когато първо се обхожда лявото поддърво, след това коренът, а накрая дясното поддърво, се получава смесено обхождане. При обхождане ляво поддърво, дясно поддърво, корен се получава последващо обхождане.

\begin{verbatim}
void preorder(Node* tree) {
  if (tree == nullptr) {
    return;
  }

  visit(tree->data);
  preorder(tree->left);
  preorder(tree->right);
}

void inorder(Node* tree) {
  if (tree == nullptr) {
    return;
  }

  inorder(tree->left);
  visit(tree->data);
  inorder(tree->right);
}

void postorder(Node* tree) {
  if (tree == nullptr) {
    return;
  }

  postorder(tree->left);
  postorder(tree->right);
  visit(tree->data);
}
\end{verbatim}

\subsection{Представяния на двойно дърво}\label{s:18.6.3}

Използват се три начина за представяне на двойно дърво:

\begin{enumerate}
    \item свързано представяне;
    \item верижно представяне чрез масиви;
    \item представяне чрез списък на бащите.
\end{enumerate}

При свързаното представяне всеки възел съдържа стойност и два указателя. Този начин е естествен за рекурсивните операции.

При верижното представяне се използват три масива $a[N]$, $b[N]$ и $c[N]$, където $N$ е броят на върховете. Върховете са номерирани от $0$ до $N-1$. В $a[i]$ се съхранява стойността на $i$-тия връх, в $b[i]$ -- индексът на левия наследник, а в $c[i]$ -- индексът на десния наследник. Ако съответният наследник липсва, се записва $-1$. Отделно се пази индексът на корена.

При представяне чрез списък на бащите се използва един масив $p[N]$. В $p[i]$ се съхранява индексът на единствения баща на върха $i$. За корена се записва $-1$.

\begin{supp}[Бележка]
За общо двойно дърво източникът посочва сложност $O(n)$ за търсене, а на друго място посочва $O(\log n)$ за добавяне. Тези оценки зависят от конкретната структура и от начина на търсене. При обикновено двойно дърво без свойство за подреждане търсенето в най-общия случай е линейно. Сложността $O(\log n)$ е характерна за операции в дърво с подходяща височина, например балансирано двойно подредено дърво.
\end{supp}


\subsection{Клас за двойно кореново дърво}\label{s:18.6.4}

Класът притежава корена и отговаря за освобождаването на цялото дърво.
Инвариант: \verb|root| е \verb|nullptr| (празно дърво) или сочи възел, чиито
поддървета са непресичащи се --- никой възел не е достижим по два различни пътя,
тоест структурата е дърво, а не общ граф.

\begin{verbatim}
#include <iostream>

class BinaryTree {
  struct Node {
    int data;
    Node* left;
    Node* right;
  };

  Node* root;

  static void destroy(Node* v) {
    if (v == nullptr) return;
    destroy(v->left);
    destroy(v->right);
    delete v;
  }

  static Node* clone(const Node* v) {
    if (v == nullptr) return nullptr;
    return new Node{v->data, clone(v->left), clone(v->right)};
  }

  static int heightOf(const Node* v) {
    if (v == nullptr) return 0;
    int l = heightOf(v->left);
    int r = heightOf(v->right);
    return 1 + (l > r ? l : r);
  }

  static int countOf(const Node* v) {
    if (v == nullptr) return 0;
    return 1 + countOf(v->left) + countOf(v->right);
  }

  static void inorder(const Node* v) {
    if (v == nullptr) return;
    inorder(v->left);
    std::cout << v->data << ' ';
    inorder(v->right);
  }

public:
  BinaryTree() : root(nullptr) {}

  // Рекурсивното определение, изразено като конструктор.
  BinaryTree(int value,
             const BinaryTree& left, const BinaryTree& right)
      : root(new Node{value,
                      clone(left.root), clone(right.root)}) {}

  BinaryTree(const BinaryTree& other)
      : root(clone(other.root)) {}

  BinaryTree& operator=(const BinaryTree& other) {
    if (this != &other) {
      Node* copied = clone(other.root);   // първо копираме,
      destroy(root);                      // после освобождаваме
      root = copied;
    }
    return *this;
  }

  ~BinaryTree() { destroy(root); }

  bool empty() const { return root == nullptr; }
  int size() const { return countOf(root); }
  int height() const { return heightOf(root); }
  void printInorder() const { inorder(root); }
};
\end{verbatim}

Конструкторът с три аргумента е точно рекурсивната дефиниция: стойност, ляво
поддърво, дясно поддърво. Затова
\begin{verbatim}
BinaryTree leaf4(4, BinaryTree(), BinaryTree());
BinaryTree leaf12(12, BinaryTree(), BinaryTree());
BinaryTree t(10, leaf4, leaf12);
\end{verbatim}
строи дървото с корен $10$ и наследници $4$ и $12$, без нито един \verb|new| в
кода на потребителя.

В \verb|operator=| копирането предхожда освобождаването: ако \verb|clone| хвърли
изключение поради недостиг на памет, обектът остава непроменен вместо да остане
с освободен корен.

\section{Двойно подредено дърво}\label{s:18.7}

\begin{defn}[Двойно подредено дърво]
Двойно подредено дърво от тип $T$ е двойно дърво, за което върховете на лявото поддърво са по-малки от корена, върховете на дясното поддърво са по-големи от корена, а лявото и дясното поддърво също са двойно подредени дървета от тип $T$.
\end{defn}

Празното двойно дърво се приема за двойно подредено дърво. В източника е посочено, че не се включват елементи със същата стойност.

Например за дърво с корен $10$, ляво поддърво с корен $4$ и дясно поддърво с корен $12$ са валидни само поддървета, които спазват същото правило за подреждане.

\subsection{Търсене}\label{s:18.7.1}

Търсенето на елемент $x$ започва от корена:

\begin{enumerate}
    \item ако дървото е празно, елементът не е намерен;
    \item ако $x$ съвпада със стойността на корена, елементът е намерен;
    \item ако $x$ е по-малък от корена, търсенето продължава в лявото поддърво;
    \item ако $x$ е по-голям от корена, търсенето продължава в дясното поддърво.
\end{enumerate}

\begin{verbatim}
Node* search(Node* tree, int value) {
  if (tree == nullptr || tree->data == value) {
    return tree;
  }

  if (value < tree->data) {
    return search(tree->left, value);
  }

  return search(tree->right, value);
}
\end{verbatim}

Източникът посочва сложност $O(\log n)$ за търсене. Тази оценка се получава при дърво с логаритмична височина. При небалансирано дърво височината може да бъде линейна и тогава търсенето е $O(n)$.

\subsection{Включване}\label{s:18.7.2}

Включването също започва от корена. Ако дървото е празно, се създава нов възел с празни ляво и дясно поддърво. Ако стойността е по-малка от корена, включването продължава в лявото поддърво. Ако е по-голяма, продължава в дясното. При същата стойност включването не се извършва.

\begin{verbatim}
Node* insert(Node* tree, int value) {
  if (tree == nullptr) {
    Node* item = new Node;
    item->data = value;
    item->left = nullptr;
    item->right = nullptr;
    return item;
  }

  if (value < tree->data) {
    tree->left = insert(tree->left, value);
  } else if (value > tree->data) {
    tree->right = insert(tree->right, value);
  }

  return tree;
}
\end{verbatim}

Посочената в източника сложност е $O(\log n)$ при дърво с логаритмична височина. При небалансирано дърво включването може да бъде $O(n)$.

\subsection{Изключване}\label{s:18.7.3}

При изключване първо се търси елементът по правилото за подреденото дърво. Ако елементът не съществува, изключване не се извършва.

Разграничават се следните случаи:

\begin{enumerate}
    \item Изтриваният връх е лист. Той може да бъде премахнат, като съответният указател на неговия баща стане \verb|nullptr|.
    \item Връхът има само ляво поддърво. На негово място се поставя лявото поддърво.
    \item Връхът има само дясно поддърво. На негово място се поставя дясното поддърво.
    \item Връхът има и ляво, и дясно поддърво. Избира се най-левият връх на дясното поддърво. Неговата стойност се поставя на мястото на изтриваната стойност, след което избраният връх се изключва от дясното поддърво.
\end{enumerate}

При последния случай избраният връх е най-малкият елемент в дясното поддърво и следователно запазва свойството за подреждане.

\begin{verbatim}
Node* removeMin(Node* tree, int& value) {
  if (tree->left == nullptr) {
    value = tree->data;
    Node* right = tree->right;
    delete tree;
    return right;
  }

  tree->left = removeMin(tree->left, value);
  return tree;
}

Node* remove(Node* tree, int value) {
  if (tree == nullptr) {
    return nullptr;
  }

  if (value < tree->data) {
    tree->left = remove(tree->left, value);
    return tree;
  }

  if (value > tree->data) {
    tree->right = remove(tree->right, value);
    return tree;
  }

  if (tree->left == nullptr) {
    Node* right = tree->right;
    delete tree;
    return right;
  }

  if (tree->right == nullptr) {
    Node* left = tree->left;
    delete tree;
    return left;
  }

  int successor;
  tree->right = removeMin(tree->right, successor);
  tree->data = successor;
  return tree;
}
\end{verbatim}

Както при търсенето и включването, източникът посочва сложност $O(\log n)$ за изключване. Тя се отнася до дърво с логаритмична височина. В общия случай сложността се определя от височината на дървото.

\subsection{Обхождане на двойно подредено дърво}\label{s:18.7.4}

При смесено обхождане -- ляво поддърво, корен, дясно поддърво -- елементите на двойното подредено дърво се получават във възходящ ред.

При обхождане корен, дясно поддърво, ляво поддърво елементите се получават в низходящ ред. Това следва от правилото, че всички елементи в лявото поддърво са по-малки от корена, а всички в дясното са по-големи.


\subsection{Клас за двойно подредено дърво}\label{s:18.7.5}

Инвариант: за всеки възел $v$ всички стойности в лявото поддърво са строго
по-малки от \verb|v->data|, а всички в дясното --- строго по-големи; равни
стойности не се съхраняват. Смесеното обхождане на класа затова връща
нарастваща редица.

\begin{verbatim}
class SearchTree {
  struct Node {
    int data;
    Node* left;
    Node* right;
  };

  Node* root;

  static void destroy(Node* v) {
    if (v == nullptr) return;
    destroy(v->left);
    destroy(v->right);
    delete v;
  }

  static Node* clone(const Node* v) {
    if (v == nullptr) return nullptr;
    return new Node{v->data, clone(v->left), clone(v->right)};
  }

  static Node* insert(Node* v, int x) {
    if (v == nullptr) return new Node{x, nullptr, nullptr};
    if (x < v->data) v->left = insert(v->left, x);
    else if (x > v->data) v->right = insert(v->right, x);
    return v;             // x == v->data: нищо не се променя
  }

  static Node* minNode(Node* v) {
    while (v->left != nullptr) v = v->left;
    return v;
  }

  static Node* remove(Node* v, int x) {
    if (v == nullptr) return nullptr;
    if (x < v->data) { v->left = remove(v->left, x); return v; }
    if (x > v->data) { v->right = remove(v->right, x); return v; }

    if (v->left == nullptr) {      // нула или един наследник
      Node* r = v->right;
      delete v;
      return r;
    }
    if (v->right == nullptr) {
      Node* l = v->left;
      delete v;
      return l;
    }
    // два наследника: най-малкият вдясно
    Node* s = minNode(v->right);
    v->data = s->data;
    v->right = remove(v->right, s->data);
    return v;
  }

public:
  SearchTree() : root(nullptr) {}
  SearchTree(const SearchTree& other)
      : root(clone(other.root)) {}

  SearchTree& operator=(const SearchTree& other) {
    if (this != &other) {
      Node* copied = clone(other.root);
      destroy(root);
      root = copied;
    }
    return *this;
  }

  ~SearchTree() { destroy(root); }

  bool contains(int x) const {
    const Node* v = root;
    while (v != nullptr) {
      if (x == v->data) return true;
      v = (x < v->data) ? v->left : v->right;
    }
    return false;
  }

  void insert(int x) { root = insert(root, x); }
  void remove(int x) { root = remove(root, x); }
};
\end{verbatim}

Изключването на връх с два наследника заменя стойността му с най-малката в
дясното поддърво и после премахва \emph{нея} --- този възел по построение няма
ляв наследник, затова попада в лесния случай. Заменящата стойност е по-голяма от
всичко вляво и по-малка от всичко останало вдясно, следователно инвариантът се
запазва.

И трите операции --- \verb|contains|, \verb|insert| и \verb|remove| --- слизат
по един-единствен път от корена, затова струват $O(h)$, където $h$ е височината.
Това е $O(\log n)$ само при височина от порядъка на $\log n$; при включване на
вече наредени стойности дървото се изражда в списък и трите операции стават
$O(n)$.

\section{Сравнение на структурите}\label{s:18.8}

Стекът и опашката са линейни структури, но се различават по реда на достъп. При стека се използва LIFO, а при опашката -- FIFO. Списъкът е по-обща линейна структура, защото допуска включване и изключване от произволно място.

Дървото е разклонена структура. При двойното дърво всеки връх има най-много ляв и десен наследник. Физическото представяне чрез указатели следва естествено рекурсивното определение на дървото, докато масивните представяния използват индекси вместо адреси.

\subsection{Сложност на операциите}\label{s:18.8.1}

Таблицата обобщава цената на основните операции. Навсякъде $n$ е броят на
съхранените елементи, а $h$ --- височината на дървото. Оценките са в
\emph{най-лошия} случай, освен изрично отбелязаните като амортизирани. Приема се,
че сравнението и копирането на един елемент струват $O(1)$; при елементи с
по-скъпо копиране цената на преоразмеряването и на \verb|clone| се умножава
съответно.

\begin{center}
\small
\begin{tabular}{llccc}
\toprule
Структура & Реализация & Включване & Изключване & Намиране\\
\midrule
Стек   & статична (масив)      & $O(1)$                    & $O(1)$                    & $O(n)$\\
Стек   & динамична (удвояване) & $O(1)^{1}$                & $O(1)$                    & $O(n)$\\
Стек   & свързана              & $O(1)$                    & $O(1)$                    & $O(n)$\\
\midrule
Опашка & статична (циклична)   & $O(1)$                    & $O(1)$                    & $O(n)$\\
Опашка & динамична             & $O(1)^{1}$                & $O(1)$                    & $O(n)$\\
Опашка & свързана              & $O(1)$                    & $O(1)$                    & $O(n)$\\
\midrule
Списък & едносвързан           & $O(1)^{2}$                & $O(n)^{3}$                & $O(n)$\\
Списък & двусвързан            & $O(1)^{2}$                & $O(1)^{4}$                & $O(n)$\\
Списък & последователен        & $O(n)$                    & $O(n)$                    & $O(n)^{5}$\\
\midrule
Двойно дърво & свързана        & $O(1)^{6}$                & $O(1)^{6}$                & $O(n)$\\
ДПД          & свързана        & $O(h)$                    & $O(h)$                    & $O(h)$\\
\bottomrule
\end{tabular}
\end{center}

\begin{center}
\small
\begin{tabular}{@{}ll@{}}
$^{1}$ амортизирано & $^{2}$ в края; при двусвързан и в началото\\
$^{3}$ по стойност --- търси се предшественикът & $^{4}$ при вече известен възел\\
$^{5}$ $O(\log n)$ при сортиран масив & $^{6}$ при вече известен баща\\
\end{tabular}
\end{center}

Единственият ред, който заблуждава, ако се чете без уговорката, е последният:
$O(h)$ е $\Theta(\log n)$ при балансирано дърво и $\Theta(n)$ при изродено.
Балансирането не се разглежда в този въпрос, затова оценките за ДПД се дават
чрез $h$, а не чрез $\log n$.

\section{Изпитен фокус}\label{s:18.9}

\begin{itemize}
    \item Да се дефинира структура от данни и да се различат логическо описание и физическо представяне.
    \item Да се обяснят стекът и правилото LIFO.
    \item Да се опишат операциите включване, изключване, достъп до върха и търсене в стек.
    \item Да се различат трите реализации, назовани от конспекта: статична (масив с фиксиран капацитет), динамична (масив с удвояване, $O(1)$ амортизирано включване) и свързана (възли с указатели).
    \item Да се сравнят последователното и свързаното представяне на стек.
    \item Да се обяснят опашката и правилото FIFO.
    \item Да се опишат началото и краят на опашката и цикличното последователно представяне.
    \item Да се обясни защо при свързана опашка са нужни указатели към началото и края.
    \item Да се дефинира списък и да се сравнят едносвързаното, двусвързаното и последователното представяне.
    \item Да се покажат връзките в едносвързан и двусвързан списък.
    \item Да се дефинира дърво и двойно дърво рекурсивно.
    \item Да се обяснят корен, наследник, листо, вътрешен връх, ляво и дясно поддърво.
    \item Да се опишат свързаното представяне, верижното представяне чрез масиви и списъкът на бащите.
    \item Да се обясни обхождането и редът на действията при предварително, смесено и последващо обхождане.
    \item Да се дефинира двойно подредено дърво.
    \item Да се възпроизведат алгоритмите за търсене, включване и изключване в двойно подредено дърво.
    \item Да се обясни изключването на връх с нула, един или два наследника.
    \item Да се посочи, че оценките $O(\log n)$ за операциите в двойно подредено дърво предполагат логаритмична височина, докато в общия случай височината може да доведе до линейна сложност.
    \item \textbf{Да се напише клас} за всяка от петте структури --- списък, стек, опашка, двойно кореново дърво и двойно подредено дърво --- с инвариант, конструктор, основни операции и деструктор. Конспектът иска „дефиниране на клас“, а не само свободни функции над отворена структура.
    \item Да се обясни защо копиращ конструктор, операция за присвояване и деструктор вървят заедно при клас, който притежава динамична памет, и какво се случва, ако някой от тях липсва.
    \item Да се възпроизведе таблицата със сложностите (\S18.8.1) заедно с допусканията: най-лош случай, $O(1)$ на елемент за сравнение и копиране, и $O(h)$ вместо $O(\log n)$ за ДПД.
\end{itemize}


\chapter{Функционално програмиране. Обща характеристика. Модели на оценяване. Функции от по-висок ред.}
%% Drafted from the mapped corpus by scripts/draft_topics.py.
% Model: gpt-5.6-luna; source characters: 23156.
\section{Функционално програмиране}

\subsection{Обща характеристика на функционалния стил}

Функционалното програмиране е стил на програмиране, в който основните логически единици са функциите. Функцията е програмна единица, която прилага определено преобразуване върху аргументи и връща резултат. Програмата се изгражда чрез дефиниране, използване и комбиниране на функции.

Основното действие във функционалната програма е обръщението към функция. Основният начин за разделяне на програмата на части е дефинирането на име и свързването му с израз, който изчислява стойността на функцията. Основният начин за комбиниране е суперпозицията на функции: резултатът от една функция се подава като аргумент на друга.

Функциите във функционалния стил могат:

\begin{itemize}
    \item да имат параметри;
    \item да се прилагат върху аргументи;
    \item да приемат други функции като аргументи;
    \item да връщат функции като резултат;
    \item да бъдат дефинирани рекурсивно, тоест чрез обръщения към самите себе си.
\end{itemize}

Функционалният стил е представител на декларативния стил на програмиране. При декларативния подход програмата описва какво трябва да бъде полученият резултат, вместо подробно да задава последователност от промени в паметта. Във функционалната програма обикновено липсват традиционни процедурни елементи като присвояване, цикли, заделяне и промяна на клетки в паметта, както и управление на изпълнението чрез прескачане.

Програмите във функционалния стил могат да се разглеждат като списък от дефиниции на функции или като списък от равенства, които описват зависимостите между стойностите. Изпълнението се свежда до оценяване на изрази и прилагане на функции.

\begin{defn}
\emph{Функционално програмиране} е стил на програмиране, при който програмите се изграждат основно чрез дефиниране, прилагане и композиция на функции, а изчисляването се представя като оценяване на изрази.
\end{defn}

Съществен компонент на функционалната програма е липсата на странични ефекти в чистия функционален език. Когато изразът е чист и се оценява в един и същ контекст, той дава една и съща стойност. Това позволява програмата да се разглежда като описание на зависимости между стойности, а не като поредица от изменения на общо състояние.

\subsection{Изрази и стойности в Scheme}

За илюстрация на функционалните понятия ще бъде използван синтаксисът на Scheme. Изразът е атом или комбинация от изрази.

Атомарните изрази не се декомпозират повече. Типични атомарни изрази в Scheme са:

\begin{itemize}
    \item булевите константи \texttt{\#t} и \texttt{\#f};
    \item числовите константи, например \texttt{15}, \texttt{2/3} и \texttt{-1.23};
    \item знаковите константи, например \texttt{\#\textbackslash a} и \texttt{\#\textbackslash newline};
    \item низовите константи, например \texttt{"Scheme"} и \texttt{"hi"};
    \item идентификаторите, тоест имената, които могат да бъдат свързани със стойности;
    \item вградените функции, например \texttt{+}, \texttt{-}, \texttt{*}, \texttt{/} и \texttt{sqrt}.
\end{itemize}

Комбинация в Scheme е израз, конструиран като списък от изрази, заградени в скоби:

\begin{verbatim}
(<expression0> <expression1> ... <expressionn>)
\end{verbatim}

Първият израз, означен като \texttt{<expression0>}, се нарича оператор. Очаква се неговата оценка да бъде функция. Останалите изрази са операнди. Стойността на комбинацията се получава, като функцията, зададена чрез оператора, се приложи върху стойностите на операндите.

Например:

\begin{verbatim}
(+ 2 3)
\end{verbatim}

е комбинация, чийто оператор е \texttt{+}, а операндите са \texttt{2} и \texttt{3}. Оценката на комбинацията е резултатът от прилагането на функцията за събиране върху двете числа.

Правилата за оценяване на някои основни изрази са следните:

\begin{itemize}
    \item оценката на число е самото число;
    \item оценката на низ е самият низ;
    \item оценката на вграден оператор е реализацията на този оператор;
    \item оценката на име е стойността, свързана с името в текущата среда.
\end{itemize}

\begin{defn}
\emph{Среда} е контекстът, който съдържа свързвания между имена и стойности. При оценяване на име се търси стойността, свързана с това име в текущата среда.
\end{defn}

\subsection{Дефиниране и използване на функции}

Дефинирането на функция представлява свързване на име с израз, който описва изчисляването на резултата. Дефиницията е средство за абстракция, защото позволява сложното изчисление да бъде скрито зад име и използвано многократно.

Дефиниране на име чрез израз в Scheme:

\begin{verbatim}
(define <name> <expression>)
\end{verbatim}

Например:

\begin{verbatim}
(define size 2)
\end{verbatim}

свързва името \texttt{size} със стойността \texttt{2}.

Функция с формални параметри може да се дефинира чрез:

\begin{verbatim}
(define (<name> <formal-parameters>) <body>)
\end{verbatim}

Пример:

\begin{verbatim}
(define (square x)
  (* x x))
\end{verbatim}

Тук \texttt{square} е името на функцията, \texttt{x} е формален параметър, а \texttt{(* x x)} е тялото. При обръщението

\begin{verbatim}
(square 5)
\end{verbatim}

функцията се прилага върху аргумента \texttt{5} и резултатът е \texttt{25}.

\begin{defn}
\emph{Формален параметър} е име, използвано в дефиницията на функцията, което представя стойността на съответния аргумент при конкретно приложение.
\end{defn}

\begin{defn}
\emph{Аргумент} е изразът или стойността, подадени при конкретно обръщение към функция.
\end{defn}

Формалните параметри и действителните аргументи се съпоставят по позиция. При прилагането на функцията формалните параметри получават стойностите на съответните аргументи, след което се оценява тялото на функцията.

Рекурсията е основен метод за дефиниране на функции. При рекурсивна дефиниция функцията се обръща към самата себе си. Източниците посочват рекурсията като характерен начин за дефиниране във функционалното програмиране, без да налагат конкретна рекурсивна задача или пример.

\subsection{Условни форми и локални дефиниции}

Условната форма \texttt{if} има синтаксис:

\begin{verbatim}
(if <condition> <expression1> <expression2>)
\end{verbatim}

Първо се оценява условието. Ако неговата стойност е \texttt{\#t}, се оценява и връща \texttt{<expression1>}. В противен случай се оценява и връща \texttt{<expression2>}.

Тъй като се оценява само избраното разклонение, \texttt{if} е специална форма, а не обикновена функция. Специалните форми са конструкции, които са изключения от общото правило за оценяване на комбинациите.

За многозначен избор се използва \texttt{cond}:

\begin{verbatim}
(cond
  (<condition1> <expression1>)
  ...
  (<conditionn> <expressionn>)
  (else <expressionm>))
\end{verbatim}

Условията се оценяват последователно. При първото условие със стойност \texttt{\#t} се оценява съответният израз и неговата стойност става стойност на цялата форма. Ако нито едно условие не е изпълнено, се избира клонът \texttt{else}.

Локални имена могат да се създават чрез специалната форма \texttt{let}:

\begin{verbatim}
(let ((<symbol1> <expression1>)
      ...
      (<symboln> <expressionn>))
  <body>)
\end{verbatim}

Оценката на \texttt{let} в среда $E$ се описва по следния начин:

\begin{enumerate}
    \item създава се нова среда $E_1$, която е разширение на $E$;
    \item всеки израз $\texttt{<expression}_i\texttt{>}$ се оценява в $E$;
    \item получената стойност се свързва със съответното име $\texttt{<symbol}_i\texttt{>}$ в $E_1$;
    \item тялото се оценява в $E_1$ и неговата стойност се връща като стойност на \texttt{let}.
\end{enumerate}

Символично:

\[
E_1=E\cup
\{\texttt{<symbol}_i\texttt{>} \mapsto
\operatorname{eval}(\texttt{<expression}_i\texttt{>},E)
\mid 1\leq i\leq n\},
\]

а резултатът е:

\[
\operatorname{eval}(\texttt{<body>},E_1).
\]

В Haskell локални дефиниции могат да се задават чрез \texttt{where}:

\begin{verbatim}
<function-definition>
  where
    <definition1>
    ...
    <definitionn>
\end{verbatim}

Областта на действие на тези дефиниции е ограничена до дефиницията на функцията. Дефинициите в \texttt{where} могат да бъдат взаимно рекурсивни.

\section{Оценяване на изрази}

\subsection{Общо правило за оценяване на комбинация}

Общото правило за оценяване на обикновена комбинация се състои от следните стъпки:

\begin{enumerate}
    \item оценяват се подизразите на комбинацията;
    \item оценката на най-левия подизраз се разглежда като функция;
    \item тази функция се прилага върху оценките на останалите подизрази.
\end{enumerate}

Например:

\begin{verbatim}
(* (+ 2 (* 4 6)) (+ 3 5 7))
\end{verbatim}

се оценява чрез последователността:

\begin{verbatim}
(* (+ 2 24) 15)
(* 26 15)
390
\end{verbatim}

Тук общото правило се прилага към вложените комбинации \texttt{(* 4 6)}, \texttt{(+ 2 (* 4 6))}, \texttt{(+ 3 5 7)} и външната комбинация.

\begin{keydefn}
При обикновена комбинация първо се оценяват операторът и операндите, а след това функцията, получена като стойност на оператора, се прилага върху стойностите на операндите.
\end{keydefn}

Конструкциите \texttt{define}, \texttt{if} и \texttt{cond} не следват непосредствено това правило. Те са специални форми и имат собствени правила за оценяване.

\subsection{Оценяване на приложение на функция}

При комбинация, чийто оператор е функция, се оценяват елементите на комбинацията и се прилага процедурата, която е стойност на оператора, върху стойностите на операндите.

При дефиницията:

\begin{verbatim}
(define (square x)
  (* x x))
\end{verbatim}

специалната форма \texttt{define} създава свързване между името \texttt{square} и зададената функция. Това свързване се добавя в текущата среда. Функцията може по-късно да бъде приложена чрез:

\begin{verbatim}
(square 4)
\end{verbatim}

Моделът на заместването представя прилагането на функцията чрез заместване на формалните параметри със съответните аргументи. За \texttt{(square 4)} тялото

\begin{verbatim}
(* x x)
\end{verbatim}

се разглежда след заместването като

\begin{verbatim}
(* 4 4)
\end{verbatim}

и се оценява до \texttt{16}. Оценката на последния израз от тялото е оценката на обръщението към функцията.

Този модел е концептуален: той обяснява как аргументите се свързват с формалните параметри. Реалната реализация може да използва друга вътрешна техника, но поведението се описва чрез същите зависимости.

\section{Модели на оценяване}

При модела на заместването са възможни два основни подхода за оценяване на аргументите: апликативен и нормален.

\subsection{Апликативно оценяване}

Апликативното оценяване се нарича още строго оценяване или оценяване по стойност (\textit{call-by-value}). При него първо се оценяват операндите, след което функцията се прилага върху получените стойности.

Схематично:

\[
\text{оценяване на операндите}
\longrightarrow
\text{получаване на стойности}
\longrightarrow
\text{прилагане на функцията}.
\]

Ако имаме приложение от вида

\begin{verbatim}
(f e1 e2)
\end{verbatim}

апликативният подход първо оценява \texttt{e1} и \texttt{e2}, а след това прилага \texttt{f} върху получените резултати.

Предимството на този подход, посочено в източниците, е, че преди прилагането на функцията всеки аргумент е сведен до една стойност. Това е пестеливо откъм памет, тъй като не се налага да се съхраняват неоценени изрази. Аргументът е, че се извършва оценяване дори когато стойността му в крайна сметка не е необходима на тялото на функцията.

\subsection{Нормално оценяване}

Нормалното оценяване се нарича още лениво оценяване или оценяване по име (\textit{call-by-name}). При него аргументите не се оценяват предварително. Вместо това имената на аргументите се заместват в тялото на функцията и изразите се оценяват, когато това стане необходимо.

Схематично:

\[
\text{заместване в тялото}
\longrightarrow
\text{достигане до примитивни операции}
\longrightarrow
\text{оценяване на необходимите изрази}.
\]

Следователно нормалният подход първо разгъва функцията и чак след това оценява аргументите, които участват в получения израз.

Нормалното оценяване може да избегне оценяване на аргумент, който не се използва. То обаче може да доведе до многократно оценяване на един и същ израз, ако той се появява повече от веднъж след заместването. Така се пести време за предварително оценяване само когато изразът не е необходим, но може да се изразходва повече време при многократно пресмятане. Същевременно е необходимо да се съхраняват или повторно да се разгръщат неоценени изрази, поради което подходът е по-непестелив откъм памет.

\begin{keydefn}
При апликативното оценяване аргументите се оценяват преди прилагането на функцията. При нормалното оценяване аргументите се заместват и се оценяват само когато са необходими.
\end{keydefn}

\begin{remark}
Една и съща програма може да се държи различно при апликативен и нормален подход, особено когато функцията не използва някой аргумент или когато оценяването на аргумент би довело до проблем. Източниците посочват също, че в стандартните езици има лениво поведение в отделни конструкции, например при конюнкция, когато след лъжлив операнд останалите операнди не се пресмятат.
\end{remark}

\begin{supp}[Уточнение за терминологията]
В учебната литература изразът ``лениво оценяване'' често се използва по-широко от ``call-by-name'' и може да включва запомняне на вече изчислена стойност. Разглежданият тук материал представя нормалния подход чрез заместване и акцентира върху възможното многократно оценяване на един и същ израз.
\end{supp}

\section{Функции от по-висок ред}

\subsection{Определение}

\begin{defn}
\emph{Функция от по-висок ред} е функция, която приема функция като параметър или връща функция като резултат.
\end{defn}

Функциите от по-висок ред са естествено следствие от идеята, че функциите са стойности, с които програмата може да работи. Те могат да се подават на други функции, да се връщат като резултат и да се използват за абстрахиране на повтарящи се схеми на изчисление.

Функция, която приема друга функция като параметър, може да прилага подадената функция върху дадена стойност. Например предикатът за неподвижна точка може да бъде записан в Scheme като:

\begin{verbatim}
(define (fixed-point f x)
  (= (f x) x))
\end{verbatim}

Тук \texttt{f} е параметър, чиято стойност трябва да бъде функция, а \texttt{x} е стойност, върху която \texttt{f} се прилага.

Практически примери за функции от по-висок ред са \texttt{map}, \texttt{filter}, \texttt{foldl}, \texttt{foldr} и \texttt{apply}. Функцията \texttt{map} прилага дадена функция върху всеки елемент на списък и връща нов списък:

\begin{verbatim}
(map (lambda (x) (* x x)) '(1 2 3))
\end{verbatim}

Резултатът е:

\begin{verbatim}
'(1 4 9)
\end{verbatim}

Функцията \texttt{filter} приема предикат и списък и връща елементите, за които предикатът е истина:

\begin{verbatim}
(filter even? '(1 2 3 4 5))
\end{verbatim}

Резултатът е:

\begin{verbatim}
'(2 4)
\end{verbatim}

Функциите \texttt{foldr} и \texttt{foldl} свиват списък чрез подадена функция и начална стойност. За пример:

\begin{verbatim}
(foldr - 0 '(1 2 3))
\end{verbatim}

се интерпретира като:

\[
1-(2-(3-0)).
\]

При \texttt{foldl} обработката се извършва отляво надясно чрез акумулиране на междинния резултат:

\begin{verbatim}
(foldl - 0 '(1 2 3))
\end{verbatim}

В Racket функцията за \texttt{foldl} получава първо текущия елемент и после акумулатора, затова изразът се интерпретира като:

\[
3-(2-(1-0))=2.
\]

Този ред на аргументите е различен от често използваната Haskell конвенция, при която акумулаторът е първи.

Функцията \texttt{apply} приема функция и списък от аргументи, като използва елементите на списъка като отделни аргументи:

\begin{verbatim}
(apply + '(1 2 3))
\end{verbatim}

е еквивалентно на:

\begin{verbatim}
(+ 1 2 3)
\end{verbatim}

и дава резултат \texttt{6}.

\subsection{Къринг}

В Haskell функция с няколко аргумента може да се разглежда като функция на един аргумент, която връща функция, очакваща следващия аргумент. Това представяне се нарича къринг.

Ако функцията има тип:

\[
t_1\to t_2\to t_3,
\]

тя приема аргумент от тип $t_1$ и връща функция от тип:

\[
t_2\to t_3.
\]

Следователно частичното прилагане на функцията създава нова функция. Например:

\begin{verbatim}
div50 x = div 50 x
\end{verbatim}

може да се разглежда като дефиниране на:

\begin{verbatim}
div50 = div 50
\end{verbatim}

и след това:

\begin{verbatim}
div50 4
\end{verbatim}

дава \texttt{12}.

\begin{defn}
\emph{Къринг} е представяне на функция с няколко аргумента като последователност от функции, всяка от които приема един аргумент и връща функция или краен резултат.
\end{defn}

\subsection{Анонимни функции}

Анонимна, или ламбда, функция е функция без зададено име. Тя може да бъде конструирана на мястото, където е необходима, особено като аргумент на функция от по-висок ред.

Синтаксисът в Scheme е:

\begin{verbatim}
(lambda (<parameters>) <body>)
\end{verbatim}

Пример:

\begin{verbatim}
(lambda (x) (+ x 3))
\end{verbatim}

оценява до функционален обект с параметър \texttt{x} и тяло \texttt{(+ x 3)}. Този обект може да бъде приложен непосредствено:

\begin{verbatim}
((lambda (x) (+ x 3)) 5)
\end{verbatim}

Резултатът е \texttt{8}.

Анонимните функции позволяват на програмиста да задава поведението на функция от по-висок ред без предварително именуване на помощна функция. В предишния пример ламбда-функцията е подадена директно на \texttt{map}.

\subsection{Функции, които връщат функции}

Функция от по-висок ред може да връща нова функция като резултат. Пример е функцията \texttt{twice}, която приема функция \texttt{f} и връща функция, прилагаща \texttt{f} два пъти:

\begin{verbatim}
(define (twice f)
  (lambda (x) (f (f x))))
\end{verbatim}

Обръщението:

\begin{verbatim}
(twice square)
\end{verbatim}

връща функция. Прилагането на получения резултат върху \texttt{3} дава:

\begin{verbatim}
((twice square) 3)
\end{verbatim}

Това е еквивалентно на прилагане на \texttt{square} два пъти:

\[
(3^2)^2=81.
\]

Друг пример е функция, която създава функция за прибавяне на фиксирано число:

\begin{verbatim}
(define (n+ n)
  (lambda (i) (+ i n)))
\end{verbatim}

Така може да се дефинира:

\begin{verbatim}
(define 1+ (n+ 1))
\end{verbatim}

Полученото \texttt{1+} е функция, която прибавя фиксирана стойност към своя аргумент.

\begin{supp}[Бележка към примерната стойност]
При дефиницията \texttt{(define 1+ (n+ 1))} обръщението \texttt{(1+ 4)} според самото тяло \texttt{(+ i n)} дава \texttt{5}. Среща се и изписване на друга резултатна стойност в учебния материал; тя не следва от показаната дефиниция, затова тук е запазена зависимостта, зададена от кода.
\end{supp}

\section{Изпитен фокус}

\begin{itemize}
    \item Да се дефинира функционалното програмиране и да се опишат функциите като основни програмни единици.
    \item Да се обясни декларативният характер на функционалния стил, композицията на функции и ролята на рекурсията.
    \item Да се различат атом, комбинация, оператор и операнд в Scheme.
    \item Да се възпроизведат синтаксисът и смисълът на \texttt{define}, \texttt{if}, \texttt{cond}, \texttt{let} и локалните дефиниции чрез \texttt{where}.
    \item Да се изложи общото правило за оценяване на комбинация и да се проследи примерът с вложено събиране и умножение.
    \item Да се обясни моделът на заместването при прилагане на функция.
    \item Да се сравнят апликативното оценяване \textit{call-by-value} и нормалното оценяване \textit{call-by-name} по ред на оценяване, време и памет.
    \item Да се дефинира функция от по-висок ред и да се дадат примери с функции като параметри и резултати.
    \item Да се обяснят кърингът, частичното прилагане и анонимните ламбда-функции.
    \item Да се проследи примерът \texttt{twice}, който връща функция, прилагаща подадена функция два пъти.
\end{itemize}


\chapter{Функционално програмиране. Списъци. Потоци и отложено оценяване.}
% AUTO-GENERATED by scripts/build_topics.py — do not edit.
% Edit topics/bodies/topic_20.tex or topics/preamble.tex instead.
%% Placeholder. Replace with the written chapter.
\textit{Източници:} PDF \texttt{20.pdf}, снимки \texttt{20та тема}, \texttt{Задачи/FP.txt}

\subsection*{Анотация от конспекта}
Функционално програмиране. Списъци. Потоци и отложено оценяване.


\chapter{Синтаксис и семантика на предикатното смятане от първи ред.}
%% Drafted from the mapped corpus by scripts/draft_topics.py.
% Model: gpt-5.6-terra; source characters: 275279.
\begin{konspekt}
\kreq{Дефиниции (синтаксис)}{език на предикатното смятане от първи ред, терм, формула, област на действие на квантор, свободна и свързана променлива, затворена формула --- \kref{21.1}.}
\kreq{Дефиниции (семантика)}{структура и оценка на индивидните променливи --- \kref{21.2}.}
\kreq{Дефиниции (семантика)}{вярност на формула в структура при оценка, тъждествена вярност, предикатна тавтология --- \kref{21.3}.}
\kreq{Доказателство}{верността при оценка не зависи от стойностите на променливите, които не участват свободно във формулата --- \kref{21.3}.}
\kreq{Формулировка}{лема за субституциите в \emph{общия} й вид --- за произволна формула и допустима субституция --- \kref{21.4}.}
\kreq{Доказателство}{лемата за субституциите в \emph{частния случай} на безкванторна формула. Само този случай се иска с доказателство --- \kref{21.4}.}
\kreq{Типове задачи}{определимост на свойства в дадена структура; изпълнимост на множество формули чрез посочване на модел --- \kref{21.5}.}
\kreq{Литература}{[13], [19], [37].}
\end{konspekt}

\section{Език и синтаксис}\label{s:21.1}

\subsection{Предикатен език от първи ред}\label{s:21.1.1}

\begin{defn}
Езикът на предикатното смятане от първи ред $\mathcal L$ има логическа и нелогическа част.

Логическата част съдържа изброимо множество от индивидуални променливи
\[
\mathrm{Var}=\{x,y,z,x_0,x_1,\ldots\},
\]
логическите връзки $\neg,\wedge,\vee,\Rightarrow,\Leftrightarrow$, кванторите $\forall,\exists$ и помощните символи $(,)$ и запетая. По избор езикът може да съдържа и формален символ за равенство $\doteq$.

Нелогическата част се състои от:
\begin{itemize}
\item множество $\mathrm{Const}_{\mathcal L}$ от индивидуални константи;
\item множество $\mathrm{Func}_{\mathcal L}$ от функционални символи;
\item множество $\mathrm{Pred}_{\mathcal L}$ от предикатни символи.
\end{itemize}
На всеки функционален и предикатен символ се съпоставя положително естествено число, наречено \emph{арност}. Ако $f$ е $n$-местен функционален символ, пишем $\#(f)=n$; аналогично за предикатен символ $p$.
\end{defn}

\begin{remark}
Езикът е от \emph{първи ред}, защото кванторите се отнасят само за елементи на универсума. Не се квантифицира пряко по функции, предикати или множества от елементи.
\end{remark}

\subsection{Термове}\label{s:21.1.2}

\begin{defn}
Множеството на термовете на $\mathcal L$ се дефинира индуктивно:
\begin{enumerate}
\item всяка индивидуална променлива е терм;
\item всяка индивидуална константа е терм;
\item ако $f\in\mathrm{Func}_{\mathcal L}$, $\#(f)=n$, и $\tau_1,\ldots,\tau_n$ са термове, то
\[
f(\tau_1,\ldots,\tau_n)
\]
е терм.
\end{enumerate}
\end{defn}

\begin{example}
Ако $c$ е константа, $f$ е едноместен, а $g$ е двуместен функционален символ, то
\[
x,\qquad c,\qquad f(x),\qquad g(f(x),c)
\]
са термове.
\end{example}

За терм $\tau$ с $\operatorname{Var}[\tau]$ означаваме множеството от променливите, които участват в него:
\[
\operatorname{Var}[x]=\{x\},\qquad
\operatorname{Var}[c]=\varnothing,
\]
\[
\operatorname{Var}[f(\tau_1,\ldots,\tau_n)]
 =\bigcup_{i=1}^{n}\operatorname{Var}[\tau_i].
\]
\begin{defn}[Затворен терм]
Терм без променливи се нарича \emph{затворен} или \emph{основен} терм.
\end{defn}

\subsection{Формули}\label{s:21.1.3}

\begin{defn}
Ако $p\in\mathrm{Pred}_{\mathcal L}$ е $n$-местен и $\tau_1,\ldots,\tau_n$ са термове, то
\[
p(\tau_1,\ldots,\tau_n)
\]
се нарича \emph{атомарна формула}. Ако езикът е с формално равенство, то и
\[
\tau_1\doteq\tau_2
\]
е атомарна формула.
\end{defn}

\begin{defn}
Множеството на формулите на $\mathcal L$ се задава индуктивно:
\begin{enumerate}
\item атомарните формули са формули;
\item ако $\varphi$ и $\psi$ са формули, то $\neg\varphi$,
\[
(\varphi\wedge\psi),\quad
(\varphi\vee\psi),\quad
(\varphi\Rightarrow\psi),\quad
(\varphi\Leftrightarrow\psi)
\]
са формули;
\item ако $x$ е индивидуална променлива и $\varphi$ е формула, то
\[
\forall x\,\varphi,\qquad \exists x\,\varphi
\]
са формули.
\end{enumerate}
\end{defn}

\begin{defn}
Формула $\psi$ е \emph{подформула} на формула $\varphi$, ако $\psi$ се получава като част от синтактичното построение на $\varphi$.
\end{defn}

\subsection{Област на действие; свободни и свързани променливи}\label{s:21.1.4}

\begin{defn}
В подформула от вида
\[
\forall x\,\psi\qquad\text{или}\qquad\exists x\,\psi
\]
формулата $\psi$ се нарича \emph{област на действие} на съответния квантор.
\end{defn}

\begin{defn}
Едно участие на променливата $x$ във формула е \emph{свързано}, ако лежи в областта на действие на квантор по $x$. В противен случай участието е \emph{свободно}.
\end{defn}

Най-близкият квантор по $x$, чиято област на действие съдържа дадено участие, определя дали това участие е свързано. Една и съща променлива може да има както свободни, така и свързани участия.

\begin{example}
Във формулата
\[
p(x)\wedge\forall x\,q(x,y)
\]
първото участие на $x$ е свободно, а второто е свързано. Всички участия на $y$ са свободни.
\end{example}

С $\operatorname{Var}^{\mathrm{free}}[\varphi]$ означаваме множеството на свободните променливи на $\varphi$, а с $\operatorname{Var}^{\mathrm{bd}}[\varphi]$ --- множеството на променливите с поне едно свързано участие. Важни рекурсивни равенства са:
\[
\operatorname{Var}^{\mathrm{free}}[p(\tau_1,\ldots,\tau_n)]
=\bigcup_{i=1}^{n}\operatorname{Var}[\tau_i],
\]
\[
\operatorname{Var}^{\mathrm{free}}[\neg\varphi]
=\operatorname{Var}^{\mathrm{free}}[\varphi],
\]
\[
\operatorname{Var}^{\mathrm{free}}[\varphi\circ\psi]
=\operatorname{Var}^{\mathrm{free}}[\varphi]\cup
\operatorname{Var}^{\mathrm{free}}[\psi],
\qquad
\circ\in\{\wedge,\vee,\Rightarrow,\Leftrightarrow\},
\]
\[
\operatorname{Var}^{\mathrm{free}}[\forall x\,\varphi]
=\operatorname{Var}^{\mathrm{free}}[\varphi]\setminus\{x\},
\qquad
\operatorname{Var}^{\mathrm{free}}[\exists x\,\varphi]
=\operatorname{Var}^{\mathrm{free}}[\varphi]\setminus\{x\}.
\]

\begin{defn}
Формула $\varphi$ се нарича \emph{затворена}, ако
\[
\operatorname{Var}^{\mathrm{free}}[\varphi]=\varnothing.
\]
\end{defn}

\section{Структури, оценки и стойности на термове}\label{s:21.2}

\begin{defn}
Структура за езика $\mathcal L$ е двойка
\[
\mathcal A=\langle A,I\rangle,
\]
където $A\neq\varnothing$ е \emph{универсум} или \emph{носител} на структурата, а $I$ е интерпретация на нелогическите символи:
\begin{itemize}
\item за всяка константа $c$ е зададен елемент $c^{\mathcal A}\in A$;
\item за всеки $n$-местен функционален символ $f$ е зададена функция
\[
f^{\mathcal A}:A^n\to A;
\]
\item за всеки $n$-местен предикатен символ $p$ е зададена релация
\[
p^{\mathcal A}\subseteq A^n.
\]
\end{itemize}
\end{defn}

\begin{remark}
Функционалният символ $f$ е синтактичен знак, докато $f^{\mathcal A}$ е функция върху универсумa. Аналогично $p$ е знак, а $p^{\mathcal A}$ е релация.
\end{remark}

\begin{example}
За език с константа $0$, двуместен функционален символ $+$ и двуместен предикатен символ $<$ естествените числа могат да се разглеждат като структура
\[
\mathcal N=\langle\mathbb N,0^{\mathcal N},+^{\mathcal N},<^{\mathcal N}\rangle.
\]
\end{example}

\begin{defn}
\emph{Оценка на индивидуалните променливи} в структура $\mathcal A$ е функция
\[
v:\mathrm{Var}\to A.
\]
За $a\in A$ с $v_x^a$ означаваме оценката, модифицирана в $x$ с $a$:
\[
v_x^a(y)=
\begin{cases}
a,&y=x,\\
v(y),&y\ne x.
\end{cases}
\]
\end{defn}

Стойността на терм $\tau$ в $\mathcal A$ при оценка $v$, означавана с
\[
\llbracket\tau\rrbracket^{\mathcal A}_v,
\]
се определя рекурсивно:
\[
\llbracket x\rrbracket^{\mathcal A}_v=v(x),
\qquad
\llbracket c\rrbracket^{\mathcal A}_v=c^{\mathcal A},
\]
\[
\llbracket f(\tau_1,\ldots,\tau_n)\rrbracket^{\mathcal A}_v
=
f^{\mathcal A}\bigl(
\llbracket\tau_1\rrbracket^{\mathcal A}_v,\ldots,
\llbracket\tau_n\rrbracket^{\mathcal A}_v
\bigr).
\]

\begin{prop}
Нека $\tau$ е терм, а $v_1,v_2$ са оценки в $\mathcal A$. Ако
\[
v_1(x)=v_2(x)
\quad\text{за всяко }x\in\operatorname{Var}[\tau],
\]
то
\[
\llbracket\tau\rrbracket^{\mathcal A}_{v_1}
=
\llbracket\tau\rrbracket^{\mathcal A}_{v_2}.
\]
\end{prop}

\begin{proof}
Доказателството е индукция по построението на $\tau$. За променлива твърдението следва от условието; за константа стойността не зависи от оценката. При
\[
\tau=f(\tau_1,\ldots,\tau_n)
\]
оценките съвпадат върху променливите на всеки $\tau_i$. Индукционната хипотеза дава равенство на стойностите на всички аргументи, а после и равенство на стойностите след прилагане на $f^{\mathcal A}$.
\end{proof}

\begin{cor}
Стойността на затворен терм не зависи от оценката.
\end{cor}

\begin{proof}
Множеството от променливи на затворен терм е празно. Всякакви две оценки съвпадат върху това множество. Предходното твърдение за зависимостта само от променливите на терма тогава дава еднакви стойности при двете оценки.
\end{proof}

\section{Семантика на формулите}\label{s:21.3}

\subsection{Вярност при оценка}\label{s:21.3.1}

За формула $\varphi$, структура $\mathcal A$ и оценка $v$ записваме
\[
\mathcal A\models\varphi[v]
\]
и четем: „$\varphi$ е вярна в $\mathcal A$ при оценка $v$“. Релацията се определя рекурсивно.

За атомарна формула:
\[
\mathcal A\models p(\tau_1,\ldots,\tau_n)[v]
\quad\Longleftrightarrow\quad
\bigl(
\llbracket\tau_1\rrbracket^{\mathcal A}_v,\ldots,
\llbracket\tau_n\rrbracket^{\mathcal A}_v
\bigr)\in p^{\mathcal A}.
\]
При наличие на формално равенство:
\[
\mathcal A\models(\tau_1\doteq\tau_2)[v]
\quad\Longleftrightarrow\quad
\llbracket\tau_1\rrbracket^{\mathcal A}_v
=
\llbracket\tau_2\rrbracket^{\mathcal A}_v.
\]

За логическите връзки се използва класическата двузначна семантика:
\[
\mathcal A\models\neg\varphi[v]
\quad\Longleftrightarrow\quad
\mathcal A\not\models\varphi[v],
\]
\[
\mathcal A\models(\varphi\wedge\psi)[v]
\quad\Longleftrightarrow\quad
\mathcal A\models\varphi[v]\ \text{и}\
\mathcal A\models\psi[v],
\]
\[
\mathcal A\models(\varphi\vee\psi)[v]
\quad\Longleftrightarrow\quad
\mathcal A\models\varphi[v]\ \text{или}\
\mathcal A\models\psi[v],
\]
\[
\mathcal A\models(\varphi\Rightarrow\psi)[v]
\quad\Longleftrightarrow\quad
\mathcal A\not\models\varphi[v]\ \text{или}\
\mathcal A\models\psi[v].
\]

За кванторите:
\[
\mathcal A\models\forall x\,\varphi[v]
\quad\Longleftrightarrow\quad
\text{за всяко }a\in A:
\mathcal A\models\varphi[v_x^a],
\]
\[
\mathcal A\models\exists x\,\varphi[v]
\quad\Longleftrightarrow\quad
\text{съществува }a\in A:
\mathcal A\models\varphi[v_x^a].
\]

\begin{example}
Нека $p^{\mathcal A}\subseteq A$ е едноместна релация. Тогава
\[
\mathcal A\models\forall x\,p(x)[v]
\]
означава, че $p^{\mathcal A}=A$, а
\[
\mathcal A\models\exists x\,p(x)[v]
\]
означава, че $p^{\mathcal A}\ne\varnothing$.
\end{example}

\subsection{Зависимост само от свободните променливи}\label{s:21.3.2}

\begin{keythm}
Нека $\varphi$ е формула, а $v_1,v_2$ са оценки в структурата $\mathcal A$. Ако
\[
v_1(x)=v_2(x)
\quad\text{за всяко }
x\in\operatorname{Var}^{\mathrm{free}}[\varphi],
\]
то
\[
\mathcal A\models\varphi[v_1]
\quad\Longleftrightarrow\quad
\mathcal A\models\varphi[v_2].
\]
\end{keythm}

\begin{proof}
Доказваме с индукция по построението на $\varphi$.

Нека $\varphi=p(\tau_1,\ldots,\tau_n)$. Всички променливи на $\tau_i$ са свободни във формулата. По твърдението за термовете
\[
\llbracket\tau_i\rrbracket^{\mathcal A}_{v_1}
=
\llbracket\tau_i\rrbracket^{\mathcal A}_{v_2}
\qquad (1\leq i\leq n).
\]
Следователно едната наредена $n$-торка принадлежи на $p^{\mathcal A}$ точно тогава, когато принадлежи другата. Случаят с формално равенство е аналогичен.

При отрицание и при двоичните логически връзки твърдението следва непосредствено от индукционната хипотеза и класическата семантика.

Нека
\[
\varphi=\forall x\,\psi.
\]
Тогава свободните променливи на $\varphi$ са свободните променливи на $\psi$, различни от $x$. За произволно $a\in A$ оценките $v_{1\,x}^a$ и $v_{2\,x}^a$ съвпадат върху всички свободни променливи на $\psi$: за $x$ и двете дават $a$, а за останалите следва от условието. По индукционната хипотеза
\[
\mathcal A\models\psi[v_{1\,x}^a]
\Longleftrightarrow
\mathcal A\models\psi[v_{2\,x}^a].
\]
Следователно всеобщото квантифициране дава еднаква стойност. Случаят $\exists x\,\psi$ е аналогичен.
\end{proof}

\begin{cor}
Ако $\varphi$ е затворена формула, нейната вярност в дадена структура не зависи от оценката.
\end{cor}

\begin{proof}
За затворена формула $\operatorname{FV}(\varphi)=\varnothing$. Всякакви две оценки съвпадат върху свободните й променливи. По доказаната теорема за независимост от останалите променливи истинностната стойност е една и съща.
\end{proof}

\subsection{Тъждествена вярност и предикатни тавтологии}\label{s:21.3.3}

\begin{defn}
Формулата $\varphi$ е \emph{тъждествено вярна в структурата} $\mathcal A$, ако
\[
\mathcal A\models\varphi[v]
\]
за всяка оценка $v$ в $\mathcal A$. Пишем
\[
\mathcal A\models\varphi.
\]
\end{defn}

\begin{defn}
Формулата $\varphi$ е \emph{предикатна тавтология} или \emph{общовалидна}, ако е тъждествено вярна във всяка структура за езика:
\[
\models\varphi.
\]
\end{defn}

\begin{remark}
За затворена формула е достатъчно да се провери една произволна оценка: по предходното следствие всички оценки дават една и съща истинностна стойност.
\end{remark}

\section{Субституции}\label{s:21.4}

\subsection{Субституция на термове}\label{s:21.4.1}

\begin{defn}
Субституция е съпоставяне, което на всяка индивидуална променлива задава терм. Често се задава крайно чрез
\[
s=\left\{\frac{x_1}{t_1},\ldots,\frac{x_n}{t_n}\right\},
\]
където $x_1,\ldots,x_n$ са различни променливи. За терм $\tau$ с $\tau s$ означаваме терма, получен при едновременно заместване на съответните променливи.
\end{defn}

\begin{lem}[Лема за субституциите на термове]
Нека $\tau(x_1,\ldots,x_n)$ е терм, $t_1,\ldots,t_n$ са термове, а $v$ е оценка в $\mathcal A$. Ако $\omega$ е оценка, зададена чрез
\[
\omega(x_i)=\llbracket t_i\rrbracket^{\mathcal A}_v
\quad (1\leq i\leq n),
\]
и съвпадаща с $v$ върху останалите променливи, то
\[
\llbracket\tau[x_1/t_1,\ldots,x_n/t_n]\rrbracket^{\mathcal A}_v
=
\llbracket\tau\rrbracket^{\mathcal A}_\omega.
\]
\end{lem}

\begin{proof}
Доказателството е индукция по построението на $\tau$. За $\tau=x_i$ равенството следва от дефиницията на $\omega$; за променлива, която не се заменя, и за константа е непосредствено. При
\[
\tau=f(\tau_1,\ldots,\tau_k)
\]
се прилага индукционната хипотеза към всеки аргумент и после дефиницията на стойността на функционален терм.
\end{proof}

\subsection{Допустимост на заместването във формули}\label{s:21.4.2}

При замяна във формула се заменят само свободните участия на съответната променлива. Трябва да се избягва \emph{захващане} на променлива.

\begin{defn}
Променлива $y$ е опасна за формула $\varphi$ при субституция $s$, ако $y$ участва в терма $s(x)$ за някоя свободна променлива $x$ на $\varphi$, за която $s(x)\ne x$.
\end{defn}

Заместването е допустимо, когато свободно участие, заменено с терм, не попада в област на действие на квантор по променлива, участваща в този терм.

\begin{example}
Нека
\[
\varphi=\forall y\,(q(x,y)\Rightarrow p(y)).
\]
Замяната на свободното $x$ с $f(y)$ не е допустима: полученото
\[
\forall y\,(q(f(y),y)\Rightarrow p(y))
\]
променя ролята на $y$, защото то става свързано от квантора.
\end{example}

\begin{defn}[Допустима субституция]
Субституцията $s$ е \emph{допустима} за формулата $\varphi$, ако за всяка
свободна променлива $x$ на $\varphi$ с $s(x)\ne x$ никое свободно участие на $x$
в $\varphi$ не попада в областта на действие на квантор по променлива, която
участва в терма $s(x)$. Тоест никоя променлива на заместващия терм не се
свързва при заместването.
\end{defn}

Конспектът иска твърдението за стойността при замяна на променливи с термове да
бъде \emph{формулирано} в общия си вид и да бъде \emph{доказано} в частния случай
на безкванторна формула. Затова тук стоят и двете: общата формулировка, която
трябва да може да се възпроизведе точно, и подробната индукция за безкванторния
случай.

\begin{keythm}[Лема за субституциите, общ вид]
Нека $\varphi$ е произволна формула, $s$ е субституция, \emph{допустима} за
$\varphi$, $v$ е оценка в структурата $\mathcal A$, а $\omega$ е оценката
\[
\omega(x)=\llbracket xs\rrbracket^{\mathcal A}_v
\]
за всяка променлива $x$. Тогава
\[
\mathcal A\models(\varphi s)[v]
\quad\Longleftrightarrow\quad
\mathcal A\models\varphi[\omega].
\]
\end{keythm}

\begin{remark}
Изискването за допустимост не може да се пропусне. За
$\varphi=\exists y\,(x\ne y)$ и $s(x)=y$ лявата страна дава
$\exists y\,(y\ne y)$, което е невярно във всяка структура, докато дясната е
вярна във всяка структура с поне два елемента. Точно това е захващането на
променлива от предходния пример.
\end{remark}

\begin{proof}
Индукция по построението на $\varphi$. Атомарните случаи и случаите с
$\neg,\wedge,\vee,\Rightarrow,\Leftrightarrow$ не използват допустимостта и са
изписани подробно в теоремата за безкванторния случай по-долу; повторени тук,
те дават базата и пропозиционалните стъпки.

Остава случаят с квантор. Нека $\varphi=\forall y\,\psi$ и нека $s_y$ е
субституцията, която съвпада с $s$ извън $y$ и полага $s_y(y)=y$: заместват се
само свободните участия, а $y$ е свързана в $\varphi$. Тогава
$\varphi s=\forall y\,(\psi s_y)$ и $s_y$ е допустима за $\psi$. По дефиниция на
семантиката
\[
\mathcal A\models(\varphi s)[v]
\quad\Longleftrightarrow\quad
\text{за всяко }a\in A:\ \mathcal A\models(\psi s_y)[v_y^a].
\]
Прилагаме индукционната хипотеза за $\psi$, $s_y$ и оценката $v_y^a$. Тя дава
\[
\mathcal A\models(\psi s_y)[v_y^a]
\quad\Longleftrightarrow\quad
\mathcal A\models\psi[\omega_a],
\qquad
\omega_a(x)=\llbracket xs_y\rrbracket^{\mathcal A}_{v_y^a}.
\]
Сравняваме $\omega_a$ с $\omega_y^a$ върху свободните променливи на $\psi$.
За $x=y$ имаме $ys_y=y$, тоест $\omega_a(y)=a=\omega_y^a(y)$. За свободна
променлива $x\ne y$ на $\psi$, която е свободна и в $\varphi$, имаме
$xs_y=xs$. Ако $s(x)=x$, стойността е $v(x)$ и в двата случая. Ако
$s(x)\ne x$, допустимостта на $s$ за $\varphi$ казва, че $y$ не участва в
терма $s(x)$; тогава по твърдението за зависимостта на стойността на терм само
от участващите в него променливи
\[
\llbracket s(x)\rrbracket^{\mathcal A}_{v_y^a}
=
\llbracket s(x)\rrbracket^{\mathcal A}_{v}
=\omega(x)=\omega_y^a(x).
\]
Следователно $\omega_a$ и $\omega_y^a$ съвпадат върху
$\operatorname{Var}^{\mathrm{free}}[\psi]$ и по теоремата за зависимост само от
свободните променливи
\[
\mathcal A\models\psi[\omega_a]
\quad\Longleftrightarrow\quad
\mathcal A\models\psi[\omega_y^a].
\]
Събирайки еквивалентностите за всяко $a\in A$, получаваме
\[
\mathcal A\models(\varphi s)[v]
\ \Longleftrightarrow\
\text{за всяко }a\in A:\ \mathcal A\models\psi[\omega_y^a]
\ \Longleftrightarrow\
\mathcal A\models\forall y\,\psi\,[\omega].
\]
Случаят $\varphi=\exists y\,\psi$ се получава по същия начин, като „за всяко
$a$“ се замени със „съществува $a$“.
\end{proof}

\begin{keythm}[Лема за субституциите, безкванторен случай]
Нека $\varphi$ е безкванторна формула, $s$ е субституция, $v$ е оценка в $\mathcal A$, а $\omega$ е оценката
\[
\omega(x)=\llbracket xs\rrbracket^{\mathcal A}_v
\]
за всяка променлива $x$. Тогава
\[
\mathcal A\models(\varphi s)[v]
\quad\Longleftrightarrow\quad
\mathcal A\models\varphi[\omega].
\]
\end{keythm}

Това е \emph{частният случай}, чието доказателство конспектът изисква. При
безкванторна $\varphi$ всяка субституция е допустима, защото няма квантор, който
да свърже променлива, затова хипотезата за допустимост отпада.

\begin{proof}
Доказваме с индукция по построението на безкванторната формула $\varphi$.

Нека
\[
\varphi=p(\tau_1,\ldots,\tau_n).
\]
По лемата за субституциите на термове:
\[
\llbracket\tau_i s\rrbracket^{\mathcal A}_v
=
\llbracket\tau_i\rrbracket^{\mathcal A}_\omega
\qquad (1\leq i\leq n).
\]
Следователно
\[
\begin{aligned}
\mathcal A\models(\varphi s)[v]
&\Longleftrightarrow
\bigl(
\llbracket\tau_1s\rrbracket^{\mathcal A}_v,\ldots,
\llbracket\tau_ns\rrbracket^{\mathcal A}_v
\bigr)\in p^{\mathcal A}\\
&\Longleftrightarrow
\bigl(
\llbracket\tau_1\rrbracket^{\mathcal A}_\omega,\ldots,
\llbracket\tau_n\rrbracket^{\mathcal A}_\omega
\bigr)\in p^{\mathcal A}\\
&\Longleftrightarrow
\mathcal A\models\varphi[\omega].
\end{aligned}
\]
Случаят с равенство е аналогичен.

Нека $\varphi=\neg\psi$. По индукционната хипотеза
\[
\mathcal A\models(\psi s)[v]
\Longleftrightarrow
\mathcal A\models\psi[\omega].
\]
След отрицание получаваме търсената еквивалентност за $\neg\psi$.

Нека
\[
\varphi=(\psi\wedge\chi).
\]
Тогава
\[
(\psi\wedge\chi)s=(\psi s\wedge\chi s).
\]
По индукционната хипотеза:
\[
\mathcal A\models(\psi s)[v]\Longleftrightarrow\mathcal A\models\psi[\omega],
\]
\[
\mathcal A\models(\chi s)[v]\Longleftrightarrow\mathcal A\models\chi[\omega].
\]
Следователно
\[
\mathcal A\models((\psi\wedge\chi)s)[v]
\Longleftrightarrow
\mathcal A\models(\psi\wedge\chi)[\omega].
\]
Случаите за $\vee$, $\Rightarrow$ и $\Leftrightarrow$ са аналогични. Понеже $\varphi$ е безкванторна, случаи с квантори няма.
\end{proof}

\section{Определимост и изпълнимост}\label{s:21.5}

\subsection{Определими свойства}\label{s:21.5.1}

Нека $\varphi[x_1,\ldots,x_n]$ означава, че всички свободни променливи на $\varphi$ са сред $x_1,\ldots,x_n$. Формулата определя в структурата $\mathcal A$ множеството
\[
D_\varphi=
\left\{
(a_1,\ldots,a_n)\in A^n
\mid
\mathcal A\models\varphi[v]
\text{ за оценка }v(x_i)=a_i
\right\}.
\]
Това е коректно определено именно защото истинността на формулата зависи само от свободните й променливи.

\begin{example}
В структурата
\[
\mathcal N=\langle\mathbb N,0,+,<\rangle
\]
формулата
\[
\exists z\,(x=z+z)
\]
определя множеството на четните естествени числа.
\end{example}

\subsection{Изпълнимост на множества от формули}\label{s:21.5.2}

\begin{defn}
Множество $\Gamma$ от формули е \emph{изпълнимо в структурата} $\mathcal A$, ако съществува оценка $v$, за която
\[
\mathcal A\models\psi[v]
\qquad\text{за всяка }\psi\in\Gamma.
\]
\end{defn}

\begin{defn}
Структурата $\mathcal A$ е \emph{модел} на $\Gamma$, ако
\[
\mathcal A\models\psi
\qquad\text{за всяка }\psi\in\Gamma.
\]
Пишем
\[
\mathcal A\models\Gamma.
\]
\end{defn}

\begin{example}
Нека езикът съдържа едноместен предикат $P$ и двуместен предикат $R$. Множеството
\[
\Gamma=
\left\{
\forall x\,\exists y\,R(x,y),\
\forall x\,\bigl(P(x)\Rightarrow\exists y\,R(y,x)\bigr),\
\exists x\,P(x)
\right\}
\]
е изпълнимо. Модел е структура с универсум $A=\{a\}$, за която
\[
P^{\mathcal A}=\{a\},
\qquad
R^{\mathcal A}=\{(a,a)\}.
\]
Действително, всяка квантифицирана променлива може да получи стойност $a$, а всички формули от $\Gamma$ са верни.
\end{example}

\begin{supp}[Бележка]
При формули със свободни променливи трябва да се различават:
\[
\text{„има една оценка, която удовлетворява всички формули“}
\]
и
\[
\text{„всички формули са тъждествено верни в структурата“}.
\]
Второто е по-силно условие. За затворени формули двете понятия съвпадат, защото истинността не зависи от оценката.
\end{supp}

\section{Изпитен фокус}\label{s:21.6}

\begin{itemize}
\item Да се дефинират език от първи ред, арност, терм, атомарна формула и предикатна формула.
\item Да се обяснят област на действие на квантор, свободно и свързано участие, свободна променлива и затворена формула.
\item Да се дефинират структура, универсум, интерпретация, оценка и модифицирана оценка $v_x^a$.
\item Да се даде рекурсивната семантика на термовете и формулите, особено на $\forall$ и $\exists$.
\item Да се различават вярност при оценка $\mathcal A\models\varphi[v]$, тъждествена вярност $\mathcal A\models\varphi$ и предикатна тавтология $\models\varphi$.
\item Да се възпроизведе доказателствената идея, че стойността на формула зависи само от свободните й променливи: индукция по построението на формулата; при кванторите се сравняват модифицираните оценки.
\item Да се формулира лемата за субституциите на термове.
\item Да се формулира лемата за субституциите в \emph{общия} й вид --- за произволна формула $\varphi$ и субституция $s$, \emph{допустима} за $\varphi$: $\mathcal A\models(\varphi s)[v]\Leftrightarrow\mathcal A\models\varphi[\omega]$, където $\omega(x)=\llbracket xs\rrbracket^{\mathcal A}_v$. Конспектът иска формулировката в общия вид.
\item Да се докаже лемата за субституциите в частния случай на \emph{безкванторна} формула чрез индукция по построението й. Само този частен случай се изисква с доказателство.
\item Да се обясни защо е нужна допустимост на заместването и да се даде пример за захващане на променлива.
\item Да се решават задачи за определимост чрез формула и за изпълнимост чрез явно посочване на структура, интерпретации и при нужда оценка.
\end{itemize}


\chapter{Изводимост и компютърно генериране на доказателства.}
%% Drafted from the mapped corpus by scripts/draft_topics.py.
% Model: gpt-5.6-terra; source characters: 269101.
\section{Основни понятия}

Методът на резолюциите е доказателствен метод, който свежда проверката за логическо следване до проверка за неизпълнимост. Идеята е проста: за да се докаже, че от множество формули $\Gamma$ следва формулата $\varphi$, разглеждаме множеството
\[
\Gamma\cup\{\neg\varphi\}.
\]
Ако то е неизпълнимо, то няма структура, в която всички формули от $\Gamma$ са верни, а $\varphi$ е невярна; следователно $\Gamma\models\varphi$. Резолюцията търси синтактично свидетелство за тази неизпълнимост: извеждане на празния дизюнкт.

\begin{defn}
Литерал е атомарна формула или отрицание на атомарна формула. Ако $L$ е литерал, то противоположният му литерал ще означаваме с $\overline L$:
\[
\overline{P(t_1,\ldots,t_n)}=\neg P(t_1,\ldots,t_n),
\qquad
\overline{\neg P(t_1,\ldots,t_n)}=P(t_1,\ldots,t_n).
\]
\end{defn}

\begin{defn}
Елементарна дизюнкция е безкванторна формула, изградена от литерали само с дизюнкция. Формула е в \emph{конюнктивна нормална форма} (КНФ), ако е конюнкция от елементарни дизюнкции.
\end{defn}

В резолютивния метод не е необходимо да се пазят редът и повторенията на литералите в една дизюнкция. Затова дизюнкциите се представят като множества.

\begin{defn}
Дизюнкт е крайно множество от литерали. Допуска се и празният дизюнкт, означаван с $\Box$ или $\varnothing$.
\end{defn}

Например формулата
\[
P(x)\vee\neg Q(f(x))\vee P(x)
\]
съответства на дизюнкта
\[
\{P(x),\neg Q(f(x))\}.
\]
Празният дизюнкт няма литерали и съответства на лъжа: той е неверен при всяка оценка.

\begin{defn}
Нека $\mathcal M$ е структура, а $v$ е оценка на индивидуалните променливи. Дизюнктът $D$ е \emph{верен в $\mathcal M$ при $v$}, ако поне един негов литерал е верен:
\[
\mathcal M\models_v D
\quad\Longleftrightarrow\quad
(\exists L\in D)\ \mathcal M\models_v L.
\]
Той е неверен в $\mathcal M$ при $v$, ако всички негови литерали са неверни.
\end{defn}

\begin{defn}
Дизюнктът $D$ е \emph{тъждествено верен в структурата} $\mathcal M$, ако
\[
(\forall v)\ \mathcal M\models_v D.
\]
Множество от дизюнкти $\Gamma$ е \emph{изпълнимо}, ако съществува структура $\mathcal M$, в която всеки дизюнкт от $\Gamma$ е тъждествено верен:
\[
\mathcal M\models\Gamma
\quad\Longleftrightarrow\quad
(\forall D\in\Gamma)\,(\forall v)\ \mathcal M\models_v D.
\]
Тогава $\mathcal M$ се нарича модел на $\Gamma$.
\end{defn}

Важно е да се различават дизюнктът $\{p,\neg p\}$ и множеството от формули $\{p,\neg p\}$. Първият е тавтологичен дизюнкт, понеже съдържа противоположни литерали. Второто множество е неизпълнимо, понеже изисква едновременно $p$ и $\neg p$ да са верни.

\begin{remark}
Ако формула е приведена в КНФ
\[
C_1\wedge\cdots\wedge C_m,
\]
то на нея съответства множеството от дизюнкти $\{D_1,\ldots,D_m\}$, където $D_i$ е множеството от литералите в $C_i$. Изпълнимостта на формулата е еквивалентна на изпълнимостта на съответното множество от дизюнкти.
\end{remark}

\section{Резолюция за дизюнкти}

\subsection{Същинска резолвента}

Най-напред разглеждаме случая, в който два дизюнкта съдържат буквално противоположни литерали.

\begin{defn}
Нека $D_1$ и $D_2$ са дизюнкти, $L\in D_1$ и $\overline L\in D_2$. Дизюнктът
\[
R_L(D_1,D_2)
=
(D_1\setminus\{L\})\cup(D_2\setminus\{\overline L\})
\]
се нарича \emph{непосредствена резолвента} на $D_1$ и $D_2$ относно литерала $L$.
\end{defn}

Например:
\[
\{p,q\},\qquad \{\neg p,r\}
\]
имат резолвента
\[
\{q,r\}.
\]
При резолюция се премахва точно една двойка противоположни литерали. Ако два дизюнкта съдържат повече от една противоположна двойка, не се премахват всички двойки едновременно.

\begin{example}
За дизюнктите
\[
D_1=\{p,q,r\},
\qquad
D_2=\{\neg p,\neg q,r\}
\]
резолвентите са
\[
\{q,\neg q,r\}
\quad\text{и}\quad
\{p,\neg p,r\},
\]
но не и $\{r\}$.
\end{example}

\begin{lem}[Резолютивна лема]
Нека $D$ е резолвента на $D_1$ и $D_2$. За всяка булева интерпретация $I$ е вярно:
\[
I\models D_1\ \text{и}\ I\models D_2
\quad\Longrightarrow\quad
I\models D.
\]
\end{lem}

\begin{proof}
Нека
\[
D=(D_1\setminus\{L\})\cup(D_2\setminus\{\overline L\}).
\]
Разглеждаме два случая.

Ако $I\models L$, то $I\not\models\overline L$. Понеже $I\models D_2$, някой литерал от $D_2$, различен от $\overline L$, трябва да е верен. Той принадлежи на $D$, следователно $I\models D$.

Ако $I\not\models L$, то понеже $I\models D_1$, някой литерал от $D_1$, различен от $L$, е верен. Той принадлежи на $D$, следователно $I\models D$.
\end{proof}

\subsection{Резолютивен извод}

\begin{defn}
Нека $\Gamma$ е множество от дизюнкти. \emph{Резолютивен извод} на дизюнкта $\delta$ от $\Gamma$ е крайна редица
\[
\delta_0,\delta_1,\ldots,\delta_n,
\qquad \delta_n=\delta,
\]
такава че за всеки $\delta_i$ е изпълнено поне едно от условията:

\begin{enumerate}
\item $\delta_i\in\Gamma$;
\item $\delta_i$ е резолвента на два дизюнкта, срещащи се по-рано в редицата.
\end{enumerate}

Казваме, че $\delta$ се извежда от $\Gamma$ с резолюция и пишем
\[
\Gamma\vdash_r\delta.
\]
\end{defn}

\begin{keythm}[Коректност на резолютивната изводимост]
Ако
\[
\Gamma\vdash_r\Box,
\]
то $\Gamma$ е неизпълнимо.
\end{keythm}

\begin{proof}
Нека
\[
\delta_0,\delta_1,\ldots,\delta_n=\Box
\]
е резолютивен извод от $\Gamma$. Ще докажем по индукция по $i$, че всеки модел на $\Gamma$ удовлетворява $\delta_i$.

Нека $\mathcal M\models\Gamma$. Ако $\delta_i\in\Gamma$, то по дефиниция
\[
\mathcal M\models\delta_i.
\]
Ако $\delta_i$ е резолвента на два по-ранни дизюнкта $\delta_j$ и $\delta_k$, то по индукционната хипотеза
\[
\mathcal M\models\delta_j
\quad\text{и}\quad
\mathcal M\models\delta_k.
\]
От резолютивната лема следва
\[
\mathcal M\models\delta_i.
\]

Следователно всеки модел на $\Gamma$ би трябвало да удовлетворява и $\Box$. Но празният дизюнкт няма верен литерал и е неверен при всяка интерпретация. Следователно $\Gamma$ няма модел.
\end{proof}

\begin{example}
Нека
\[
\Gamma=
\bigl\{
\{p,q\},
\{\neg p,q\},
\{p,\neg q\},
\{\neg p,\neg q\}
\bigr\}.
\]
Получаваме извода
\[
\begin{array}{rll}
1.&\{p,\neg q\} & \in\Gamma,\\
2.&\{p,q\} & \in\Gamma,\\
3.&\{p\} & \text{резолвента на 1 и 2 относно }q,\\
4.&\{\neg p,\neg q\} & \in\Gamma,\\
5.&\{\neg p,q\} & \in\Gamma,\\
6.&\{\neg p\} & \text{резолвента на 4 и 5 относно }q,\\
7.&\Box & \text{резолвента на 3 и 6 относно }p.
\end{array}
\]
Следователно $\Gamma$ е неизпълнимо.
\end{example}

\section{Предикатна либерална резолюция}

При предикатните дизюнкти противоположните литерали невинаги са синтактично еднакви. Например $P(x)$ и $\neg P(f(y))$ могат да станат противоположни след подходяща субституция. Затова е необходима унификация.

\subsection{Субституции и унификация}

\begin{defn}
Субституция е крайно множество
\[
\sigma=
\left\{
\frac{x_1}{t_1},\ldots,\frac{x_n}{t_n}
\right\},
\]
където $x_1,\ldots,x_n$ са различни индивидуални променливи, а $t_1,\ldots,t_n$ са термове. Резултатът от прилагането на $\sigma$ върху дизюнкта $\varepsilon$ се означава с $\varepsilon\sigma$.
\end{defn}

\begin{defn}
Субституцията $\sigma$ е унификатор на термовете $t_1,\ldots,t_m$, ако
\[
t_1\sigma=\cdots=t_m\sigma.
\]
Тя е най-общ унификатор, ако всеки друг унификатор се получава от нея чрез последваща субституция.
\end{defn}

При унификацията трябва да се спазва проверката за участие: не може да се замества $x$ с терм, в който участва самото $x$. Например уравнението
\[
x=f(x)
\]
няма допустим унификатор.

\begin{example}
Литералите
\[
P(x,f(y))
\quad\text{и}\quad
\neg P(f(z),f(a))
\]
се унифицират чрез
\[
\sigma=
\left\{
\frac{x}{f(z)},
\frac{y}{a}
\right\}.
\]
След прилагане на $\sigma$ получаваме противоположните литерали
\[
P(f(z),f(a))
\quad\text{и}\quad
\neg P(f(z),f(a)).
\]
\end{example}

\subsection{Предикатна либерална резолвента}

\begin{defn}
Нека $\varepsilon_1$ и $\varepsilon_2$ са дизюнкти. Първо техните променливи се преименуват така, че двата дизюнкта да нямат общи променливи. Нека
\[
P(t_1,\ldots,t_n)\in\varepsilon_1,
\qquad
\neg P(s_1,\ldots,s_n)\in\varepsilon_2,
\]
а $\sigma$ е най-общ унификатор на съответните аргументи. Тогава дизюнктът
\[
\left(
(\varepsilon_1\setminus\{P(t_1,\ldots,t_n)\})
\cup
(\varepsilon_2\setminus\{\neg P(s_1,\ldots,s_n)\})
\right)\sigma
\]
се нарича \emph{предикатна либерална резолвента} на $\varepsilon_1$ и $\varepsilon_2$.
\end{defn}

Терминът ``либерална'' подчертава, че преди резолюцията литералите могат да се направят противоположни чрез унификация, а не е необходимо да са буквално противоположни още в изходните дизюнкти.

\begin{example}
Нека
\[
\varepsilon_1=\{P(x),Q(x)\},
\qquad
\varepsilon_2=\{\neg P(f(y)),R(y)\}.
\]
За $P(x)$ и $\neg P(f(y))$ най-общ унификатор е
\[
\sigma=\left\{\frac{x}{f(y)}\right\}.
\]
Следователно либералната резолвента е
\[
\{Q(f(y)),R(y)\}.
\]
\end{example}

\begin{defn}
Нека $\Gamma$ е множество от предикатни дизюнкти. \emph{Предикатен резолютивен извод} на $\delta$ от $\Gamma$ е крайна редица от дизюнкти, в която всеки член е или елемент на $\Gamma$, или предикатна либерална резолвента на два по-ранни члена. Отново пишем
\[
\Gamma\vdash_r\delta.
\]
\end{defn}

\begin{supp}[Бележка]
В някои представяния предикатната резолюция се описва чрез дизюнкти с ограничения, които записват равенствата между аргументите на резолвираните литерали. След елиминиране на ограниченията чрез алгоритъм за унификация се получава същата резолвента с приложена субституция. Тук използваме непосредственото описание чрез най-общ унификатор.
\end{supp}

\subsection{Коректност и пълнота}

\begin{keythm}[Коректност на предикатната резолютивна изводимост]
Ако
\[
\Gamma\vdash_r\Box,
\]
то множеството от предикатни дизюнкти $\Gamma$ е неизпълнимо.
\end{keythm}

Доказателствената идея е същата като при същинската резолюция. Ако дадена структура удовлетворява двата родителски дизюнкта при всички оценки, то след прилагане на унифициращата субституция тя удовлетворява и резолвентата. Следователно всеки дизюнкт в резолютивния извод е логическо следствие на изходното множество. Появата на $\Box$ противоречи на съществуването на модел.

\begin{keythm}[Пълнота на резолютивната изводимост]
За множество $\Gamma$ от първоредови дизюнкти в език без формално равенство е в сила:
\[
\Gamma\ \text{е неизпълнимо}
\quad\Longleftrightarrow\quad
\Gamma\vdash_r\Box.
\]
\end{keythm}

При същинската резолюция пълнотата може да се изрази чрез резолютивното затваряне. Нека
\[
R(S)=S\cup
\{D\mid D\text{ е резолвента на два дизюнкта от }S\},
\]
\[
S_0=S,\qquad S_{n+1}=R(S_n),\qquad
S^*=\bigcup_{n\geq 0}S_n.
\]
Тогава
\[
S\vdash_rD
\quad\Longleftrightarrow\quad
D\in S^*,
\]
а в частност
\[
S\text{ е неизпълнимо}
\quad\Longleftrightarrow\quad
\Box\in S^*
\quad\Longleftrightarrow\quad
S\vdash_r\Box.
\]

За предикатния случай пълнотата е свързана с теоремата на Ербран. След сколемизация и преминаване към затворени универсални формули се разглеждат техните затворени частни случаи. Неизпълнимостта на първоредовото множество се свежда до булева неизпълнимост на подходящо крайно множество от такива случаи, за което резолюцията е пълна.

Ако равенството е част от логическия език със стандартната си семантика, обикновената резолюция сама не е достатъчна за тази формулировка на пълнотата. Необходими са подходящи аксиоми за равенство или специализирани правила като параметодулация/суперпозиция.

\begin{remark}
Пълнотата не означава, че всяка стратегия за избор на двойки дизюнкти непременно ще намери празния дизюнкт за краен брой стъпки. Теоремата твърди, че при неизпълнимо множество съществува краен резолютивен извод на $\Box$.
\end{remark}

\section{Компютърно генериране на доказателства}

Резолюцията е особено подходяща за автоматизация, защото всяка стъпка има краен и проверим вид: посочват се два родителски дизюнкта, избрани литерали и унификатор. Така доказателството може да се пази като дърво или като линейна последователност от стъпки.

Общият процес има следните етапи:

\begin{enumerate}
\item Формулира се задачата като проверка на неизпълнимост, обикновено чрез добавяне на отрицанието на целта.
\item Отстраняват се импликациите и еквивалентностите.
\item Отрицанията се свеждат непосредствено пред атомарни формули.
\item Формулата се привежда в пренексна форма и се сколемизира, когато има квантори за съществуване.
\item Отпадат универсалните квантори и матрицата се привежда в КНФ.
\item Генерират се резолвенти чрез унификация.
\item При получаване на $\Box$ се извежда неизпълнимостта на изходното множество.
\end{enumerate}

\begin{example}[Предикатно доказване на неизпълнимост]
Нека са дадени формулите
\[
\forall x\,\bigl(Student(x)\Rightarrow Learns(x)\bigr),
\]
\[
Student(ivan),
\]
\[
\neg Learns(ivan).
\]
След преминаване към дизюнкти получаваме:
\[
D_1=\{\neg Student(x),Learns(x)\},
\qquad
D_2=\{Student(ivan)\},
\qquad
D_3=\{\neg Learns(ivan)\}.
\]
Резолюцията на $D_1$ и $D_2$ използва унификатора
\[
\sigma=\left\{\frac{x}{ivan}\right\}
\]
и дава
\[
D_4=\{Learns(ivan)\}.
\]
След това $D_4$ и $D_3$ дават
\[
\Box.
\]
Следователно трите изходни формули са неизпълними като множество.
\end{example}

\section{Хорнови дизюнкти и Пролог}

\begin{defn}
Хорнов дизюнкт е дизюнкт с най-много един положителен литерал.
\end{defn}

Най-често се разглеждат следните три вида:

\[
\{P\}
\]
е факт;

\[
\{P,\neg Q_1,\ldots,\neg Q_n\}
\]
е правило, което се чете като
\[
Q_1\wedge\cdots\wedge Q_n\Rightarrow P;
\]

\[
\{\neg Q_1,\ldots,\neg Q_n\}
\]
е цел, която съответства на въпроса дали могат да се докажат едновременно $Q_1,\ldots,Q_n$.

Хорновите правила и факти образуват хорнова програма. Резолюцията за хорнови дизюнкти е основата на логическото програмиране и на езика Пролог. Вместо да се генерират произволни резолвенти, обикновено се работи с цел и с правило, чиято глава се унифицира с избран подцелеви атом.

\begin{example}[Дефиниране на предикати на Пролог]
Следната програма описва предиката \verb|predok(X,Y)|: ``\verb|X| е предшественик на \verb|Y|''.

\begin{verbatim}
roditel(ivan, maria).
roditel(maria, petar).

predok(X, Y) :- roditel(X, Y).
predok(X, Y) :- roditel(X, Z), predok(Z, Y).
\end{verbatim}

Фактът \verb|roditel(ivan,maria).| съответства на едноелементния дизюнкт
\[
\{roditel(ivan,maria)\}.
\]
Правилото
\begin{verbatim}
predok(X, Y) :- roditel(X, Z), predok(Z, Y).
\end{verbatim}
съответства на хорновия дизюнкт
\[
\{
predok(X,Y),
\neg roditel(X,Z),
\neg predok(Z,Y)
\}.
\]
\end{example}

\begin{example}[Резолютивно доказване на цел]
Нека целта е
\begin{verbatim}
?- predok(ivan, petar).
\end{verbatim}
Тя се представя с отрицателния дизюнкт
\[
\{\neg predok(ivan,petar)\}.
\]
Резолюция с рекурсивното правило дава нова цел
\[
\{
\neg roditel(ivan,Z),
\neg predok(Z,petar)
\}.
\]
Фактът $roditel(ivan,maria)$ унифицира $Z$ с $maria$, откъдето получаваме
\[
\{\neg predok(maria,petar)\}.
\]
Резолюция с основното правило за $predok$ дава
\[
\{\neg roditel(maria,petar)\}.
\]
Този литерал се резолвира с факта $roditel(maria,petar)$ и се получава $\Box$. Следователно целта е следствие на програмата.
\end{example}

\begin{remark}
В Пролог обикновено се използва SLD-резолюция: на всяка стъпка се избира подцел, унифицира се с главата на правило или факт и се заменя с тялото на правилото. Натрупаните унификатори дават отговора за променливите в заявката.
\end{remark}

\section{Изпитен фокус}

\begin{itemize}
\item Да се дефинират: литерал, елементарна дизюнкция, КНФ, дизюнкт, верен дизюнкт в структура при оценка, тъждествено верен дизюнкт и изпълнимо множество от дизюнкти.
\item Да се знае, че $\Box$ е празният дизюнкт и е неверен при всяка интерпретация.
\item Да се дефинира непосредствена резолвента:
\[
R_L(D_1,D_2)=
(D_1\setminus\{L\})\cup(D_2\setminus\{\overline L\}).
\]
\item Да се дефинират субституция, унификатор и най-общ унификатор.
\item Да се дефинира предикатна либерална резолвента чрез преименуване на променливите, унификация на противоположни атоми и прилагане на унификатора към остатъка от дизюнктите.
\item Да се дефинира резолютивен извод и означението $\Gamma\vdash_r\delta$.
\item Да се формулират теоремите за коректност и пълнота:
\[
\Gamma\vdash_r\Box\Rightarrow\Gamma\text{ е неизпълнимо},
\qquad
\Gamma\text{ е неизпълнимо}\Rightarrow\Gamma\vdash_r\Box.
\]
Пълнотата в този вид се отнася за първоредов език без формално равенство; при равенство са нужни допълнителни аксиоми или правила.
\item Да се възпроизведе доказателството за коректност: индукция по стъпките на извода, използваща резолютивната лема; празният дизюнкт не може да бъде удовлетворен.
\item Да се обясни ролята на Ербрановите частни случаи и унификацията при предикатната резолюция.
\item Да се решава задача за превеждане на предикатни формули в дизюнкти и за извеждане на $\Box$.
\item Да се разпознават факти, правила и цели като хорнови дизюнкти и да се обясни връзката им с Пролог и SLD-резолюцията.
\end{itemize}


\chapter{Бази от данни. Релационен модел на данните.}
%% Drafted from the mapped corpus by scripts/draft_topics.py.
% Model: gpt-5.6-luna; source characters: 14340.
\section{Бази от данни}

\subsection{Предназначение на базите от данни}

База от данни е организирана колекция от данни, предназначена за съхраняване, обработване и извличане на информация. Данните се описват чрез определена структура, а операциите върху тях се изпълняват посредством специализиран език за работа с бази от данни.

Релационната база от данни представлява колекция от релации. В практическа реализация релациите обикновено се представят като таблици, но математическото им разглеждане се основава на множества от кортежи. Релационният модел дава възможност данните да бъдат описани чрез формална схема и да бъдат обработвани чрез операции от релационната алгебра или чрез заявки на SQL.

Проектирането на релационна база от данни включва следните основни стъпки:

\begin{enumerate}
    \item определяне на данните, които ще се съхраняват;
    \item определяне на връзките между обектите;
    \item определяне на ключовете и свързващите колони;
    \item определяне на ограниченията върху обектите и връзките между тях;
    \item отстраняване на евентуални недостатъци и излишества;
    \item реализиране на базата от данни.
\end{enumerate}

Тези стъпки водят от описание на предметната област към конкретна структура от релации и атрибути.

\section{Релационен модел на данните}

\subsection{Математическа основа}

Релационният модел се основава на математическото понятие за $n$-членна релация. Една $n$-членна релация е множество от елементи, всеки от които се състои от $n$ компоненти. Тези елементи се наричат $n$-торки или кортежи.

Чрез една релация се моделира даден клас от обекти, а всеки кортеж от релацията представя конкретен обект от този клас. Например релация за филми може да описва класа от всички филми, като всеки неин кортеж представя един конкретен филм.

\begin{defn}
Домейн е именувано множество, разглеждано като множество от допустими стойности на дадена величина.
\end{defn}

Домейнът определя какви стойности могат да се използват за даден атрибут. Например атрибутът \texttt{title} може да има домейн от низове, а атрибутът \texttt{year} може да има домейн от цели числа:

\[
\texttt{title}:\texttt{string}, \qquad
\texttt{year}:\texttt{integer}.
\]

\begin{defn}
В контекста на базите от данни релация е таблица, чиито редове са кортежи, а стойностите в една колона принадлежат на домейна, асоцииран със съответния атрибут.
\end{defn}

Релацията е множество от уникални кортежи. Следователно в математическия релационен модел един и същ кортеж не се повтаря.

\begin{defn}
Атрибут е име на колона в релацията, с което е свързан определен домейн.
\end{defn}

Например релацията за филми може да има атрибути \texttt{title}, \texttt{year}, \texttt{length} и \texttt{filmType}. Съответните домейни могат да бъдат низове, цели числа и други прости типове.

\begin{defn}
Кортеж е ред от релацията, съдържащ конкретна стойност за всеки атрибут.
\end{defn}

Ако атрибутите са подредени като
\[
(\texttt{title},\texttt{year},\texttt{length},\texttt{filmType}),
\]
примерен кортеж е
\[
(\texttt{Star\ wars},1977,124,\texttt{Color}).
\]

Компонентите на всеки кортеж принадлежат на домейните на съответните атрибути. Кортежите следват предварително определената подредба на атрибутите от схемата на релацията.

Релационният модел изисква стойностите на компонентите да бъдат атомарни. Това означава, че във всяка позиция на кортежа се съдържа една проста стойност, например цяло число или низ. Сложни стойности като списъци и масиви не се разглеждат като един компонент на кортежа в основния релационен модел.

\subsection{Схема и екземпляр на релация}

\begin{defn}
Схема на релация е името на релацията заедно с множеството от нейните атрибути.
\end{defn}

Схемата се записва чрез името на релацията, последвано от имената на атрибутите в скоби:

\[
\textit{Movies}(\textit{title},\textit{year},\textit{length},\textit{filmType}).
\]

В математическото описание атрибутите образуват множество. При записването на схемата обаче се фиксира стандартна подредба, за да бъде еднозначно определена позицията на всеки компонент в кортежа.

\begin{defn}
Екземпляр на релация е множеството от кортежи, които в даден момент се съдържат в релацията.
\end{defn}

Схемата описва структурата на релацията, докато екземплярът описва нейното текущо съдържание. Схемата обикновено се променя по-рядко, а екземплярът може да се изменя при добавяне, промяна или изтриване на данни.

\begin{defn}
Кардиналност на екземпляр на релация е броят на кортежите в този екземпляр.
\end{defn}

Ако релацията \textit{Movies} съдържа $120$ кортежа, нейната кардиналност за съответния момент е $120$.

\begin{defn}
Схема на релационна база от данни е съвкупността от схемите на всички релации в базата от данни.
\end{defn}

Например една база от данни може да има следната схема:

\[
\begin{aligned}
\textit{Movie}(&\textit{Title},\textit{Year},\textit{Length}),\\
\textit{StarsIn}(&\textit{MovieTitle},\textit{MovieYear},\textit{StarName}).
\end{aligned}
\]

Подчертаването се използва, за да се означат атрибути, които в дадения пример участват в идентифицирането на кортежите. Така в \textit{Movie} заглавието, годината и дължината са показани като определящи атрибути, а в \textit{StarsIn} са показани \textit{MovieTitle} и \textit{StarName}. Връзката между двете релации се изразява чрез съответните атрибути за филм и година.

\section{Операции върху релационна база от данни}

Операциите върху релационната база от данни могат да бъдат разгледани според това дали описват структурата, обработват данните или управляват достъпа и запазването на промените.

\subsection{DDL операции}

Data Definition Language, или DDL, е подезик за описание и промяна на схемата на базата от данни. Основните DDL команди са \texttt{CREATE}, \texttt{ALTER} и \texttt{DROP}.

Пример за създаване на таблица е:

\begin{verbatim}
CREATE TABLE movies (
  id INTEGER PRIMARY KEY,
  title CHAR(50) NOT NULL,
  length INTEGER NULL
);
\end{verbatim}

Тук се създава таблица \texttt{movies} с три колони. Колоната \texttt{id} е от тип \texttt{INTEGER} и е обявена като първичен ключ. Колоната \texttt{title} съдържа низ с дължина до 50 символа и е обявена като задължителна чрез \texttt{NOT NULL}. За \texttt{length} е разрешена празна стойност чрез \texttt{NULL}.

Промяна на схемата може да се извърши чрез \texttt{ALTER TABLE}:

\begin{verbatim}
ALTER TABLE movies ADD release_date DATE;
ALTER TABLE movies DROP COLUMN length;
\end{verbatim}

Първата команда добавя атрибута \texttt{release\_date}, а втората премахва атрибута \texttt{length}. Изтриването на цялата таблица се извършва чрез:

\begin{verbatim}
DROP TABLE movies;
\end{verbatim}

При изпълнение на DDL операция промените върху схемата се прилагат незабавно.

\subsection{DML операции}

Data Manipulation Language, или DML, е подезик за работа със съдържанието на релациите. Основните DML команди са \texttt{SELECT}, \texttt{INSERT}, \texttt{UPDATE} и \texttt{DELETE}.

Командата \texttt{SELECT} извлича данни:

\begin{verbatim}
SELECT title FROM movies
WHERE length IS NOT NULL AND length < 120
ORDER BY release_date DESC
\end{verbatim}

Заявката извлича заглавията на филмите, чиято дължина не е празна и е по-малка от $120$. Резултатът се подрежда по \texttt{release\_date} в низходящ ред.

Добавяне на кортеж се извършва с \texttt{INSERT}:

\begin{verbatim}
INSERT INTO movies (title, length, release_date)
VALUES ('Star wars', 124, '09-04-1984')
\end{verbatim}

Промяна на вече съществуващи данни се извършва с \texttt{UPDATE}:

\begin{verbatim}
UPDATE movies
SET title = 'Star Wars'
WHERE title = 'Star wars'
\end{verbatim}

Изтриване на кортежи се извършва чрез \texttt{DELETE}:

\begin{verbatim}
DELETE FROM movies
WHERE length < 120
\end{verbatim}

След изпълнение на DML операция промените трябва да бъдат потвърдени с \texttt{COMMIT}, за да станат постоянни.

\subsection{Заявки към релационната база от данни}

Заявка е описание на желан резултат, получен чрез обработване на една или повече релации. В SQL заявките най-често се използва командата \texttt{SELECT}, която може да включва избиране на колони, условие за подбор на редове и подреждане на резултата.

В релационната алгебра заявката се представя чрез алгебричен израз. Всеки оператор приема релация или релации и връща нова релация. По този начин сложните заявки могат да се изграждат чрез комбиниране на по-прости операции.

\section{Релационна алгебра}

\begin{defn}
Релационната алгебра е междинен език за изчисляване на заявки, който задава правила за обработване на релации чрез алгебрични изрази.
\end{defn}

Операциите в релационната алгебра се изпълняват върху релации. В ядрото на релационната алгебра релациите се разглеждат като множества от кортежи. Разширената релационна алгебра разглежда релациите като мултимножества от кортежи, при които могат да се срещат повторения.

Основните класове операции са:

\begin{itemize}
    \item операции върху множества;
    \item операции за отсяване на части от релациите;
    \item операции за комбиниране на кортежи от две релации;
    \item операция за преименуване.
\end{itemize}

Основните операции са селекция, проекция, обединение, разлика, декартово произведение и преименуване. Сечение, тета-съединение и естествено съединение могат да се изразят чрез основните операции и поради това се разглеждат като допълнителни.

\subsection{Обединение}

Обединението на две релации $R$ и $S$ се означава с

\[
R\cup S.
\]

Операцията е приложима, когато релациите имат съвместими схеми: еднакъв брой атрибути, които си съответстват. Резултатът е релация със същата схема, съдържаща всички кортежи, които принадлежат на $R$ или на $S$.

Обединението е бинарна операция, комутативна и асоциативна:

\[
R\cup S=S\cup R,
\]

\[
(R\cup S)\cup T=R\cup(S\cup T).
\]

Повтарящите се кортежи не се включват многократно, тъй като релациите в ядрото на релационната алгебра са множества.

\subsection{Разлика}

Разликата на $R$ и $S$ се означава с

\[
R-S.
\]

Както при обединението, схемите трябва да бъдат съвместими. Резултатът съдържа всички кортежи, които са в $R$, но не са в $S$:

\[
R-S=\{t\mid t\in R \text{ и } t\notin S\}.
\]

Разликата не е комутативна. В общия случай

\[
R-S\neq S-R.
\]

Редът на операндите има значение: първата релация определя множеството, от което се изваждат кортежи.

\subsection{Сечение}

Сечението на $R$ и $S$ се означава с

\[
R\cap S.
\]

Операцията се прилага върху релации със съвместими схеми и връща релация, съдържаща само кортежите, общи за двете релации:

\[
R\cap S=\{t\mid t\in R \text{ и } t\in S\}.
\]

Сечението е комутативно и асоциативно:

\[
R\cap S=S\cap R,
\]

\[
(R\cap S)\cap T=R\cap(S\cap T).
\]

То може да бъде изразено чрез основните операции:

\[
R\cap S=R-(R-S).
\]

\subsection{Проекция}

Проекцията е унарна операция, която избира определени атрибути от релация. Означава се с

\[
\pi_{A_1,A_2,\ldots,A_k}(R),
\]

където $A_1,A_2,\ldots,A_k$ е списък от атрибути.

Резултатът съдържа само посочените колони. Ако след премахването на останалите колони се получат еднакви кортежи, повторенията се премахват. Например

\[
\pi_{\textit{Title},\textit{Year},\textit{Length}}(\textit{Movie})
\]

избира трите посочени атрибута от релацията \textit{Movie}.

Проекцията променя схемата, като намалява броя на атрибутите. Затова тя често се описва като вертикално отсяване на релацията.

\subsection{Селекция}

Селекцията, наричана още рестрикция, е унарна операция, която избира кортежите, удовлетворяващи определено условие. Означава се с

\[
\sigma_C(R),
\]

където $C$ е предикат.

Предикатът може да сравнява два атрибута или атрибут с константа. Допустимите оператори са

\[
=,\quad \neq,\quad <,\quad >,\quad \leq,\quad \geq.
\]

Пример:

\[
\sigma_{\textit{Year}<1980}(\textit{Movie}).
\]

Резултатът съдържа само филмите, чиято година е по-малка от $1980$. Схемата остава същата, но броят на кортежите може да намалее. Поради това селекцията се разглежда като хоризонтално отсяване на релацията.

Селекцията е комутативна при последователно прилагане:

\[
\sigma_{C_1}(\sigma_{C_2}(R))
=
\sigma_{C_2}(\sigma_{C_1}(R)).
\]

\subsection{Декартово произведение}

Декартовото произведение на $R$ и $S$ се означава с

\[
R\times S.
\]

Резултатът съдържа всички двойки от кортежи, при които първият кортеж е от $R$, а вторият е от $S$. Схемата на резултата е обединение на схемите на $R$ и $S$.

Ако $R$ има $m$ кортежа, а $S$ има $n$ кортежа, резултатът съдържа всички възможни комбинации между тях. При еднакви имена на атрибути се използва квалифицирано име от вида

\[
\textit{имеНаРелация}.\textit{имеНаАтрибут}.
\]

Например могат да се използват \textit{Movie.Year} и \textit{StarsIn.MovieYear}.

Декартовото произведение е бинарна операция, комутативна и асоциативна според разглеждането на съответните атрибути и подредбата на схемата.

\subsection{Тета-съединение}

Тета-съединението на $R$ и $S$ при условие $C$ се означава с

\[
R\bowtie_C S.
\]

То се получава чрез декартово произведение, последвано от селекция:

\[
R\bowtie_C S=\sigma_C(R\times S).
\]

Резултатът има схема, получена чрез обединяване на схемите на $R$ и $S$, и съдържа само онези двойки кортежи, които удовлетворяват условието $C$.

Ако условието включва само съвпадение между атрибути, тета-съединението се нарича еквисъединение. Например условието

\[
\textit{Movie.Year}=\textit{StarsIn.MovieYear}
\]

свързва кортежи, за които стойностите на съответните атрибути съвпадат.

\subsection{Естествено съединение}

Естественото съединение се означава с

\[
R\bowtie S.
\]

То свързва релациите по всички атрибути с еднакви имена. Резултатът съдържа кортежите от декартовото произведение, за които стойностите на едноименните атрибути съвпадат. Повтарящите се колони се премахват автоматично.

Естественото съединение може да бъде представено чрез тета-съединение и проекция:

\[
R\bowtie S
=
\pi_L\left(\sigma_C(R\times S)\right),
\]

където $C$ е условието за равенство по едноименните атрибути, а $L$ съдържа всички атрибути на $R$ и онези атрибути на $S$, които не присъстват в $R$.

Ако общите атрибути са $A_1,\ldots,A_n$, условието е от вида

\[
C=(R.A_1=S.A_1)\mathbin{\mathrm{AND}}\cdots
\mathbin{\mathrm{AND}}(R.A_n=S.A_n).
\]

\subsection{Преименуване}

Преименуването се означава с оператора $\rho$. То променя името на релацията и имената на нейните атрибути, без да променя кортежите.

Записът

\[
R_1:=\rho_{R_1(A_1,\ldots,A_n)}(R_2)
\]

означава, че $R_1$ е релация с атрибути $A_1,\ldots,A_n$ и със същите кортежи като $R_2$.

Преименуването е полезно, когато една и съща релация участва повече от веднъж в алгебричен израз или когато имената на атрибутите трябва да бъдат разграничени.

\subsection{Класификация и приоритет}

Основните операции на релационната алгебра са:

\begin{itemize}
    \item обединение;
    \item разлика;
    \item декартово произведение;
    \item проекция;
    \item селекция;
    \item преименуване.
\end{itemize}

Допълнителните операции са:

\begin{itemize}
    \item сечение;
    \item тета-съединение;
    \item естествено съединение.
\end{itemize}

Те могат да се получат от основните операции чрез следните зависимости:

\[
R\cap S=R-(R-S),
\]

\[
R\bowtie_C S=\sigma_C(R\times S),
\]

\[
R\bowtie S=\pi_L\left(\sigma_C(R\times S)\right).
\]

Приоритетът на операторите в алгебричните изрази е:

\begin{enumerate}
    \item унарни оператори -- селекция, проекция и преименуване;
    \item декартово произведение и съединение;
    \item сечение;
    \item обединение и разлика;
    \item изрази в скоби.
\end{enumerate}

При нееднозначност се използват скоби.

\section{Примерни задачи}

\subsection{Съставяне на SQL заявки}

Нека е дадена таблицата \texttt{movies} със следните колони: \texttt{id}, \texttt{title}, \texttt{length} и \texttt{release\_date}.

Задача: да се изведат заглавията на филмите с ненулева дължина, по-малка от $120$, подредени по дата на издаване в низходящ ред.

\begin{verbatim}
SELECT title FROM movies
WHERE length IS NOT NULL AND length < 120
ORDER BY release_date DESC
\end{verbatim}

Задача: да се добави филм.

\begin{verbatim}
INSERT INTO movies (title, length, release_date)
VALUES ('Star wars', 124, '09-04-1984')
\end{verbatim}

Задача: да се промени заглавието на филма.

\begin{verbatim}
UPDATE movies
SET title = 'Star Wars'
WHERE title = 'Star wars'
\end{verbatim}

Задача: да се изтрият филмите с дължина под $120$.

\begin{verbatim}
DELETE FROM movies
WHERE length < 120
\end{verbatim}

След промени чрез \texttt{INSERT}, \texttt{UPDATE} или \texttt{DELETE} се изпълнява:

\begin{verbatim}
COMMIT
\end{verbatim}

\subsection{Пример за DDL}

\begin{verbatim}
CREATE TABLE movies (
  id INTEGER PRIMARY KEY,
  title CHAR(50) NOT NULL,
  length INTEGER NULL
);
\end{verbatim}

Добавяне на атрибут:

\begin{verbatim}
ALTER TABLE movies ADD release_date DATE;
\end{verbatim}

Премахване на атрибут:

\begin{verbatim}
ALTER TABLE movies DROP COLUMN length;
\end{verbatim}

Изтриване на таблицата:

\begin{verbatim}
DROP TABLE movies;
\end{verbatim}

\subsection{Тригери}

Тригерите са свързани с автоматично изпълнение на действие при настъпване на определено събитие върху данни. В предоставения материал не са зададени конкретен синтаксис, събития или примерна реализация на тригер. Поради това тук не се фиксира конкретен SQL вариант за тригер, тъй като синтаксисът зависи от използваната система за управление на бази от данни.

\section{Изпитен фокус}

За устно изложение трябва да могат да бъдат възпроизведени:

\begin{itemize}
    \item определенията за домейн, релация, атрибут, кортеж, схема и екземпляр;
    \item разликата между схема на релация и текущ екземпляр на релация;
    \item определението за схема на релационна база от данни;
    \item изискването компонентите на кортежите да бъдат атомарни;
    \item стъпките при проектиране и реализиране на релационна база от данни;
    \item предназначението на DDL и командите \texttt{CREATE}, \texttt{ALTER}, \texttt{DROP};
    \item предназначението на DML и командите \texttt{SELECT}, \texttt{INSERT}, \texttt{UPDATE}, \texttt{DELETE};
    \item ролята на \texttt{COMMIT} след DML операции;
    \item дефинициите и условията за приложимост на обединение, разлика и сечение;
    \item разликата между селекция и проекция;
    \item определението за декартово произведение;
    \item изразяването на тета-съединение чрез декартово произведение и селекция;
    \item особеността на естественото съединение да свързва по едноименни атрибути и да премахва повтарящите се колони;
    \item ролята на преименуването;
    \item класификацията на основни и допълнителни операции;
    \item зависимостите
    \[
    R\cap S=R-(R-S),\qquad
    R\bowtie_C S=\sigma_C(R\times S);
    \]
    \item приоритетът на операторите в релационната алгебра;
    \item съставянето на примерни SQL заявки за извличане, добавяне, промяна и изтриване на данни;
    \item примерите за DDL операции и предназначението на тригерите.
\end{itemize}


\chapter{Бази от данни. Нормални форми.}
% AUTO-GENERATED by scripts/build_topics.py — do not edit.
% Edit topics/bodies/topic_24.tex or topics/preamble.tex instead.
%% Placeholder. Replace with the written chapter.
\textit{Източници:} PDF \texttt{24.pdf}, снимки \texttt{24та тема}

\subsection*{Анотация от конспекта}
Бази от данни. Нормални форми. Нормални форми. Проектиране схемите на релационните бази от данни. Аномалии, ограничения, ключове. Функционални зависимости, аксиоми на Армстронг. Първа, втора, трета нормална форма, нормална форма на Бойс-Код. Многозначни зависимости; аксиоми на функционалните и многозначните зависимости; съединение без загуба; четвърта нормална форма. Примерни задачи: Привеждане на схема на базата от данни (при зададени функционални зависимости) към зададена нормална форма. Литература [32].


\chapter{Изкуствен интелект. Пространство на състоянията.}
% AUTO-GENERATED by scripts/build_topics.py — do not edit.
% Edit topics/bodies/topic_25.tex or topics/preamble.tex instead.
%% Placeholder. Replace with the written chapter.
\textit{Източници:} PDF \texttt{25.pdf}, снимки \texttt{25та тема}, \texttt{Задачи/AI.txt}

\subsection*{Анотация от конспекта}
Изкуствен интелект: Пространство на състоянията – определение, характеристики на състоянията, търсене в пространство на състояния Локално търсещи алгоритми: • Локално търсене в лъч; • Изкачване на хълмове; • Симулирано втвърдяване; • Генетични алгоритми. Задачи за удовлетворяване на ограничения: определение, търсене с възврат, намаляване на конфликтите. Избор на следващ ход при игра за двама играчи: мини-макс процедура, алфа-бета отсичане. Литература: [41].


\chapter{Съвременни софтуерни технологии.}
% AUTO-GENERATED by scripts/build_topics.py — do not edit.
% Edit topics/bodies/topic_26.tex or topics/preamble.tex instead.
%% Placeholder. Replace with the written chapter.
\textit{Източници:} PDF \texttt{26.pdf}, снимки \texttt{26та тема}, \texttt{ST.pdf}

\subsection*{Анотация от конспекта}
Съвременни софтуерни технологии. Софтуерен продукт и процес. Основни фази на софтуерните процеси. Модел на софтуерен процес – модел на водопада, модел на бързата разработка, еволюционни модели, прототипен модел, спираловиден модел. Сравнителна характеристика на моделите на софтуерни процеси. Гъвкави методи - Extreme Programming (XP) и SCRUM. Изисквания към софтуерните системи. Функционални и нефункционални изисквания. Анализ и проектиране на софтуерните изисквания. Езици за описание. UML. Верификация и валидация на софтуера. Управление на качеството. Литература [9], [39], [44].


\chapter{Архитектури на софтуерни системи.}
% AUTO-GENERATED by scripts/build_topics.py — do not edit.
% Edit topics/bodies/topic_27.tex or topics/preamble.tex instead.
%% Placeholder. Replace with the written chapter.
\textit{Източници:} PDF \texttt{27.pdf}, снимки \texttt{27ма тема}, \texttt{SA.pdf}, \texttt{SA-SI.pdf}

\subsection*{Анотация от конспекта}
Архитектури на софтуерни системи. a. Понятие за софтуерна архитектура. Структури и изгледи (structures and views) на архитектурата. Необходимост от софтуерни архитектури. Влияние на архитектурата върху проекта и организацията. b. Изисквания към качеството (нефункционални изисквания) на системата. c. Архитектурни стилове. Многослоен стил, Модел-изглед контролер (MVC), Поточни (Pipe-and-filter) архитектури. Архитектура ориентирана към услуги. d. Проектиране на софтуерната архитектура. Процес за проектиране. Избор на подходящи структури. Последователност на създаване на архитектурата. Тактики (архитектурни решения) за постигане на желаните качествени показатели. e. Документиране на софтуерната архитектура. Предназначение и основни елементи в документацията на архитектурата. Литература: [30], [52].


\part{Математика и приложения}

\chapter{Уравнения на права в равнината.}
%% Drafted from the mapped corpus by scripts/draft_topics.py.
% Model: gpt-5.6-terra; source characters: 29744.
\section{Параметрични уравнения на права}

Нека в афинната равнина е фиксирана афинна координатна система
$K=Oxy$. Права $g$ се определя еднозначно от точка $M_0$ от нея и
ненулев вектор $\vec p$, колинеарен на правата. Векторът $\vec p$ се
нарича направляващ вектор на $g$.

\begin{defn}
Нека $M_0\in g$ и $\vec p\neq\vec 0$, $\vec p\parallel g$. Векторното
параметрично уравнение на правата $g$ е
\[
 \vec r=\vec r_0+s\vec p,\qquad s\in\mathbb R,
\]
където $\vec r=\overrightarrow{OM}$ е радиус-векторът на текущата точка
$M$, а $\vec r_0=\overrightarrow{OM_0}$ е радиус-векторът на фиксираната
точка $M_0$.
\end{defn}

Действително,
\[
 M\in g
 \quad\Longleftrightarrow\quad
 \overrightarrow{M_0M}\parallel\vec p
 \quad\Longleftrightarrow\quad
 \exists s\in\mathbb R:\ \overrightarrow{M_0M}=s\vec p.
\]
Понеже
\[
 \overrightarrow{M_0M}
 =\overrightarrow{OM}-\overrightarrow{OM_0}
 =\vec r-\vec r_0,
\]
получаваме векторното параметрично уравнение.

\begin{defn}
Нека $M_0(x_0,y_0)$ и $\vec p=(p_1,p_2)\neq(0,0)$. Координатните
(скаларни) параметрични уравнения на правата през $M_0$, колинеарна на
$\vec p$, са
\[
 g:\quad
 \begin{cases}
 x=x_0+sp_1,\\
 y=y_0+sp_2,
 \end{cases}
 \qquad s\in\mathbb R.
\]
\end{defn}

Параметърът $s$ задава положението на точката върху правата. При $s=0$
се получава началната точка $M_0$. Различни параметризации могат да
задават една и съща права: може да се избере друга точка от правата или
ненулев колинеарен на $\vec p$ вектор.

\begin{example}
Правата през $M_0(2,-1)$ с направляващ вектор
$\vec p=(3,2)$ има уравнения
\[
 \vec r=(2,-1)+s(3,2),\qquad s\in\mathbb R,
\]
съответно
\[
 \begin{cases}
 x=2+3s,\\
 y=-1+2s,
 \end{cases}
 \qquad s\in\mathbb R.
\]
Например при $s=1$ се получава точката $(5,1)$.
\end{example}

\section{Общо уравнение на права}

\begin{keythm}[Теорема за общото уравнение на права]
Всяка права в равнината има уравнение от вида
\[
 Ax+By+C=0,\qquad (A,B)\neq(0,0).
\]
Обратно, всяко уравнение от този вид определя точно една права в
равнината.
\end{keythm}

\begin{proof}
Нека правата $g$ минава през $M_0(x_0,y_0)$ и има направляващ вектор
$\vec p=(p_1,p_2)\neq(0,0)$. Точка $M(x,y)$ лежи върху $g$ тогава и само
тогава, когато векторите
\[
 \overrightarrow{M_0M}=(x-x_0,y-y_0)
 \quad\text{и}\quad
 \vec p=(p_1,p_2)
\]
са колинеарни. Координатното условие за колинеарност е
\[
 \begin{vmatrix}
 x-x_0 & y-y_0\\
 p_1 & p_2
 \end{vmatrix}=0.
\]
Следователно
\[
 p_2(x-x_0)-p_1(y-y_0)=0,
\]
тоест
\[
 p_2x-p_1y-p_2x_0+p_1y_0=0.
\]
При полагането
\[
 A=p_2,\qquad B=-p_1,\qquad C=-p_2x_0+p_1y_0
\]
получаваме уравнение от търсения вид. Понеже $\vec p\neq\vec0$, то
$(A,B)=(p_2,-p_1)\neq(0,0)$.

Обратно, нека е дадено
\[
 Ax+By+C=0,\qquad (A,B)\neq(0,0).
\]
Нека $M_0(x_0,y_0)$ е негова точка. Разглеждаме вектора
\[
 \vec p=(-B,A)\neq\vec0.
\]
Съществува единствена права през $M_0$, колинеарна на $\vec p$. Нейните
параметрични уравнения са
\[
 x=x_0-sB,\qquad y=y_0+sA,\qquad s\in\mathbb R.
\]
Оттук
\[
 A(x-x_0)+B(y-y_0)=0.
\]
Тъй като $Ax_0+By_0+C=0$, последното уравнение е еквивалентно на
$Ax+By+C=0$. Следователно даденото уравнение задава точно тази права.
\end{proof}

\begin{defn}
Уравнението
\[
 g:\ Ax+By+C=0,\qquad (A,B)\neq(0,0),
\]
се нарича общо уравнение на правата $g$ спрямо координатната система
$K=Oxy$.
\end{defn}

Коефициентите на общото уравнение дават два важни вектора:
\[
 \vec n=(A,B)\perp g,\qquad
 \vec p=(-B,A)\parallel g.
\]
Векторът $\vec n$ се нарича нормален вектор на правата, а
$\vec p$ е един от нейните направляващи вектори.

\begin{prop}
Нека
\[
 g:\ Ax+By+C=0,\qquad (A,B)\neq(0,0).
\]
Векторът $\vec q=(q_1,q_2)$ е колинеарен на $g$ тогава и само тогава,
когато
\[
 Aq_1+Bq_2=0.
\]
\end{prop}

\begin{proof}
Направляващ вектор на $g$ е $\vec p=(-B,A)$. Условието
$\vec q\parallel\vec p$ е
\[
 \begin{vmatrix}
 -B & A\\
 q_1 & q_2
 \end{vmatrix}=0,
\]
което е равносилно на $Aq_1+Bq_2=0$.
\end{proof}

\begin{example}
Да се намери общото уравнение на правата през $M_0(1,2)$ с
направляващ вектор $\vec p=(4,-3)$.

От доказателството на теоремата получаваме
\[
 A=p_2=-3,\qquad B=-p_1=-4,
\]
а
\[
 C=-p_2x_0+p_1y_0=3+8=11.
\]
Следователно
\[
 -3x-4y+11=0,
\]
или еквивалентно
\[
 3x+4y-11=0.
\]
\end{example}

\section{Взаимно положение на две прави}

Нека са дадени правите
\[
 g_1:\ A_1x+B_1y+C_1=0,\qquad
 g_2:\ A_2x+B_2y+C_2=0,
\]
където
\[
 (A_i,B_i)\neq(0,0),\qquad i=1,2.
\]
Техните общи точки са решенията на системата
\[
 \begin{cases}
 A_1x+B_1y+C_1=0,\\
 A_2x+B_2y+C_2=0.
 \end{cases}
\]
Определителят на системата по неизвестните е
\[
 \Delta=
 \begin{vmatrix}
 A_1&B_1\\
 A_2&B_2
 \end{vmatrix}
 =A_1B_2-A_2B_1.
\]

\begin{keythm}
За взаимното положение на $g_1$ и $g_2$ са в сила следните критерии.

\begin{enumerate}
 \item Правите се пресичат в точно една точка тогава и само тогава,
 когато
 \[
 \Delta\neq0.
 \]

 \item Правите съвпадат тогава и само тогава, когато съществува
 $k\neq0$, за което
 \[
 A_1=kA_2,\qquad B_1=kB_2,\qquad C_1=kC_2.
 \]

 \item Правите са различни успоредни тогава и само тогава, когато
 съществува $k\neq0$, за което
 \[
 A_1=kA_2,\qquad B_1=kB_2,\qquad C_1\neq kC_2.
 \]
\end{enumerate}
\end{keythm}

\begin{proof}
Ако $\Delta\neq0$, системата има единствено решение, следователно
правите имат точно една обща точка.

Ако $\Delta=0$, двойките $(A_1,B_1)$ и $(A_2,B_2)$ са пропорционални.
Тогава или и свободните членове са пропорционални със същия коефициент,
което прави двете уравнения еквивалентни, или не са. В първия случай
правите съвпадат, а във втория системата е несъвместима и правите са
различни успоредни.
\end{proof}

\begin{remark}
Записите с отношения, например
\[
 \frac{A_1}{A_2}=\frac{B_1}{B_2},
\]
са удобни само когато съответните знаменатели са различни от нула.
Критериите чрез общ ненулев множител $k$ или чрез определителя
$\Delta$ нямат това ограничение.
\end{remark}

\begin{supp}[Бележка]
Среща се неточен критерий, според който пресичането се описва с
неравенства между всички отношения на коефициентите. Точният критерий
за пресичане е
\[
 A_1B_2-A_2B_1\neq0.
\]
Свободните членове $C_1,C_2$ определят положението на пресечната точка,
но не влияят на това дали две прави с различни направления се
пресичат.
\end{supp}

\begin{example}
За правите
\[
 g_1:\ 2x-y+1=0,\qquad
 g_2:\ 4x-2y-3=0
\]
е изпълнено
\[
 \Delta=
 \begin{vmatrix}
 2&-1\\
 4&-2
 \end{vmatrix}=0.
\]
Първите два коефициента на $g_2$ са два пъти съответните коефициенти на
$g_1$, но свободният член $-3$ не е два пъти $1$. Следователно правите
са различни успоредни.
\end{example}

\section{Декартово уравнение}

Оттук нататък при метричните понятия приемаме, че $K=Oxy$ е
ортонормирана координатна система.

\begin{defn}
Нека
\[
 g:\ Ax+By+C=0,\qquad B\neq0.
\]
След изразяване на $y$ получаваме
\[
 y=-\frac ABx-\frac CB.
\]
Ако
\[
 k=-\frac AB,\qquad n=-\frac CB,
\]
то уравнението
\[
 g:\ y=kx+n
\]
се нарича декартово уравнение на правата $g$.
\end{defn}

Коефициентът $k$ се нарича ъглов коефициент на правата. Ако
ориентираният ъгъл между положителната посока на $Ox$ и правата е
$\alpha$, то
\[
 k=\tan\alpha.
\]
Един направляващ вектор на правата $y=kx+n$ е $(1,k)$.

\begin{prop}
Права има декартово уравнение $y=kx+n$ тогава и само тогава, когато не
е успоредна на координатната ос $Oy$.
\end{prop}

\begin{proof}
Декартовото уравнение се получава точно когато $B\neq0$. Ако $B=0$,
общото уравнение има вид $Ax+C=0$, тоест описва права, успоредна на
$Oy$. Обратно, когато $B\neq0$, $y$ се изразява еднозначно чрез $x$.
\end{proof}

\begin{remark}
Правата $y=kx+n$ пресича оста $Oy$ в точката $(0,n)$. Тя пресича
оста $Ox$, ако $k\neq0$, в точката
\[
 \left(-\frac nk,0\right).
\]
\end{remark}

\begin{supp}[Бележка]
В някои материали точката $(0,n)$ е означена като пресечна точка с
$Ox$. Това е несъответствие: при уравнение $y=kx+n$ тази точка лежи на
оста $Oy$.
\end{supp}

\begin{example}
Правата
\[
 3x-2y+6=0
\]
има декартово уравнение
\[
 y=\frac32x+3.
\]
Тя има ъглов коефициент $k=\frac32$ и пресича $Oy$ в точката $(0,3)$.
\end{example}

\section{Нормално уравнение и разстояние от точка до права}

Нека
\[
 g:\ Ax+By+C=0,\qquad (A,B)\neq(0,0).
\]
Нормален вектор на правата е $\vec n=(A,B)$, а неговата дължина е
\[
 |\vec n|=\sqrt{A^2+B^2}.
\]
Следователно един от двата единични нормални вектора е
\[
 \vec n_1=
 \left(
 \frac{A}{\sqrt{A^2+B^2}},
 \frac{B}{\sqrt{A^2+B^2}}
 \right).
\]

\begin{defn}
Нормално уравнение на правата $g$ се нарича общо уравнение на същата
права, при което коефициентите пред $x$ и $y$ са координати на единичен
нормален вектор. Нормалните уравнения на $g$ са
\[
 \frac{Ax+By+C}{\sqrt{A^2+B^2}}=0
 \qquad\text{и}\qquad
 -\frac{Ax+By+C}{\sqrt{A^2+B^2}}=0.
\]
\end{defn}

Следователно всяка права има точно две нормални уравнения. Те се
получават чрез двата противоположно насочени единични нормални вектора.

\begin{keythm}[Разстояние от точка до права]
Нека
\[
 g:\ Ax+By+C=0,\qquad (A,B)\neq(0,0),
\]
и $P(x_1,y_1)$ е точка в равнината. Тогава разстоянието от $P$ до $g$
е
\[
 d(P,g)=
 \frac{|Ax_1+By_1+C|}{\sqrt{A^2+B^2}}.
\]
\end{keythm}

\begin{proof}
Нека нормалното уравнение на $g$ е
\[
 A_1x+B_1y+C_1=0,
\]
където $A_1^2+B_1^2=1$. Нека $H(x_H,y_H)$ е ортогоналната проекция на
$P$ върху $g$. Тогава $\overrightarrow{HP}$ е колинеарен на единичния
нормален вектор $(A_1,B_1)$, тоест за точно едно $\delta\in\mathbb R$
е изпълнено
\[
 \overrightarrow{HP}=\delta(A_1,B_1).
\]
Следователно
\[
 x_H=x_1-\delta A_1,\qquad
 y_H=y_1-\delta B_1.
\]
Тъй като $H\in g$, имаме
\[
 A_1x_H+B_1y_H+C_1=0.
\]
След заместване и използване на $A_1^2+B_1^2=1$ получаваме
\[
 \delta=A_1x_1+B_1y_1+C_1.
\]
Числото $\delta$ е ориентираното разстояние, а търсеното разстояние е
неговата абсолютна стойност. След заместване на $A_1,B_1,C_1$ се
получава посочената формула.
\end{proof}

\begin{remark}
Ако уравнението на правата вече е нормално, тоест
\[
 A^2+B^2=1,
\]
формулата се опростява до
\[
 d(P,g)=|Ax_1+By_1+C|.
\]
Без абсолютна стойност изразът
\[
 \delta(P,g)=\frac{Ax_1+By_1+C}{\sqrt{A^2+B^2}}
\]
е ориентираното разстояние, определено от избрания единичен нормален
вектор.
\end{remark}

\begin{example}
Да се намери разстоянието от $P(1,-2)$ до правата
\[
 g:\ 3x+4y-5=0.
\]
Изчисляваме
\[
 d(P,g)=
 \frac{|3\cdot1+4\cdot(-2)-5|}{\sqrt{3^2+4^2}}
 =\frac{|-10|}{5}=2.
\]
\end{example}

\section{Полуравнини}

Всяка права разделя равнината на две части, наречени отворени
полуравнини. Нека
\[
 g:\ Ax+By+C=0,\qquad (A,B)\neq(0,0),
\]
и въведем линейната функция
\[
 l(x,y)=Ax+By+C.
\]
За точките от правата е изпълнено $l(x,y)=0$. За всяка точка извън
правата стойността на $l(x,y)$ е или положителна, или отрицателна.

\begin{keythm}
Нека $M_1(x_1,y_1)$ и $M_2(x_2,y_2)$ не лежат на правата $g$. Те лежат
в различни полуравнини спрямо $g$ тогава и само тогава, когато
\[
 l(x_1,y_1)\,l(x_2,y_2)<0.
\]
\end{keythm}

\begin{proof}
Нека първо $M_1$ и $M_2$ са в различни полуравнини. Тогава отсечката
$M_1M_2$ пресича $g$ в точка $M_0(x_0,y_0),$ която е вътрешна за
отсечката. Следователно съществува $k<0$, за което
\[
 \overrightarrow{M_0M_1}=k\overrightarrow{M_0M_2}.
\]
Оттук
\[
 x_0=\frac{x_1-kx_2}{1-k},\qquad
 y_0=\frac{y_1-ky_2}{1-k}.
\]
Тъй като $M_0\in g$, след заместване в $l(x_0,y_0)=0$ получаваме
\[
 l(x_1,y_1)-k\,l(x_2,y_2)=0.
\]
Следователно
\[
 k=\frac{l(x_1,y_1)}{l(x_2,y_2)}<0,
\]
откъдето
\[
 l(x_1,y_1)l(x_2,y_2)<0.
\]

Обратно, нека
\[
 l(x_1,y_1)l(x_2,y_2)<0.
\]
Тогава
\[
 k=\frac{l(x_1,y_1)}{l(x_2,y_2)}<0.
\]
Точката
\[
 M_0\left(
 \frac{x_1-kx_2}{1-k},
 \frac{y_1-ky_2}{1-k}
 \right)
\]
удовлетворява $l(x_0,y_0)=0$, следователно лежи на $g$. Освен това
\[
 \overrightarrow{M_0M_1}
 =k\overrightarrow{M_0M_2},\qquad k<0.
\]
Векторите са противопосочни, значи $M_0$ е вътрешна точка на
отсечката $M_1M_2$. Следователно $M_1$ и $M_2$ са в различни
полуравнини.
\end{proof}

\begin{defn}
Двете отворени полуравнини с граница $g$ се задават аналитично чрез
\[
 \lambda_g=\{M(x,y):Ax+By+C>0\}
\]
и
\[
 \overline{\lambda}_g=\{M(x,y):Ax+By+C<0\}.
\]
\end{defn}

Знакът на $l(x,y)$ е постоянен във всяка от двете полуравнини. Това
дава практичен начин за проверка от коя страна на права се намира
дадена точка.

\begin{example}
Нека
\[
 g:\ 2x-y-1=0.
\]
За точките $P(0,0)$ и $Q(2,0)$ имаме
\[
 l(0,0)=-1,\qquad l(2,0)=3.
\]
Техният произведение е отрицателно, следователно $P$ и $Q$ лежат в
различни полуравнини спрямо $g$.
\end{example}

\section{Изпитен фокус}

\begin{itemize}
 \item Да се дефинира права чрез точка $M_0$ и ненулев направляващ
 вектор $\vec p$, както и да се запише векторно-параметричното
 уравнение
 \[
 \vec r=\vec r_0+s\vec p.
 \]

 \item Да се преминава към координатни параметрични уравнения
 \[
 x=x_0+sp_1,\qquad y=y_0+sp_2.
 \]

 \item Да се формулира и докаже теоремата за общото уравнение
 \[
 Ax+By+C=0,\qquad (A,B)\neq(0,0),
 \]
 в двете посоки, чрез условието за колинеарност и вектора
 $(-B,A)$.

 \item Да се разпознават нормален вектор $(A,B)$ и направляващ вектор
 $(-B,A)$ на права с общо уравнение.

 \item Да се определя взаимното положение на две прави чрез
 \[
 \Delta=A_1B_2-A_2B_1
 \]
 и чрез пропорционалността на тройките коефициенти.

 \item Да се извежда декартовото уравнение
 \[
 y=kx+n,\qquad k=-\frac AB,\qquad n=-\frac CB,
 \]
 при $B\neq0$, и да се обясни защо правите, успоредни на $Oy$, нямат
 такава форма.

 \item Да се нормализира общо уравнение на права и да се запишат двете
 нормални уравнения.

 \item Да се възпроизвежда формулата
 \[
 d(P,g)=\frac{|Ax_1+By_1+C|}{\sqrt{A^2+B^2}}
 \]
 и идеята за доказателството чрез ортогоналната проекция на точката
 върху правата.

 \item Да се описват полуравнините чрез знака на
 $l(x,y)=Ax+By+C$ и да се прилага критерият
 \[
 l(x_1,y_1)l(x_2,y_2)<0
 \]
 за точки, разположени в различни полуравнини.
\end{itemize}


\chapter{Уравнения на права и равнина в пространството.}
% AUTO-GENERATED by scripts/build_topics.py — do not edit.
% Edit topics/bodies/topic_29.tex or topics/preamble.tex instead.
%% Placeholder. Replace with the written chapter.
\textit{Източници:} PDF \texttt{32.pdf}, снимки \texttt{33та тема}, \texttt{Задачи/2.Равнина в пространството\_ДИ.pdf}

\begin{supp}[Бележка към картографирането]
shares 32.pdf with topic 28
\end{supp}

\subsection*{Анотация от конспекта}
Уравнения на права и равнина в пространството. Векторно-параметрично уравнение на равнина. Координатни (скаларни) параметрични уравнения на равнина в пространството. Теорема за общо уравнение на равнина. Взаимни положения между две равнини. Нормално уравнение на равнина. Разстояние от точка до равнина. Полупространства. Уравнения на права в пространството. Литература: [20].


\chapter{Симетрични оператори в крайномерни евклидови пространства. Теорема за диагонализация.}
%% Drafted from the mapped corpus by scripts/draft_topics.py.
% Model: gpt-5.6-terra; source characters: 11878.
\begin{konspekt}
\kreq{Официално заглавие}{„Симетрични оператори в крайномерни евклидови пространства. \emph{Основни свойства.} Теорема за диагонализация.“}
\kreq{Дефиниция}{симетричен оператор --- \kr{30.2}; матрицата му спрямо ортонормиран базис --- \kref{30.2.1}.}
\kreq{Доказателство}{всички характеристични корени на симетричен оператор са реални --- \kref{30.3.1}.}
\kreq{Доказателство}{всеки два собствени вектора за различни собствени стойности са ортогонални --- \kref{30.3.2}.}
\kreq{Лема с доказателство}{инвариантност на ортогоналното допълнение --- използва се в теоремата за диагонализация --- \kref{30.4}.}
\kreq{Теорема с доказателство}{съществува ортонормиран базис, в който матрицата на симетричен оператор е диагонална --- \kref{30.5}.}
\kprereq{\kr{30.1} събира понятията от линейната алгебра --- евклидово пространство, скаларно произведение, ортонормиран базис, линеен оператор и матрицата му. Анотацията започва от \kr{30.2}.}
\kreq{Литература}{[18].}
\end{konspekt}

\section{Предварителни понятия}\label{s:30.1}

\begin{supp}[Предварителен материал]
Анотацията към Въпрос 30 започва от определението за симетричен оператор
(\S30.2). Този раздел събира понятията от линейната алгебра, които се използват
по-нататък --- евклидово пространство, скаларно произведение, ортонормиран
базис, линеен оператор и матрицата му. Той е предпоставка, а не част от
изпитния въпрос; при повторение може да се прегледа бегло.
\end{supp}

Понятията в този раздел са от курса по линейна алгебра. Изброяваме ги, защото всяко от тях се използва по-нататък в главата: без тях нито определението за симетричен оператор, нито доказателството на теоремата за диагонализация са напълно записани.

\begin{defn}[Скаларно произведение и евклидово пространство]
Нека $E$ е векторно пространство над $\mathbb{R}$. Изображението
\[
(\cdot,\cdot):E\times E\to\mathbb{R}
\]
се нарича \emph{скаларно произведение}, ако за всички $x,y,z\in E$ и $\alpha\in\mathbb{R}$ са изпълнени:

\begin{itemize}
\item симетричност: $(x,y)=(y,x)$;
\item адитивност и хомогенност по първия аргумент:
\[
(x+z,y)=(x,y)+(z,y),
\qquad
(\alpha x,y)=\alpha(x,y);
\]
\item положителна определеност: $(x,x)\geq 0$, като $(x,x)=0$ точно тогава, когато $x=0$.
\end{itemize}

От симетричността следва същото и по втория аргумент, тоест скаларното произведение е \emph{билинейно}. Двойката $\bigl(E,(\cdot,\cdot)\bigr)$ се нарича \emph{евклидово пространство}; то е крайномерно, ако $\dim E<\infty$.
\end{defn}

Стандартният пример е $\mathbb{R}^n$ със скаларно произведение
\[
(x,y)=\sum_{i=1}^n x_iy_i=x^ty,
\]
което ще наричаме \emph{стандартно евклидово скаларно произведение}.

\begin{defn}[Норма и ортогоналност]
Нека $E$ е евклидово пространство. \emph{Норма} на вектора $x$ е числото
\[
\|x\|=\sqrt{(x,x)},
\]
което е определено, защото $(x,x)\geq0$. Вектор с $\|x\|=1$ се нарича \emph{единичен}; за $x\ne0$ векторът $x/\|x\|$ е единичен и се получава чрез \emph{нормализиране} на $x$.

Векторите $x$ и $y$ са \emph{ортогонални}, което се записва $x\perp y$, ако
\[
(x,y)=0.
\]
Две подмножества $A,B\subseteq E$ са ортогонални, $A\perp B$, ако $x\perp y$ за всички $x\in A$ и $y\in B$.
\end{defn}

\begin{defn}[Ортонормиран базис]
Системата вектори $f_1,\ldots,f_k$ е \emph{ортогонална}, ако $(f_i,f_j)=0$ при $i\ne j$, и \emph{ортонормирана}, ако освен това всички вектори са единични, тоест
\[
(f_i,f_j)=\delta_{ij}=
\begin{cases}
1, & i=j,\\
0, & i\ne j.
\end{cases}
\]
\emph{Ортонормиран базис} на $E$ е базис, който е ортонормирана система.
\end{defn}

Ортонормираният базис е удобен именно защото свежда скаларното произведение до координатно пресмятане: ако $e_1,\ldots,e_n$ е ортонормиран базис и
\[
x=\sum_{i=1}^n x_ie_i,
\qquad
y=\sum_{j=1}^n y_je_j,
\]
то от билинейността
\[
(x,y)=\sum_{i=1}^n\sum_{j=1}^n x_iy_j(e_i,e_j)=\sum_{i=1}^n x_iy_i=x^ty .
\]
В частност $x_i=(x,e_i)$: координатите на вектор спрямо ортонормиран базис са скаларните му произведения с базисните вектори.

\begin{prop}[Съществуване на ортонормиран базис]
Всяко крайномерно евклидово пространство $E$ с $\dim E=n\geq1$ притежава ортонормиран базис.
\end{prop}

\begin{proof}
Нека $v_1,\ldots,v_n$ е произволен базис на $E$. По метода на Грам--Шмит строим индуктивно
\[
w_1=v_1,
\qquad
w_k=v_k-\sum_{i=1}^{k-1}\frac{(v_k,w_i)}{(w_i,w_i)}\,w_i,
\qquad k=2,\ldots,n.
\]
Допускаме индуктивно, че $w_1,\ldots,w_{k-1}$ са ненулеви, попарно ортогонални и
\[
\mathcal{L}(w_1,\ldots,w_{k-1})=\mathcal{L}(v_1,\ldots,v_{k-1}).
\]
Тогава дробите са определени, защото $(w_i,w_i)>0$ за ненулево $w_i$. За $j<k$ пресмятаме
\[
(w_k,w_j)
=(v_k,w_j)-\sum_{i=1}^{k-1}\frac{(v_k,w_i)}{(w_i,w_i)}(w_i,w_j)
=(v_k,w_j)-(v_k,w_j)=0,
\]
тъй като в сумата оцелява само събираемото с $i=j$. Освен това $w_k\ne0$: иначе $v_k$ би лежал в $\mathcal{L}(w_1,\ldots,w_{k-1})=\mathcal{L}(v_1,\ldots,v_{k-1})$, което противоречи на линейната независимост на базиса. От построението следва и равенството на обвивките за $k$.

Накрая полагаме $e_i=w_i/\|w_i\|$. Системата $e_1,\ldots,e_n$ е ортонормирана. Всяка ортонормирана система е линейно независима: ако $\sum_i\alpha_ie_i=0$, скаларното умножение с $e_j$ дава $\alpha_j=0$. Значи $n$ такива вектора образуват базис на $E$.
\end{proof}

\begin{defn}[Линейна обвивка и ортогонално допълнение]
С $\mathcal{L}(v_1,\ldots,v_k)$ означаваме \emph{линейната обвивка} на векторите $v_1,\ldots,v_k$, тоест множеството от всички техни линейни комбинации; това е най-малкото подпространство, което ги съдържа. В частност за $u\ne0$
\[
\mathcal{L}(u)=\{\alpha u\mid\alpha\in\mathbb{R}\}
\]
е правата, породена от $u$.

\emph{Ортогонално допълнение} на подпространството $U\subseteq E$ е множеството
\[
U^\perp=\{x\in E\mid (x,u)=0\ \text{за всяко }u\in U\}.
\]
То е подпространство на $E$, защото условието е линейно по $x$.
\end{defn}

\begin{prop}[Ортогонално разлагане]
Нека $E$ е крайномерно евклидово пространство и $U\subseteq E$ е подпространство. Тогава
\[
E=U\oplus U^\perp,
\qquad
\dim U^\perp=\dim E-\dim U.
\]
\end{prop}

\begin{proof}
Ако $U=\{0\}$, твърдението е очевидно, защото тогава $U^\perp=E$. Нека $\dim U=k\geq1$ и нека $u_1,\ldots,u_k$ е ортонормиран базис на $U$; такъв съществува по предходното твърдение, приложено към $U$ със същото скаларно произведение.

За произволно $x\in E$ полагаме
\[
p=\sum_{i=1}^k (x,u_i)u_i\in U,
\qquad
r=x-p.
\]
За всяко $j$ имаме
\[
(r,u_j)=(x,u_j)-\sum_{i=1}^k (x,u_i)(u_i,u_j)=(x,u_j)-(x,u_j)=0,
\]
защото базисът е ортонормиран. Понеже $r$ е ортогонален на всички $u_j$, по линейност той е ортогонален на всяка тяхна линейна комбинация, тоест $r\in U^\perp$. Следователно $x=p+r$ и
\[
E=U+U^\perp.
\]

Сумата е директна: ако $x\in U\cap U^\perp$, то $x$ е ортогонален на себе си, тоест $(x,x)=0$, откъдето $x=0$ поради положителната определеност. Оттук следва и равенството за размерностите.
\end{proof}

\begin{defn}[Собствена стойност, собствен вектор, собствено подпространство]
Нека $\varphi:E\to E$ е линеен оператор. Числото $\lambda\in\mathbb{R}$ е \emph{собствена стойност} на $\varphi$, ако съществува вектор $x\ne0$, за който
\[
\varphi(x)=\lambda x.
\]
Всеки такъв $x$ се нарича \emph{собствен вектор} за $\lambda$. Множеството
\[
E_\lambda=\{x\in E\mid \varphi(x)=\lambda x\}
\]
е подпространство на $E$ и се нарича \emph{собствено подпространство} за $\lambda$; то се състои от собствените вектори за $\lambda$ и нулевия вектор.
\end{defn}

\begin{defn}[Характеристичен полином и характеристичен корен]
Нека $A\in M_{n\times n}(\mathbb{R})$. Полиномът
\[
\chi_A(\lambda)=\det(A-\lambda I)
\]
от степен $n$ се нарича \emph{характеристичен полином} на $A$. Неговите корени в $\mathbb{C}$ се наричат \emph{характеристични корени}.

Подобните матрици имат един и същ характеристичен полином, затова матриците на един и същ оператор спрямо различни базиси дават един и същ полином. Той се нарича характеристичен полином на оператора $\varphi$.
\end{defn}

Връзката между двете понятия е следната. Ако $A$ е матрицата на $\varphi$ спрямо някакъв базис, то за $\lambda\in\mathbb{R}$
\[
\varphi(x)=\lambda x\ \text{за някое }x\ne0
\iff
(A-\lambda I)X=0\ \text{има ненулево решение}
\iff
\det(A-\lambda I)=0 .
\]
Следователно \emph{собствените стойности на $\varphi$ са точно реалните характеристични корени}. Затова твърдението, че всички характеристични корени на симетричен оператор са реални, е по същество твърдение за пълнотата на спектъра му.

\begin{remark}
По основната теорема на алгебрата всеки полином от степен $n\geq1$ с комплексни коефициенти има корен в $\mathbb{C}$. Затова всеки оператор в пространство с $\dim E\geq1$ има поне един характеристичен корен, но той не е задължително реален; именно това е препятствието, което симетричността премахва.
\end{remark}

\begin{defn}[Инвариантно подпространство и ограничение]
Подпространството $U\subseteq E$ е \emph{инвариантно} относно линейния оператор $\varphi:E\to E$, ако
\[
\varphi(U)\subseteq U.
\]
Тогава \emph{ограничението}
\[
\varphi|_U:U\to U,
\qquad
\varphi|_U(x)=\varphi(x),
\]
е коректно определен линеен оператор в $U$.
\end{defn}

\begin{defn}[Ортогонална матрица]
Матрицата $Q\in M_{n\times n}(\mathbb{R})$ се нарича \emph{ортогонална}, ако
\[
Q^tQ=I,
\]
тоест $Q^{-1}=Q^t$. Еквивалентно, стълбовете на $Q$ образуват ортонормирана система спрямо стандартното скаларно произведение в $\mathbb{R}^n$.
\end{defn}

\section{Симетрични оператори}\label{s:30.2}

Оттук нататък $E$ е крайномерно евклидово пространство над $\mathbb{R}$ със скаларно произведение $(\cdot,\cdot)$ и $\dim E=n$.

\begin{defn}
Линейният оператор $\varphi:E\to E$ се нарича \emph{симетричен}, ако за всеки два вектора $x,y\in E$ е изпълнено
\[
(\varphi(x),y)=(x,\varphi(y)).
\]
\end{defn}

Това определение изразява съгласуваност между действието на оператора и скаларното произведение. При симетричен оператор прехвърлянето на оператора от първия към втория аргумент на скаларното произведение не променя стойността му.

\begin{remark}
В по-общия случай на унитарно пространство над $\mathbb{C}$ аналогичното понятие е \emph{ермитов оператор}. Тогава условието е
\[
\langle\varphi(x),y\rangle=\langle x,\varphi(y)\rangle.
\]
В настоящата глава разглеждаме евклидови пространства над $\mathbb{R}$, поради което използваме термина \emph{симетричен оператор}.
\end{remark}

\subsection{Матрица на симетричен оператор}\label{s:30.2.1}

Нека $e_1,\ldots,e_n$ е ортонормиран базис на $E$ и
\[
A=(a_{ij})\in M_{n\times n}(\mathbb{R})
\]
е матрицата на $\varphi$ спрямо този базис. По определение координатите на $\varphi(e_j)$ са елементите от $j$-тия стълб на $A$, тоест
\[
\varphi(e_j)=\sum_{i=1}^n a_{ij}e_i.
\]

\begin{defn}
Матрицата $A=(a_{ij})\in M_{n\times n}(\mathbb{R})$ се нарича \emph{симетрична}, ако
\[
A^t=A.
\]
Еквивалентно,
\[
a_{ij}=a_{ji}\qquad (1\leq i,j\leq n).
\]
\end{defn}

\begin{prop}
Нека $\varphi:E\to E$ е линеен оператор и $e_1,\ldots,e_n$ е ортонормиран базис на $E$. Тогава следните условия са еквивалентни:

\begin{enumerate}
\item $\varphi$ е симетричен оператор;
\item матрицата на $\varphi$ спрямо базиса $e_1,\ldots,e_n$ е симетрична.
\end{enumerate}
\end{prop}

\begin{proof}
Нека първо $\varphi$ е симетричен оператор и $A=(a_{ij})$ е неговата матрица спрямо дадения ортонормиран базис. За произволни $i,j\in\{1,\ldots,n\}$ имаме
\[
a_{ji}=(\varphi(e_i),e_j).
\]
Тъй като $\varphi$ е симетричен,
\[
(\varphi(e_i),e_j)=(e_i,\varphi(e_j)).
\]
От ортонормираността на базиса следва
\[
(e_i,\varphi(e_j))=a_{ij}.
\]
Следователно
\[
a_{ji}=a_{ij}
\]
за всички $i,j$, откъдето $A^t=A$.

Обратно, нека матрицата $A$ е симетрична. Тогава
\[
(\varphi(e_i),e_j)=a_{ji}=a_{ij}=(e_i,\varphi(e_j))
\]
за всички базисни вектори $e_i,e_j$.

Нека сега
\[
x=\sum_{i=1}^n x_i e_i,\qquad
y=\sum_{j=1}^n y_j e_j.
\]
Използвайки линейността на $\varphi$ и билинейността на скаларното произведение, получаваме
\[
\begin{aligned}
(\varphi(x),y)
&=
\left(
\varphi\left(\sum_{i=1}^n x_i e_i\right),
\sum_{j=1}^n y_j e_j
\right) \\
&=
\sum_{i=1}^n\sum_{j=1}^n
x_i y_j(\varphi(e_i),e_j) \\
&=
\sum_{i=1}^n\sum_{j=1}^n
x_i y_j(e_i,\varphi(e_j)) \\
&=
\left(
\sum_{i=1}^n x_i e_i,
\varphi\left(\sum_{j=1}^n y_j e_j\right)
\right) \\
&=(x,\varphi(y)).
\end{aligned}
\]
Следователно $\varphi$ е симетричен оператор.
\end{proof}

\begin{cor}
Линеен оператор в крайномерно евклидово пространство е симетричен точно тогава, когато матрицата му спрямо произволен ортонормиран базис е симетрична.
\end{cor}

\begin{proof}
В ортонормиран базис скаларното произведение има вид $x^ty$. Ако $A$ е матрицата на оператора, условието за симетричност е $(Ax)^ty=x^tAy$ за всички $x,y$, тоест $x^t(A^t-A)y=0$. Избирането на координатни вектори за $x,y$ показва, че това е еквивалентно на $A^t=A$. Аргументът важи във всеки ортонормиран базис.
\end{proof}

\begin{supp}[Бележка]
За матрици с комплексни елементи условието
\[
\overline{A}^{\,t}=A
\]
описва ермитова, а не непременно симетрична матрица. В реалния случай $\overline{A}=A$, така че това условие се свежда до $A^t=A$.
\end{supp}

\section{Собствени стойности и собствени вектори}\label{s:30.3}

Нека $\varphi:E\to E$ е симетричен оператор. Основните му спектрални свойства са, че характеристичните му корени са реални и че собствените подпространства за различни собствени стойности са ортогонални.

\subsection{Реалност на характеристичните корени}\label{s:30.3.1}

\begin{thm}
Всички характеристични корени на симетричен оператор в крайномерно евклидово пространство са реални числа.
\end{thm}

\begin{proof}
Нека $A$ е матрицата на $\varphi$ спрямо ортонормиран базис на $E$. Съгласно предходното твърдение $A$ е реална симетрична матрица.

Нека $\lambda\in\mathbb{C}$ е характеристичен корен на $A$. Тогава съществува ненулев комплексен стълб-вектор
\[
z=
\begin{pmatrix}
z_1\\
\vdots\\
z_n
\end{pmatrix}
\in\mathbb{C}^n
\]
такъв, че
\[
Az=\lambda z.
\]
Умножаваме отляво с ред-вектора
\[
\overline{z}^{\,t}=
(\overline{z_1},\ldots,\overline{z_n}).
\]
Получаваме
\[
\overline{z}^{\,t}Az
=
\lambda\,\overline{z}^{\,t}z
=
\lambda\sum_{i=1}^n |z_i|^2.
\]
Нека
\[
q=\overline{z}^{\,t}Az.
\]
Тъй като $A$ е реална симетрична матрица и $q$ е скалар,
\[
\overline q
=z^tA\overline z
=(z^tA\overline z)^t
=\overline z^{\,t}A^tz
=\overline z^{\,t}Az
=q.
\]
Следователно $q$ е реално число. От друга страна, комплексното спрягане на равенството $Az=\lambda z$ дава
\[
A\overline{z}=\overline{\lambda}\,\overline{z}.
\]
Следователно
\[
q=\overline q
=z^tA\overline z
=
\overline{\lambda}\sum_{i=1}^n |z_i|^2.
\]
Значи
\[
\lambda\sum_{i=1}^n |z_i|^2
=
\overline{\lambda}\sum_{i=1}^n |z_i|^2.
\]
Понеже $z\ne0$, то
\[
\sum_{i=1}^n |z_i|^2>0.
\]
Следователно
\[
\lambda=\overline{\lambda},
\]
което означава, че $\lambda\in\mathbb{R}$.
\end{proof}

\begin{remark}
Доказателството разглежда характеристичните корени като комплексни числа. Това е необходимо, защото характеристичният полином на реална матрица може да има комплексни корени. За симетрична матрица доказаното свойство показва, че всички тези корени всъщност са реални.
\end{remark}

\begin{cor}
Всеки симетричен оператор $\varphi:E\to E$ има собствена стойност, ако $\dim E\geq1$.
\end{cor}

\begin{proof}
Характеристичният полином на матрицата на $\varphi$ има комплексен корен. По теоремата всички негови корени са реални, следователно този корен е реална собствена стойност на $\varphi$.
\end{proof}

\subsection{Ортогоналност на собствени вектори}\label{s:30.3.2}

\begin{prop}
Нека $\varphi:E\to E$ е симетричен оператор. Ако $u,v\in E\setminus\{0\}$ са собствени вектори, съответстващи съответно на различните собствени стойности $\lambda$ и $\mu$, то
\[
u\perp v.
\]
\end{prop}

\begin{proof}
Имаме
\[
\varphi(u)=\lambda u,\qquad
\varphi(v)=\mu v,
\]
където $\lambda\ne\mu$. От симетричността на $\varphi$ следва
\[
(\varphi(u),v)=(u,\varphi(v)).
\]
След заместване получаваме
\[
(\lambda u,v)=(u,\mu v),
\]
тоест
\[
\lambda(u,v)=\mu(u,v).
\]
Следователно
\[
(\lambda-\mu)(u,v)=0.
\]
Понеже $\lambda\ne\mu$, имаме
\[
(u,v)=0.
\]
Следователно $u\perp v$.
\end{proof}

\begin{cor}
Собствените подпространства на симетричен оператор, съответстващи на различни собствени стойности, са ортогонални.
\end{cor}

\begin{proof}
Нека $E_\lambda$ и $E_\mu$ са собствените подпространства за $\lambda$ и $\mu$, където $\lambda\ne\mu$. Всеки ненулев вектор от $E_\lambda$ е собствен вектор за $\lambda$, а всеки ненулев вектор от $E_\mu$ е собствен вектор за $\mu$. По предходното твърдение всеки такъв чифт вектори е ортогонален. Следователно
\[
E_\lambda\perp E_\mu.
\]
\end{proof}

\begin{remark}
Условието за различност на собствените стойности е съществено. Два собствени вектора, отговарящи на една и съща собствена стойност, не са задължително ортогонални. В собственото подпространство обаче може да се избере ортонормиран базис.
\end{remark}

\section{Инвариантност на ортогоналното допълнение}\label{s:30.4}

За доказателството на теоремата за диагонализация е необходимо свойството, че ортогоналното допълнение на собствен вектор е инвариантно относно симетричния оператор.

\begin{lem}
Нека $\varphi:E\to E$ е симетричен оператор и $u\in E\setminus\{0\}$ е собствен вектор:
\[
\varphi(u)=\lambda u.
\]
Тогава ортогоналното допълнение
\[
\mathcal{L}(u)^\perp
\]
е инвариантно относно $\varphi$, тоест
\[
\varphi\bigl(\mathcal{L}(u)^\perp\bigr)
\subseteq
\mathcal{L}(u)^\perp.
\]
\end{lem}

\begin{proof}
Нека $x\in\mathcal{L}(u)^\perp$. Тогава
\[
(x,u)=0.
\]
Ще докажем, че $\varphi(x)\in\mathcal{L}(u)^\perp$. От симетричността на $\varphi$ следва
\[
(\varphi(x),u)=(x,\varphi(u)).
\]
Тъй като $\varphi(u)=\lambda u$, получаваме
\[
(\varphi(x),u)=(x,\lambda u)=\lambda(x,u)=0.
\]
Следователно
\[
\varphi(x)\perp u,
\]
тоест
\[
\varphi(x)\in\mathcal{L}(u)^\perp.
\]
\end{proof}

\begin{lem}
При условията на предходната лема ограничението
\[
\varphi\big|_{\mathcal{L}(u)^\perp}:
\mathcal{L}(u)^\perp\to\mathcal{L}(u)^\perp
\]
е симетричен оператор.
\end{lem}

\begin{proof}
Нека $x,y\in\mathcal{L}(u)^\perp$. Понеже $\varphi$ е симетричен в $E$, имаме
\[
(\varphi(x),y)=(x,\varphi(y)).
\]
Съгласно предходната лема $\varphi(x)$ и $\varphi(y)$ принадлежат на $\mathcal{L}(u)^\perp$, следователно това равенство е именно условието ограничението на $\varphi$ върху това подпространство да е симетрично.
\end{proof}

\section{Теорема за диагонализация}\label{s:30.5}

\begin{keythm}[Теорема за диагонализация на симетричен оператор]
Нека $E$ е крайномерно евклидово пространство и $\varphi:E\to E$ е симетричен оператор. Тогава съществува ортонормиран базис
\[
e_1,\ldots,e_n
\]
на $E$, съставен от собствени вектори на $\varphi$.

Еквивалентно, съществува ортонормиран базис на $E$, спрямо който матрицата на $\varphi$ е диагонална:
\[
[\varphi]_{(e_1,\ldots,e_n)}
=
\begin{pmatrix}
\lambda_1 & 0 & \cdots & 0\\
0 & \lambda_2 & \cdots & 0\\
\vdots & \vdots & \ddots & \vdots\\
0 & 0 & \cdots & \lambda_n
\end{pmatrix},
\]
където $\lambda_1,\ldots,\lambda_n\in\mathbb{R}$ са собствени стойности на $\varphi$, като някои от тях могат да съвпадат.
\end{keythm}

\begin{proof}
Ще докажем твърдението с индукция по
\[
n=\dim E.
\]

При $n=1$ твърдението е очевидно. Ако $e_1$ е единичен вектор, то $E=\mathcal{L}(e_1)$ и поради линейността на $\varphi$ съществува $\lambda_1\in\mathbb{R}$, за което
\[
\varphi(e_1)=\lambda_1e_1.
\]
Следователно $e_1$ е ортонормиран базис от собствени вектори.

Нека твърдението е вярно за всички евклидови пространства с размерност $n-1$, и нека сега $\dim E=n$.

По доказаното по-горе симетричният оператор $\varphi$ има реална собствена стойност $\lambda_1$. Нека
\[
u_1\in E\setminus\{0\}
\]
е съответен собствен вектор:
\[
\varphi(u_1)=\lambda_1u_1.
\]
Нормализираме $u_1$ и полагаме
\[
e_1=\frac{u_1}{\|u_1\|}.
\]
Тогава $\|e_1\|=1$ и
\[
\varphi(e_1)=\lambda_1e_1.
\]

Нека
\[
U=\mathcal{L}(e_1)^\perp.
\]
По твърдението за ортогоналното разлагане $E=\mathcal{L}(e_1)\oplus U$, а тъй като $\dim\mathcal{L}(e_1)=1$, това е подпространство с размерност $n-1$. От лемата следва, че $U$ е инвариантно относно $\varphi$. Следователно ограничението
\[
\varphi|_U:U\to U
\]
е коректно определен линеен оператор. Освен това то е симетрично.

По индукционното предположение за оператора $\varphi|_U$ съществува ортонормиран базис
\[
e_2,\ldots,e_n
\]
на $U$, съставен от собствени вектори на ограничението. Следователно съществуват реални числа $\lambda_2,\ldots,\lambda_n$, за които
\[
\varphi(e_i)=\lambda_i e_i,
\qquad i=2,\ldots,n.
\]

Тъй като разлагането
\[
E=\mathcal{L}(e_1)\oplus U
\]
е ортогонално --- всеки вектор от $U$ е ортогонален на $e_1$ по определението за $U$ --- системата
\[
e_1,e_2,\ldots,e_n
\]
е ортонормиран базис на $E$. Всеки негов вектор е собствен вектор на $\varphi$. Следователно матрицата на $\varphi$ спрямо този базис е
\[
\begin{pmatrix}
\lambda_1 & 0 & \cdots & 0\\
0 & \lambda_2 & \cdots & 0\\
\vdots & \vdots & \ddots & \vdots\\
0 & 0 & \cdots & \lambda_n
\end{pmatrix}.
\]
Така индукционната стъпка е доказана.
\end{proof}

\begin{cor}
Всеки симетричен оператор в крайномерно евклидово пространство е диагонализируем спрямо ортонормиран базис.
\end{cor}

\begin{proof}
По теоремата за диагонализация съществува ортонормиран базис от собствени вектори $e_1,\ldots,e_n$. Ако $\varphi(e_i)=\lambda_i e_i$, $i$-тата колона на матрицата в този базис има единствена възможна ненулева координата $\lambda_i$ на ред $i$. Следователно матрицата е $\operatorname{diag}(\lambda_1,\ldots,\lambda_n)$.
\end{proof}

\begin{remark}
Теоремата не твърди само съществуване на базис, в който матрицата е диагонална. Същественото допълнително свойство е, че базисът може да бъде избран ортонормиран. Именно това е специфичната особеност на симетричните оператори.
\end{remark}

\subsection{Връзка със симетричните матрици}\label{s:30.5.1}

Нека $A\in M_{n\times n}(\mathbb{R})$ е симетрична матрица. Разглеждаме линейния оператор
\[
\varphi:\mathbb{R}^n\to\mathbb{R}^n,
\qquad
\varphi(x)=Ax,
\]
спрямо стандартното евклидово скаларно произведение. Матрицата на $\varphi$ спрямо каноничния ортонормиран базис е $A$, следователно $\varphi$ е симетричен оператор.

От теоремата за диагонализация следва, че съществува ортонормиран базис от собствени вектори на $A$. Ако $Q$ е матрицата, чиито стълбове са тези базисни вектори, то $Q$ е ортогонална матрица и в този базис матрицата на оператора е диагонална. Следователно симетричната матрица може да бъде приведена до диагонален вид чрез преход към ортонормиран базис.

\begin{example}
Нека
\[
A=
\begin{pmatrix}
2 & 1\\
1 & 2
\end{pmatrix}.
\]
Матрицата е симетрична, защото
\[
A^t=
\begin{pmatrix}
2 & 1\\
1 & 2
\end{pmatrix}
=A.
\]
Характеристичният полином е
\[
\det(A-\lambda I)
=
\begin{vmatrix}
2-\lambda & 1\\
1 & 2-\lambda
\end{vmatrix}
=
(2-\lambda)^2-1.
\]
Следователно характеристичните корени са
\[
\lambda_1=3,\qquad \lambda_2=1.
\]

За $\lambda_1=3$ получаваме собствен вектор
\[
u_1=
\begin{pmatrix}
1\\
1
\end{pmatrix},
\]
а за $\lambda_2=1$ --- собствен вектор
\[
u_2=
\begin{pmatrix}
1\\
-1
\end{pmatrix}.
\]
Техният скаларен продукт е
\[
(u_1,u_2)=1\cdot1+1\cdot(-1)=0,
\]
което илюстрира ортогоналността на собствени вектори за различни собствени стойности.

След нормализиране получаваме ортонормирани собствени вектори
\[
e_1=\frac{1}{\sqrt{2}}
\begin{pmatrix}
1\\
1
\end{pmatrix},
\qquad
e_2=\frac{1}{\sqrt{2}}
\begin{pmatrix}
1\\
-1
\end{pmatrix}.
\]
Спрямо ортонормирания базис $e_1,e_2$ матрицата на оператора е
\[
\begin{pmatrix}
3 & 0\\
0 & 1
\end{pmatrix}.
\]
\end{example}

\section{Обобщение на свойствата}\label{s:30.6}

Нека $\varphi:E\to E$ е симетричен оператор в крайномерно евклидово пространство. Тогава са изпълнени следните основни свойства:

\begin{enumerate}
\item За всеки $x,y\in E$ е изпълнено
\[
(\varphi(x),y)=(x,\varphi(y)).
\]

\item Матрицата на $\varphi$ спрямо всеки ортонормиран базис е реална симетрична матрица.

\item Всички характеристични корени, съответно всички собствени стойности, на $\varphi$ са реални.

\item Ако $u$ и $v$ са собствени вектори за различни собствени стойности, то
\[
(u,v)=0.
\]

\item Ако $u$ е собствен вектор, то подпространството
\[
\mathcal{L}(u)^\perp
\]
е инвариантно относно $\varphi$.

\item Съществува ортонормиран базис на $E$, съставен от собствени вектори на $\varphi$.

\item Спрямо този ортонормиран базис матрицата на $\varphi$ е диагонална, а елементите по главния диагонал са реалните собствени стойности на оператора.
\end{enumerate}

\section{Изпитен фокус}\label{s:30.7}

\begin{itemize}
\item Да се даде определението за симетричен оператор:
\[
(\varphi(x),y)=(x,\varphi(y))
\qquad
\text{за всички }x,y\in E.
\]

\item Да се знае, че матрицата на симетричен оператор спрямо ортонормиран базис е симетрична, и да се възпроизведе доказателството чрез равенството
\[
a_{ji}=(\varphi(e_i),e_j)=(e_i,\varphi(e_j))=a_{ij}.
\]

\item Да се знае и обратното: симетрична матрица спрямо ортонормиран базис определя симетричен оператор.

\item Да се формулира и докаже, че всички характеристични корени на симетричен оператор са реални, като се използва равенството
\[
\lambda\sum_{i=1}^n|z_i|^2
=
\overline{\lambda}\sum_{i=1}^n|z_i|^2.
\]

\item Да се формулира и докаже ортогоналността на собствени вектори за различни собствени стойности:
\[
(\lambda-\mu)(u,v)=0.
\]

\item Да се знае, че ако $u$ е собствен вектор, то $\mathcal{L}(u)^\perp$ е инвариантно относно симетричния оператор, защото
\[
(\varphi(x),u)=(x,\varphi(u))=\lambda(x,u)=0.
\]

\item Да се формулира теоремата за диагонализация: всеки симетричен оператор в крайномерно евклидово пространство има ортонормиран базис от собствени вектори.

\item Да се възпроизведе доказателствената идея за теоремата: индукция по $\dim E$, избор и нормализиране на собствен вектор, преминаване към неговото ортогонално допълнение и прилагане на индукционното предположение върху ограничението на оператора.
\end{itemize}


\chapter{Симетрична и алтернативна група. Теорема на Кейли. Теорема за хомоморфизмите.}
%% Drafted from the mapped corpus by scripts/draft_topics.py.
% Model: gpt-5.6-terra; source characters: 20622.
\section{Симетрична група и цикличен запис}

\begin{defn}
Нека $M$ е множество. Множеството
\[
\operatorname{Sym}(M)=\{f\colon M\to M\mid f\text{ е биекция}\}
\]
е група относно композицията на изображенията. Тя се нарича \emph{симетрична група на множеството $M$}. Неутралният ѝ елемент е тъждественото изображение
\[
\operatorname{Id}_M\colon M\to M,\qquad \operatorname{Id}_M(x)=x.
\]
\end{defn}

Ако $|M|=n$, то след избиране на номерация $M=\{1,\ldots,n\}$ симетричната група се бележи със $S_n$ и се нарича \emph{симетрична група от степен $n$}. Нейните елементи се наричат \emph{пермутации}. Всяка пермутация е биекция на $\{1,\ldots,n\}$, следователно се определя еднозначно от редицата
\[
\sigma(1),\sigma(2),\ldots,\sigma(n),
\]
която е пренареждане на числата $1,\ldots,n$. Поради това
\[
|S_n|=n!.
\]

\begin{defn}
За $\sigma\in S_n$ множеството
\[
\operatorname{Supp}(\sigma)=\{i\in\{1,\ldots,n\}\mid \sigma(i)\ne i\}
\]
се нарича \emph{носител} на $\sigma$.
\end{defn}

\begin{defn}
Пермутацията $\xi\in S_n$ е \emph{цикъл с дължина $k$}, ако има $k$ различни числа
\[
i_1,\ldots,i_k
\]
такива, че
\[
\xi(i_s)=i_{s+1}\quad (1\le s<k),\qquad \xi(i_k)=i_1,
\]
а всички останали числа се оставят неподвижни. Тя се записва като
\[
\xi=(i_1,i_2,\ldots,i_k).
\]
\end{defn}

Един и същ цикъл има много записи, получени чрез циклично преместване:
\[
(i_1,i_2,\ldots,i_k)
=
(i_2,\ldots,i_k,i_1).
\]
Редът на елементите в цикличния запис обаче е съществен: обикновено
\[
(i_1,i_2,\ldots,i_k)\ne(i_k,\ldots,i_2,i_1).
\]

\begin{defn}
Два цикъла се наричат \emph{независими}, ако носителите им не се пресичат.
\end{defn}

Независимите цикли комутират, тъй като всеки от тях оставя неподвижни всички елементи от носителя на другия. Следователно редът на множителите в произведение на независими цикли няма значение.

\begin{thm}
Всяка нетъждествена пермутация $\sigma\in S_n\setminus\{\operatorname{Id}\}$ се представя като произведение на независими цикли. Представянето е единствено с точност до реда на множителите и до цикличното преместване в записа на всеки отделен цикъл.
\end{thm}

\begin{proof}
Избираме $i_1\in\operatorname{Supp}(\sigma)$ и разглеждаме последователността
\[
i_1,\ \sigma(i_1),\ \sigma^2(i_1),\ldots.
\]
Тъй като множеството $\{1,\ldots,n\}$ е крайно, някой член трябва да се повтори. Първото повторение е именно $i_1$: ако $\sigma(i_r)=i_s$ за $1\le s<r$, то от инективността на $\sigma$ следва $i_{r-1}=i_{s-1}$, което противоречи на избора на първо повторение, освен когато $s=1$. Получава се цикъл
\[
(i_1,\sigma(i_1),\ldots,\sigma^{k-1}(i_1)).
\]

Множеството от елементи на този цикъл се запазва от $\sigma$, а върху допълнението му $\sigma$ отново е пермутация. Повтаряйки процедурата върху останалите неподвижни още елементи от носителя, получаваме разлагане на $\sigma$ в независими цикли.

За единствеността нека имаме две разлагания в независими цикли. Ако $i$ принадлежи на цикъл от първото разлагане, то в другото разлагане има точно един цикъл, който мести $i$. И в двете разлагания последователните образи
\[
i,\ \sigma(i),\ \sigma^2(i),\ldots
\]
са едни и същи; следователно двата цикъла съвпадат. След премахването им разсъждението се прилага за останалите цикли.
\end{proof}

\begin{example}
Нека $\sigma\in S_8$ е зададена чрез
\[
\sigma(1)=4,\quad \sigma(4)=7,\quad \sigma(7)=1,
\]
\[
\sigma(2)=6,\quad \sigma(6)=2,\quad \sigma(3)=3,\quad \sigma(5)=5,\quad \sigma(8)=8.
\]
Тогава
\[
\sigma=(1,4,7)(2,6).
\]
Фиксираните точки $3$, $5$ и $8$ обикновено не се записват.
\end{example}

\section{Спрягане в $S_n$}

\begin{defn}
Елементите $a$ и $b$ на група $G$ се наричат \emph{спрегнати}, ако съществува $c\in G$, за което
\[
b=cac^{-1}.
\]
\end{defn}

В симетричната група спрягането има особено ясен смисъл: то просто преименува елементите, участващи в цикъла.

\begin{lem}
За $\rho\in S_n$ и цикъл $(i_1,\ldots,i_k)\in S_n$ е изпълнено
\[
\rho(i_1,\ldots,i_k)\rho^{-1}
=
\bigl(\rho(i_1),\ldots,\rho(i_k)\bigr).
\]
\end{lem}

\begin{proof}
Нека $1\le s\le k$, като индексите в цикъла се разглеждат по модул $k$. Тогава
\[
\bigl[\rho(i_1,\ldots,i_k)\rho^{-1}\bigr](\rho(i_s))
=
\rho(i_1,\ldots,i_k)(i_s)
=
\rho(i_{s+1}).
\]
Това е точно действието на цикъла
\[
\bigl(\rho(i_1),\ldots,\rho(i_k)\bigr).
\]
Ако $j$ не участва в първоначалния цикъл, то $\rho(j)$ не участва в получения цикъл и и двете страни оставят $\rho(j)$ неподвижно.
\end{proof}

Ако
\[
\sigma=\alpha_1\alpha_2\cdots\alpha_r
\]
е разлагане в независими цикли, тогава
\[
\rho\sigma\rho^{-1}
=
(\rho\alpha_1\rho^{-1})
(\rho\alpha_2\rho^{-1})
\cdots
(\rho\alpha_r\rho^{-1}).
\]
Следователно при спрягане се запазват дължините и броят на независимите цикли.

\begin{defn}
Две пермутации от $S_n$ имат \emph{еднакъв цикличен строеж}, ако в разлаганията си в независими цикли имат еднакъв брой цикли от всяка дължина.
\end{defn}

\begin{prop}
Две пермутации $\sigma,\tau\in S_n$ са спрегнати в $S_n$ тогава и само тогава, когато имат еднакъв цикличен строеж.
\end{prop}

\begin{proof}
Ако $\tau=\rho\sigma\rho^{-1}$, лемата показва, че всеки цикъл на $\sigma$ се преобразува в цикъл със същата дължина. Значи цикличният строеж се запазва.

Обратно, нека $\sigma$ и $\tau$ имат еднакъв цикличен строеж. Сдвояваме циклите им с еднаква дължина и дефинираме $\rho$ така, че да изпраща последователните елементи на всеки цикъл на $\sigma$ в последователните елементи на съответния цикъл на $\tau$. Върху фиксираните точки $\rho$ се избира произволно като биекция към фиксираните точки на $\tau$. Тогава лемата дава
\[
\tau=\rho\sigma\rho^{-1}.
\]
\end{proof}

\section{Транспозиции, четност и алтернативна група}

\begin{defn}
Цикъл с дължина $2$ се нарича \emph{транспозиция}. Той има вида
\[
(i,j)
\]
и разменя $i$ и $j$, като оставя всички останали елементи неподвижни.
\end{defn}

За всяка транспозиция е изпълнено
\[
(i,j)^{-1}=(i,j),\qquad (i,j)^2=\operatorname{Id}.
\]

\begin{prop}
Всеки цикъл с дължина $k\ge2$ се представя като произведение на $k-1$ транспозиции:
\[
(i_1,i_2,\ldots,i_k)
=
(i_1,i_k)(i_1,i_{k-1})\cdots(i_1,i_3)(i_1,i_2).
\]
\end{prop}

\begin{proof}
Прилагаме множителите отдясно наляво. Елементът $i_1$ се изпраща в $i_2$ от първата транспозиция и следващите го фиксират. За $2\le s<k$ елементът $i_s$ остава неподвижен до $(i_1,i_s)$, която го изпраща в $i_1$; следващата $(i_1,i_{s+1})$ го изпраща в $i_{s+1}$, където остава. Елементът $i_k$ се изпраща в $i_1$ от последната транспозиция. Всички останали елементи са фиксирани. Това е точно действието на цикъла.
\end{proof}



\begin{cor}
Всяка пермутация от $S_n$ се представя като произведение на транспозиции.
\end{cor}

\begin{proof}
Разлагаме пермутацията на независими цикли по доказаната теорема. Заменяме всеки цикъл с дължина поне $2$ с произведението от транспозиции от предходното твърдение. Неподвижните точки не изискват множител; тъждествената пермутация се представя с празното произведение.
\end{proof}

Наистина, първо разлагаме пермутацията на независими цикли, след което всеки цикъл заменяме с произведението от транспозиции, дадено по-горе.

\begin{thm}
Ако
\[
\tau_1\tau_2\cdots\tau_k=\operatorname{Id},
\]
където $\tau_1,\ldots,\tau_k$ са транспозиции, то $k$ е четно число. Следователно ако една и съща пермутация е представена като произведение съответно на $k$ и на $m$ транспозиции, то
\[
k\equiv m\pmod 2.
\]
\end{thm}

\begin{proof}
На пермутация $\sigma$ съпоставяме пермутационната матрица $P_\sigma$, която изпраща $e_i$ в $e_{\sigma(i)}$. Тогава $P_{\sigma\rho}=P_\sigma P_\rho$. Матрицата на транспозиция се получава от единичната с размяна на две колони, затова детерминантата й е $-1$. От произведението, равно на тъждествената пермутация, получаваме $1=(-1)^k$, следователно $k$ е четно. Ако два записа на една пермутация имат дължини $k,m$, умножаваме първия по обратния на втория. Получаваме запис на тъждествената с $k+m$ транспозиции, така че $k+m$ е четно, тоест $k\equiv m\pmod2$.
\end{proof}

Теоремата позволява да се въведе четността на пермутация независимо от избраното разлагане на транспозиции.

\begin{defn}
Пермутацията $\sigma\in S_n$ се нарича \emph{четна}, ако в произволно нейно представяне като произведение на транспозиции броят на транспозициите е четен. Тя се нарича \emph{нечетна}, ако този брой е нечетен.
\end{defn}

Тъждествената пермутация е четна. Произведението на две четни пермутации е четно; обратната на четна пермутация също е четна. Аналогично, произведението на четна и нечетна пермутация е нечетно, а произведението на две нечетни пермутации е четно.

\begin{remark}
Нека $\sigma\in S_n$ и разгледаме редицата
\[
\sigma(1),\ldots,\sigma(n).
\]
Двойката $\sigma(p),\sigma(q)$ образува \emph{инверсия}, ако
\[
1\le p<q\le n
\quad\text{и}\quad
\sigma(p)>\sigma(q).
\]
Четността на пермутацията съвпада с четността на броя на нейните инверсии.
\end{remark}

\begin{defn}
За $n\ge2$ множеството
\[
A_n=\{\sigma\in S_n\mid \sigma\text{ е четна}\}
\]
се нарича \emph{алтернативна група от степен $n$}.
\end{defn}

\begin{thm}
За всяко $n\ge2$ множеството $A_n$ е нормална подгрупа на $S_n$ с индекс $2$. В частност
\[
A_n\triangleleft S_n,
\qquad
[S_n:A_n]=2,
\qquad
|A_n|=\frac{n!}{2}.
\]
\end{thm}

\begin{proof}
Критерият за подгрупа показва, че $A_n\le S_n$: ако $\sigma,\tau\in A_n$, то $\sigma\tau^{-1}$ е произведение на четен брой транспозиции и следователно принадлежи на $A_n$.

Нека $(i,j)$ е фиксирана транспозиция. Произведението $(i,j)\sigma$ е нечетно за всяка $\sigma\in A_n$, така че
\[
(i,j)A_n\subseteq S_n\setminus A_n.
\]
Обратно, ако $\tau$ е нечетна, то $(i,j)\tau$ е четна. Следователно
\[
\tau=(i,j)\bigl((i,j)\tau\bigr)\in(i,j)A_n.
\]
Така
\[
S_n\setminus A_n=(i,j)A_n.
\]
Получаваме разлагането на $S_n$ в два леви съседни класа:
\[
S_n=A_n\mathbin{\dot\cup}(i,j)A_n.
\]
Следователно индексът е $2$. Подгрупа с индекс $2$ е нормална, а от теоремата на Лагранж следва
\[
|A_n|=\frac{|S_n|}{[S_n:A_n]}=\frac{n!}{2}.
\]
\end{proof}

\section{Хомоморфизми, ядро и образ}

\begin{defn}
Нека $(G_1,\cdot)$ и $(G_2,\ast)$ са групи. Отображението
\[
\varphi\colon G_1\to G_2
\]
се нарича \emph{хомоморфизъм на групи}, ако за всички $a,b\in G_1$ е изпълнено
\[
\varphi(a\cdot b)=\varphi(a)\ast\varphi(b).
\]
\end{defn}

Ако груповите операции са означени еднакво, обичайно се пише просто
\[
\varphi(ab)=\varphi(a)\varphi(b).
\]

\begin{prop}
Нека $\varphi\colon G_1\to G_2$ е хомоморфизъм, а $e_1$ и $e_2$ са неутралните елементи на $G_1$ и $G_2$. Тогава:
\[
\varphi(e_1)=e_2;
\]
\[
\varphi(a^{-1})=\varphi(a)^{-1}
\qquad\text{за всяко }a\in G_1.
\]
\end{prop}

\begin{proof}
От $e_1e_1=e_1$ получаваме
\[
\varphi(e_1)=\varphi(e_1)\varphi(e_1).
\]
Умножаваме отляво с $\varphi(e_1)^{-1}$ и получаваме $\varphi(e_1)=e_2$.

Освен това
\[
e_2=\varphi(e_1)=\varphi(aa^{-1})
=\varphi(a)\varphi(a^{-1}),
\]
поради което $\varphi(a^{-1})$ е обратен на $\varphi(a)$.
\end{proof}

\begin{defn}
За хомоморфизъм $\varphi\colon G_1\to G_2$ множествата
\[
\ker\varphi
=
\{a\in G_1\mid \varphi(a)=e_2\}
\]
и
\[
\operatorname{Im}\varphi
=
\{\varphi(a)\mid a\in G_1\}
\]
се наричат съответно \emph{ядро} и \emph{образ} на $\varphi$.
\end{defn}

\begin{prop}
Ако $\varphi\colon G_1\to G_2$ е хомоморфизъм, то
\[
\ker\varphi\le G_1,
\qquad
\operatorname{Im}\varphi\le G_2.
\]
\end{prop}

\begin{proof}
За $a,b\in\ker\varphi$ имаме
\[
\varphi(ab^{-1})
=
\varphi(a)\varphi(b)^{-1}
=
e_2e_2^{-1}
=
e_2.
\]
Следователно $ab^{-1}\in\ker\varphi$ и критерият за подгрупа доказва първото твърдение.

Ако $x=\varphi(a)$ и $y=\varphi(b)$ принадлежат на $\operatorname{Im}\varphi$, то
\[
xy^{-1}
=
\varphi(a)\varphi(b)^{-1}
=
\varphi(ab^{-1})\in\operatorname{Im}\varphi.
\]
Следователно и образът е подгрупа.
\end{proof}

За подгрупа $H\le G$ левият съседен клас с представител $a\in G$ е
\[
aH=\{ah\mid h\in H\}.
\]
Ако $H$ е нормална подгрупа, пишем $H\triangleleft G$ и левите и десните съседни класове съвпадат. Тогава множеството
\[
G/H=\{aH\mid a\in G\}
\]
е факторгрупа относно операцията
\[
(aH)(bH)=abH.
\]

\begin{prop}
Нека $\varphi\colon G_1\to G_2$ е хомоморфизъм и $H=\ker\varphi$. Тогава:
\[
\varphi(a)=\varphi(b)
\quad\Longleftrightarrow\quad
aH=bH
\]
за всички $a,b\in G_1$.
\end{prop}

\begin{proof}
Имаме
\[
\varphi(a)=\varphi(b)
\quad\Longleftrightarrow\quad
\varphi(a^{-1}b)=e_2
\quad\Longleftrightarrow\quad
a^{-1}b\in H.
\]
Последното условие е еквивалентно на $b\in aH$, а то от своя страна е еквивалентно на $aH=bH$.
\end{proof}

\section{Теорема за хомоморфизмите}

\begin{keythm}[Теорема за хомоморфизмите на групи]
Нека
\[
\varphi\colon G_1\to G_2
\]
е хомоморфизъм. Тогава
\[
\ker\varphi\triangleleft G_1
\]
и $\varphi$ индуцира изоморфизъм
\[
\overline{\varphi}\colon G_1/\ker\varphi\longrightarrow\operatorname{Im}\varphi,
\qquad
\overline{\varphi}(a\ker\varphi)=\varphi(a).
\]
За естествения хомоморфизъм
\[
\pi\colon G_1\to G_1/\ker\varphi,
\qquad
\pi(a)=a\ker\varphi,
\]
е изпълнено
\[
\overline{\varphi}\circ\pi=\varphi.
\]
Следователно
\[
G_1/\ker\varphi\cong\operatorname{Im}\varphi.
\]
\end{keythm}

\begin{proof}
За $a\in G_1$ и $h\in\ker\varphi$ имаме
\[
\varphi(aha^{-1})
=
\varphi(a)\varphi(h)\varphi(a)^{-1}
=
\varphi(a)e_2\varphi(a)^{-1}
=
e_2.
\]
Следователно $aha^{-1}\in\ker\varphi$, тоест $\ker\varphi\triangleleft G_1$.

Дефинираме
\[
\overline{\varphi}(a\ker\varphi)=\varphi(a).
\]
Това е коректно дефинирано, защото от
\[
a\ker\varphi=b\ker\varphi
\]
следва $\varphi(a)=\varphi(b)$.

За $a,b\in G_1$:
\[
\overline{\varphi}\bigl((a\ker\varphi)(b\ker\varphi)\bigr)
=
\overline{\varphi}(ab\ker\varphi)
=
\varphi(ab)
=
\varphi(a)\varphi(b),
\]
тоест $\overline{\varphi}$ е хомоморфизъм.

Той е сюрективен по дефиницията на образа: всеки елемент на $\operatorname{Im}\varphi$ има вид $\varphi(a)$ и е образ на $a\ker\varphi$. Ако
\[
\overline{\varphi}(a\ker\varphi)
=
\overline{\varphi}(b\ker\varphi),
\]
то $\varphi(a)=\varphi(b)$, откъдето $a\ker\varphi=b\ker\varphi$. Значи $\overline{\varphi}$ е и инективен.

Накрая:
\[
(\overline{\varphi}\circ\pi)(a)
=
\overline{\varphi}(a\ker\varphi)
=
\varphi(a).
\]
\end{proof}

\begin{example}
Нека
\[
\varphi\colon \mathbb{Z}\to\mathbb{Z}_m,
\qquad
\varphi(k)=\overline{k},
\]
където $\mathbb{Z}$ е адитивната група на целите числа, а $\mathbb{Z}_m$ е адитивната група по модул $m$. Тогава
\[
\ker\varphi=m\mathbb{Z},
\qquad
\operatorname{Im}\varphi=\mathbb{Z}_m.
\]
Теоремата за хомоморфизмите дава
\[
\mathbb{Z}/m\mathbb{Z}\cong\mathbb{Z}_m.
\]
\end{example}

\section{Теорема на Кейли}

\begin{keythm}[Теорема на Кейли]
Всяка крайна група $G$ от ред $n$ е изоморфна на подгрупа на симетричната група $S_n$.
\end{keythm}

\begin{proof}
Разглеждаме действието на $G$ върху самата себе си чрез леви транслации. За всяко $a\in G$ дефинираме
\[
\lambda_a\colon G\to G,
\qquad
\lambda_a(x)=ax.
\]
Отображението $\lambda_a$ е биекция, защото
\[
\lambda_a^{-1}=\lambda_{a^{-1}}.
\]
Следователно $\lambda_a\in\operatorname{Sym}(G)$.

Нека
\[
\Phi\colon G\to\operatorname{Sym}(G),
\qquad
\Phi(a)=\lambda_a.
\]
За $a,b,x\in G$ е изпълнено
\[
(\Phi(a)\circ\Phi(b))(x)
=
\Phi(a)(bx)
=
a(bx)
=
(ab)x
=
\Phi(ab)(x).
\]
Следователно $\Phi$ е хомоморфизъм.

Неговото ядро е тривиално. Действително, ако $\Phi(a)=\operatorname{Id}_G$, то
\[
ab=b
\qquad\text{за всяко }b\in G.
\]
За $b=e$ получаваме $a=e$. Значи
\[
\ker\Phi=\{e\}.
\]
По теоремата за хомоморфизмите
\[
G/\{e\}\cong\operatorname{Im}\Phi.
\]
Но $G/\{e\}\cong G$, а $\operatorname{Im}\Phi$ е подгрупа на $\operatorname{Sym}(G)$. Тъй като $|G|=n$, имаме
\[
\operatorname{Sym}(G)\cong S_n.
\]
Следователно $G$ е изоморфна на подгрупа на $S_n$.
\end{proof}

\begin{supp}[Бележка]
По-общото твърдение, доказано чрез левите транслации, е че всяка група $G$ е изоморфна на подгрупа на $\operatorname{Sym}(G)$. Формулировката с $S_n$ е специализацията за крайна група с $n$ елемента.
\end{supp}

\section{Изпитен фокус}

\begin{itemize}
\item Да се дадат определенията за $\operatorname{Sym}(M)$, $S_n$, носител на пермутация, цикъл, независими цикли, транспозиция, четна и нечетна пермутация, $A_n$, хомоморфизъм, ядро, образ, нормална подгрупа и факторгрупа.

\item Да се обясни алгоритъмът за разлагане на пермутация в независими цикли: започва се от елемент от носителя, проследяват се последователните му образи до завръщане в началния елемент и процедурата се повтаря върху останалите елементи.

\item Да се възпроизведе формулата
\[
(i_1,\ldots,i_k)
=
(i_1,i_k)(i_1,i_{k-1})\cdots(i_1,i_2)
\]
и изводът, че всяка пермутация е произведение на транспозиции.

\item Да се знае, че броят на транспозициите в две разлагания на една и съща пермутация има една и съща четност, и да се използва това за определяне на четни и нечетни пермутации.

\item Да се формулира и докаже, че
\[
A_n\triangleleft S_n,\qquad [S_n:A_n]=2,\qquad |A_n|=\frac{n!}{2}.
\]

\item Да се използва формулата за спрягане на цикъл:
\[
\rho(i_1,\ldots,i_k)\rho^{-1}
=
(\rho(i_1),\ldots,\rho(i_k)).
\]
Да се обясни критерият: две пермутации в $S_n$ са спрегнати точно когато имат еднакъв цикличен строеж.

\item Да се докажат основните свойства на хомоморфизма:
\[
\varphi(e_1)=e_2,
\qquad
\varphi(a^{-1})=\varphi(a)^{-1},
\qquad
\ker\varphi\le G_1,
\qquad
\operatorname{Im}\varphi\le G_2.
\]

\item Да се формулира теоремата за хомоморфизмите:
\[
G_1/\ker\varphi\cong\operatorname{Im}\varphi,
\]
като се обясни защо отображението
\[
a\ker\varphi\longmapsto\varphi(a)
\]
е коректно дефинирано, хомоморфизъм, сюрективно и инективно.

\item Да се възпроизведе доказателствената идея на теоремата на Кейли: групата действа върху себе си чрез леви транслации $x\mapsto ax$; получава се хомоморфизъм към симетрична група с тривиално ядро.
\end{itemize}


\chapter{Теорема на Ферма. Теореми за средните стойности (Рол, Лагранж, Коши). Формула на Тейлър.}
%% Drafted from the mapped corpus by scripts/draft_topics.py.
% Model: gpt-5.6-terra; source characters: 15413.
\begin{konspekt}
\kreq{Дефиниция}{локален екстремум на функция на една променлива --- \kref{32.1}.}
\kreq{Формулировка и доказателство}{необходимо условие за локален екстремум при диференцируеми функции --- теорема на Ферма --- \kref{32.1}.}
\kreq{Доказателства}{теоремите на Рол, Лагранж и Коши, заедно с $g(a)\ne g(b)$ при предположенията на Коши --- \kref{32.2}.}
\kreq{Извеждане}{формула на Тейлър с остатъчен член във формата на Лагранж --- \kref{32.3}.}
\kextra{теоремата на Вайерщрас се \emph{използва без доказателство} при теоремата на Рол --- това е изрично разрешено от анотацията.}
\kreq{Литература}{[6], [8], [10].}
\end{konspekt}

\section{Локални екстремуми и теорема на Ферма}\label{s:32.1}

\begin{defn}[Локален екстремум]
Нека $f:D\to\mathbb{R}$, където $D\subseteq\mathbb{R}$, и нека $x_0\in D$.

Казваме, че $f$ има \emph{локален минимум} в точката $x_0$, ако съществува $\delta>0$, за което
\[
(x_0-\delta,x_0+\delta)\subseteq D
\]
и
\[
f(x_0)\leq f(x)
\quad\text{за всяко }x\in(x_0-\delta,x_0+\delta).
\]
Аналогично, $f$ има \emph{локален максимум} в $x_0$, ако съществува $\delta>0$, за което
\[
(x_0-\delta,x_0+\delta)\subseteq D
\]
и
\[
f(x_0)\geq f(x)
\quad\text{за всяко }x\in(x_0-\delta,x_0+\delta).
\]
Точка, в която функцията има локален минимум или локален максимум, се нарича \emph{точка на локален екстремум}.
\end{defn}

\begin{keythm}[Теорема на Ферма --- необходимо условие за локален екстремум]
Нека $f:D\to\mathbb{R}$ и нека $x_0$ е точка на локален екстремум на $f$. Ако $f$ е диференцируема в $x_0$, то
\[
f'(x_0)=0.
\]
\end{keythm}

\begin{proof}
Нека първо $x_0$ е точка на локален минимум. Тогава съществува $\delta>0$, такова че
\[
f(x_0)\leq f(x)
\]
за всяко $x\in(x_0-\delta,x_0+\delta)$.

Ако $x\in(x_0,x_0+\delta)$, то $x-x_0>0$ и $f(x)-f(x_0)\geq 0$. Следователно
\[
\frac{f(x)-f(x_0)}{x-x_0}\geq 0.
\]
При преминаване към граница отдясно получаваме
\[
f'(x_0)\geq 0.
\]

Ако $x\in(x_0-\delta,x_0)$, то $x-x_0<0$, докато отново $f(x)-f(x_0)\geq 0$. Затова
\[
\frac{f(x)-f(x_0)}{x-x_0}\leq 0.
\]
При преминаване към граница отляво следва
\[
f'(x_0)\leq 0.
\]
Тъй като производната съществува, левият и десният предел съвпадат. Следователно едновременно е изпълнено
\[
f'(x_0)\geq 0
\qquad\text{и}\qquad
f'(x_0)\leq 0,
\]
тоест $f'(x_0)=0$.

Ако $x_0$ е точка на локален максимум на $f$, то тя е точка на локален минимум на $-f$. От вече доказаното
\[
(-f)'(x_0)=0,
\]
откъдето $-f'(x_0)=0$, т.е. $f'(x_0)=0$.
\end{proof}

\begin{remark}
Теоремата на Ферма дава необходимо, но не достатъчно условие за локален екстремум. Условието $f'(x_0)=0$ само посочва възможна точка на екстремум.
\end{remark}

\section{Теореми за средните стойности}\label{s:32.2}

В този раздел нека $a<b$ и $f$ е непрекъсната в затворения интервал $[a,b]$ и диференцируема в отворения интервал $(a,b)$.

За доказателството на първата теорема използваме теоремата на Вайерщрас: всяка непрекъсната функция върху краен затворен интервал достига най-голямата и най-малката си стойност.

\subsection{Теорема на Рол}\label{s:32.2.1}

\begin{keythm}[Теорема на Рол]
Ако
\[
f(a)=f(b),
\]
то съществува точка $c\in(a,b)$, за която
\[
f'(c)=0.
\]
\end{keythm}

\begin{proof}
Понеже $f$ е непрекъсната в $[a,b]$, от теоремата на Вайерщрас следва, че тя достига своя минимум и максимум в някои точки
\[
x_{\min},x_{\max}\in[a,b].
\]

Ако поне една от точките $x_{\min}$ и $x_{\max}$ принадлежи на $(a,b)$, то в нея функцията има локален екстремум. Функцията е диференцируема в тази вътрешна точка, следователно по теоремата на Ферма производната в нея е нула. Получаваме търсената точка $c$.

Остава случаят, когато и минимумът, и максимумът се достигат само в краищата на интервала. Тъй като
\[
f(a)=f(b),
\]
минималната и максималната стойност са равни. Следователно функцията е константна в $[a,b]$. Тогава
\[
f'(x)=0
\quad\text{за всяко }x\in(a,b),
\]
и всяка точка от $(a,b)$ може да бъде избрана за $c$.
\end{proof}

\begin{supp}[Бележка]
В едно от предоставените изложения при разглеждането на вътрешна точка на минимум е записано $f(x_{\min})=0$. От теоремата на Ферма следва не равенство за стойността на функцията, а
\[
f'(x_{\min})=0.
\]
Именно това равенство се използва в доказателството на теоремата на Рол.
\end{supp}

\begin{remark}
Геометрично теоремата на Рол гласи, че ако графиката на функцията започва и завършва на една и съща височина, то между тези две точки има точка с хоризонтална допирателна.
\end{remark}

\subsection{Теорема на Лагранж}\label{s:32.2.2}

\begin{keythm}[Теорема на Лагранж за крайните нараствания]
Съществува точка $c\in(a,b)$, такава че
\[
f(b)-f(a)=f'(c)(b-a).
\]
Еквивалентно,
\[
f'(c)=\frac{f(b)-f(a)}{b-a}.
\]
\end{keythm}

\begin{proof}
Да разгледаме функцията
\[
h(x)=f(x)-kx,
\]
където константата $k$ ще изберем така, че да е изпълнено условието на теоремата на Рол:
\[
h(a)=h(b).
\]
Това равенство е еквивалентно на
\[
f(a)-ka=f(b)-kb,
\]
или, тъй като $a\ne b$,
\[
k=\frac{f(b)-f(a)}{b-a}.
\]

Функцията $h$ е непрекъсната в $[a,b]$ и диференцируема в $(a,b)$, защото $f$ има същите свойства. При избраното $k$ е изпълнено $h(a)=h(b)$. По теоремата на Рол съществува $c\in(a,b)$, за което
\[
h'(c)=0.
\]
Но
\[
h'(x)=f'(x)-k.
\]
Следователно
\[
0=h'(c)=f'(c)-k,
\]
тоест
\[
f'(c)=k=\frac{f(b)-f(a)}{b-a}.
\]
След умножение по $b-a$ получаваме
\[
f(b)-f(a)=f'(c)(b-a).
\]
\end{proof}

\begin{remark}
Числото
\[
\frac{f(b)-f(a)}{b-a}
\]
е средната скорост на изменение на функцията върху $[a,b]$, а $f'(c)$ е моментната скорост на изменение в подходяща вътрешна точка $c$.
\end{remark}

\subsection{Теорема на Коши}\label{s:32.2.3}

\begin{keythm}[Обобщена теорема на Лагранж --- теорема на Коши]
Нека $f$ и $g$ са непрекъснати в $[a,b]$ и диференцируеми в $(a,b)$. Нека освен това
\[
g'(x)\ne 0
\quad\text{за всяко }x\in(a,b).
\]
Тогава съществува $c\in(a,b)$, такова че
\[
\frac{f(b)-f(a)}{g(b)-g(a)}
=
\frac{f'(c)}{g'(c)}.
\]
\end{keythm}

\begin{proof}
\emph{Стъпка 1: знаменателят е ненулев.} Да допуснем противното, тоест $g(a)=g(b)$. Тогава $g$ е непрекъсната в $[a,b]$, диференцируема в $(a,b)$ и приема равни стойности в двата края. По теоремата на Рол съществува $\xi\in(a,b)$, за което
\[
g'(\xi)=0.
\]
Това противоречи на условието $g'(x)\ne 0$ за всяко $x\in(a,b)$. Следователно
\[
g(b)-g(a)\ne 0,
\]
така че лявата страна на доказваното равенство е определена.

\emph{Стъпка 2: спомагателната функция.} Полагаме
\[
H(x)=\bigl(f(b)-f(a)\bigr)\bigl(g(x)-g(a)\bigr)-\bigl(g(b)-g(a)\bigr)\bigl(f(x)-f(a)\bigr).
\]
Числата $f(b)-f(a)$ и $g(b)-g(a)$ са константи, затова $H$ е линейна комбинация на $f$ и $g$ с добавена константа. Значи $H$ е непрекъсната в $[a,b]$ и диференцируема в $(a,b)$, защото $f$ и $g$ имат тези свойства.

\emph{Стъпка 3: равни стойности в краищата.} Пресмятаме двете стойности:
\[
H(a)=\bigl(f(b)-f(a)\bigr)\underbrace{\bigl(g(a)-g(a)\bigr)}_{=0}-\bigl(g(b)-g(a)\bigr)\underbrace{\bigl(f(a)-f(a)\bigr)}_{=0}=0,
\]
\[
H(b)=\bigl(f(b)-f(a)\bigr)\bigl(g(b)-g(a)\bigr)-\bigl(g(b)-g(a)\bigr)\bigl(f(b)-f(a)\bigr)=0,
\]
тъй като двете произведения в $H(b)$ съвпадат. Следователно
\[
H(a)=H(b)=0.
\]

\emph{Стъпка 4: прилагане на теоремата на Рол.} Функцията $H$ удовлетворява трите условия на Рол, затова съществува $c\in(a,b)$, за което
\[
H'(c)=0.
\]

\emph{Стъпка 5: разписване на производната.} Диференцираме $H$, като константните множители излизат пред производната, а производните на $g(x)-g(a)$ и $f(x)-f(a)$ са съответно $g'(x)$ и $f'(x)$:
\[
H'(x)=\bigl(f(b)-f(a)\bigr)g'(x)-\bigl(g(b)-g(a)\bigr)f'(x).
\]
Полагаме $x=c$ и използваме $H'(c)=0$:
\[
0=\bigl(f(b)-f(a)\bigr)g'(c)-\bigl(g(b)-g(a)\bigr)f'(c),
\]
тоест
\[
\bigl(f(b)-f(a)\bigr)g'(c)=\bigl(g(b)-g(a)\bigr)f'(c).
\]

\emph{Стъпка 6: деление.} По условие $g'(c)\ne 0$, а по стъпка 1 $g(b)-g(a)\ne 0$. Следователно произведението
\[
\bigl(g(b)-g(a)\bigr)g'(c)\ne 0
\]
и можем да разделим на него двете страни на последното равенство:
\[
\frac{\bigl(f(b)-f(a)\bigr)g'(c)}{\bigl(g(b)-g(a)\bigr)g'(c)}
=
\frac{\bigl(g(b)-g(a)\bigr)f'(c)}{\bigl(g(b)-g(a)\bigr)g'(c)}.
\]
Отляво се съкращава $g'(c)$, а отдясно --- $g(b)-g(a)$, откъдето
\[
\frac{f(b)-f(a)}{g(b)-g(a)}
=
\frac{f'(c)}{g'(c)}.
\]
\end{proof}

\begin{prop}
При условията на теоремата на Коши е изпълнено
\[
g(a)\ne g(b).
\]
\end{prop}

\begin{proof}
Да допуснем противното:
\[
g(a)=g(b).
\]
Тогава функцията $g$ удовлетворява условията на теоремата на Рол. Следователно съществува $c\in(a,b)$, за което
\[
g'(c)=0.
\]
Това противоречи на предположението, че $g'(x)\ne 0$ за всяко $x\in(a,b)$. Следователно
\[
g(a)\ne g(b).
\]
\end{proof}

\begin{proof}[Второ доказателство на теоремата на Коши]
От предходното твърдение знаменателят $g(b)-g(a)$ е различен от нула. Нека
\[
k=\frac{f(b)-f(a)}{g(b)-g(a)}
\]
и дефинираме
\[
h(x)=f(x)-k g(x).
\]
Като линейна комбинация на $f$ и $g$, функцията $h$ е непрекъсната в $[a,b]$ и диференцируема в $(a,b)$.

От избора на $k$ следва
\[
k\bigl(g(b)-g(a)\bigr)=f(b)-f(a).
\]
Следователно
\[
f(a)-kg(a)=f(b)-kg(b),
\]
тоест
\[
h(a)=h(b).
\]
По теоремата на Рол съществува $c\in(a,b)$, за което
\[
h'(c)=0.
\]
Тъй като
\[
h'(x)=f'(x)-k g'(x),
\]
получаваме
\[
f'(c)-k g'(c)=0.
\]
Понеже $g'(c)\ne 0$, можем да разделим на $g'(c)$:
\[
\frac{f'(c)}{g'(c)}=k.
\]
От определението на $k$ следва
\[
\frac{f'(c)}{g'(c)}
=
\frac{f(b)-f(a)}{g(b)-g(a)}.
\]
\end{proof}

\begin{cor}
Теоремата на Лагранж е частен случай на теоремата на Коши при
\[
g(x)=x.
\]
\end{cor}

\begin{proof}
За $g(x)=x$ имаме $g'(x)=1$ и $g(b)-g(a)=b-a$. Формулата на Коши приема вида
\[
\frac{f(b)-f(a)}{b-a}=f'(c).
\]
\end{proof}

\section{Формула на Тейлър}\label{s:32.3}

\subsection{Полином на Тейлър}\label{s:32.3.1}

\begin{defn}[Полином на Тейлър]
Нека функцията $f$ притежава производни до ред $n$ в точката $x_0$. Полиномът
\[
T_n(x)=\sum_{k=0}^{n}\frac{f^{(k)}(x_0)}{k!}(x-x_0)^k
\]
се нарича \emph{полином на Тейлър от ред $n$} на функцията $f$ около точката $x_0$.

Разписан, той има вида
\[
T_n(x)=f(x_0)+f'(x_0)(x-x_0)
+\frac{f''(x_0)}{2!}(x-x_0)^2+\cdots
+\frac{f^{(n)}(x_0)}{n!}(x-x_0)^n.
\]
\end{defn}

\begin{lem}
Ако $f$ притежава производни до ред $n$ в $x_0$, то
\[
T_n^{(i)}(x_0)=f^{(i)}(x_0),
\qquad i=0,1,\ldots,n.
\]
\end{lem}

\begin{proof}
При диференциране $i$ пъти всеки член
\[
\frac{f^{(k)}(x_0)}{k!}(x-x_0)^k
\]
с $k<i$ изчезва. За $k\geq i$ се получава множител
\[
k(k-1)\cdots(k-i+1)(x-x_0)^{k-i}.
\]
При $x=x_0$ всички членове с $k>i$ са нула, защото съдържат положителна степен на $x-x_0$. Остава само членът с $k=i$, който е равен на
\[
f^{(i)}(x_0).
\]
Следователно
\[
T_n^{(i)}(x_0)=f^{(i)}(x_0).
\]
\end{proof}

\subsection{Остатъчен член във формата на Лагранж}\label{s:32.3.2}

Полиномът на Тейлър приближава $f$ около $x_0$; разликата между функцията и приближението получава собствено име.

\begin{defn}[Остатъчен член]
Нека $f$ притежава производни до ред $n$ в $x_0$ и нека $T_n$ е нейният полином на Тейлър около $x_0$. Функцията
\[
R_n(x)=f(x)-T_n(x)
\]
се нарича \emph{остатъчен член от ред $n$}. По построение
\[
f(x)=T_n(x)+R_n(x).
\]
\end{defn}

От лемата по-горе $T_n^{(i)}(x_0)=f^{(i)}(x_0)$ за $i=0,1,\ldots,n$, следователно
\[
R_n^{(i)}(x_0)=0,\qquad i=0,1,\ldots,n.
\]
Остатъкът и първите му $n$ производни се анулират в $x_0$; въпросът е с каква скорост той расте при отдалечаване от $x_0$. Теоремата по-долу дава точен отговор чрез стойност на $(n+1)$-вата производна във вътрешна точка.

\begin{keythm}[Формула на Тейлър с остатъчен член на Лагранж]
Нека са дадени цяло число $n\geq 0$, точка $x_0\in\mathbb{R}$ и число $\delta>0$. Нека функцията
\[
f:(x_0-\delta,x_0+\delta)\to\mathbb{R}
\]
притежава производни до ред $n+1$ включително във \emph{всяка} точка на интервала, тоест $f^{(n+1)}(t)$ съществува за всяко $t\in(x_0-\delta,x_0+\delta)$.

Тогава за всяко
\[
x\in(x_0-\delta,x_0+\delta),\qquad x\ne x_0,
\]
съществува точка $c$, разположена \emph{строго} между $x_0$ и $x$, тоест
\[
c\in(x_0,x)\ \text{ при } x>x_0,
\qquad
c\in(x,x_0)\ \text{ при } x<x_0,
\]
за която
\[
f(x)
=
\sum_{k=0}^{n}\frac{f^{(k)}(x_0)}{k!}(x-x_0)^k
+
\frac{f^{(n+1)}(c)}{(n+1)!}(x-x_0)^{n+1}.
\]
Еквивалентно, съществува $\theta\in(0,1)$, за което $c=x_0+\theta(x-x_0)$.

При $x=x_0$ равенството е изпълнено тривиално: сумата се свежда до $f(x_0)$, а остатъчният член е нула.
\end{keythm}

Условието е върху отворен интервал около $x_0$, а не върху затворен: така за всяко $x$ от интервала целият затворен интервал с краища $x_0$ и $x$ лежи вътре в областта, върху която $f^{(n+1)}$ съществува. От съществуването на $f^{(k+1)}$ следва непрекъснатост на $f^{(k)}$, затова $f,f',\ldots,f^{(n)}$ са непрекъснати в целия интервал и отделно предположение за това не е нужно.

\begin{proof}
Случаят $x=x_0$ е непосредствен: всички степени $(x-x_0)^k$ при $k\geq 1$ са нула, така че и двете страни са равни на $f(x_0)$.

Нека по-нататък $x\ne x_0$ е фиксирано. Означаваме с
\[
I=[\min(x_0,x),\ \max(x_0,x)]
\]
затворения интервал с краища $x_0$ и $x$. Тъй като и двете точки лежат в $(x_0-\delta,x_0+\delta)$, а този интервал е изпъкнал, имаме $I\subseteq(x_0-\delta,x_0+\delta)$. Достатъчно е да докажем, че
\[
R_n(x)=\frac{f^{(n+1)}(c)}{(n+1)!}(x-x_0)^{n+1}
\]
за подходяща вътрешна точка $c$ на $I$, защото $f(x)=T_n(x)+R_n(x)$.

\emph{Стъпка 1: спомагателните функции.} При фиксирано $x$ разглеждаме функциите на променливата $t$
\[
\varphi(t)
=
f(x)-\sum_{k=0}^{n}\frac{f^{(k)}(t)}{k!}(x-t)^k
=
f(x)-f(t)-\frac{f'(t)}{1!}(x-t)-\cdots-\frac{f^{(n)}(t)}{n!}(x-t)^n
\]
и
\[
\psi(t)=(x-t)^{n+1}.
\]
Обърнете внимание, че тук $x$ е константа, а $t$ пробягва $I$; в $\varphi$ производните на $f$ се пресмятат в $t$, а не в $x_0$.

Функцията $\varphi$ е крайна сума от произведения на $f^{(k)}(t)$ с полиноми на $t$. За $k\leq n$ производната $f^{(k)}$ е диференцируема в целия интервал $(x_0-\delta,x_0+\delta)$, защото $f^{(k+1)}$ съществува там. Следователно $\varphi$ е диференцируема, а значи и непрекъсната, в отворено множество, съдържащо $I$. Същото важи очевидно и за полинома $\psi$. В частност $\varphi$ и $\psi$ са непрекъснати в $I$ и диференцируеми във вътрешността му.

\emph{Стъпка 2: стойности в краищата.} При $t=x$ всички степени $(x-t)^k$ с $k\geq 1$ се анулират, така че от сумата остава само членът с $k=0$:
\[
\varphi(x)=f(x)-f(x)=0.
\]
При $t=x_0$ сумата е точно полиномът на Тейлър:
\[
\varphi(x_0)
=
f(x)-\sum_{k=0}^{n}\frac{f^{(k)}(x_0)}{k!}(x-x_0)^k
=
f(x)-T_n(x)
=
R_n(x).
\]
Аналогично
\[
\psi(x)=0,
\qquad
\psi(x_0)=(x-x_0)^{n+1}\ne 0.
\]

\emph{Стъпка 3: разписване на $\varphi'$.} Диференцираме почленно, като за $k\geq 1$ използваме правилото за произведение:
\[
\frac{d}{dt}\left[\frac{f^{(k)}(t)}{k!}(x-t)^k\right]
=
\frac{f^{(k+1)}(t)}{k!}(x-t)^k
-
\frac{f^{(k)}(t)}{k!}\,k\,(x-t)^{k-1},
\]
където вторият член идва от $\frac{d}{dt}(x-t)^k=-k(x-t)^{k-1}$. За $k=0$ вторият член липсва, тъй като $(x-t)^0=1$ е константа. Понеже $\dfrac{k}{k!}=\dfrac{1}{(k-1)!}$, получаваме
\[
\varphi'(t)
=
-\sum_{k=0}^{n}\frac{f^{(k+1)}(t)}{k!}(x-t)^k
+
\sum_{k=1}^{n}\frac{f^{(k)}(t)}{(k-1)!}(x-t)^{k-1}.
\]
Във втората сума полагаме $j=k-1$ и тя се превръща в
\[
\sum_{j=0}^{n-1}\frac{f^{(j+1)}(t)}{j!}(x-t)^{j},
\]
тоест в първата сума без нейния последен член. Двете суми се унищожават член по член и остава
\[
\varphi'(t)
=
-\frac{f^{(n+1)}(t)}{n!}(x-t)^n.
\]
Това е причината да се избере точно тази спомагателна функция: сумата е телескопична и нейната производна съдържа само $(n+1)$-вата производна на $f$.

За $\psi$ имаме направо
\[
\psi'(t)=-(n+1)(x-t)^n.
\]
За всяко $t$ от вътрешността на $I$ е $t\ne x$, следователно $(x-t)^n\ne 0$ и
\[
\psi'(t)\ne 0.
\]

\emph{Стъпка 4: прилагане на теоремата на Коши.} Функциите $\varphi$ и $\psi$ са непрекъснати в $I$, диференцируеми във вътрешността му и $\psi'$ не се анулира там. По теоремата на Коши съществува вътрешна точка $c$ на $I$, тоест точка строго между $x_0$ и $x$, за която
\[
\frac{\varphi(x)-\varphi(x_0)}
{\psi(x)-\psi(x_0)}
=
\frac{\varphi'(c)}{\psi'(c)}.
\]
Теоремата на Коши е формулирана за интервал $[a,b]$ с $a<b$. Ако $x<x_0$, тя се прилага с $a=x$ и $b=x_0$; лявата страна не се променя, защото при размяна на краищата числителят и знаменателят сменят знака си едновременно.

\emph{Стъпка 5: заместване.} Използваме стойностите от стъпка 2 и производните от стъпка 3:
\[
\frac{0-R_n(x)}{0-(x-x_0)^{n+1}}
=
\frac{
-\dfrac{f^{(n+1)}(c)}{n!}(x-c)^n
}{
-(n+1)(x-c)^n
}.
\]
Отляво знаците се съкращават. Отдясно $c$ е вътрешна точка, затова $c\ne x$ и $(x-c)^n\ne 0$, така че този множител може да бъде съкратен; освен това
\[
\frac{1}{n!\,(n+1)}=\frac{1}{(n+1)!}.
\]
Получаваме
\[
\frac{R_n(x)}{(x-x_0)^{n+1}}
=
\frac{f^{(n+1)}(c)}{(n+1)!},
\]
а след умножение по $(x-x_0)^{n+1}\ne 0$
\[
R_n(x)
=
\frac{f^{(n+1)}(c)}{(n+1)!}(x-x_0)^{n+1}.
\]
Тъй като $f(x)=T_n(x)+R_n(x)$, това е точно формулата на Тейлър с остатъчен член на Лагранж.
\end{proof}

\begin{remark}
Точката $c$ в остатъчния член зависи както от $x$, така и от $n$; тя не е дадена явно, а само съществува. Затова често се означава с $\xi(x)$:
\[
f(x)
=
f(x_0)+f'(x_0)(x-x_0)+\cdots+
\frac{f^{(n)}(x_0)}{n!}(x-x_0)^n
+
\frac{f^{(n+1)}(\xi(x))}{(n+1)!}(x-x_0)^{n+1}.
\]
Винаги $\xi(x)$ е между $x_0$ и $x$.
\end{remark}

\begin{example}
За $n=1$ формулата на Тейлър приема вида
\[
f(x)=f(x_0)+f'(x_0)(x-x_0)
+\frac{f''(c)}{2}(x-x_0)^2,
\]
където $c$ е между $x_0$ и $x$.
\end{example}

\section{Изпитен фокус}\label{s:32.4}

\begin{itemize}
\item Да се дефинират локален минимум, локален максимум и локален екстремум, като се подчертае съществуването на околност на точката, съдържаща се в дефиниционната област.

\item Да се формулира и докаже теоремата на Ферма чрез знака на отношението
\[
\frac{f(x)-f(x_0)}{x-x_0}
\]
отляво и отдясно на точка на локален минимум; случаят на локален максимум се свежда до функцията $-f$.

\item Да се формулира теоремата на Рол и да се възпроизведе доказателството чрез теоремата на Вайерщрас и теоремата на Ферма. Трябва да се разгледа и случаят, когато минимумът и максимумът се достигат само в краищата.

\item Да се формулира теоремата на Лагранж и да се изведе от Рол чрез спомагателната функция
\[
h(x)=f(x)-kx,
\qquad
k=\frac{f(b)-f(a)}{b-a}.
\]

\item Да се формулира теоремата на Коши, включително условието
\[
g'(x)\ne 0\quad\text{за }x\in(a,b).
\]
Да се докаже предварително, че то влече $g(a)\ne g(b)$, чрез противоречие с теоремата на Рол.

\item Да се докаже теоремата на Коши чрез функцията
\[
h(x)=f(x)-k g(x),
\qquad
k=\frac{f(b)-f(a)}{g(b)-g(a)}.
\]

\item Да се знае полиномът на Тейлър
\[
T_n(x)=\sum_{k=0}^{n}\frac{f^{(k)}(x_0)}{k!}(x-x_0)^k
\]
и свойството, че неговите производни до ред $n$ в $x_0$ съвпадат със съответните производни на $f$.

\item Да се формулира формулата на Тейлър с остатъчен член във формата на Лагранж: предположението е за съществуване на $f^{(n+1)}$ във всяка точка на отворен интервал около $x_0$, а точката $c$ е строго между $x_0$ и $x$.

\item Да се възпроизведе доказателството: фиксира се $x$, въвеждат се
\[
\varphi(t)=f(x)-\sum_{k=0}^{n}\frac{f^{(k)}(t)}{k!}(x-t)^k,
\qquad
\psi(t)=(x-t)^{n+1},
\]
пресмятат се $\varphi(x)=0$, $\varphi(x_0)=R_n(x)$, $\psi(x)=0$, $\psi(x_0)=(x-x_0)^{n+1}$, разписва се телескопичната производна
\[
\varphi'(t)=-\frac{f^{(n+1)}(t)}{n!}(x-t)^n,
\]
и към $\varphi$ и $\psi$ се прилага теоремата на Коши.
\end{itemize}


\chapter{Определен интеграл. Суми на Дарбу. Интегруемост. Теорема на Нютон-Лайбниц.}
%% Drafted from the mapped corpus by scripts/draft_topics.py.
% Model: gpt-5.6-terra; source characters: 18188.
\section{Определен интеграл и суми на Дарбу}

Нека $f\colon [a,b]\to\mathbb{R}$, където $a<b$. Целта на определения интеграл е да се постави на строга основа идеята за натрупана величина: например лице под графика, изминат път при зададена скорост или общ принос на непрекъснато изменяща се функция.

\subsection{Разбиване на интервал}

\begin{defn}
\label{def:partition}
\emph{Разбиване} на интервала $[a,b]$ наричаме всяка крайна система от точки
\[
\tau\colon\quad a=x_0<x_1<\cdots<x_n=b.
\]
Точките $x_0,x_1,\ldots,x_n$ се наричат \emph{делящи точки} на разбиването. С разбиването са свързани подинтервалите
\[
[x_{i-1},x_i],\qquad i=1,\ldots,n,
\]
чийто съюз е $[a,b]$ и два от които или не се пресичат, или имат обща само крайна точка.
\end{defn}

\begin{defn}
\emph{Диаметър} на разбиването $\tau$ се нарича числото
\[
d(\tau)=\max_{1\leq i\leq n}(x_i-x_{i-1}).
\]
\end{defn}

Колкото по-малък е диаметърът, толкова по-фино е разбит интервалът. При изследване на интегруемостта ще използваме разбивания, за които $d(\tau)$ е достатъчно малък.

\subsection{Малки и големи суми на Дарбу}

Нека функцията $f$ е ограничена върху $[a,b]$ и нека
\[
\tau\colon\quad a=x_0<x_1<\cdots<x_n=b
\]
е разбиване. За всеки подинтервал полагаме
\[
m_i=\inf\{f(x)\mid x\in[x_{i-1},x_i]\},
\qquad
M_i=\sup\{f(x)\mid x\in[x_{i-1},x_i]\}.
\]
Тъй като $f$ е ограничена, числата $m_i$ и $M_i$ са крайни.

\begin{defn}
\emph{Малка сума на Дарбу} на $f$ по разбиването $\tau$ се нарича числото
\[
s_f(\tau)=\sum_{i=1}^{n}m_i(x_i-x_{i-1}).
\]
\emph{Голяма сума на Дарбу} на $f$ по разбиването $\tau$ се нарича числото
\[
S_f(\tau)=\sum_{i=1}^{n}M_i(x_i-x_{i-1}).
\]
\end{defn}

Геометрично малката сума се получава, ако над всеки подинтервал се построи правоъгълник с височина $m_i$, а голямата --- с височина $M_i$. Следователно малката сума дава долна, а голямата --- горна оценка на търсената натрупана величина.

\[
\begin{array}{c}
\text{над }[x_{i-1},x_i]\text{ височината на долния правоъгълник е }m_i,
\\[1mm]
\text{а височината на горния правоъгълник е }M_i:
\\[2mm]
\begin{array}{c}
\text{\rule{35mm}{0.4pt}}\\[-1mm]
\hspace{4mm}x_{i-1}\hspace{20mm}x_i
\end{array}
\end{array}
\]

По-точно, при добавяне на делящи точки долните правоъгълници могат да се повишат, а горните --- да се понижат.

\begin{lem}
\label{lem:refinement}
Нека $\tau^*$ се получава от разбиването $\tau$ чрез добавяне на една или повече нови делящи точки. Тогава
\[
s_f(\tau^*)\geq s_f(\tau),
\qquad
S_f(\tau^*)\leq S_f(\tau).
\]
\end{lem}

\begin{proof}
Достатъчно е да разгледаме добавянето на една нова точка, тъй като общият случай следва чрез многократно прилагане на същото разсъждение.

Нека $x^*\in(x_{i-1},x_i)$ е добавена точка. Единствено приносът на интервала $[x_{i-1},x_i]$ се променя. Нека
\[
m_i'=\inf_{x\in[x_{i-1},x^*]}f(x),
\qquad
m_i''=\inf_{x\in[x^*,x_i]}f(x).
\]
Понеже
\[
[x_{i-1},x^*]\subseteq[x_{i-1},x_i],
\qquad
[x^*,x_i]\subseteq[x_{i-1},x_i],
\]
имаме
\[
m_i'\geq m_i,
\qquad
m_i''\geq m_i.
\]
Следователно
\begin{align*}
s_f(\tau^*)-s_f(\tau)
&=m_i'(x^*-x_{i-1})+m_i''(x_i-x^*)-m_i(x_i-x_{i-1})\\
&\geq m_i(x^*-x_{i-1})+m_i(x_i-x^*)-m_i(x_i-x_{i-1})\\
&=0.
\end{align*}
Значи $s_f(\tau^*)\geq s_f(\tau)$.

За големите суми разсъждението е аналогично. Ако $M_i'$ и $M_i''$ са супремумите на функцията върху двата нови подинтервала, то
\[
M_i'\leq M_i,\qquad M_i''\leq M_i.
\]
Затова
\[
S_f(\tau^*)\leq S_f(\tau).
\]
\end{proof}

\begin{supp}[Бележка]
В част от предоставените записи означенията за горен и долен интеграл са разменени. Съгласно определенията чрез суми на Дарбу горният интеграл се задава чрез инфимум на \emph{големите} суми, а долният --- чрез супремум на \emph{малките} суми. Това е и единственото съгласувано с неравенството $s_f(\tau)\leq S_f(\tau)$ означение.
\end{supp}

\subsection{Горен и долен интеграл на Дарбу}

За всяко разбиване $\tau$ е изпълнено
\[
s_f(\tau)\leq S_f(\tau).
\]
Освен това всяка малка сума е не по-голяма от всяка голяма сума. Действително, ако $\tau_1$ и $\tau_2$ са две разбивания, разглеждаме общото им доразбиване $\tau_1\cup\tau_2$. От Лема \ref{lem:refinement} следва
\[
s_f(\tau_1)\leq s_f(\tau_1\cup\tau_2)\leq S_f(\tau_1\cup\tau_2)\leq S_f(\tau_2).
\]

\begin{defn}
\emph{Долен интеграл на Дарбу} на ограничената функция $f$ върху $[a,b]$ наричаме
\[
\underline{\int_a^b} f
=
\sup\{s_f(\tau)\mid \tau\text{ е разбиване на }[a,b]\}.
\]
\emph{Горен интеграл на Дарбу} наричаме
\[
\overline{\int_a^b} f
=
\inf\{S_f(\tau)\mid \tau\text{ е разбиване на }[a,b]\}.
\]
\end{defn}

От разгледаното неравенство между малките и големите суми следва, че
\[
\underline{\int_a^b} f\leq\overline{\int_a^b} f.
\]

\begin{defn}
Ограничената функция $f\colon[a,b]\to\mathbb{R}$ се нарича \emph{интегруема по Риман} върху $[a,b]$, ако
\[
\underline{\int_a^b} f=\overline{\int_a^b} f.
\]
Общата стойност се нарича \emph{риманов определен интеграл} на $f$ върху $[a,b]$ и се означава с
\[
\int_a^b f(x)\,dx.
\]
\end{defn}

Така определението чрез Дарбу формализира идеята, че долните и горните стъпаловидни оценки могат да бъдат направени произволно близки.

\section{Критерий за интегруемост}

\begin{thm}[Критерий на Дарбу]
\label{thm:darbourcriterion}
Нека $f\colon[a,b]\to\mathbb{R}$ е ограничена. Тогава $f$ е интегруема по Риман тогава и само тогава, когато за всяко $\varepsilon>0$ съществува разбиване $\tau$ на $[a,b]$, за което
\[
S_f(\tau)-s_f(\tau)<\varepsilon.
\]
\end{thm}

\begin{proof}
Нека първо $f$ е интегруема и
\[
I=\int_a^b f(x)\,dx
=
\underline{\int_a^b}f
=
\overline{\int_a^b}f.
\]
Нека $\varepsilon>0$. От свойството на супремума съществува разбиване $\tau_1$, за което
\[
s_f(\tau_1)>I-\frac{\varepsilon}{2}.
\]
От свойството на инфимума съществува разбиване $\tau_2$, за което
\[
S_f(\tau_2)<I+\frac{\varepsilon}{2}.
\]
Нека $\tau=\tau_1\cup\tau_2$ е общото доразбиване. Според Лема \ref{lem:refinement},
\[
s_f(\tau)\geq s_f(\tau_1),
\qquad
S_f(\tau)\leq S_f(\tau_2).
\]
Следователно
\[
S_f(\tau)-s_f(\tau)
\leq S_f(\tau_2)-s_f(\tau_1)
<
\left(I+\frac{\varepsilon}{2}\right)
-
\left(I-\frac{\varepsilon}{2}\right)
=
\varepsilon.
\]

Обратно, нека за всяко $\varepsilon>0$ има разбиване $\tau$, за което
\[
S_f(\tau)-s_f(\tau)<\varepsilon.
\]
За всяко разбиване е изпълнено
\[
s_f(\tau)\leq\underline{\int_a^b}f
\leq\overline{\int_a^b}f
\leq S_f(\tau).
\]
Затова
\[
0\leq
\overline{\int_a^b}f-\underline{\int_a^b}f
\leq S_f(\tau)-s_f(\tau)
<\varepsilon.
\]
Това е вярно за всяко $\varepsilon>0$, откъдето
\[
\overline{\int_a^b}f=\underline{\int_a^b}f.
\]
Следователно $f$ е интегруема по Риман.
\end{proof}

\section{Интегруемост на непрекъснатите функции}

Ще използваме без доказателство теоремата на Кантор: всяка непрекъсната функция върху краен затворен интервал е равномерно непрекъсната. Ще използваме също, че непрекъсната функция върху $[a,b]$ е ограничена.

\begin{thm}
Всяка непрекъсната функция $f\colon[a,b]\to\mathbb{R}$ е интегруема по Риман.
\end{thm}

\begin{proof}
Нека $\varepsilon>0$. Понеже $f$ е равномерно непрекъсната върху $[a,b]$, съществува $\delta>0$, такова че за всички $x',x''\in[a,b]$ от
\[
|x'-x''|<\delta
\]
следва
\[
|f(x')-f(x'')|<\frac{\varepsilon}{b-a}.
\]

Избираме разбиване
\[
\tau\colon\quad a=x_0<x_1<\cdots<x_n=b,
\qquad
d(\tau)<\delta.
\]
За произволен индекс $i$ и произволни точки
\[
x',x''\in[x_{i-1},x_i]
\]
имаме
\[
|x'-x''|\leq x_i-x_{i-1}\leq d(\tau)<\delta.
\]
Следователно
\[
|f(x')-f(x'')|<\frac{\varepsilon}{b-a}.
\]
Оттук осцилацията на $f$ върху всеки подинтервал удовлетворява
\[
M_i-m_i\leq\frac{\varepsilon}{b-a}.
\]
Тогава
\begin{align*}
S_f(\tau)-s_f(\tau)
&=\sum_{i=1}^{n}(M_i-m_i)(x_i-x_{i-1})\\
&\leq
\frac{\varepsilon}{b-a}
\sum_{i=1}^{n}(x_i-x_{i-1})\\
&=
\frac{\varepsilon}{b-a}(b-a)
=
\varepsilon.
\end{align*}
По Теорема \ref{thm:darbourcriterion} функцията $f$ е интегруема по Риман.
\end{proof}

\section{Основни свойства на римановия интеграл}

Следващите свойства се използват без доказателство. Нека разглежданите функции са интегруеми по Риман върху съответните интервали.

\begin{prop}[Линейност]
Ако $f,g\colon[a,b]\to\mathbb{R}$ са интегруеми и $\lambda,\mu\in\mathbb{R}$, то $\lambda f+\mu g$ е интегруема и
\[
\int_a^b\bigl(\lambda f(x)+\mu g(x)\bigr)\,dx
=
\lambda\int_a^b f(x)\,dx
+
\mu\int_a^b g(x)\,dx.
\]
В частност,
\[
\int_a^b(f(x)+g(x))\,dx
=
\int_a^b f(x)\,dx+\int_a^b g(x)\,dx,
\]
\[
\int_a^b\lambda f(x)\,dx
=
\lambda\int_a^b f(x)\,dx.
\]
\end{prop}

\begin{prop}[Адитивност по интервала]
Ако $c\in(a,b)$, то $f$ е интегруема върху $[a,b]$ тогава и само тогава, когато е интегруема върху $[a,c]$ и $[c,b]$. В този случай
\[
\int_a^b f(x)\,dx
=
\int_a^c f(x)\,dx+\int_c^b f(x)\,dx.
\]
\end{prop}

\begin{prop}[Смяна на границите]
За интегруема функция $f$ е изпълнено
\[
\int_b^a f(x)\,dx=-\int_a^b f(x)\,dx.
\]
\end{prop}

\begin{prop}[Позитивност и интегриране на неравенства]
Ако $f(x)\geq0$ за всяко $x\in[a,b]$, то
\[
\int_a^b f(x)\,dx\geq0.
\]
По-общо, ако $f(x)\leq g(x)$ за всяко $x\in[a,b]$, то
\[
\int_a^b f(x)\,dx\leq\int_a^b g(x)\,dx.
\]
\end{prop}

\begin{cor}
Ако $f$ е интегруема върху $[a,b]$, то $|f|$ също е интегруема и
\[
\left|\int_a^b f(x)\,dx\right|
\leq
\int_a^b |f(x)|\,dx.
\]
\end{cor}

\begin{example}
За константната функция $f(x)=k$ е изпълнено
\[
\int_a^b k\,dx=k(b-a).
\]
\end{example}

\section{Теорема за средната стойност за интеграли}

\begin{thm}[Теорема за средната стойност]
\label{thm:meanintegral}
Нека $f\colon[a,b]\to\mathbb{R}$ е непрекъсната. Тогава съществува $c\in[a,b]$, такова че
\[
\int_a^b f(x)\,dx=f(c)(b-a).
\]
\end{thm}

\begin{proof}
Понеже $f$ е непрекъсната върху затворения интервал $[a,b]$, тя приема най-малка и най-голяма стойност. Нека
\[
m=\min_{x\in[a,b]}f(x),
\qquad
M=\max_{x\in[a,b]}f(x).
\]
Тогава
\[
m\leq f(x)\leq M,
\qquad x\in[a,b].
\]
Чрез свойството за интегриране на неравенства получаваме
\[
\int_a^b m\,dx
\leq
\int_a^b f(x)\,dx
\leq
\int_a^b M\,dx.
\]
Следователно
\[
m(b-a)
\leq
\int_a^b f(x)\,dx
\leq
M(b-a).
\]
Тъй като $b-a>0$,
\[
m\leq
\frac{1}{b-a}\int_a^b f(x)\,dx
\leq M.
\]
Непрекъснатата функция приема всички стойности между своя минимум и максимум. Следователно съществува $c\in[a,b]$, за което
\[
f(c)=\frac{1}{b-a}\int_a^b f(x)\,dx.
\]
Умножаваме по $b-a$ и получаваме
\[
\int_a^b f(x)\,dx=f(c)(b-a).
\]
\end{proof}

Числото
\[
\frac{1}{b-a}\int_a^b f(x)\,dx
\]
се нарича средна стойност на функцията $f$ върху $[a,b]$. Теоремата показва, че за непрекъсната функция тази средна стойност действително се достига като стойност $f(c)$ в поне една точка от интервала.

\section{Теорема на Нютон--Лайбниц}

\begin{keythm}[Теорема на Нютон--Лайбниц]
\label{thm:newtonleibniz}
Нека $f\colon[a,b]\to\mathbb{R}$ е непрекъсната и
\[
F(x)=\int_a^x f(t)\,dt,
\qquad x\in[a,b].
\]
Тогава $F$ е диференцируема във всяка вътрешна точка на $[a,b]$ и
\[
F'(x)=f(x).
\]
В краищата $a$ и $b$ същото равенство се разбира чрез съответната едностранна производна.
\end{keythm}

\begin{proof}
Нека $x_0\in[a,b]$ е фиксирана точка и $h\ne0$ е такова, че
\[
x_0+h\in[a,b].
\]
От адитивността на интеграла следва
\[
F(x_0+h)-F(x_0)
=
\int_{x_0}^{x_0+h}f(t)\,dt.
\]
Следователно
\[
\frac{F(x_0+h)-F(x_0)}{h}
=
\frac{1}{h}\int_{x_0}^{x_0+h}f(t)\,dt.
\]
Изваждаме $f(x_0)$ и използваме, че
\[
f(x_0)=\frac{1}{h}\int_{x_0}^{x_0+h}f(x_0)\,dt.
\]
Получаваме
\[
\frac{F(x_0+h)-F(x_0)}{h}-f(x_0)
=
\frac{1}{h}\int_{x_0}^{x_0+h}
\bigl(f(t)-f(x_0)\bigr)\,dt.
\]

Нека $\varepsilon>0$. От непрекъснатостта на $f$ в $x_0$ съществува $\delta>0$, такова че
\[
|t-x_0|<\delta
\quad\Longrightarrow\quad
|f(t)-f(x_0)|<\varepsilon.
\]
Ако $|h|<\delta$, то всяка точка $t$ между $x_0$ и $x_0+h$ удовлетворява
\[
|t-x_0|\leq |h|<\delta.
\]

При $h>0$, използвайки неравенството за интеграла, имаме
\begin{align*}
\left|
\frac{F(x_0+h)-F(x_0)}{h}-f(x_0)
\right|
&\leq
\frac{1}{h}
\int_{x_0}^{x_0+h}|f(t)-f(x_0)|\,dt\\
&\leq
\frac{1}{h}\int_{x_0}^{x_0+h}\varepsilon\,dt
=
\varepsilon.
\end{align*}
При $h<0$ аналогично,
\begin{align*}
\left|
\frac{F(x_0+h)-F(x_0)}{h}-f(x_0)
\right|
&\leq
\frac{1}{|h|}
\int_{x_0+h}^{x_0}|f(t)-f(x_0)|\,dt\\
&\leq
\frac{1}{|h|}
\int_{x_0+h}^{x_0}\varepsilon\,dt
=
\varepsilon.
\end{align*}

Следователно
\[
\lim_{h\to0}
\frac{F(x_0+h)-F(x_0)}{h}
=
f(x_0).
\]
Значи $F'(x_0)=f(x_0)$. Тъй като $x_0$ е произволна точка, твърдението е доказано.
\end{proof}

Теоремата показва, че функцията
\[
x\longmapsto\int_a^x f(t)\,dt
\]
е примитивна на непрекъснатата функция $f$.

\subsection{Изчисляване на определени интеграли}

\begin{cor}[Формула на Нютон--Лайбниц]
Нека $f\colon[a,b]\to\mathbb{R}$ е непрекъсната и $\Phi$ е примитивна на $f$ върху $[a,b]$, тоест
\[
\Phi'(x)=f(x).
\]
Тогава
\[
\int_a^b f(x)\,dx=\Phi(b)-\Phi(a).
\]
Обичайното означение е
\[
\int_a^b f(x)\,dx=\left.\Phi(x)\right|_a^b.
\]
\end{cor}

\begin{proof}
По Теорема \ref{thm:newtonleibniz} функцията
\[
F(x)=\int_a^x f(t)\,dt
\]
също е примитивна на $f$. Следователно за някоя константа $C$ е изпълнено
\[
\Phi(x)=F(x)+C.
\]
При $x=a$ получаваме
\[
\Phi(a)=F(a)+C=C,
\]
защото
\[
F(a)=\int_a^a f(t)\,dt=0.
\]
Следователно
\[
\Phi(x)=\int_a^x f(t)\,dt+\Phi(a),
\]
или
\[
\int_a^x f(t)\,dt=\Phi(x)-\Phi(a).
\]
При $x=b$ получаваме търсената формула.
\end{proof}

\begin{example}
Да се изчисли
\[
\int_0^1 3x^2\,dx.
\]
Примитивна на $3x^2$ е $\Phi(x)=x^3$. Следователно
\[
\int_0^1 3x^2\,dx
=
\left.x^3\right|_0^1
=
1^3-0^3
=
1.
\]
\end{example}

Практическият алгоритъм за пресмятане на определен интеграл на непрекъсната функция е следният:

\begin{enumerate}
\item Намира се примитивна функция $\Phi$ на подинтегралната функция $f$.
\item Изчисляват се стойностите $\Phi(a)$ и $\Phi(b)$.
\item Прилага се формулата
\[
\int_a^b f(x)\,dx=\Phi(b)-\Phi(a).
\]
\end{enumerate}

\section{Изпитен фокус}

\begin{itemize}
\item Да се дефинират разбиване на $[a,b]$, делящи точки и диаметър $d(\tau)$.
\item Да се дефинират $m_i$, $M_i$, малка сума $s_f(\tau)$ и голяма сума $S_f(\tau)$ на Дарбу.
\item Да се възпроизведе доказателството, че при доразбиване малките суми не намаляват, а големите не нарастват.
\item Да се дефинират долен и горен интеграл на Дарбу и интегруемост по Риман чрез равенството им.
\item Да се докаже критерият на Дарбу:
\[
f\text{ е интегруема}
\Longleftrightarrow
\forall\varepsilon>0\ \exists\tau:
S_f(\tau)-s_f(\tau)<\varepsilon.
\]
\item Да се обясни доказателството за интегруемост на непрекъсната функция: теорема на Кантор, избор на разбиване с малък диаметър и оценка на $S_f(\tau)-s_f(\tau)$.
\item Да се изброят линейност, адитивност, позитивност, интегриране на неравенства, смяна на границите и неравенството с $|f|$.
\item Да се докаже теоремата за средната стойност за интеграли чрез минимум, максимум, интегриране на неравенства и свойството за междинните стойности.
\item Да се докаже теоремата на Нютон--Лайбниц чрез разликата
\[
\frac{1}{h}\int_x^{x+h}\bigl(f(t)-f(x)\bigr)\,dt
\]
и непрекъснатостта на $f$.
\item Да се използва формулата
\[
\int_a^b f(x)\,dx=\Phi(b)-\Phi(a)
\]
за изчисляване на определен интеграл чрез примитивна функция.
\end{itemize}


\chapter{Итерационни методи за решаване на нелинейни уравнения.}
% AUTO-GENERATED by scripts/build_topics.py — do not edit.
% Edit topics/bodies/topic_34.tex or topics/preamble.tex instead.
%% Placeholder. Replace with the written chapter.
\textit{Източници:} PDF \texttt{33.pdf}, снимки \texttt{34та тема}

\subsection*{Анотация от конспекта}
Итерационни методи за решаване на нелинейни уравнения. Да се дефинира понятието неподвижна точка на изображението  и да се докаже, че ако  е непрекъснато изображение на интервала [a, b] в себе си, то  има поне една неподвижна точка в [a, b]. Да се покаже, че решаването на уравнението f(x) = 0 може да се сведе към намиране на неподвижна точка. Да се дефинира понятието свиващо изображение и да се докаже, че ако  е непрекъснато изображение на интервала [a, b] в себе си и е свиващо с константа на Липшиц q< 1, то: а) уравнението x= (x) има единствен корен  в [a, b]; б) редицата \{xn\} от последователни приближения (при произволно x0[a, b] и xn+1 = (xn), n = 0, 1, 2, ..., клони към  при n→, като |xn – |  (b – a)qn, за всяко n. Да се получи като следствие, че ако  е корен на уравнението x= (x) и  има непрекъсната производна в околност U на , за която |′(ξ)| < 1, то при достатъчно добро начално приближение x0 итерационният процес, породен от , е сходящ със скоростта на геометрична прогресия. Да се дефинира понятието ред на сходимост. Да се дадат геометрична илюстрация, формула за последователните приближения и ред на сходимост при: метод на хордите, метод на секущите и метод на Нютон. Да се докаже, че при метода на хордите сходимостта е със скоростта на геометричната прогресия (при условие, че коренът е отделен в достатъчно малък интервал). Литература: [1], [2], [16].


\chapter{Случайни величини с дискретни разпределения – биномно, геометрично, Поасоново.}
%% Drafted from the mapped corpus by scripts/draft_topics.py.
% Model: gpt-5.6-terra; source characters: 16734.
\section{Дискретни случайни величини и разпределения}

\begin{defn}[Случайна величина]
Върху вероятностно пространство $(\Omega,\mathcal A,\mathbb P)$ случайна величина е измерима функция $X:\Omega\to\mathbb R$, тоест $\{\omega:X(\omega)\le x\}\in\mathcal A$ за всяко $x\in\mathbb R$. Ако $\omega\in\Omega$ е елементарен изход, то стойността $X(\omega)$ зависи от реализирания изход $\omega$.
\end{defn}

\begin{defn}
Случайната величина $X$ се нарича \emph{дискретна}, ако множеството от стойностите, които тя приема, е крайно или изброимо.
\end{defn}

Нека $H_1,H_2,\ldots$ е разлагане на $\Omega$ и нека $x_1,x_2,\ldots$ са различни реални числа. Дискретната случайна величина може да се представи като
\[
X(\omega)=\sum_j x_j I_{H_j}(\omega),
\]
където
\[
I_{H_j}(\omega)=
\begin{cases}
1, & \omega\in H_j,\\
0, & \omega\notin H_j
\end{cases}
\]
е индикаторът на множеството $H_j$.

\begin{defn}
Нека $X$ е дискретна случайна величина, която приема стойностите
\[
x_1,x_2,\ldots .
\]
\emph{Дискретно разпределение на $X$} наричаме съответствието
\[
\begin{array}{c|ccccc}
X & x_1 & x_2 & \cdots & x_j & \cdots\\
\hline
\mathbb{P} & p_1 & p_2 & \cdots & p_j & \cdots
\end{array},
\]
където
\[
p_j=\mathbb{P}(X=x_j).
\]
\end{defn}

Вероятностите в таблицата трябва да удовлетворяват две основни условия:
\[
p_j\geq 0\quad\text{за всяко }j,
\qquad
\sum_j p_j=1.
\]
Първото условие изразява неотрицателността на вероятността, а второто --- нормираността: случайната величина трябва да приеме някоя от възможните си стойности с вероятност $1$.

По-общо, вероятността е функция, дефинирана върху $\sigma$-алгебра $\mathcal{A}$ от събития. За всяко $A\in\mathcal{A}$ е изпълнено
\[
\mathbb{P}(A)\geq 0,
\qquad
\mathbb{P}(\Omega)=1.
\]
Ако $A$ и $B$ са несъвместими събития, то
\[
\mathbb{P}(A\cup B)=\mathbb{P}(A)+\mathbb{P}(B).
\]
От свойствата на вероятността следва и монотонността:
\[
A\subseteq B\quad\Longrightarrow\quad \mathbb{P}(A)\leq \mathbb{P}(B).
\]

\begin{defn}[Математическо очакване]
За дискретна случайна величина с вероятности $p_j=\mathbb{P}(X=x_j)$ математическото очакване е
\[
\mathbb{E}X=\sum_j x_jp_j,
\]
когато редът $\sum_j |x_j|p_j$ е сходящ.
\end{defn} \begin{defn}[Дисперсия]
Дисперсията се определя чрез
\[
\mathbb{D}X=\mathbb{E}\bigl[(X-\mathbb{E}X)^2\bigr]
=\sum_j (x_j-\mathbb{E}X)^2p_j.
\]
При $\mathbb E X^2<\infty$ разкриването на квадрата и линейността на очакването дават равенството
\[
\mathbb{D}X=\mathbb{E}X^2-(\mathbb{E}X)^2.
\]
\end{defn}

\section{Пораждаща функция на дискретна случайна величина}

\begin{defn}
Нека $X$ е неотрицателна целочислена дискретна случайна величина, за която
\[
\mathbb{P}(X=k)=p_k,\qquad k=0,1,2,\ldots .
\]
\emph{Пораждаща функция} на $X$ се нарича функцията
\[
G_X(s)=\mathbb{E}s^X=\sum_{k=0}^{\infty}p_ks^k,
\]
за онези стойности на $s$, за които редът е сходящ.
\end{defn}

Коефициентът пред $s^k$ в реда за $G_X(s)$ е именно вероятността
\[
\mathbb{P}(X=k).
\]
Следователно пораждащата функция кодира цялото дискретно разпределение.

\begin{prop}
Нека $X$ е неотрицателна целочислена случайна величина. За пораждащата й функция са в сила следните свойства, като производните при $1$ са леви граници и за формулата за дисперсията се изисква $\mathbb E X^2<\infty$:
\[
G_X(1)=1,
\]
\[
G_X'(1)=\mathbb{E}X,
\]
а при съществуваща втора производна
\[
\mathbb{D}X=G_X''(1)+G_X'(1)-\bigl(G_X'(1)\bigr)^2.
\]
\end{prop}

\begin{proof}
Действително, от нормираността на разпределението непосредствено се получава
\[
G_X(1)=\sum_{k=0}^{\infty}p_k=1.
\]
Първата производна има вида
\[
G_X'(s)=\sum_{k=1}^{\infty}kp_ks^{k-1},
\]
За $0\le s<1$ почленното диференциране е допустимо вътре в радиуса на сходимост. При $s\uparrow1$ монотонната сходимост на неотрицателните членове дава граничната производна, откъдето следва формулата за математическото очакване. Аналогично,
\[
G_X''(1)=\sum_{k=0}^{\infty}k(k-1)p_k.
\]
Тъй като
\[
k^2=k(k-1)+k,
\]
получаваме
\[
\mathbb{E}X^2=G_X''(1)+G_X'(1),
\]
а след това и формулата за дисперсията.
\end{proof}

Пораждащите функции са удобни, защото чрез едно и също пресмятане дават както вероятностите, така и основните моменти на разпределението.

\section{Биномно разпределение}

\subsection{Схема на Бернули}

Разглеждаме опит с два възможни изхода: успех и неуспех. Нека вероятността за успех е $p$, а вероятността за неуспех е
\[
q=1-p.
\]
Случайна величина $Y$, която приема стойност $1$ при успех и стойност $0$ при неуспех, има разпределение на Бернули. За нея
\[
\mathbb{P}(Y=1)=p,
\qquad
\mathbb{P}(Y=0)=q.
\]
Следователно
\[
\mathbb{E}Y=p,
\qquad
\mathbb{D}Y=pq.
\]

Биномното разпределение възниква при фиксиран брой независими опити на Бернули, когато вероятността за успех е една и съща във всеки опит.

\begin{defn}
Нека са извършени $n$ независими опита на Бернули с вероятност за успех $p$. Ако $X$ е броят на успехите в тези $n$ опита, казваме, че $X$ има \emph{биномно разпределение} с параметри $n$ и $p$ и пишем
\[
X\sim\mathrm{Bin}(n,p).
\]
\end{defn}

Случайната величина $X$ приема стойностите
\[
0,1,\ldots,n.
\]
За да се получат точно $k$ успеха, трябва да се изберат $k$ от $n$-те позиции, в които успехът да настъпи. За всяка такава подредба вероятността е
\[
p^kq^{n-k}.
\]
Броят на възможните избори е $\binom{n}{k}$, затова
\[
\mathbb{P}(X=k)=\binom{n}{k}p^kq^{n-k},
\qquad k=0,1,\ldots,n.
\]

\begin{prop}
Вероятностите на биномното разпределение са нормировани:
\[
\sum_{k=0}^{n}\binom{n}{k}p^kq^{n-k}=1.
\]
\end{prop}

\begin{proof}
По биномната формула на Нютон
\[
\sum_{k=0}^{n}\binom{n}{k}p^kq^{n-k}=(p+q)^n.
\]
Тъй като $p+q=1$, получаваме
\[
(p+q)^n=1.
\]
\end{proof}

\begin{example}
Нека една монета има вероятност $p$ да се падне ези и се хвърля $n$ пъти независимо. Ако $X$ е броят на падналите се ези, то
\[
X\sim\mathrm{Bin}(n,p).
\]
При правилна монета $p=\frac12$.
\end{example}

\begin{example}
Правилен зар се хвърля $5$ пъти. Нека $X$ е броят на получените шестици. Всеки опит има успех с вероятност
\[
p=\frac16,
\]
следователно
\[
X\sim\mathrm{Bin}\left(5,\frac16\right).
\]
Например вероятността да се получат точно две шестици е
\[
\mathbb{P}(X=2)
=
\binom52\left(\frac16\right)^2
\left(\frac56\right)^3.
\]
\end{example}

\subsection{Пораждаща функция, математическо очакване и дисперсия}

За $X\sim\mathrm{Bin}(n,p)$ пораждащата функция е
\[
G_X(s)
=
\sum_{k=0}^{n}\binom{n}{k}p^kq^{n-k}s^k.
\]
Като обединим $p^k$ и $s^k$, получаваме
\[
G_X(s)
=
\sum_{k=0}^{n}\binom{n}{k}(ps)^kq^{n-k}.
\]
Отново по биномната формула на Нютон:
\[
G_X(s)=(q+ps)^n.
\]

Първата производна е
\[
G_X'(s)=np(q+ps)^{n-1}.
\]
Следователно
\[
\mathbb{E}X=G_X'(1)
=np(p+q)^{n-1}=np.
\]

Втората производна е
\[
G_X''(s)=n(n-1)p^2(q+ps)^{n-2}.
\]
При $s=1$:
\[
G_X''(1)=n(n-1)p^2.
\]
Затова
\[
\begin{aligned}
\mathbb{D}X
&=G_X''(1)+G_X'(1)-\bigl(G_X'(1)\bigr)^2\\
&=n(n-1)p^2+np-n^2p^2\\
&=np(1-p)=npq.
\end{aligned}
\]

\begin{keythm}
Ако
\[
X\sim\mathrm{Bin}(n,p),
\]
то
\[
\mathbb{E}X=np,
\qquad
\mathbb{D}X=npq=np(1-p).
\]
\end{keythm}

\begin{proof}
Същите формули следват и от представянето
\[
X=X_1+\cdots+X_n,
\]
където $X_i$ са независими случайни величини на Бернули с параметър $p$. Тогава
\[
\mathbb{E}X
=
\mathbb{E}X_1+\cdots+\mathbb{E}X_n
=np,
\]
а по независимостта
\[
\mathbb{D}X
=
\mathbb{D}X_1+\cdots+\mathbb{D}X_n
=npq.
\]
\end{proof}

\section{Геометрично разпределение}

\subsection{Дефиниция и интерпретация}

Геометричното разпределение описва повтаряне на независими опити на Бернули до настъпването на първия успех. Броят на опитите не е предварително ограничен.

\begin{defn}
Нека се извършват независими опити на Бернули с вероятност за успех $p$, където
\[
0<p<1.
\]
Нека $X$ е броят на неуспехите преди първия успех. Тогава $X$ има \emph{геометрично разпределение} с параметър $p$, означавано с
\[
X\sim\mathrm{Geo}(p),
\]
ако
\[
\mathbb{P}(X=k)=q^kp,
\qquad k=0,1,2,\ldots,
\]
където $q=1-p$.
\end{defn}

Събитието $\{X=k\}$ означава, че първите $k$ опита са неуспешни, а следващият опит е успешен. Поради независимостта вероятността му е
\[
q^kp.
\]

\begin{prop}
Вероятностите на геометричното разпределение са нормировани:
\[
\sum_{k=0}^{\infty}q^kp=1.
\]
\end{prop}

\begin{proof}
Тъй като $0<q<1$, редът е геометрична прогресия:
\[
\sum_{k=0}^{\infty}q^kp
=
p\sum_{k=0}^{\infty}q^k
=
p\frac{1}{1-q}.
\]
Понеже $1-q=p$, следва
\[
p\frac{1}{1-q}=\frac{p}{p}=1.
\]
\end{proof}

\begin{example}
Хвърляме правилен зар, докато се падне шестица. Ако $X$ е броят на неуспешните хвърляния преди първата шестица, то
\[
X\sim\mathrm{Geo}\left(\frac16\right).
\]
Например събитието $X=4$ означава, че са се получили четири нешестици, последвани от шестица, и има вероятност
\[
\mathbb{P}(X=4)
=
\left(\frac56\right)^4\frac16.
\]
\end{example}

\begin{supp}[Бележка]
В някои формулировки случайна величина се дефинира като броя на \emph{опитите} до първия успех. Тогава тя приема стойности $1,2,\ldots$ и се различава с единица от използваната тук величина. В настоящата глава следваме дефиницията, при която $X$ брои неуспехите преди първия успех; затова стойностите са $0,1,2,\ldots$ и
\[
\mathbb{P}(X=k)=q^kp.
\]
\end{supp}

\subsection{Пораждаща функция, математическо очакване и дисперсия}

За $X\sim\mathrm{Geo}(p)$ имаме
\[
\begin{aligned}
G_X(s)
&=\sum_{k=0}^{\infty}\mathbb{P}(X=k)s^k\\
&=\sum_{k=0}^{\infty}q^kps^k\\
&=p\sum_{k=0}^{\infty}(qs)^k.
\end{aligned}
\]
Като използваме сумата на геометрична прогресия, получаваме
\[
G_X(s)=\frac{p}{1-qs}.
\]

Първата производна е
\[
G_X'(s)=\frac{pq}{(1-qs)^2}.
\]
Следователно
\[
\mathbb{E}X
=
G_X'(1)
=
\frac{pq}{(1-q)^2}
=
\frac{pq}{p^2}
=
\frac{q}{p}.
\]

Втората производна е
\[
G_X''(s)=\frac{2pq^2}{(1-qs)^3}.
\]
Следователно
\[
G_X''(1)=\frac{2pq^2}{(1-q)^3}
=\frac{2q^2}{p^2}.
\]
Прилагаме формулата за дисперсията:
\[
\begin{aligned}
\mathbb{D}X
&=G_X''(1)+G_X'(1)-\bigl(G_X'(1)\bigr)^2\\
&=\frac{2q^2}{p^2}+\frac{q}{p}-\frac{q^2}{p^2}\\
&=\frac{q^2+pq}{p^2}\\
&=\frac{q(p+q)}{p^2}\\
&=\frac{q}{p^2}.
\end{aligned}
\]

\begin{keythm}
Ако
\[
X\sim\mathrm{Geo}(p),
\]
където $X$ е броят неуспехи преди първия успех, то
\[
\mathbb{E}X=\frac{q}{p},
\qquad
\mathbb{D}X=\frac{q}{p^2}.
\]
\end{keythm}

\begin{proof}
За $0<p\le1$ и $q=1-p$ пораждащата функция е $G(s)=p/(1-qs)$. Получаваме $G'(1)=pq/(1-q)^2=q/p$ и $G''(1)=2pq^2/(1-q)^3=2q^2/p^2$. По формулата за дисперсията
\[\mathbb D X=\frac{2q^2}{p^2}+\frac qp-\frac{q^2}{p^2}=\frac{q(q+p)}{p^2}=\frac q{p^2}.\]
При $p=1$ величината е тъждествено нула и формулите също са валидни.
\end{proof}

\section{Поасоново разпределение}

\subsection{Дефиниция и задачи, в които възниква}

Поасоновото разпределение се използва за броя на редки събития при голям брой възможности за настъпването им. То възниква като граничен случай на биномното разпределение, когато броят на опитите е голям, вероятността за успех във всеки опит е малка, а средният брой успехи остава краен.

\begin{defn}
Казваме, че случайната величина $X$ има \emph{Поасоново разпределение} с параметър $\lambda>0$, ако приема стойностите
\[
0,1,2,\ldots
\]
с вероятности
\[
\mathbb{P}(X=k)=\frac{e^{-\lambda}\lambda^k}{k!},
\qquad k=0,1,2,\ldots .
\]
Означаваме
\[
X\sim\mathrm{Po}(\lambda).
\]
\end{defn}

Параметърът $\lambda$ е средният брой настъпвания на разглежданото събитие.

\begin{example}
Дисплей има $1\,000\,000$ пиксела. Вероятността един пиксел да изгори е
\[
p=\frac{1}{1\,000\,000}.
\]
При $n=1\,000\,000$ възможности средният брой изгорели пиксели е
\[
\lambda=np=1.
\]
Броят $X$ на изгорелите пиксели може да се моделира с Поасоново разпределение с параметър $\lambda=1$. Например
\[
\mathbb{P}(X=2)
=
\frac{1^2e^{-1}}{2!}
=
\frac{1}{2e}.
\]
\end{example}

\begin{prop}
Вероятностите на Поасоновото разпределение са нормировани:
\[
\sum_{k=0}^{\infty}\frac{e^{-\lambda}\lambda^k}{k!}=1.
\]
\end{prop}

\begin{proof}
Използваме реда за експоненциалната функция:
\[
e^\lambda=\sum_{k=0}^{\infty}\frac{\lambda^k}{k!}.
\]
Тогава
\[
\sum_{k=0}^{\infty}\frac{e^{-\lambda}\lambda^k}{k!}
=
e^{-\lambda}
\sum_{k=0}^{\infty}\frac{\lambda^k}{k!}
=
e^{-\lambda}e^\lambda
=
1.
\]
\end{proof}

\subsection{Връзка с биномното разпределение}

Нека
\[
X\sim\mathrm{Bin}(n,p).
\]
Когато $n$ е голямо, $p$ е малко и
\[
np\longrightarrow\lambda,
\]
биномната вероятност
\[
\mathbb{P}(X=k)=\binom{n}{k}p^k(1-p)^{n-k}
\]
се стреми към
\[
\frac{e^{-\lambda}\lambda^k}{k!},
\qquad k=0,1,2,\ldots .
\]
Това обяснява защо Поасоновото разпределение е подходящ модел за редки успехи сред много на брой опити.

\subsection{Пораждаща функция, математическо очакване и дисперсия}

Нека
\[
X\sim\mathrm{Po}(\lambda).
\]
Тогава
\[
\begin{aligned}
G_X(s)
&=\sum_{k=0}^{\infty}\mathbb{P}(X=k)s^k\\
&=\sum_{k=0}^{\infty}
\frac{e^{-\lambda}\lambda^k}{k!}s^k\\
&=e^{-\lambda}\sum_{k=0}^{\infty}
\frac{(\lambda s)^k}{k!}.
\end{aligned}
\]
От реда за експоненциалната функция следва
\[
G_X(s)=e^{-\lambda}e^{\lambda s}
=e^{\lambda(s-1)}.
\]

Първата производна е
\[
G_X'(s)=\lambda e^{\lambda(s-1)}.
\]
Следователно
\[
\mathbb{E}X=G_X'(1)=\lambda.
\]

Втората производна е
\[
G_X''(s)=\lambda^2e^{\lambda(s-1)}.
\]
Следователно
\[
G_X''(1)=\lambda^2.
\]
Тогава
\[
\begin{aligned}
\mathbb{D}X
&=G_X''(1)+G_X'(1)-\bigl(G_X'(1)\bigr)^2\\
&=\lambda^2+\lambda-\lambda^2\\
&=\lambda.
\end{aligned}
\]

\begin{keythm}
Ако
\[
X\sim\mathrm{Po}(\lambda),
\]
то
\[
\mathbb{E}X=\lambda,
\qquad
\mathbb{D}X=\lambda.
\]
\end{keythm}

\begin{proof}
Пораждащата функция е $G(s)=e^{\lambda(s-1)}$. Следователно $G'(1)=\lambda$ и $G''(1)=\lambda^2$. От доказаните формули за моментите следва $\mathbb E X=\lambda$ и $\mathbb D X=\lambda^2+\lambda-\lambda^2=\lambda$.
\end{proof}

Равенството
\[
\mathbb{E}X=\mathbb{D}X=\lambda
\]
е характерна особеност на Поасоновото разпределение.

\section{Сравнение на разглежданите разпределения}

Трите разпределения описват различни въпроси, свързани с опити с успех и неуспех.

При биномното разпределение броят опити е фиксиран, а се броят успехите:
\[
X\sim\mathrm{Bin}(n,p),
\qquad
\mathbb{P}(X=k)=\binom{n}{k}p^kq^{n-k}.
\]
То е подходящо за въпроса: ``Колко успеха ще има в точно $n$ независими опита?''

При геометричното разпределение опитите продължават до първия успех, а се броят неуспехите преди него:
\[
X\sim\mathrm{Geo}(p),
\qquad
\mathbb{P}(X=k)=q^kp.
\]
То е подходящо за въпроса: ``Колко неуспеха ще има преди първия успех?''

При Поасоновото разпределение се брои броят на редки събития:
\[
X\sim\mathrm{Po}(\lambda),
\qquad
\mathbb{P}(X=k)=\frac{e^{-\lambda}\lambda^k}{k!}.
\]
То е подходящо, когато има много възможности за настъпване на събитие, всяка с малка вероятност, и средният брой настъпвания е $\lambda$.

Основните формули са:
\[
\begin{array}{c|c|c|c}
\text{Разпределение} & \text{Параметри} & \mathbb{E}X & \mathbb{D}X\\
\hline
\mathrm{Bin}(n,p) & n,\ p,\ q=1-p & np & npq\\
\mathrm{Geo}(p) & p,\ q=1-p & \dfrac{q}{p} & \dfrac{q}{p^2}\\
\mathrm{Po}(\lambda) & \lambda>0 & \lambda & \lambda
\end{array}
\]

Съответните пораждащи функции са:
\[
G_X(s)=
\begin{cases}
(q+ps)^n, & X\sim\mathrm{Bin}(n,p),\\[4pt]
\dfrac{p}{1-qs}, & X\sim\mathrm{Geo}(p),\\[10pt]
e^{\lambda(s-1)}, & X\sim\mathrm{Po}(\lambda).
\end{cases}
\]

\section{Изпитен фокус}

\begin{itemize}
\item Да се дефинира дискретна случайна величина и да се представи нейното разпределение чрез стойности $x_j$ и вероятности $p_j=\mathbb{P}(X=x_j)$.

\item Да се посочат условията
\[
p_j\geq0,
\qquad
\sum_jp_j=1,
\]
като израз на неотрицателността и нормираността на вероятностите.

\item Да се дефинира пораждащата функция
\[
G_X(s)=\sum_{k\geq0}\mathbb{P}(X=k)s^k
\]
и да се използват формулите
\[
G_X(1)=1,\qquad
\mathbb{E}X=G_X'(1),\qquad
\mathbb{D}X=G_X''(1)+G_X'(1)-\bigl(G_X'(1)\bigr)^2.
\]

\item Да се дефинира биномното разпределение като брой успехи в $n$ независими опита на Бернули:
\[
\mathbb{P}(X=k)=\binom{n}{k}p^kq^{n-k}.
\]
Да се изведе
\[
G_X(s)=(q+ps)^n,
\qquad
\mathbb{E}X=np,
\qquad
\mathbb{D}X=npq.
\]

\item Да се обясни геометричното разпределение като брой неуспехи преди първия успех:
\[
\mathbb{P}(X=k)=q^kp.
\]
Да се изведе
\[
G_X(s)=\frac{p}{1-qs},
\qquad
\mathbb{E}X=\frac{q}{p},
\qquad
\mathbb{D}X=\frac{q}{p^2}.
\]

\item Да се дефинира Поасоновото разпределение:
\[
\mathbb{P}(X=k)=\frac{e^{-\lambda}\lambda^k}{k!}.
\]
Да се изведе
\[
G_X(s)=e^{\lambda(s-1)},
\qquad
\mathbb{E}X=\lambda,
\qquad
\mathbb{D}X=\lambda.
\]

\item Да се разпознават типичните задачи: фиксиран брой опити за биномното разпределение, чакане до първи успех за геометричното разпределение и редки събития при много възможности за Поасоновото разпределение.

\item Да се възпроизведат идеите за проверка на нормираността: биномната формула на Нютон при биномното разпределение, сумата на геометрична прогресия при геометричното разпределение и редът за $e^x$ при Поасоновото разпределение.
\end{itemize}

\end{document}
